% !TeX root = ../Book2.tex
\BookPart{振荡积分、乘子与限制问题}{b2:part:05}
% 本编导读（章级标题，不计入编号）
\BookPartGuide

一个模为一的指数因子可以使积分很小，原因是不同位置的贡献相互抵消。相位变化越快，抵消通常越强；驻点与曲率则决定哪些位置留下主要贡献。本编先研究这种几何怎样产生衰减，再用它分析频率乘法和曲面限制。选读的多线性章进一步比较不同方向之间的相互作用。

\ProseParagraphBreak
主要依赖图为
\[
\begin{tikzcd}[column sep=1.2em,row sep=large,font=\small]
 \text{第\ref{b2:ch:21}章：振荡积分}
    \arrow[r] \arrow[dr]
 &\text{第\ref{b2:ch:23}章：限制问题}
    \arrow[d] \arrow[dr]&\\
 \text{第\ref{b2:ch:22}章：乘子}
    \arrow[r]
 &\text{第\ref{b2:ch:24}章：Bochner--Riesz}
 &\text{\chapref{b2:ch:25}：多线性}
\end{tikzcd}
\]
\chapref{b2:ch:22}用第五编的奇异积分与平方函数研究频率乘法算子，把符号的光滑性转为空间范数估计。\chapref{b2:ch:21}到\chapref{b2:ch:24}的箭头表示核渐近；\chapref{b2:ch:23}到\chapref{b2:ch:24}表示频率薄壳估计。从混合指数的薄壳估计走向同指数乘子界，还要控制不同薄壳的叠加。

\ProseParagraphBreak
\chapref{b2:ch:25}从方向横截的模型走向一般多线性限制与解耦，其中深层定理以证明概要介绍。计算径向核时，可回查\appref{b2:app:A}的\cref{b2:thm:A2-sphere-fourier,b2:thm:A2-bessel-asymptotic,b2:thm:A2-bochner-riesz}，它们分别给出球面 Fourier 公式、固定实阶 Bessel 渐近及球内幂权核。进一步的应用有两个方向：\appref{b2:app:E}研究随时间的色散，\appref{b2:app:G}用相位追踪奇性的位置和方向。

\ProseParagraphBreak
本编的比较读物包括 Elias M. Stein 的\cite{Stein1993Harmonic}与 Christopher D. Sogge 的\cite{Sogge2017Fourier}；多线性和解耦的正式来源为\cite{Bennett2006Multilinear,BourgainDemeter2015Proof}。遇到一个新的估计，先核对变量维数、输入输出指数、曲率或横截下界及允许的尺度损失，再判断它能否嵌入当前证明。

% !TeX root = ../../Book2.tex
\chapter{振荡积分}
\label{b2:ch:21}
\BookChapterNumber{b2:ch:21}
\begin{chapterabstract}
振荡积分的大小取决于哪些部分能够相互抵消。相位没有驻点时，分部积分把抵消转为衰减，但边界和振幅导数必须一并控制；相位有驻点时，局部尺度、Hesse 矩阵与振幅消失阶决定主要贡献。本章从这些标量问题出发，再用 \(TT^*\) 方法处理允许输入自由振荡的算子，区分普通 Hesse 矩阵、混合 Hesse 矩阵和相位驻点方程的三种作用。
\end{chapterabstract}

\begin{learninggoals}\begin{itemize}
\item\label{b2:item:21-goal-1}区分端点项、驻点贡献和非驻相余项，说明支撑与振幅导数如何进入估计常数。
\item\label{b2:item:21-goal-2}由次水平集和分部积分理解 van der Corput 衰减，辨别有限型、单调性与参数尺度的作用。
\item\label{b2:item:21-goal-3}通过 Morse 坐标和复 Gaussian 计算驻相系数，核对行列式、符号、Jacobi 因子及参数余项。
\item\label{b2:item:21-goal-4}用混合导数下界、\(TT^*\) 与 Schur 估计证明局部 \(L^2\) 衰减，并判断局部化和参数一致性的适用范围。
\end{itemize}\end{learninggoals}

积分 \(\int e^{2\pi\mathrm i\lambda\phi}\cdot a\) 中的指数因子始终具有模一。对它直接取绝对值，只能得到振幅的总量，无法看到 \(\lambda\) 增大后的抵消。一个可行的起点是观察相位的一阶导数：它不为零时，可以把指数写成一个导数；它等于零时，则应先找出附近的自然尺度。由此产生非驻相法与驻相法，但两种方法都需要知道振幅在哪里非零。

\ProseParagraphBreak

本章先在一维把边界项完整保留下来，再证明 van der Corput 引理。随后由复 Gaussian 与光滑换元得到多维驻相展开，并讨论临界点碰撞、振幅随尺度变化和曲面曲率。一般上界与渐近首项承担不同任务：一个统一上界未必给出非零主项，而几个各自非零的驻点贡献也可能相互抵消。

\ProseParagraphBreak

进入\secref{b2:sec:2104}后，对象从一个固定振幅的积分变成作用于整个 \(L^2\) 单位球的算子。输入可以选择自己的相位，不能把标量驻相公式直接用于每一个输入。\(TT^*\) 把问题改写成两个输出点之间的相位差，混合 Hesse 矩阵由此出现。该矩阵在一个点可逆，只先保证足够小的邻域内有统一估计；把小片相加与直接在整个支撑上使用同一个梯度下界，是不同的论证。

\ProseParagraphBreak

分部积分从相位导数中提取衰减，Morse 换元则把非退化驻点附近的相位化成二次型；两种方法分别处理远离驻点和接近驻点的区域。选读的\secref{b2:sec:2105}进一步用相位记录位置与频率之间的对应关系，介绍 Fourier 积分算子。前七道习题围绕主线安排；最后一道用实际波包检验半波传播的符号与尺度。

\ProseParagraphBreak

Elias M. Stein 的《Harmonic Analysis: Real-Variable Methods, Orthogonality, and Oscillatory Integrals》\cite{Stein1993Harmonic}可用于继续阅读振荡积分及其在调和分析中的应用。本章保留指数中的 \(2\pi\)，使用其他归一化的资料时，应同时换算驻相系数和传播参数。

% !TeX root = ../../Book2.tex
\section{一维振荡积分}
\label{b2:sec:2101}

正积分的大小由振幅总量决定，振荡积分还受到相位变化的制约。对有限区间 \([\alpha,\beta]\)、实值相位 \(\phi\) 和复值振幅 \(a\)，记
\[
 I_\lambda(a,\phi)
 \coloneqq\int_\alpha^\beta
 e^{2\pi\mathrm i\lambda\phi(t)}\cdot a(t)\,\mathrm dt,
 \quad \lambda>0.
\]
本章主要考察 \(\lambda\) 很大的情形；本节的一阶和高阶导数估计对任意正 \(\lambda\) 都成立。这一点在把导数下界吸收到参数后尤其有用，因为新的有效参数未必大于一。

\ProseParagraphBreak

直接取绝对值只能得到 \(|I_\lambda|\leq\|a\|_1\)，完全看不到振荡。要获得衰减，须比较邻近位置的相位：导数远离零时，相邻波动可以通过分部积分抵消；导数接近零的区域则不能这样处理。van der Corput 方法把这两个区域分开，并用高阶导数控制后一部分的长度。

\subsection{分部积分}
\label{b2:subsec:210101}

先在足够光滑的情形中辨认边界项。设 \(\phi\in C^2([\alpha,\beta];\mathbb R)\)、\(|\phi'|\geq c>0\)，并设 \(a:[\alpha,\beta]\to\mathbb C\) 绝对连续。此时 \(a/\phi'\) 也绝对连续，由
\[
 e^{2\pi\mathrm i\lambda\phi(t)}
 =\frac1{2\pi\mathrm i\lambda\phi'(t)}
   \frac{\mathrm d}{\mathrm dt}e^{2\pi\mathrm i\lambda\phi(t)}
\]
得到
\[
 \begin{aligned}
 I_\lambda(a,\phi)
 &=\left[
 \frac{e^{2\pi\mathrm i\lambda\phi(t)}\cdot a(t)}
      {2\pi\mathrm i\lambda\phi'(t)}
 \right]_\alpha^\beta\\
 &\hspace{1em}
 -\frac1{2\pi\mathrm i\lambda}
 \int_\alpha^\beta e^{2\pi\mathrm i\lambda\phi(t)}
       \cdot\left(\frac a{\phi'}\right)'(t)\,\mathrm dt.
 \end{aligned}
\]
商的导数是 \(a'/\phi'-a\phi''/(\phi')^2\)，所以
\[
 |I_\lambda(a,\phi)|
 \leq\frac1{2\pi\lambda}
 \left[
 \frac{|a(\alpha)|+|a(\beta)|+\|a'\|_1}{c}
 +\frac1{c^2}\int_\alpha^\beta|a(t)|\cdot|\phi''(t)|\,\mathrm dt
 \right].
\]
这个公式清楚地区分了振幅变化、相位导数变化和两个端点的贡献。相位本身很大并不提供下界；给 \(\phi\) 加上一个任意大的常数，只会改变整体复相位。

\ProseParagraphBreak

振幅在端点消失时，边界项为零；否则不能将其省略。例如
\[
 \int_0^1e^{2\pi\mathrm i\lambda t}\,\mathrm dt
 =\frac{e^{2\pi\mathrm i\lambda}-1}{2\pi\mathrm i\lambda}.
\]
沿 \(\lambda=n+1/2\) 的序列，其模为 \(1/(\pi\lambda)\)。这里相位没有驻点，振幅甚至是常数，但仍只有一阶衰减。将一个光滑函数硬截断在闭区间上，与使振幅光滑地消失在端点附近，是两种不同的操作。

\subsection{非驻相}
\label{b2:subsec:210102}

若振幅有紧支撑，导数转移不会产生边界项，可以反复使用刚才的恒等式。为便于后面处理参数族，下面把估计常数所需的振幅范数也写出来。

\Needspace{6\baselineskip}
\begin{proposition}[非驻相衰减]\label{b2:prop:21-nonstationary-one}
设 \(\alpha<\beta\)，\(a\in C_c^\infty((\alpha,\beta);\mathbb C)\)，\(\phi\in C^\infty((\alpha,\beta);\mathbb R)\)。设 \(V\) 是 \(\operatorname{supp}a\) 的一个邻域，在 \(V\) 上有 \(|\phi'|\geq c>0\)。对任意整数 \(N\geq1\)、任意 \(\lambda>0\)，
\[
 \left|\int_\alpha^\beta e^{2\pi\mathrm i\lambda\phi(t)}
                \cdot a(t)\,\mathrm dt\right|
 \leq C_N\lambda^{-N}\sum_{j=0}^N\|a^{(j)}\|_1.
\]
若在 \(V\) 上的 \(\phi',\ldots,\phi^{(N+1)}\) 绝对值不超过 \(M\)，则 \(C_N\) 只依赖 \(N,c^{-1},M\)，不依赖 \(\lambda\) 或振幅。所有 \(L^1\) 范数均在 \((\alpha,\beta)\) 上取。
\end{proposition}
\begin{proof}
令 \(a_0=a\)，并在 \(V\) 上递归定义
\[
 a_{j+1}=-\left(\frac{a_j}{\phi'}\right)'.
\]
因 \(a_j\) 的支撑始终包含于 \(\operatorname{supp}a\)，可将这些函数在 \(V\) 外延为零，得到光滑紧支撑函数。每次积分分部都没有边界项，因此
\[
 I_\lambda(a,\phi)
 =(2\pi\mathrm i\lambda)^{-N}
   \int_\alpha^\beta e^{2\pi\mathrm i\lambda\phi(t)}
          \cdot a_N(t)\,\mathrm dt.
\]
递归中的每次求导至多增加一阶振幅导数；系数则由 \(\phi'\) 的倒数与相位的各阶导数相乘、相加得到。归纳可写成
\[
 a_N=\sum_{j=0}^N c_{N,j}\cdot a^{(j)},
\]
其中每个 \(c_{N,j}\) 只含至多 \(N+1\) 阶的相位导数及 \((\phi')^{-1}\) 的有限次幂。题设给这些系数一致上界，故
\(\|a_N\|_1\leq C\sum_{j=0}^N\|a^{(j)}\|_1\)。
取绝对值即得结论。
\end{proof}

这里“非驻相”指相位一阶导数在振幅支撑附近有正下界。相位在别处可以有驻点，因为那里没有振幅。反之，一个参数族中的驻点即使尚未进入支撑，只要到支撑的距离趋于零，相位导数下界就可能变小，估计常数也会相应增大。

\subsection{van der Corput 引理}
\label{b2:subsec:210103}

高阶导数下界允许一阶导数消失，因此不能在整个区间上使用倒导数。先建立只需要 \(C^1\) 正则性的一阶估计。单调导数未必经典可微，直接写出 \(\phi''\) 并积分并不合法；改用相位本身作变量，可以避开这一问题。

\Needspace{6\baselineskip}
\begin{lemma}[一阶导数估计]\label{b2:lem:21-monotone-first-derivative}
设 \(\alpha<\beta\)，\(\phi\in C^1([\alpha,\beta];\mathbb R)\)，\(|\phi'|\geq c>0\)。若连续函数 \(1/\phi'\) 在该区间的总变差为有限数 \(V\)，则对任意 \(\lambda>0\) 与子区间 \([u,v]\subset[\alpha,\beta]\)，
\[
 \left|\int_u^v e^{2\pi\mathrm i\lambda\phi(t)}\,\mathrm dt\right|
 \leq\frac{2c^{-1}+V}{2\pi\lambda}.
\]
这里总变差是所有有限分割上
\(\sum_j|q(t_j)-q(t_{j-1})|\) 的上确界，其中 \(q=1/\phi'\)。
特别地，若 \(\phi'\) 单调，则上述积分不超过 \(3/(2\pi\lambda c)\)。
\end{lemma}
\begin{proof}
导数连续且不为零，所以它符号恒定，\(\phi\) 严格单调。相位换元 \(s=\phi(t)\) 将积分写成
\[
 \int_{\phi(u)}^{\phi(v)}
       e^{2\pi\mathrm i\lambda s}\cdot q_*(s)\,\mathrm ds,
 \quad
 q_*(s)=\frac1{\phi'(\phi^{-1}(s))}.
\]
若上下限倒置，先反转积分方向，绝对值不变。函数 \(q_*\) 连续，\(\|q_*\|_\infty\leq c^{-1}\)，且沿单调换元的总变差不超过 \(V\)。

在相位区间上取网格大小趋于零的分割，用 \(q_*\) 在分割点的值作分段线性插值 \(q_n\)。连续性保证 \(q_n\) 一致收敛到 \(q_*\)；插值不增大上确界范数，而且
\[
 \int|q_n'(s)|\,\mathrm ds
 =\sum_j|q_*(s_j)-q_*(s_{j-1})|\leq V.
\]
每个 \(q_n\) 绝对连续，普通积分分部给出
\[
 \left|\int e^{2\pi\mathrm i\lambda s}
                 \cdot q_n(s)\,\mathrm ds\right|
 \leq\frac{2\|q_n\|_\infty+\|q_n'\|_1}{2\pi\lambda}
 \leq\frac{2c^{-1}+V}{2\pi\lambda}.
\]
在固定的有限相位区间上，积分随一致逼近收敛，得到 \(q_*\) 的结论。若 \(\phi'\) 单调，倒数也单调，并且符号固定、绝对值至多为 \(c^{-1}\)，故其总变差不超过 \(c^{-1}\)。
\end{proof}

第二个准备事实说明：一个高阶导数始终不小的函数，不可能在很长的集合上保持很小。所控制的是集合的总长度，不要求这个集合本身是一个区间。

\Needspace{6\baselineskip}
\begin{lemma}[次水平集的长度]\label{b2:lem:21-sublevel-measure}
设 \(r\geq1\) 为整数，\(\alpha<\beta\)，\(h\in C^r([\alpha,\beta];\mathbb R)\)，且 \(|h^{(r)}|\geq c>0\)。对任意 \(\delta>0\)，有
\[
 \left|\{\,t\in[\alpha,\beta]:|h(t)|\leq\delta\,\}\right|
 \leq 2r\left(\frac{\delta}{c}\right)^{1/r}.
\]
\end{lemma}
\begin{proof}
记左侧集合为 \(E\)，其测度为 \(L\)。\(L=0\) 时无需证明。设 \(L>0\)，则 \(E\) 是非空紧集。函数
\[
 F(t)=|E\cap[\alpha,t]|
\]
连续且为 \(1\)-Lipschitz 函数，从 \(0\) 增至 \(L\)。取 \(t_0=\min E\)、\(t_r=\max E\)。对 \(1\leq j\leq r-1\)，令 \(t_j\) 为满足 \(F(t)\geq jL/r\) 的最小点。连续性给出 \(F(t_j)=jL/r\)；该点必在 \(E\) 中，否则其邻域不交 \(E\)，\(F\) 在此邻域为常数，与最小性矛盾。端点也分别满足 \(F(t_0)=0\)、\(F(t_r)=L\)。因此，对 \(0\leq i<j\leq r\)，
\[
 t_j-t_i\geq F(t_j)-F(t_i)=\frac{(j-i)L}{r}.
\]
这些点既落在 \(E\) 中，又具有可计算的间距。

用这 \(r+1\) 个点作 Lagrange 插值多项式
\[
 P(t)=\sum_{j=0}^r h(t_j)
       \prod_{\substack{0\leq\ell\leq r\\\ell\ne j}}
       \frac{t-t_\ell}{t_j-t_\ell}.
\]
它的 \(t^r\) 系数为
\[
 A=\sum_{j=0}^r
       \frac{h(t_j)}
       {\prod_{\ell\ne j}(t_j-t_\ell)}.
\]
因为 \(|h(t_j)|\leq\delta\)，且
\[
 \left|\prod_{\ell\ne j}(t_j-t_\ell)\right|
 \geq (L/r)^r\,j!(r-j)!,
\]
所以
\[
 |A|
 \leq\delta\left(\frac rL\right)^r
       \sum_{j=0}^r\frac1{j!(r-j)!}
 =\frac{\delta}{r!}\left(\frac{2r}{L}\right)^r.
\]
函数 \(h-P\) 在 \(t_0,\ldots,t_r\) 为零。先在每两个相邻零点之间应用 Rolle 定理，得到至少 \(r\) 个一阶导数零点；如此重复 \(r\) 次，得到某个 \(\theta\in(t_0,t_r)\)，使
\[
 0=(h-P)^{(r)}(\theta)
   =h^{(r)}(\theta)-r!A.
\]
于是 \(c\leq r!|A|\leq\delta(2r/L)^r\)，整理即得所述长度界。
\end{proof}

现在可以证明\term[van-der-corput-yinli]{van der Corput 引理}[van der Corput lemma]。一阶情形需要控制倒导数的变化；二阶及以上情形则由最高阶导数下界自动限制低阶导数反复转向的次数。

\Needspace{6\baselineskip}
\begin{theorem}[van der Corput 引理]\label{b2:thm:21-vdc}
设 \(\alpha<\beta\)，\(k\geq1\) 为整数，\(c>0\)、\(\lambda>0\)，\(\phi\in C^k([\alpha,\beta];\mathbb R)\)，\(a:[\alpha,\beta]\to\mathbb C\) 绝对连续。若 \(k=1\)，假设 \(|\phi'|\geq c\) 且 \(\phi'\) 单调；若 \(k\geq2\)，假设 \(|\phi^{(k)}|\geq c\)。则
\[
 \left|\int_\alpha^\beta
 e^{2\pi\mathrm i\lambda\phi(t)}\cdot a(t)\,\mathrm dt\right|
 \leq C_k(\lambda c)^{-1/k}
       \bigl(\|a\|_\infty+\|a'\|_1\bigr),
\]
其中 \(C_k\) 只依赖 \(k\)，不依赖区间长度、相位、振幅、\(\lambda\) 或 \(c\)。
\end{theorem}
\begin{proof}
先取振幅恒为一。\(k=1\) 时直接用\cref{b2:lem:21-monotone-first-derivative}。以下设 \(k\geq2\)，取任意阈值 \(\delta>0\)，将区间分为
\[
 E_\delta=\{\,t:|\phi'(t)|\leq\delta\,\}
 \quad\text{和}\quad
 [\alpha,\beta]\setminus E_\delta.
\]
对 \(h=\phi'\)、\(r=k-1\) 应用\cref{b2:lem:21-sublevel-measure}，得到
\[
 \left|\int_{E_\delta}e^{2\pi\mathrm i\lambda\phi(t)}
                         \,\mathrm dt\right|
 \leq |E_\delta|
 \leq 2(k-1)\left(\frac{\delta}{c}\right)^{1/(k-1)}.
\]

补集上的积分使用一阶估计，但须先控制其单调区间数。若 \(k=2\)，\(\phi''\) 连续且不为零，因而符号固定，\(\phi'\) 在整个区间单调。若 \(k>2\)，\(\phi''\) 至多有 \(k-2\) 个互异零点：否则从 \(k-1\) 个零点出发，连续应用 \(k-2\) 次 Rolle 定理，会产生一个 \(\phi^{(k)}\) 的零点，矛盾。因此可将区间分成至多 \(k-1\) 个 \(\phi'\) 单调的区间。在每一片内，\(\phi'>\delta\) 和 \(\phi'<-\delta\) 各至多占一个区间。补集于是由至多 \(2(k-1)\) 个区间组成；端点的开闭不影响积分。在每个这样的区间上，先对内部闭子区间使用一阶引理，再令端点趋向边界，得到
\[
 \left|\int_{[\alpha,\beta]\setminus E_\delta}
             e^{2\pi\mathrm i\lambda\phi(t)}\,\mathrm dt\right|
 \leq \frac{C_k}{\lambda\delta}.
\]
两部分合并为
\[
 \left|\int_\alpha^\beta
          e^{2\pi\mathrm i\lambda\phi(t)}\,\mathrm dt\right|
 \leq C_k\left[
       \left(\frac{\delta}{c}\right)^{1/(k-1)}
       +\frac1{\lambda\delta}\right].
\]
选
\(\delta=c^{1/k}\lambda^{-(k-1)/k}\)，两项都等于
\((\lambda c)^{-1/k}\) 的常数倍。因此该估计适用于全部正的 \(\lambda,c\)，包括有效参数 \(\lambda c<1\) 的情形。

最后处理振幅。上述结论在任意子区间上仍成立，因为导数假设在限制后不变。令
\[
 F(t)=\int_\alpha^t e^{2\pi\mathrm i\lambda\phi(s)}\,\mathrm ds.
\]
则 \(\|F\|_\infty\leq C_k(\lambda c)^{-1/k}\)。函数 \(a\) 绝对连续，积分分部给出
\[
 I_\lambda(a,\phi)
 =a(\beta)F(\beta)-\int_\alpha^\beta a'(t)F(t)\,\mathrm dt,
\]
从而得到所述上界。
\end{proof}

一阶中的单调性可以替换为受控的总变差，但不能仅保留 \(|\phi'|\geq c\)。沿用一阶引理，并作同样的振幅积分分部，可得
\[
 |I_\lambda(a,\phi)|
 \leq \frac{2c^{-1}+\operatorname{Var}(1/\phi')}{2\pi\lambda}
       \bigl(\|a\|_\infty+\|a'\|_1\bigr).
\]
因此关键不是导数是否有零点这一件事，还包括倒导数能否反复变化而积累很大的总变差。

\ProseParagraphBreak

下面的光滑相位族给出真实的失败例子。对正整数 \(N\)，令
\[
 q_N(s)=\frac{2+\cos(2\pi Ns)}3,
 \quad
 T_N(s)=\int_0^s q_N(v)\,\mathrm dv,
 \quad 0\leq s\leq1.
\]
有 \(1/3\leq q_N\leq1\)、\(T_N(0)=0\)、\(T_N(1)=2/3\)，所以 \(T_N\) 有光滑反函数
\(\phi_N:[0,2/3]\to[0,1]\)，并且
\[
 1\leq\phi_N'(t)
 =\frac1{q_N(\phi_N(t))}\leq3.
\]
然而在固定区间上取常数振幅、取 \(\lambda=N\)，换元 \(t=T_N(s)\) 给出
\[
 \begin{aligned}
 \int_0^{2/3}e^{2\pi\mathrm iN\phi_N(t)}\,\mathrm dt
 &=\int_0^1e^{2\pi\mathrm iNs}\cdot
                  \frac{2+\cos(2\pi Ns)}3\,\mathrm ds\\
 &=\frac16.
 \end{aligned}
\]
最后一步把余弦写成两个复指数；所有非零整数频率的全周期积分为零，只留下常数项。故不存在只依赖导数下界的统一 \(C/N\) 估计。此例的倒导数总变差恰为 \(4N/3\)，也解释了为什么变差版本没有产生矛盾。

\subsection{衰减估计}
\label{b2:subsec:210104}

导数下界 \(c\) 与高频参数始终以有效参数 \(\lambda c\) 一起出现。即使 \(\lambda\geq1\)，当相位接近退化时也可能有 \(\lambda c<1\)；上一定理仍然成立，只是此时其右侧可能不如平凡界 \(\|a\|_1\)。实际应用可取两者的较小值。不能先只证明参数大于一的定理，再不加说明地将它用于 \(\lambda c<1\)。

\ProseParagraphBreak

相位 \(t^k\) 展示了幂次的尺度来源。换元 \(t=\lambda^{-1/k}u\) 后，
\[
 \int e^{2\pi\mathrm i\lambda t^k}\cdot a(t)\,\mathrm dt
 =\lambda^{-1/k}
   \int e^{2\pi\mathrm i u^k}\cdot a(\lambda^{-1/k}u)\,\mathrm du.
\]
缩放确定了观察驻点的自然邻域。要取得非零渐近首项，还需要控制扩张积分区间上的振荡尾部，并排除对称性抵消。二次情形的非零主项将在\secref{b2:sec:2102}中算出，从而把自然尺度与实际衰减阶数联系起来。

\ProseParagraphBreak

有限阶导数条件也可局部使用。若 \(\phi\) 在紧支撑附近光滑，且每一点至少有某个 \(1\leq j\leq k\) 满足 \(|\phi^{(j)}|\geq c>0\)，则连续性给出有限个邻域，使每一片上某个固定阶导数有下界 \(c/2\)。乘上光滑分割后，\(j\geq2\) 的片用 van der Corput 引理；\(j=1\) 的片用紧支撑非驻相估计，无须额外假设该片导数单调。对 \(\lambda\geq1\)，有限项相加得到 \(O(\lambda^{-1/k})\)。覆盖、分割函数导数、相位的有限阶导数与下界都参与常数，不能仅记录一个抽象的“有限型”条件。

\ProseParagraphBreak

这种统一上界在驻点碰撞时比逐点二次展开更稳定。取
\[
 I(\lambda,s)=\int_{\mathbb R}
 e^{2\pi\mathrm i\lambda(t^3-st)}\cdot a(t)\,\mathrm dt,
 \quad a\in C_c^\infty(\mathbb R;\mathbb C).
\]
三阶导数恒为六，因此 van der Corput 引理给出
\(|I(\lambda,s)|\leq C_a\lambda^{-1/3}\)，其中 \(C_a\) 与 \(s\in\mathbb R\) 无关。若 \(s>0\)，两个驻点是 \(\pm\sqrt{s/3}\)，各点的二阶导数大小为 \(2\sqrt{3s}\)。后面的二次驻相公式因而会出现 \(s^{-1/4}\lambda^{-1/2}\) 的系数，不能在 \(s\downarrow0\) 时保持统一。

\ProseParagraphBreak

换元 \(t=\lambda^{-1/3}u\) 后，指数中的相位成为
\(u^3-(s\lambda^{2/3})u\)。因此碰撞区域的自然参数是 \(s\lambda^{2/3}\)，相应范围是 \(s\asymp\lambda^{-2/3}\)。在此范围，驻点间距 \(\sqrt{s}\) 与三次振荡尺度 \(\lambda^{-1/3}\) 相当，两个局部二次模型已不能统一分离。对于固定的小正 \(s\)，二次驻点的局部宽度约为 \(\lambda^{-1/2}s^{-1/4}\)；要求它远小于 \(\sqrt{s}\)，也恰好得到 \(s\lambda^{2/3}\gg1\)。

\ProseParagraphBreak

\cref{b2:fig:21-critical-collision}画的是相位导数，而不是原相位；其横轴交点直接显示驻点的数目和碰撞。统一三阶估计无需写出 Airy 函数，尺度比较已经说明了分别使用二次近似的范围。

\begin{figure}[htbp]
\centering
\begin{tikzpicture}[x=1cm,y=.6cm]
\foreach \shift/\offset/\paneltitle in {-3.2/.65/{\(s<0\)},0/0/{\(s=0\)},3.2/-.65/{\(s>0\)}}{
\begin{scope}[shift={(\shift,0)}]
\draw[->] (-1.25,0)--(1.35,0) node[right]{\(t\)};
\draw[AnalysisBlue,thick,domain=-1.05:1.05,samples=45] plot(\x,{1.3*\x*\x+\offset});
\node at (0,2.35){\paneltitle};
\end{scope}}
\fill (0,0) circle(1.3pt);
\fill (2.493,0) circle(1.3pt);
\fill (3.907,0) circle(1.3pt);
\node[below,font=\small] at (2.493,-.95){\(-\sqrt{s/3}\)};
\node[below,font=\small] at (3.907,-.95){\(\sqrt{s/3}\)};
\node at (0,-2.7){碰撞尺度：\(\sqrt{s}\asymp\lambda^{-1/3}\)，即 \(s\lambda^{2/3}\asymp1\)};
\end{tikzpicture}
\caption[三次相位的驻点碰撞]{\(\phi_s'(t)=3t^2-s\) 的三个符号情形，不按比例。横轴交点是驻点；当两个交点的间距降到三次振荡尺度时，不能再把它们作为两个统一非退化的二次模型。}
\label{b2:fig:21-critical-collision}
\end{figure}

非驻相区域仍需连同支撑一起判断。若 \(s<0\)，则
\(\phi_s'(t)=3t^2+|s|\geq|s|\)，固定负参数给出任意阶衰减，但其常数随 \(|s|\downarrow0\) 增长。若振幅支撑满足 \(|t|\geq\varepsilon>0\)，则对足够小的 \(|s|\)，一阶导数仍有统一正下界，可以恢复对这些参数一致的快速衰减。退化的影响由相位、支撑与参数范围共同决定。

% !TeX root = ../../Book2.tex
\section{驻相法}
\label{b2:sec:2102}

一阶导数消失后，非驻相分部积分的分母发生奇异，主要贡献便可能集中在驻点附近。若驻点 \(x_0\) 的二次项非退化，位移 \(h\) 使相位变化约为 \(\lambda h^{\mathsf T}Hh/2\)；在各特征值固定且非零时，变化为常数量级的位移尺度是 \(\lambda^{-1/2}\)。这解释了每个非退化方向为何贡献半个衰减幂次。本节先把远离驻点的部分除去，再用可逆换元将驻点附近的相位精确化为二次型，最后计算其常数与余项。

\subsection{非退化临界点}
\label{b2:subsec:210201}

\Needspace{6\baselineskip}
\begin{definition}\label{b2:def:21-critical}
设整数 \(d\geq1\)，\(U\subset\mathbb R^d\) 为开集，\(\phi\in C^2(U;\mathbb R)\)，\(x_0\in U\)。若 \(\nabla\phi(x_0)=0\)，称 \(x_0\) 为驻点；若同时
\(\det D^2\phi(x_0)\ne0\)，称其为\term[fei-tuihua-zhudian]{非退化驻点}[nondegenerate critical point]。
\end{definition}
逆函数定理用于映射 \(\nabla\phi\)，说明非退化驻点在一个小邻域内是唯一的。反过来，孤立驻点未必非退化：\(t^4\) 的原点就是例子。第二种情形的衰减阶由更高阶项决定，不能套用二次驻相公式。

\ProseParagraphBreak
设相位在紧集 \(K\) 的一个开邻域内光滑。若 \(K\) 中所有驻点非退化，则其中驻点只有有限个。否则一列互异驻点在 \(K\) 内有聚点，梯度的连续性使聚点仍为驻点，与局部孤立性矛盾。于是可以分别计算各驻点贡献，再对剩余部分使用非驻相法；不同贡献的复相位可能抵消，因此各项上界之和只是上界。

\ProseParagraphBreak

先把后一步所需的多维工具在此证明。不能直接在多个坐标中同时除以各个偏导数，因为梯度非零并不保证每一个偏导数都非零。沿梯度方向微分，则只需除以整个梯度的平方长度。

\Needspace{6\baselineskip}
\begin{proposition}[多维非驻相衰减]\label{b2:prop:21-nonstationary-multi}
设整数 \(d\geq1\)，\(U\subset\mathbb R^d\) 开，\(a\in C_c^\infty(U;\mathbb C)\)，\(\phi\in C^\infty(U;\mathbb R)\)。设在 \(K=\operatorname{supp}a\) 的一个邻域 \(W\Subset U\) 内有 \(|\nabla\phi|\geq c>0\)。对每个整数 \(N\geq1\) 及 \(\lambda\geq1\)，
\[
 \left|\int_U e^{2\pi\mathrm i\lambda\phi(x)}
                   \cdot a(x)\,\mathrm dx\right|
 \leq C_N\lambda^{-N}
       \sum_{|\alpha|\leq N}\|\partial^\alpha a\|_{L^1(U)}.
\]
常数 \(C_N\) 只依赖 \(d,N,c^{-1}\) 及 \(\phi\) 在 \(W\) 上至多 \(N+1\) 阶导数的上界。因而固定支撑、相同梯度下界与这些导数的一致界给出一致的估计。
\end{proposition}
\begin{proof}
在 \(W\) 内令
\[
 V(x)=\frac{\nabla\phi(x)}{|\nabla\phi(x)|^2},
 \quad
 Ta=-\operatorname{div}(Va).
\]
因为 \(V\cdot\nabla\phi=1\)，有
\[
 V\cdot\nabla e^{2\pi\mathrm i\lambda\phi}
 =2\pi\mathrm i\lambda e^{2\pi\mathrm i\lambda\phi}.
\]
振幅紧支于 \(W\)，散度积分没有边界项，故
\[
 \int e^{2\pi\mathrm i\lambda\phi}\cdot a
 =(2\pi\mathrm i\lambda)^{-1}
   \int e^{2\pi\mathrm i\lambda\phi}\cdot Ta.
\]
这里使用双线性积分分部，\(T\) 是 \(V\cdot\nabla\) 的形式转置；没有对振幅取复共轭。具体地，
\[
 Ta=-V\cdot\nabla a-(\operatorname{div}V)a.
\]
每个 \(T^ja\) 的支撑仍包含于 \(K\)，故可以重复 \(N\) 次，得到
\[
 \int e^{2\pi\mathrm i\lambda\phi}\cdot a
 =(2\pi\mathrm i\lambda)^{-N}
   \int e^{2\pi\mathrm i\lambda\phi}\cdot T^Na.
\]
用乘积公式归纳，\(T^Na\) 是 \(\partial^\alpha a\)、\(|\alpha|\leq N\)，乘上 \(V\) 的至多 \(N\) 阶导数所组成的有限和。对 \(V=\nabla\phi/|\nabla\phi|^2\) 求导，分母只有 \(|\nabla\phi|\) 的负幂，分子只用到 \(\phi\) 的至多 \(N+1\) 阶导数。题设因此给出这些系数的统一上界，取各项的 \(L^1\) 范数即得结论。
\end{proof}

在驻点附近选一个等于一的光滑截断后，剩余振幅仍紧支撑；若其支撑上没有驻点，紧性便给出梯度的正下界。这正是下面把远离驻点的部分吸收到任意阶余项中的依据。对有边界的积分域，必须另行保留边界项，本命题不适用于直接的示性函数截断。

\subsection{Morse 归约与相位二次型}
\label{b2:subsec:210202}

对实对称可逆矩阵 \(H\)，记其正特征值数减负特征值数为 \(\operatorname{sgn}H\)，称为这里使用的符号差。Taylor 展开只说明相位近似于一个二次型；\term[morse-yinli]{Morse 引理}[Morse lemma]给出更强的局部事实：通过改变坐标，可以把所有高阶项吸收到坐标与振幅中，使新相位精确成为二次型。

\Needspace{6\baselineskip}
\begin{lemma}[Morse 局部归约]\label{b2:lem:21-morse-normal-form}
设整数 \(d\geq1\)，\(U\subset\mathbb R^d\) 开，\(\phi\in C^\infty(U;\mathbb R)\)，\(x_0\in U\) 满足 \(\nabla\phi(x_0)=0\) 且 \(H_0=D^2\phi(x_0)\) 可逆。则存在 \(x_0\) 的邻域 \(U_0\)、原点的邻域 \(V_0\)，以及光滑微分同胚 \(u=\kappa(x):U_0\longrightarrow V_0\)，使 \(\kappa(x_0)=0\) 且
\[
 \phi(x)=\phi(x_0)+\frac12u^{\mathsf T}Su,
 \quad
 S=\operatorname{diag}(\varepsilon_1,\ldots,\varepsilon_d),
 \quad \varepsilon_j\in\{1,-1\}.
\]
其中 \(\sum_j\varepsilon_j=\operatorname{sgn}H_0\)。记逆映射为 \(x=X(u)\)，则
\[
 |\det DX(0)|=|\det H_0|^{-1/2}.
\]
\end{lemma}
\begin{proof}
先平移 \(x_0\) 到原点，再作一个固定正交换元，将 \(H_0\) 化为
\(\operatorname{diag}(\mu_1,\ldots,\mu_d)\)，其中每个 \(\mu_j\neq0\)。这两次换元的 Jacobi 因子绝对值都为一。暂用 \(x\) 表示变换后的坐标。在足够小的球内，线段 \(tx\)、\(0\leq t\leq1\)，都留在相位定义域中。带积分余项的 Taylor 公式给出精确恒等式
\[
 \phi(x)-\phi(0)=\frac12x^{\mathsf T}G(x)x,
 \quad
 G(x)=2\int_0^1(1-t)D^2\phi(tx)\,\mathrm dt.
\]
矩阵 \(G(x)\) 实对称、光滑，且 \(G(0)=H_0\)。这里把高阶项保留在矩阵函数中，并没有把它们丢进误差。

先看二维的配方。若
\[
 G(x)=
 \begin{pmatrix}p(x)&q(x)\\q(x)&r(x)\end{pmatrix},
\]
则在 \(p\neq0\) 的区域有
\[
 x^{\mathsf T}G(x)x
 =p(x)\left(x_1+\frac{q(x)}{p(x)}x_2\right)^2
  +\left(r(x)-\frac{q(x)^2}{p(x)}\right)x_2^2.
\]
在原点，两个新系数分别为 \(\mu_1,\mu_2\)，因而缩小邻域后都非零且符号不变。把它们绝对值的平方根分别乘入括号，就得到两个新的坐标表达式。系数依赖于原变量 \(x\) 并不破坏这个逐点恒等式；剩下要核验的是这些表达式合起来是否真的可作坐标。

一般维数按同一规则逐步配方。第一步以 \(p_1(x)=G_{11}(x)\) 为主元，提出
\[
 p_1(x)\left(x_1+\sum_{k>1}
                     \frac{G_{1k}(x)}{p_1(x)}x_k\right)^2,
\]
余下变量的对称系数矩阵为
\[
 G^{(2)}_{kl}(x)
 =G_{kl}(x)-\frac{G_{k1}(x)G_{1l}(x)}{p_1(x)}
 \quad(k,l>1).
\]
由于 \(G(0)\) 已经对角化，\(p_1(0)=\mu_1\)，且
\(G^{(2)}(0)=\operatorname{diag}(\mu_2,\ldots,\mu_d)\)。下一步以 \(G^{(2)}_{22}\) 为主元，重复同一配方。归纳可知，第 \(j\) 步主元 \(p_j\) 在原点等于 \(\mu_j\)，所以每一步都可以在一个更小邻域内继续。只有有限步，最后取这些邻域的公共部分，所有主元均不为零、符号固定，且全部商与平方根都光滑。

因此得到一个光滑可逆的上三角矩阵函数 \(B(x)\)，使
\[
 G(x)=B(x)^{\mathsf T}S B(x),
 \quad
 S=\operatorname{diag}\bigl(\operatorname{sgn}\mu_1,\ldots,
                            \operatorname{sgn}\mu_d\bigr).
\]
这里 \(B\) 的对角元为 \(\sqrt{|p_j|}\)，其余系数来自上述配方。令
\(u=\kappa(x)=B(x)x\)，则已经有精确的相位恒等式。还需注意：矩阵 \(B(x)\) 逐点可逆，本身不等于映射 \(x\mapsto B(x)x\) 已知可逆。真正的依据是在原点
\[
 D\kappa(0)=B(0),
\]
因为对 \(B\) 求导后乘上的 \(x\) 在原点为零。矩阵 \(B(0)\) 可逆，逆函数定理于是给出光滑局部逆映射。

在驻点处对相位恒等式求两次导数，得到
\[
 H_0=D\kappa(0)^{\mathsf T}S D\kappa(0).
\]
取行列式绝对值，并用 \(|\det S|=1\)，便有
\[
 |\det D\kappa(0)|^2=|\det H_0|,
 \quad
 |\det DX(0)|=|\det H_0|^{-1/2}.
\]
恢复原来的正交坐标后，这一绝对值公式不变。符号差由原点的对角化已经确定，故等于 \(\operatorname{sgn}H_0\)。
\end{proof}

预先对角化的作用是保证固定消元顺序中的主元不消失。仅有矩阵可逆不够：\(\left(\begin{smallmatrix}0&1\\1&0\end{smallmatrix}\right)\) 可逆，却不能直接把左上角作为第一个主元。本证明只在一个固定矩阵上选择正交坐标，随后用连续性工作；没有要求在一个矩阵族中光滑地选择所有特征向量。

\ProseParagraphBreak

例如
\[
 \phi(x_1,x_2)=\frac12\left[x_1^2+2x_1x_2+(2+x_1)x_2^2\right]
 =\frac12\left[(x_1+x_2)^2+(1+x_1)x_2^2\right].
\]
在原点附近可取 \(u_1=x_1+x_2\)、\(u_2=\sqrt{1+x_1}\,x_2\)。其在原点的导数行列式为一，故确为局部坐标，并将含三次项的相位精确化为 \((u_1^2+u_2^2)/2\)。原 Hesse 矩阵为
\(\left(\begin{smallmatrix}1&1\\1&2\end{smallmatrix}\right)\)，行列式也为一，与 Jacobi 公式相符。

\ProseParagraphBreak

换元以后振幅成为 \(a(X(u))|\det DX(u)|\)，不能只替换函数自变量而忽略体积因子。一个 Hesse 特征值变小时，对应方向的驻点尺度变宽，\(|\det H_0|^{-1/2}\) 也可能增大；这是非退化下界必须进入估计常数的原因。

\subsection{驻相引理}
\label{b2:subsec:210203}

二次振荡核没有绝对可积性，不能直接把实 Gaussian 积分中的正参数换成纯虚数。先在实部为正的参数域计算，再取边界极限，便能同时确定常数与复平方根的分支。

\Needspace{6\baselineskip}
\begin{proposition}[复 Gaussian 的 Fourier 变换]\label{b2:prop:21-complex-gaussian}
设 \(a\in\mathbb C\)、\(\operatorname{Re}a>0\)、\(\xi\in\mathbb R\)。则下列积分绝对收敛，且
\[
 F(a,\xi)=\int_{\mathbb R}e^{-\pi a x^2}e^{-2\pi\mathrm ix\xi}\,\mathrm dx
 =a^{-1/2}e^{-\pi\xi^2/a}.
\]
这里 \(a^{-1/2}=|a|^{-1/2}e^{-\mathrm i\arg(a)/2}\)，
\(-\pi/2<\arg(a)<\pi/2\)；在正实轴上取正平方根。
\end{proposition}
\begin{proof}
固定 \(\varepsilon>0\)，写 \(a=\varepsilon+\mathrm it\)。任意有限 \(t,\xi\) 区间上的参数导数由多项式乘 \(e^{-\pi\varepsilon x^2}\) 控制，故可以在积分号下求导；下面的实轴分部积分也没有无穷远边界项。积分
\(\partial_x(e^{-\pi a x^2-2\pi\mathrm ix\xi})\) 得到
\(\int x e^{-\pi a x^2-2\pi\mathrm ix\xi}\,\mathrm dx=-\mathrm i\xi F/a\)，所以
\(\partial_\xi F=-(2\pi\xi/a)F\)。解这个实变量一阶方程可得
\(F(a,\xi)=F(a,0)e^{-\pi\xi^2/a}\)。

再积分 \(\partial_x(xe^{-\pi a x^2})\)，得到
\(\int x^2e^{-\pi a x^2}\,\mathrm dx=F(a,0)/(2\pi a)\)。于是
\(\partial_tF(\varepsilon+\mathrm it,0)=-\mathrm iF/(2a)\)。实 Gaussian 公式给初值 \(F(\varepsilon,0)=\varepsilon^{-1/2}\)，上述连续分支的 \((\varepsilon+\mathrm it)^{-1/2}\) 满足同一方程和初值，故两者相等。整个计算只沿实轴作分部积分，复平方根由初值连续确定。
\end{proof}
由于 \(\operatorname{Re}(1/a)>0\)，此公式还有
\( |F(a,\xi)|\leq |a|^{-1/2}\)。在纯二次振荡的正则化中取
\(a=\varepsilon/\pi-\mathrm i\lambda\varepsilon_j\)、\(\varepsilon_j=\pm1\)，则模上界为 \(\lambda^{-1/2}\)，与正则化参数无关。高维对角二次型逐坐标相乘即可；乘上可积的 \(\widehat b\) 后，这正是下面支配收敛所需的共同控制。

先把二次型与光滑振幅的配对写成一个可绝对积分的 Fourier 表达式。这不仅给出首项，也给出每个更高阶系数及可检查的余项。

\Needspace{6\baselineskip}
\begin{proposition}[二次驻相展开]\label{b2:prop:21-quadratic-expansion}
设整数 \(d\geq1\)，\(S=\operatorname{diag}(\varepsilon_1,\ldots,\varepsilon_d)\)，其中 \(\varepsilon_j\in\{1,-1\}\)，且 \(b\in C_c^\infty(\mathbb R^d;\mathbb C)\)。采用负号 \(2\pi\) Fourier 归一化，并令
\[
 D_S=\sum_{k=1}^d\varepsilon_k\partial_{u_k}^2.
\]
对 \(\lambda\geq1\)，有精确恒等式
\begin{equation}\label{b2:eq:21-quadratic-fourier}
 \int_{\mathbb R^d}e^{\pi\mathrm i\lambda u^{\mathsf T}Su}
                 \cdot b(u)\,\mathrm du
 =\lambda^{-d/2}e^{\mathrm i\pi\operatorname{sgn}S/4}
   \int_{\mathbb R^d}e^{-\pi\mathrm i\eta^{\mathsf T}S\eta/\lambda}
                 \cdot\widehat b(\eta)\,\mathrm d\eta.
\end{equation}
对任意整数 \(N\geq1\)，因而有
\[
 \begin{aligned}
 \int e^{\pi\mathrm i\lambda u^{\mathsf T}Su}\cdot b(u)\,\mathrm du
 &=\lambda^{-d/2}e^{\mathrm i\pi\operatorname{sgn}S/4}
   \left[
   \sum_{j=0}^{N-1}\frac1{j!}
     \left(\frac{\mathrm i}{4\pi\lambda}\right)^j
     D_S^jb(0)+R_N(\lambda)
   \right],\\
 |R_N(\lambda)|
 &\leq\frac{\pi^N}{N!}\lambda^{-N}
      \int_{\mathbb R^d}|\eta|^{2N}\cdot
                        |\widehat b(\eta)|\,\mathrm d\eta.
 \end{aligned}
\]
若 \(\operatorname{supp}b\subset K\)，其中 \(K\) 为固定紧集，则对任意整数 \(M>2N+d\)，右端积分至多为
\[
 C_{K,M,N}\sum_{|\alpha|\leq M}
                    \|\partial^\alpha b\|_\infty.
\]
\end{proposition}
\begin{proof}
先在左侧乘 \(e^{-\epsilon|u|^2}\)，其中 \(\epsilon>0\)。由 Schwartz 函数的 Fourier 反演，
\[
 b(u)=\int_{\mathbb R^d}e^{2\pi\mathrm i u\cdot\eta}
                  \cdot\widehat b(\eta)\,\mathrm d\eta.
\]
阻尼 Gaussian 在 \(u\) 上可积，\(\widehat b\) 在 \(\eta\) 上可积，因此可用 Fubini 交换两个积分。内层积分分成 \(d\) 个一维积分，在\cref{b2:prop:21-complex-gaussian}中分别取
\[
 a_k=\epsilon/\pi-\mathrm i\lambda\varepsilon_k,
 \quad \xi=-\eta_k.
\]
由于 \(\operatorname{Re}(1/a_k)>0\)，每一因子的模不超过
\(|a_k|^{-1/2}\leq\lambda^{-1/2}\)。所以交换后的被积函数统一由
\(\lambda^{-d/2}|\widehat b(\eta)|\) 控制。让 \(\epsilon\downarrow0\) 时，
\[
 a_k^{-1/2}\longrightarrow
        \lambda^{-1/2}e^{\mathrm i\pi\varepsilon_k/4},
 \quad
 e^{-\pi\eta_k^2/a_k}\longrightarrow
        e^{-\pi\mathrm i\varepsilon_k\eta_k^2/\lambda}.
\]
左侧则由可积函数 \(|b|\) 控制。两侧分别使用支配收敛，便得到\cref{b2:eq:21-quadratic-fourier}。Gaussian 阻尼保证交换积分合法，可积控制则使去掉阻尼后的极限仍保留该恒等式。

对实数 \(t\)，Taylor 积分余项给出
\[
 \left|e^{\mathrm it}
       -\sum_{j=0}^{N-1}\frac{(\mathrm it)^j}{j!}\right|
 \leq\frac{|t|^N}{N!}.
\]
在\cref{b2:eq:21-quadratic-fourier}中取
\(t=-\pi\eta^{\mathsf T}S\eta/\lambda\)，并用
\(|\eta^{\mathsf T}S\eta|\leq|\eta|^2\)，得到所列余项界。各个多项式项可以积分，因为 \(\widehat b\) 是 Schwartz 函数。Fourier 求导公式给出
\[
 \int_{\mathbb R^d}(\eta^{\mathsf T}S\eta)^j
                   \cdot\widehat b(\eta)\,\mathrm d\eta
 =(-4\pi^2)^{-j}D_S^jb(0).
\]
将这一因子与 \((-\pi\mathrm i/\lambda)^j\) 相乘，正好得到
\((\mathrm i/(4\pi\lambda))^j\)，从而确定次项的符号和常数。

最后估计所需的 Fourier 矩。对 \(|\eta|\geq1\)，选择
\(|\eta_k|\geq|\eta|/\sqrt d\)，沿第 \(k\) 个坐标分部积分 \(M\) 次，得到
\[
 |\widehat b(\eta)|
 \leq C_{d,M}|\eta|^{-M}\cdot\|\partial_k^Mb\|_1.
\]
对 \(|\eta|\leq1\) 用 \(|\widehat b(\eta)|\leq\|b\|_1\)。固定紧支撑使这些 \(L^1\) 范数由 \(L^\infty\) 半范数控制，故
\[
 |\widehat b(\eta)|
 \leq C_{K,M}(1+|\eta|)^{-M}
       \sum_{|\alpha|\leq M}\|\partial^\alpha b\|_\infty.
\]
条件 \(M>2N+d\) 保证乘上 \(|\eta|^{2N}\) 后仍可积，证明完成。
\end{proof}

下面采用\term[zhuxiang-gongshi]{驻相公式}[stationary phase formula]的下述版本。指数固定为 \(e^{2\pi\mathrm i\lambda\phi}\)：相位中二次项的系数是 \(1/2\)，因此归约后的指数是 \(e^{\pi\mathrm i\lambda u^{\mathsf T}Su}\)。

\Needspace{6\baselineskip}
\begin{theorem}[驻相公式]\label{b2:thm:21-stationary-phase}
设整数 \(d\geq1\)，\(U\subset\mathbb R^d\) 开，\(a\in C_c^\infty(U;\mathbb C)\)，\(\phi\in C^\infty(U;\mathbb R)\)。假定 \(\operatorname{supp}a\) 中只有一个驻点 \(x_0\)，且 \(H=D^2\phi(x_0)\) 可逆。则当 \(\lambda\to+\infty\) 时，
\[
 \int e^{2\pi\mathrm i\lambda\phi(x)}a(x)\,\mathrm dx
 =e^{2\pi\mathrm i\lambda\phi(x_0)}
   \lambda^{-d/2}e^{\mathrm i\pi\operatorname{sgn}H/4}
   \frac{a(x_0)}{|\det H|^{1/2}}
   +O(\lambda^{-d/2-1}).
\]
\end{theorem}
\begin{proof}
由\cref{b2:lem:21-morse-normal-form}，在 \(x_0\) 附近选择 Morse 坐标 \(u=\kappa(x)\)，逆映射记为 \(X\)。取一个支撑于该坐标邻域的光滑截断 \(\chi\)，在 \(x_0\) 的更小邻域内等于一，分解
\(a=\chi a+(1-\chi)a\)。第二部分的紧支撑上没有驻点，所以梯度有正下界；\cref{b2:prop:21-nonstationary-multi}使其积分为 \(O(\lambda^{-M})\)，其中整数 \(M\) 可以任意选取。

第一部分换元后成为
\[
 e^{2\pi\mathrm i\lambda\phi(x_0)}
 \int e^{\pi\mathrm i\lambda u^{\mathsf T}Su}
                   \cdot b(u)\,\mathrm du,
 \quad
 b(u)=(\chi a)(X(u))\cdot|\det DX(u)|.
\]
由于 \(\chi a\) 的支撑紧包含于坐标邻域，\(b\) 向坐标像外延零后属于 \(C_c^\infty(\mathbb R^d)\)，没有人为造成的边界跳跃。Morse 引理给出
\[
 b(0)=a(x_0)|\det H|^{-1/2},
 \quad
 \operatorname{sgn}S=\operatorname{sgn}H.
\]
对\cref{b2:prop:21-quadratic-expansion}取 \(N=1\)，再在非驻相部分取 \(M>d/2+1\)，便得到所述主项和余项。
\end{proof}
二次模型的误差已由有限个振幅导数控制。Morse 换元把 \(b\) 写成原振幅与 Jacobi 因子的乘积，有限次链式法则又将这些导数化成相位、振幅及逆映射的有限阶导数。因此参数族的统一性需要控制支撑、Hesse 最小奇异值、所用导数，以及远离驻点部分的梯度下界；单独记录每个参数的行列式非零不够。

该公式是固定振幅、固定相位下的渐近式。若振幅支撑随 \(\lambda\) 缩小或导数增长，余项常数可能增长；此时应先把自然尺度规范化。

% 扩写定位：b2:subsec:210203
边界还会改变驻相的首项。考虑一个临界点恰好位于积分端点的模型
\[
 J(\lambda)=\int_0^\infty e^{\pi\mathrm i\lambda t^2}a(t)\,\mathrm dt,
 \quad a\in C_c^\infty([0,\infty)).
\]
此处光滑是指可延拓到零点邻域的光滑函数。若 \(a(0)\ne0\)，则
\[
 J(\lambda)=\frac12e^{\pi\mathrm i/4}a(0)\lambda^{-1/2}+O(\lambda^{-1}).
\]
取在零点附近为一的偶截断 \(\chi\in C_c^\infty(\mathbb R)\)，在半直线上分解
\[
 a(t)=a(0)\chi(t)+t b(t).
\]
分子在零点消失，商 \(b(t)=[a(t)-a(0)\chi(t)]/t\) 因而可光滑延拓到零点，并且仍紧支撑。第一项的半直线积分由偶性恰为全直线截断二次积分的一半；\cref{b2:prop:21-quadratic-expansion}给出它的主项
\(\tfrac12e^{\pi\mathrm i/4}a(0)\lambda^{-1/2}\)。
第二项令 \(v=t^2\)，成为
\[
 \frac12\int_0^\infty e^{\pi\mathrm i\lambda v}
                         \cdot b(\sqrt v)\,\mathrm dv.
\]
新振幅在零点连续，且对 \(v>0\)，其导数为
\(b'(\sqrt v)/(2\sqrt v)\)，绝对积分恰为
\(\int_0^\infty|b'(t)|\,\mathrm dt\)。因此它绝对连续，虽有可能的 \(v^{-1/2}\) 导数奇性，仍可作一阶分部积分，得到界
\[
 \frac1{2\pi\lambda}
       \left[|b(0)|+\int_0^\infty|b'(t)|\,\mathrm dt\right].
\]
这证明了余项为 \(O(\lambda^{-1})\)，也确定了首项中的一半来自端点几何，而不是 Fourier 归一化的改变。

\ProseParagraphBreak

若临界点离开端点而进入内部，其首项重新是完整的二次驻相系数；若移到区间外，则固定参数下没有临界点而只剩边界项。随着临界点在距离 \(\lambda^{-1/2}\) 的区域内移动，应使用保留端点的 Fresnel 积分，不能把内部公式和非驻相公式硬拼成一个统一展开。这个模型同时说明，振幅紧支于开集和振幅在闭区间端点非零是不同假设。

\subsection{渐近展开}
\label{b2:subsec:210204}

\Cref{b2:prop:21-quadratic-expansion}已经给出二次模型的任意阶展开。一般光滑相位先经 Morse 换元，每个系数便由新振幅 \(b\) 在原点的有限阶导数确定：第 \(j\) 项用到 \(D_S^jb(0)\)，因而需要新振幅的至多 \(2j\) 阶导数。若相位原本就是 \(x^{\mathsf T}Hx/2\)，其中 \(H\) 实对称可逆，则只需固定线性换元。这时记
\[
 D_H=\sum_{k,l=1}^d(H^{-1})_{kl}\partial_{x_k}\partial_{x_l},
\]
可将前两项直接写成
\[
 \begin{aligned}
 \int e^{\pi\mathrm i\lambda x^{\mathsf T}Hx}
                  \cdot a(x)\,\mathrm dx
 ={}&\frac{\lambda^{-d/2}e^{\pi\mathrm i\operatorname{sgn}H/4}}
             {|\det H|^{1/2}}\\
 &{}\times
 \left[a(0)+\frac{\mathrm i}{4\pi\lambda}D_Ha(0)
               +O(\lambda^{-2})\right],
 \quad a\in C_c^\infty(\mathbb R^d).
 \end{aligned}
\]
具体地，若 \(x=Tu\) 且 \(T^{\mathsf T}HT=S\)，则
\(H^{-1}=TST^{\mathsf T}\)。链式法则把 \(D_S\) 化为上面的 \(D_H\)，线性换元的 Jacobi 因子为常数 \(|\det H|^{-1/2}\)。这一步不能直接推广为非二次相位的次项公式，因为非线性换元的 Jacobi 因子也会变化。

\ProseParagraphBreak
领先系数为零并不意味着整个积分为零。例如一维二次相位配以偶振幅 \(a(t)=t^2\chi(t)\)，其中 \(\chi\in C_c^\infty(\mathbb R)\) 偶且在零点附近为一。此时 \(a(0)=0\)、\(a''(0)=2\)，所以
\[
 \int e^{\pi\mathrm i\lambda t^2}\cdot t^2\chi(t)\,\mathrm dt
 =\frac{\mathrm i e^{\pi\mathrm i/4}}{2\pi}
        \lambda^{-3/2}+O(\lambda^{-5/2}).
\]
若改用奇振幅 \(t\chi(t)\)，积分则由对称性恒为零。前者是局部消失阶带来的增益，后者是全局对称性抵消；两种机制应分清。

\ProseParagraphBreak

实际应用常把 Fourier 方向、曲面上的位置或另一个积分变量当作参数。此时逐个参数都能使用驻相还不够，还须知道邻域、余项及其参数导数可以统一选择。下面给出一个固定紧支撑的版本。

\Needspace{6\baselineskip}
\begin{proposition}[含参数的驻相展开]\label{b2:prop:21-parametric-stationary}
设整数 \(d,m\geq1\)，\(U\subset\mathbb R^d\)、\(V\subset\mathbb R^m\) 开，\(K\Subset U\) 为紧集。设
\(\phi\in C^\infty(U\times V;\mathbb R)\)、
\(a\in C^\infty(U\times V;\mathbb C)\)，且每个 \(v\in V\) 都满足
\(\operatorname{supp}_x a(\,\cdot\,,v)\subset K\)。假定对每个 \(v\in V\)，\(\phi(\,\cdot\,,v)\) 在 \(K\) 中恰有一个驻点 \(x_0(v)\)，并且
\[
 H(v)=D_x^2\phi(x_0(v),v)
\]
可逆。令
\[
 c(v)=\phi(x_0(v),v),
 \quad
 I(\lambda,v)=\int_U e^{2\pi\mathrm i\lambda\phi(x,v)}
                              \cdot a(x,v)\,\mathrm dx.
\]
则 \(x_0,c\) 光滑，存在光滑系数 \(c_j(v)\)，使对任意整数 \(N\geq1\)、
紧集 \(P\Subset V\) 及参数多重指标 \(\beta\)，
\[
 \begin{aligned}
 e^{-2\pi\mathrm i\lambda c(v)}\lambda^{d/2}I(\lambda,v)
 &=\sum_{j=0}^{N-1}\lambda^{-j}c_j(v)+R_N(\lambda,v),\\
 \sup_{v\in P}|\partial_v^\beta R_N(\lambda,v)|
 &\leq C_{P,N,\beta}\lambda^{-N}
 \quad(\lambda\geq1).
 \end{aligned}
\]
首项为
\[
 c_0(v)=e^{\pi\mathrm i\operatorname{sgn}H(v)/4}
           \frac{a(x_0(v),v)}{|\det H(v)|^{1/2}}.
\]
\end{proposition}
\begin{proof}
固定 \(\lambda\) 时，公共紧支撑保证可以在积分号下对参数任意次求导。先核验驻点随参数连续。若 \(v_n\to v\)，则 \(x_0(v_n)\in K\) 的任意收敛子列极限仍在 \(K\)，梯度的连续性使它成为参数 \(v\) 的驻点。唯一性迫使这个极限为 \(x_0(v)\)。紧性于是给出整个驻点列的收敛。在
\(\nabla_x\phi(x,v)=0\) 上应用含参数的隐函数定理，由 \(H(v)\) 可逆得到 \(x_0(v)\) 光滑，\(c(v)\) 也随之光滑。

固定 \(v_*\in V\)，先把 \(x_0(v)\) 平移到原点。选择一个固定正交矩阵将 \(H(v_*)\) 对角化，而不随 \(v\) 选择特征向量。参数接近 \(v_*\) 且空间变量足够小时，\cref{b2:lem:21-morse-normal-form}证明中的全部消元主元仍非零且符号固定。因此同一配方给出光滑依赖 \((x,v)\) 的 Morse 映射。含参数的逆函数定理使逆映射 \(X(u,v)\) 同样光滑，二次矩阵 \(S\) 在这一参数片上固定。

对紧参数集 \(P\) 取有限个这样的参数片，并在每片中选随 \(x_0(v)\) 平移的截断。必要时缩小各片，使截断后的新振幅支撑都包含于一个固定 \(u\) 紧集，而且每个预先指定阶的 \(u,v\) 导数一致有界。这来自光滑映射在稍小的紧集上的有界性。在临界点处，Hesse 最小奇异值的正下界保证逆映射的一阶导数受控；反复微分逆映射恒等式，便把更高阶导数界归结为相位的相应有限阶导数。

去掉临界值因子后，局部积分正是一个固定二次相位配光滑参数振幅 \(b(u,v)\)。对\cref{b2:prop:21-quadratic-expansion}及其余项估计应用 \(\partial_v^\beta\)，只需把 \(b\) 换成 \(\partial_v^\beta b\)。于是每片上的系数光滑，余项连同该阶参数导数均为 \(O(\lambda^{-N})\)。在重叠参数片上，系数相同：将两种展开相减，先令 \(\lambda\to\infty\) 比较常数项，再依次乘以 \(\lambda\) 比较下一项，便逐项得到唯一性。因此这些局部系数可拼成 \(V\) 上的光滑函数。

还须控制截断之外的部分。它在 \(K\times P\) 内远离驻点图像，且没有其他驻点，故全梯度 \(|\nabla_x\phi|\) 有统一正下界。对归一化后的非驻相积分求 \(\beta\) 阶参数导数，最多产生 \(|\beta|\) 个 \(\lambda\) 因子，外面另有 \(\lambda^{d/2}\)。预先将\cref{b2:prop:21-nonstationary-multi}中的积分分部次数取为大于 \(N+|\beta|+d/2\) 的整数，这一部分仍为 \(O(\lambda^{-N})\)。有限个参数片上的常数取最大值，即得结论。
\end{proof}

命题中的紧性把非退化条件转为统一下界。若参数范围不紧，或相位族不在公共邻域光滑，必须另行假定固定支撑、Hesse 的统一非退化下界、所需有限阶导数界，以及远离驻点区域的梯度下界。特别要先去掉迅速变化的临界值因子：直接微分 \(e^{2\pi\mathrm i\lambda c(v)}\) 会产生 \(\lambda\)，不能把未归一化积分的参数导数误当作与函数值具有同一衰减阶。

% 扩写定位：b2:subsec:210204
驻相公式也适用于一个临界子流形，前提是垂直于临界片的二次型非退化。在坐标已经把临界片拉直为 \(z=0\) 时，可以把切向变量 \(y\) 当作参数，再在法向变量 \(z\) 上应用刚才的统一结论。若这些坐标来自先前的一次换元，体积的 Jacobi 因子须已经吸收到下面的振幅 \(a\) 中。

\Needspace{6\baselineskip}
\begin{proposition}[非退化法向临界片]\label{b2:prop:21-clean-stationary}
设整数 \(d\geq1\)、\(0\leq r\leq d\)，\(Y\subset\mathbb R^{d-r}\) 是非空连通开集，\(Z\subset\mathbb R^r\) 是包含原点的开集。设
\(\phi\in C^\infty(Y\times Z;\mathbb R)\)、
\(a\in C_c^\infty(Y\times Z;\mathbb C)\)，且对所有 \(y\in Y\) 都有
\(\mathrm d\phi(y,0)=0\)。假定 \(\operatorname{supp}a\) 中除 \(z=0\) 外没有其他全梯度驻点。当 \(r\geq1\) 时，进一步假定法向 Hesse 矩阵
\(H(y)=\partial_z^2\phi(y,0)\) 对所有 \(y\in Y\) 可逆，其最小奇异值在振幅的 \(y\) 投影附近有统一正下界。则 \(\phi(y,0)=c\) 为常数，且当 \(\lambda\to+\infty\) 时，
\[
 \int e^{2\pi\mathrm i\lambda\phi(y,z)}a(y,z)\,\mathrm dy\,\mathrm dz
 =e^{2\pi\mathrm i\lambda c}\lambda^{-r/2}
 \int e^{\pi\mathrm i\operatorname{sgn}H(y)/4}
 \frac{a(y,0)}{|\det H(y)|^{1/2}}\,\mathrm dy
 +O(\lambda^{-r/2-1}).
\]
约定 \(\mathbb R^0=\{0\}\)，零维积分取函数在该点的值。\(r=0\) 时取空矩阵行列式为一、符号差为零，公式成为没有余项的恒等式；\(r=d\) 时没有 \(y\) 变量，公式即孤立驻点情形。
\end{proposition}
\begin{proof}
临界条件使函数 \(y\mapsto\phi(y,0)\) 的各一阶导数为零，故它局部常值；连通性使其在 \(Y\) 上为同一个常数。若 \(r=0\)，全部变量都是切向变量，相位处处为 \(c\)，直接提出指数因子即得恒等式。若 \(r=d\) 且原点在振幅支撑中，命题就是前面的驻相公式；若原点不在支撑中，则主项为零，\cref{b2:prop:21-nonstationary-multi}给出任意阶衰减。以下设 \(1\leq r<d\)。

振幅的 \(y\) 投影为紧集。在其中每一点附近，将 \(y\) 当作参数，\(\partial_z\phi(y,0)=0\) 与 \(H(y)\) 可逆保证 \(z=0\) 在一个共同的小法向邻域内是唯一的法向驻点。取有限个这样的参数片与光滑分割，并在各片内加一个支撑于该法向邻域、在零点附近为一的截断。\Cref{b2:prop:21-parametric-stationary}于是适用于每片的 \(z\) 积分，给出
\[
 e^{2\pi\mathrm i\lambda c}\lambda^{-r/2}
 e^{\pi\mathrm i\operatorname{sgn}H(y)/4}
 \frac{a_{\mathrm{loc}}(y,0)}{|\det H(y)|^{1/2}}
 +O(\lambda^{-r/2-1}),
\]
其中余项对该片的 \(y\) 一致。它们的 \(y\) 支撑均在固定紧集中，故可对 \(y\) 积分，余项保持同阶；在 \(z=0\) 上各截断之和为一，主项相加就是题中积分。

剩余振幅紧支撑且远离 \(z=0\)。题设排除了此处的全梯度驻点，紧性给出 \(|\nabla_{y,z}\phi|\) 的正下界。用\cref{b2:prop:21-nonstationary-multi}在全部 \(d\) 个变量上分部积分，可取得任意阶衰减。这里不能仅凭全梯度非零，就对每个固定 \(y\) 在 \(z\) 中分部积分：远处可能有 \(\partial_z\phi=0\) 而 \(\partial_y\phi\neq0\)。先局部处理法向驻点，再对余部使用全梯度，才覆盖了题设允许的情况。
\end{proof}
衰减幂次由真正振荡的法向维数 \(r\) 决定，而不是积分空间的总维数。例如 \(\phi(y,z)=|z|^2/2\) 与乘积振幅 \(a=b(y)c(z)\)，对 \(y\) 的积分没有任何 \(\lambda\) 衰减；只有 \(z\) 变量贡献 \(\lambda^{-r/2}\)。圆柱曲面的 Fourier 变换也体现这种秩信息：弯曲方向提供驻相增益，平直方向是否衰减取决于振幅和 Fourier 方向，不能从总维数套用球面的公式。

\ProseParagraphBreak

这一局部公式不是任意退化临界点的公式。若 Hesse 矩阵在法向也退化，Morse 坐标不再存在；若临界集有交叉或尖点，也不能选成单一光滑片。比如 \(t^4\) 在零点的衰减为 \(\lambda^{-1/4}\)，已与任何一维非退化二次型的 \(\lambda^{-1/2}\) 不同。所需的是有限型分析，而不是把行列式的零值放进分母。

\ProseParagraphBreak

非二次相位的下一项并不只取决于振幅的二阶导数。这个差别可以用一维局部计算明确看到。为固定常数，假设
\(\phi(0)=\phi'(0)=0\)、\(\phi''(0)=1\)，记
\(c_3=\phi'''(0)\)、\(c_4=\phi^{(4)}(0)\)，于是
\[
 \phi(t)=\frac{t^2}2+\frac{c_3}6t^3+\frac{c_4}{24}t^4+O(t^5).
\]
商 \(2\phi(t)/t^2\) 在零点有光滑延拓且值为一，所以在足够小的邻域取正平方根，定义
\[
 y=t\sqrt{\frac{2\phi(t)}{t^2}}
   =t+\frac{c_3}{6}t^2
       +\left(\frac{c_4}{24}-\frac{c_3^2}{72}\right)t^3+O(t^4).
\]
其导数在零点等于一，故有光滑逆映射。把待定的三次 Taylor 多项式代回上式、逐项比较，得到
\[
 t=y-\frac{c_3}6y^2+
           \left(\frac{5c_3^2}{72}-\frac{c_4}{24}\right)y^3+O(y^4).
\]
逆映射在此保持方向，因此新振幅是
\(b(y)=a(t(y))\cdot t'(y)\)。若记逆映射的二次、三次系数分别为 \(\alpha,\beta\)，则 \(b\) 的二次系数是
\(a''(0)/2+3\alpha a'(0)+3\beta a(0)\)。代入这两个系数，得到
\[
 b''(0)=a''(0)-c_3a'(0)
 +\left(\frac{5c_3^2}{12}-\frac{c_4}4\right)a(0).
\]
二次模型的展开于是给出
\[
 I_\lambda=e^{\pi\mathrm i/4}\lambda^{-1/2}
 \left[a(0)+\frac{\mathrm i}{4\pi\lambda}b''(0)
 +O(\lambda^{-2})\right],
\]
这里振幅光滑紧支于该局部片，或已用截断把其余非驻相部分分开。若未减去原来的临界值，整个右侧还须乘以 \(e^{2\pi\mathrm i\lambda\phi(0)}\)。

\ProseParagraphBreak

即使原振幅在驻点附近为常数，三次和四次相位也会改变次项，因为坐标变换的体积因子不是常数。若只把二次公式中的 \(a''(0)\) 原样用于任意相位，就会遗漏这些项。计算首项时只需变换在驻点的一阶微分；计算更高项则需要相位与变换的更高阶 Taylor 数据。这解释了为什么完整渐近展开的导数假设随所需阶数增加。

\ProseParagraphBreak

有限个驻点的贡献需要先相加，才能讨论积分的整体大小。假定支撑中只有非退化驻点 \(x_1,\ldots,x_\ell\)，选各点附近互不相交的截断，再把其余部分留给非驻相法，便有
\[
 I_\lambda=\lambda^{-d/2}\sum_{\nu=1}^{\ell}
 e^{2\pi\mathrm i\lambda\phi(x_\nu)}
 e^{\pi\mathrm i\operatorname{sgn}H_\nu/4}
 \frac{a(x_\nu)}{|\det H_\nu|^{1/2}}
 +O(\lambda^{-d/2-1}).
\]
这里每个局部首项都可能非零，但它们的和仍可能抵消。取一维相位 \(\phi(t)=t^3/3-t\)，振幅只支在 \(\pm1\) 的两个小邻域，并在两个驻点等于一。两点的相位值为 \(\mp2/3\)，二阶导数为 \(\pm2\)，故首项合成为
\[
 \sqrt2\,\lambda^{-1/2}
 \cos\left(\frac{4\pi\lambda}{3}-\frac\pi4\right).
\]
沿余弦消失的一列正参数，积分只剩 \(O(\lambda^{-3/2})\)。因此一般驻相定理提供上界与展开；若要求统一的下界，还必须排除不同驻点之间的干涉。反过来，只有一个驻点且该点振幅非零时，首项模有正下界，余项便保证充分大参数下确实具有 \(\lambda^{-d/2}\) 的大小。

\ProseParagraphBreak

首项中的行列式与体积因子也有坐标无关的含义。设 \(x=F(y)\) 是局部光滑换元，\(F(y_0)=x_0\)，且 \(B=DF(y_0)\) 可逆。在驻点处一阶相位导数为零，二阶链式法则中的变换二阶导数项消失，因而新 Hesse 矩阵为 \(B^{\mathsf T}HB\)，新振幅为 \(a(x_0)|\det B|\)。于是
\[
 \frac{a(x_0)|\det B|}
 {|\det(B^{\mathsf T}HB)|^{1/2}}
 =\frac{a(x_0)}{|\det H|^{1/2}}.
\]
实可逆合同也保持正负惯性指数，所以相位因子不变。这说明驻相系数与选择哪一套 Morse 坐标无关。若只改变相位而忘记变换振幅中的 Jacobi 因子，就会破坏这种不变性；前面的一维次项计算正是这一原则在更高 Taylor 阶上的体现。

% !TeX root = ../../Book2.tex
\section{多维振荡积分}
\label{b2:sec:2103}

多维积分中，可利用的振荡方向数未必等于变量个数。相位梯度不为零时，只需沿一个合适的向量场积分分部；非退化驻点附近，每个二次方向提供半个幂次的衰减；若某些方向完全平直，它们只贡献普通积分。本节把这些方法用于参数振幅和曲面 Fourier 变换，并说明统一估计需要怎样局部化。

\subsection{相位函数}
\label{b2:subsec:210301}

\cref{b2:prop:21-nonstationary-multi}已经建立多维非驻相估计。恢复本处所需条件：相位 \(\phi\) 是开集 \(U\subset\mathbb R^d\) 上的实光滑函数，振幅 \(a\in C_c^\infty(U;\mathbb C)\)，并且 \(|\nabla\phi|\geq c>0\) 在振幅支撑的一个邻域成立。相应的向量场及微分算子为
\[
 V(x)=\frac{\nabla\phi(x)}{|\nabla\phi(x)|^2},
 \quad
 L=\frac1{2\pi\mathrm i\lambda}V\cdot\nabla,
 \quad
 Le^{2\pi\mathrm i\lambda\phi}
 =e^{2\pi\mathrm i\lambda\phi}.
\]
在这里转移的是双线性积分中的微分算子，而不是带复共轭的 Hilbert 伴随。由散度公式，其形式转置作用于振幅为
\[
 L^{\mathrm t}a
 =-\frac1{2\pi\mathrm i\lambda}\operatorname{div}(Va)
 =-\frac1{2\pi\mathrm i\lambda}
   \bigl(V\cdot\nabla a+(\operatorname{div}V)a\bigr).
\]
因此即使原振幅局部为常数，相位的变化也会通过 \(\operatorname{div}V\) 进入下一次振幅。具体地，
\[
 \operatorname{div}V
 =\frac{\Delta\phi}{|\nabla\phi|^2}
  -2\,\frac{(D^2\phi\,\nabla\phi)\cdot\nabla\phi}
               {|\nabla\phi|^4}.
\]
梯度下界控制分母；相位的有限阶导数控制分子。反复转移 \(N\) 次后，得到 \(\lambda^{-N}\) 乘以振幅至多 \(N\) 阶导数的有限线性组合，其系数由相位至多 \(N+1\) 阶导数与梯度倒数控制。

\ProseParagraphBreak

这个描述也给出使用参数族时的核对方法。若 \(z\) 遍历某个参数集合，\(a_z\) 的支撑都在同一紧集 \(K\Subset U\) 内，所需导数的 \(L^1\) 范数一致有界，\(\phi_z\) 在同一邻域上满足统一梯度下界和相位导数界，那么任意固定次数的非驻相估计就对 \(z\) 一致。参数集合是否紧只是获得这些界的一种途径；实际使用的是这些定量条件。

\ProseParagraphBreak

紧支撑负责消除积分域边界，不能被“相位在内部没有驻点”替代。对光滑有界区域 \(\Omega\)，一次散度积分分部会出现
\[
 \begin{aligned}
 \int_\Omega e^{2\pi\mathrm i\lambda\phi}\cdot a\,\mathrm dx
 &=
 \frac1{2\pi\mathrm i\lambda}
 \int_{\partial\Omega}
 e^{2\pi\mathrm i\lambda\phi}\cdot a(V\cdot n)\,\mathrm d\sigma\\
 &\hspace{1em}
 -\frac1{2\pi\mathrm i\lambda}
 \int_\Omega e^{2\pi\mathrm i\lambda\phi}\cdot
                    \operatorname{div}(Va)\,\mathrm dx,
 \end{aligned}
\]
其中 \(n\) 是外单位法向。边界积分是否衰减，取决于相位限制在边界后的切向变化。平直立方体上的线性相位甚至可以在某些面上恒定。例如
\[
 \int_{[0,1]^d}e^{2\pi\mathrm i\lambda x_1}\,\mathrm dx
 =\frac{e^{2\pi\mathrm i\lambda}-1}{2\pi\mathrm i\lambda}.
\]
其余 \(d-1\) 个变量只提供体积一，并没有增加衰减。这与光滑紧支撑振幅的快速衰减没有矛盾。

\subsection{振幅}
\label{b2:subsec:210302}

振幅的支撑、导数和在驻点处的消失阶都参与估计。尤其当振幅依赖于 \(\lambda\) 时，不能把它当作固定函数代入渐近公式，再把增长中的导数藏在余项常数里。一个完整可算的模型是
\[
 I_\lambda
 =\int_{\mathbb R^d}e^{\pi\mathrm i\lambda|x|^2}
                    \cdot a(\lambda^\theta x)\,\mathrm dx,
 \quad
 a\in C_c^\infty(\mathbb R^d;\mathbb C),
 \quad \theta\geq0.
\]
这里 \(\lambda\geq1\)，\(\theta\) 固定。换元 \(u=\lambda^\theta x\) 得
\[
 I_\lambda=\lambda^{-d\theta}
 \int_{\mathbb R^d}e^{\pi\mathrm i\tau|u|^2}
                       \cdot a(u)\,\mathrm du,
 \quad
 \tau=\lambda^{1-2\theta}.
\]
振幅的尺度归一化后，真正需要比较的是 \(\tau\) 与一的大小。

\ProseParagraphBreak

当 \(0\leq\theta<1/2\) 时，\(\tau\to\infty\)。在原点附近用二次驻相公式，其余部分用非驻相估计，得到
\[
 I_\lambda
 =e^{\pi\mathrm i d/4}\lambda^{-d/2}
   \left[a(0)+O\bigl(\lambda^{-(1-2\theta)}\bigr)\right].
\]
因子 \(\lambda^{-d\theta}\tau^{-d/2}\) 恰好等于 \(\lambda^{-d/2}\)。此时振幅变化尺度比驻相尺度宽，首项仍只看 \(a(0)\)，但相对于首项的余项是 \(\lambda^{-(1-2\theta)}\)，已不是固定振幅时的 \(\lambda^{-1}\)。若 \(a(0)=0\)，须继续看更高阶项，不能从此式断言首项非零。

\ProseParagraphBreak

当 \(\theta=1/2\) 时，\(\tau=1\)，于是有精确等式
\[
 I_\lambda
 =\lambda^{-d/2}
   \int_{\mathbb R^d}e^{\pi\mathrm i|u|^2}
                       \cdot a(u)\,\mathrm du.
\]
系数依赖整个振幅，通常不能用 \(e^{\pi\mathrm i d/4}a(0)\) 代替。振幅与驻点邻域在同一尺度变化，归一化后的积分已没有一个趋于无穷的参数。

\ProseParagraphBreak

当 \(\theta>1/2\) 时，\(\tau\to0\)。利用
\(|e^{\mathrm it}-1|\leq|t|\) 与 \(a\) 的紧支撑，得到
\[
 I_\lambda
 =\lambda^{-d\theta}
   \left[
   \int_{\mathbb R^d}a(u)\,\mathrm du
   +O\bigl(\lambda^{1-2\theta}\bigr)
   \right].
\]
这里误差由
\(\pi\tau\int|u|^2\cdot|a(u)|\,\mathrm du\)
控制。衰减主要来自支撑体积缩小，而不是相位在该支撑上产生越来越多的波动。这三个范围也说明，同一条固定振幅驻相公式不能在支撑尺度穿过 \(\lambda^{-1/2}\) 时原样保持。

\ProseParagraphBreak

非驻相区域也有类似的导数代价。设 \(\delta\geq0\) 固定。若相位固定、梯度在固定紧集附近有正下界，而 \(a_\lambda\) 支撑于该紧集，并满足
\[
 \|\partial^\gamma a_\lambda\|_1
 \leq A_\gamma\lambda^{\delta|\gamma|},
 \quad \lambda\geq1,
\]
则转移 \(N\) 次导数后至多得到
\[
 |I_\lambda(a_\lambda,\phi)|
 \leq C_N\lambda^{-N(1-\delta)}.
\]
当 \(0\leq\delta<1\)，且每个所需阶数都有上述控制时，仍能通过增加 \(N\) 得到任意固定的负幂。若 \(\delta\geq1\)，这套估计本身不再提供衰减；例如振幅中包含 \(e^{-2\pi\mathrm i\lambda\phi}\) 时，原来的振荡可以被完全消去。

\subsection{曲率}
\label{b2:subsec:210303}

曲面 Fourier 变换把几何曲率转成相位的二阶导数。先在一个图形坐标片中写
\[
 F(y)=(y,h(y)),
 \quad y\in U\subset\mathbb R^{d-1},
 \quad h\in C^\infty(U;\mathbb R).
\]
曲面体积元为
\[
 \mathrm d\sigma(F(y))
 =J(y)\,\mathrm dy,
 \quad
 J(y)=\sqrt{1+|\nabla h(y)|^2}.
\]
这个因子来自切向基向量
\(\partial_jF=e_j+(\partial_jh)e_d\)
的 Gram 矩阵 \(G=I+\nabla h(\nabla h)^{\mathsf T}\)：其行列式为 \(1+|\nabla h|^2\)，体积元是行列式的平方根。忽略这个因子会改变曲面测度，也会改变驻相系数。

\ProseParagraphBreak

取上向单位法向
\(n=(-\nabla h,1)/J\)。
第二基本形式在此切向基下的矩阵为 \(B=D^2h/J\)，因为
\(\partial_{ij}F\cdot n=(\partial_{ij}h)/J\)。
Gaussian 曲率是形算子 \(G^{-1}B\) 的行列式，故
\[
 \mathcal K(y)
 =\frac{\det D^2h(y)}
        {\bigl(1+|\nabla h(y)|^2\bigr)^{(d+1)/2}}.
\]
法向选择可能改变曲率的符号约定，但是否为零及其绝对值不受影响。下面只用其非零性与绝对值。

\Needspace{6\baselineskip}
\begin{proposition}[非零曲率曲面密度的 Fourier 衰减]\label{b2:prop:21-curved-surface-decay}
设 \(d\geq2\) 为整数，\(\Sigma\subset\mathbb R^d\) 是光滑嵌入超曲面，\(a\in C_c^\infty(\Sigma;\mathbb C)\)，且 \(\Sigma\) 在 \(\operatorname{supp}a\) 附近没有边界。假设 \(\Sigma\) 的 Gaussian 曲率在 \(\operatorname{supp}a\) 上处处非零。定义有限复测度 \(\mu=a\,\mathrm d\sigma\)，其中 \(\mathrm d\sigma\) 为诱导曲面测度，并取
\[
 \widehat\mu(\xi)
 =\int_\Sigma e^{-2\pi\mathrm i x\cdot\xi}
                  \cdot a(x)\,\mathrm d\sigma(x).
\]
则存在仅依赖曲面在该紧支撑附近的几何及振幅的常数 \(C\)，使
\[
 |\widehat\mu(\xi)|
 \leq C(1+|\xi|)^{-(d-1)/2},
 \quad \xi\in\mathbb R^d.
\]
\end{proposition}
\begin{proof}
由支撑的紧性，可用有限个经过刚性运动得到的图形坐标片覆盖，并取其上的光滑分割。曲率连续且在支撑上不为零，故可将坐标片略微缩小，使相应紧集附近仍有非零曲率。刚性运动中的平移只给 Fourier 变换乘一个模为一的因子，正交变换只改变方向。因此只需估计每个图形片上的积分
\[
 \int_U e^{2\pi\mathrm i\rho\phi_\omega(y)}
                   \cdot b(y)\,\mathrm dy,
 \quad
 \phi_\omega(y)=-\omega'\cdot y-\omega_dh(y),
\]
其中 \(\rho=|\xi|\)、\(\omega=\xi/|\xi|\in S^{d-1}\)，
\[
 b(y)=\eta(F(y))\cdot a(F(y))\cdot J(y)
\]
是光滑紧支撑振幅，\(\eta\) 为所选分割函数。记 \(K=\operatorname{supp}b\Subset U\)。

先检查任意方向上的驻点。由
\[
 \nabla_y\phi_\omega
 =-\omega'-\omega_d\nabla h,
 \quad
 D_y^2\phi_\omega=-\omega_dD^2h,
\]
若 \(\nabla_y\phi_\omega(y)=0\)，则
\(\omega'=-\omega_d\nabla h(y)\)。
因为 \(|\omega|=1\)，必有
\[
 |\omega_d|=\frac1{J(y)}>0.
\]
这正说明 \(\omega\) 平行于曲面法向。非零曲率又使 \(D^2h(y)\) 可逆，所以驻点处 \(D_y^2\phi_\omega\) 可逆。

为了使常数对方向一致，需要同时在位置和方向上局部化。考虑紧集 \(K\times S^{d-1}\)。对其中每一对 \((y_0,\omega_0)\)，分两种情形。若相位梯度不为零，连续性给出一个乘积邻域，在其中
\(|\nabla_y\phi_\omega|\geq c_0>0\)。
若它为零，先在方向球面上取局部坐标，再由刚才的 Hesse 矩阵可逆性对方程
\(\nabla_y\phi_\omega(y)=0\)
应用含参数的隐函数定理：在足够小的方向邻域内，驻点 \(y_*(\omega)\) 光滑依赖 \(\omega\)，并在一个固定的位置邻域内唯一。进一步缩小邻域，可使其 Hesse 矩阵的最小奇异值有统一正下界，同时保证所有用到的相位导数一致有界。还可在此位置邻域内选一个共同紧集，既包含后面分割振幅的支撑，也包含这些方向下的全部驻点 \(y_*(\omega)\)。

上述乘积邻域覆盖 \(K\times S^{d-1}\)，紧性给出有限子覆盖。可在各片中取位置和方向的光滑截断之积，再除以其总和，构造该紧集附近的有限光滑分割
\(\sum_\nu\chi_\nu(y,\omega)=1\)。
有限性及紧支撑保证分割函数所需的 \(y\) 导数一致有界。于是 \(b\chi_\nu\) 仍支撑于固定紧集，而其导数界不随 \(\omega\) 改变。

对第一类片，\cref{b2:prop:21-nonstationary-multi}给出任意固定阶的 \(O(\rho^{-N})\)，常数对方向一致。对第二类片，\cref{b2:prop:21-parametric-stationary}适用于移动驻点 \(y_*(\omega)\)：固定位置邻域、统一可逆的 Hesse 矩阵、相位与振幅的有限阶导数界均已建立。因此这些片的积分一致为 \(O(\rho^{-(d-1)/2})\)。驻点即使落在某个分割振幅的零区，其首项也只会消失，不影响此上界。取 \(N\geq(d-1)/2\)，对有限片相加，得到 \(\rho\geq1\) 时所需结论。

最后，当 \(|\xi|\leq1\)，直接使用
\[
 |\widehat\mu(\xi)|
 \leq\int_\Sigma|a|\,\mathrm d\sigma<\infty.
\]
调整常数便得到整个频率空间上的估计。
\end{proof}

驻相系数还展示了曲面体积元为何必不可少。记 \(n=d-1\)，在图形片的驻点处，
\[
 |\det D_y^2\phi_\omega|
 =|\omega_d|^n\cdot|\det D^2h|
 =|\mathcal K|\cdot J^2.
\]
因此原图形振幅的几何部分满足
\[
 \frac{\eta(F(y))\cdot a(F(y))\cdot J(y)}
      {|\det D_y^2\phi_\omega(y)|^{1/2}}
 =\frac{\eta(F(y))\cdot a(F(y))}
      {|\mathcal K(y)|^{1/2}}.
\]
Jacobi 因子与 Hesse 行列式中的坐标因子恰好抵消。这与\secref{b2:sec:2102}中驻相首项的坐标不变性一致，也说明曲率下界趋于零时为何不能指望同一常数继续适用。

\ProseParagraphBreak

曲率退化必须按具体模型判断。若曲面是平面，沿法向的相位恒定，非零密度的 Fourier 变换可完全没有衰减。取整数 \(1\leq r<n\)，若局部相位恰为
\[
 \phi(u,v)=\frac{|u|^2}{2},
 \quad
 (u,v)\in\mathbb R^r\times\mathbb R^{n-r},
\]
且振幅是 \(b(u)c(v)\)，其中
\(b\in C_c^\infty(\mathbb R^r;\mathbb C)\)、
\(c\in C_c^\infty(\mathbb R^{n-r};\mathbb C)\)，则积分精确分解成
\[
 \left(\int_{\mathbb R^{n-r}}c(v)\,\mathrm dv\right)
 \left(\int_{\mathbb R^r}e^{\pi\mathrm i\lambda|u|^2}
                    \cdot b(u)\,\mathrm du\right).
\]
当 \(b(0)\int c\ne0\)，二次驻相公式给出非零的
\(\lambda^{-r/2}\) 主项。这个秩模型中，只有 \(r\) 个变量产生二次振荡；其余变量确实完全平直。圆柱曲面沿适当法向的局部模型具有这种性质。

\ProseParagraphBreak

不能把上述结论推广为“驻点 Hesse 矩阵的秩总能决定衰减”。例如对 \(m\geq2\)，相位
\[
 \phi(u,v)=\frac{|u|^2}{2}+v^{2m},
 \quad u\in\mathbb R^r,\quad v\in\mathbb R,
\]
在原点的 Hesse 矩阵仍只有秩 \(r\)，但最后一个变量并非平直。对乘积振幅，二次估计和 \(2m\) 阶 van der Corput 引理给出
\(O(\lambda^{-r/2-1/(2m)})\)。
同一个点上的二阶秩没有记录这个额外增益。一般退化曲率问题还要研究高阶项、驻点集合形状及频率方向；只有具备明确的法向非退化结构时，才可应用已证明的临界子流形公式。

\subsection{局部化}
\label{b2:subsec:210304}

局部化的目的不是简单地把积分切小，而是让所需下界和导数界在每一片中同时成立。若 \(\chi_j\) 是振幅支撑附近的有限光滑分割，则
\[
 I_\lambda(a,\phi)
 =\sum_j I_\lambda(a\chi_j,\phi).
\]
在非驻相片中记录梯度下界，在驻相片中记录 Hesse 矩阵的最小奇异值下界；分割函数导数进入新振幅。有限个固定片可以相加，但参数变化时，还必须检查能否保持同一覆盖和这些下界。曲面命题中的 \(K\times S^{d-1}\) 紧性，正是使方向分割和全部估计常数能够统一的原因。

\ProseParagraphBreak

一个方向逐渐变平的二次模型可以把这种依赖完全算清。取
\[
 J_{\lambda,\varepsilon}
 =\iint_{\mathbb R^2}
 e^{\pi\mathrm i\lambda(x^2+\varepsilon y^2)}
       \cdot a(x)b(y)\,\mathrm dx\,\mathrm dy,
 \quad \lambda\geq1,\quad\varepsilon\geq0,
\]
其中 \(a,b\in C_c^\infty(\mathbb R;\mathbb C)\) 固定。令
\[
 A(\lambda)=\int_{\mathbb R}e^{\pi\mathrm i\lambda x^2}
                           \cdot a(x)\,\mathrm dx,
 \quad
 B(\tau)=\int_{\mathbb R}e^{\pi\mathrm i\tau y^2}
                           \cdot b(y)\,\mathrm dy.
\]
Fubini 定理给出 \(J_{\lambda,\varepsilon}=A(\lambda)B(\lambda\varepsilon)\)。
对 \(A\) 用二阶 van der Corput，对 \(B\) 同时用普通积分界和二阶估计，可得统一上界
\[
 |J_{\lambda,\varepsilon}|
 \leq C_{a,b}\lambda^{-1/2}
 \begin{cases}
  1,&0\leq\lambda\varepsilon\leq1,\\
  (\lambda\varepsilon)^{-1/2},&\lambda\varepsilon\geq1.
 \end{cases}
\]
这条估计在 \(\varepsilon=0\) 也有意义。固定 \(\varepsilon>0\) 时，两个二次方向最终都产生衰减，得到
\(O(\lambda^{-1}\varepsilon^{-1/2})\)；
若 \(\lambda\varepsilon\) 很小，第二个方向尚未形成有效振荡，只能保留第一方向的 \(\lambda^{-1/2}\)。

\ProseParagraphBreak

若进一步取 \(a(0)\ne0\)、\(\int b\ne0\)，并令 \(\varepsilon=\lambda^{-2}\)，则
\[
 \lambda^{1/2}J_{\lambda,\lambda^{-2}}
 \longrightarrow
 e^{\pi\mathrm i/4}a(0)\int_{\mathbb R}b(y)\,\mathrm dy\ne0.
\]
这里第一因子用二次驻相首项，第二因子用 \(\tau=\lambda^{-1}\to0\) 时的控制收敛。因此即使每个正 \(\varepsilon\) 都给出可逆 Hesse 矩阵，也不存在对这些参数统一的 \(C\lambda^{-1}\) 界。这个例子直接检验了常数的退化，而没有改变驻点数目。后面的振荡积分算子中，混合 Hesse 矩阵的下界也必须按同样原则与有效参数一起追踪。

% !TeX root = ../../Book2.tex
\section{振荡积分算子}
\label{b2:sec:2104}

标量积分的衰减来自给定振幅下的相位抵消。算子范数却要同时控制整个 \(L^2\) 单位球；输入可以依赖于高频参数，甚至抵消相位的一部分。要找到在所有输入下仍然存在的增益，需要比较两个输出点的相位。这个比较将标量非驻相估计与 Hilbert 空间的 \(TT^*\) 方法联系起来。

\subsection{算子核}
\label{b2:subsec:210401}

设整数 \(d\ge1\)，\(U\subset\mathbb R^d_x\times\mathbb R^d_y\) 开，实值相位 \(\Phi\in C^\infty(U)\)，振幅 \(a\in C_c^\infty(U)\)。对 \(\lambda\ge1\) 定义
\[
 T_\lambda f(x)\coloneqq
 \int_{\mathbb R^d}e^{2\pi\mathrm i\lambda\Phi(x,y)}a(x,y)f(y)\,\mathrm dy.
\]
这里把振幅在 \(U\) 外延为零，只在其支撑附近使用相位。由于 \(y\) 支撑有界，Cauchy--Schwarz 不等式使积分对每个 \(x\) 都有意义；对 \(x\) 再积分，得到一个与 \(\lambda\) 无关的粗略 \(L^2\) 界。这个初步界保证伴随和算子复合已经定义，随后才研究高频衰减。

\ProseParagraphBreak
固定 \(x\) 后直接对核使用驻相法，会把 \(f\) 当作振幅的一部分，所得常数便涉及 \(f\) 的导数，而 \(L^2\) 单位球没有这种控制。更直接的障碍是分离相位
\(\Phi(x,y)=\psi(x)+\theta(y)\)。取
\(f_\lambda(y)=e^{-2\pi\mathrm i\lambda\theta(y)}g(y)\)，则
\[
 T_\lambda f_\lambda(x)
 =e^{2\pi\mathrm i\lambda\psi(x)}\int a(x,y)g(y)\,\mathrm dy.
\]
输入与输出上的指数因子都不改变 \(L^2\) 范数，因此一般不存在随 \(\lambda\) 衰减的算子界。

\ProseParagraphBreak
本节固定混合 Hesse 矩阵的行列顺序为
\[
 M(x,y)_{jk}\coloneqq
 \frac{\partial^2\Phi}{\partial y_j\partial x_k}(x,y),
 \quad 1\le j,k\le d.
\]
它把输出点的位移映成输入相位梯度的变化。分离相位的 \(M\) 恒为零，双线性相位 \(x\cdot y\) 的 \(M\) 则是单位矩阵。后一个模型在没有截断时满足
\(T_\lambda f(x)=\widehat f(-\lambda x)\)，Plancherel 给出准确的范数 \(\lambda^{-d/2}\)。混合矩阵的可逆性将使一般光滑相位在小片上具有同一衰减幂次。

\subsection{\texorpdfstring{\(TT^*\)}{TT*} 方法}
\label{b2:subsec:210402}

由伴随定义和 Fubini 定理，
\[
 T_\lambda^*g(y)=\int e^{-2\pi\mathrm i\lambda\Phi(x,y)}
                         \overline{a(x,y)}g(x)\,\mathrm dx.
\]
复合算子 \(T_\lambda T_\lambda^*\) 的核为
\[
 K_\lambda(x,z)=\int
 e^{2\pi\mathrm i\lambda(\Phi(x,y)-\Phi(z,y))}
 a(x,y)\overline{a(z,y)}\,\mathrm dy.
\]
这些换序可先对紧支撑有界输入进行。对于一般 \(L^2\) 输入，振幅各变量的紧支撑及 Cauchy--Schwarz 给出相同的绝对可积性，也可由稠密性延拓。相位差不再包含待估输入；因此非驻相法所要求的振幅导数，现在全部落在已知的 \(a\) 上。

\ProseParagraphBreak
Hilbert 空间中有
\[
 \|T_\lambda\|_{2\to2}^2
 =\|T_\lambda T_\lambda^*\|_{2\to2}.
\]
一边由算子范数的次乘性得到；另一边由
\(\|T_\lambda^*g\|_2^2
 =\langle T_\lambda T_\lambda^*g,g\rangle\)
及 \(\|T_\lambda^*\|=\|T_\lambda\|\) 得到。
在 \(x=z\) 上相位差为零，核不可能处处都小。有效的目标是在小片上证明它集中在宽度约 \(\lambda^{-1}\) 的对角邻域，再把这个邻域的体积转成算子范数。

\Needspace{6\baselineskip}
\begin{lemma}[Schur 检验]\label{b2:lem:21-schur}
设整数 \(m,n\ge1\)，\(K\) 是 \(\mathbb R^m\times\mathbb R^n\) 上的可测核，且存在有限常数 \(A,B\ge0\)，使
\[
 \mathop{\mathrm{ess\,sup}}_x\int_{\mathbb R^n}|K(x,z)|\,\mathrm dz\le A,
 \quad
 \mathop{\mathrm{ess\,sup}}_z\int_{\mathbb R^m}|K(x,z)|\,\mathrm dx\le B.
\]
则 \(Sg(x)=\int K(x,z)g(z)\,\mathrm dz\) 定义从 \(L^2(\mathbb R^n)\) 到 \(L^2(\mathbb R^m)\) 的有界算子，且
\(\|S\|_{2\to2}\le\sqrt{AB}\)。对每个 \(g\in L^2\)，该积分在几乎处处的 \(x\) 绝对收敛。
\end{lemma}
\begin{proof}
Tonelli 定理与第二个界给出
\[
 \int\!\int |K(x,z)|\cdot|g(z)|^2\,\mathrm dz\,\mathrm dx
 \le B\|g\|_2^2.
\]
因此内层积分几乎处处有限。在这些 \(x\) 上，带权 Cauchy--Schwarz 不等式给出绝对收敛和
\[
 |Sg(x)|^2
 \le\left(\int |K(x,z)|\,\mathrm dz\right)
       \left(\int |K(x,z)|\cdot|g(z)|^2\,\mathrm dz\right).
\]
对 \(x\) 积分并使用第一个界，得到
\(\|Sg\|_2^2\le AB\|g\|_2^2\)。
若最初只在有界紧支撑函数上定义 \(S\)，这个估计给出的连续延拓与上述几乎处处定义一致。
\end{proof}

\term[schur-jianyan]{Schur 检验}[Schur test]要求两个方向的积分界。对本节所需的核，若已经利用振荡证明
\[
 |K_\lambda(x,z)|\le C_N(1+\lambda|x-z|)^{-N},
\]
那么换元 \(w=\lambda(x-z)\) 给出
\[
 \int_{\mathbb R^d}(1+\lambda|x-z|)^{-N}\,\mathrm dz
 =\lambda^{-d}\int_{\mathbb R^d}(1+|w|)^{-N}\,\mathrm dw.
\]
取 \(N>d\)，右端有限，对另一个积分方向也一样。因此
\(\|T_\lambda T_\lambda^*\|_{2\to2}\le C\lambda^{-d}\)，开平方才得到 \(\lambda^{-d/2}\)。这里取核的绝对值是在非驻相估计之后；若一开始对原始振荡核取绝对值，只能回到不衰减的粗略界。

\subsection{\texorpdfstring{\(L^2\)}{L2} 估计}
\label{b2:subsec:210403}

\Needspace{6\baselineskip}
\begin{theorem}[非退化振荡算子的局部 \(L^2\) 估计]\label{b2:thm:21-oscillatory-l2}
设整数 \(d\ge1\)，\(U\subset\mathbb R^d_x\times\mathbb R^d_y\) 开，实值 \(\Phi\in C^\infty(U)\)，\((x_0,y_0)\in U\)，并设混合 Hesse 矩阵
\(M_{jk}=\partial_{y_j}\partial_{x_k}\Phi\) 在 \((x_0,y_0)\) 可逆。
则存在该点的乘积邻域 \(W\Subset U\)，使每个 \(a\in C_c^\infty(W)\) 所定义的算子
\[
 T_\lambda f(x)=\int e^{2\pi\mathrm i\lambda\Phi(x,y)}a(x,y)f(y)\,\mathrm dy
\]
满足
\[
 \|T_\lambda f\|_2\le C\lambda^{-d/2}\|f\|_2,
 \quad f\in L^2(\mathbb R^d),\quad\lambda\ge1.
\]
更一般地，若 \(a\in C_c^\infty(U)\)，且 \(M\) 在 \(\operatorname{supp}a\) 的邻域处处可逆，同一结论成立。常数可依赖于相位、振幅及其固定支撑，但与 \(f,\lambda\) 无关。
\end{theorem}
\begin{proof}
记 \(M_0=M(x_0,y_0)\)，其最小奇异值为 \(s_0>0\)。取有界凸邻域 \(X,Y\)，使
\(\overline X\times\overline Y\subset U\)，并且在其上
\(\|M(x,y)-M_0\|\le s_0/2\)。先设振幅支撑在 \(X\times Y\) 内。
当 \(x,z\in X\) 时，连接二者的线段仍在 \(X\) 中，故
\[
 \nabla_y\Phi(x,y)-\nabla_y\Phi(z,y)
 =\left[\int_0^1 M(z+t(x-z),y)\,\mathrm dt\right](x-z).
\]
括号内矩阵与 \(M_0\) 的距离不超过 \(s_0/2\)，所以相位梯度差的长度至少为 \((s_0/2)|x-z|\)。

对 \(x\ne z\) 定义归一化相位
\[
 \psi_{x,z}(y)\coloneqq
 \frac{\Phi(x,y)-\Phi(z,y)}{|x-z|}.
\]
它的梯度下界为 \(s_0/2\)。此外，对于每个多重指标 \(\beta\)，均值公式给出
\[
 \partial_y^\beta\psi_{x,z}(y)
 =\int_0^1\nabla_x\partial_y^\beta\Phi(z+t(x-z),y)
       \cdot\frac{x-z}{|x-z|}\,\mathrm dt.
\]
因此在固定紧集上，相位的所需导数一致有界，常数不随 \(x,z\) 接近而增大。振幅
\(a(x,y)\overline{a(z,y)}\) 的 \(y\) 导数及支撑体积也一致受控。
对参数 \(\mu=\lambda|x-z|\ge1\) 使用\cref{b2:prop:21-nonstationary-multi}，得到
\( |K_\lambda(x,z)|\le C_N\mu^{-N}\)。
当 \(\mu<1\) 或 \(x=z\) 时，直接用振幅的 \(L^1\) 上界。两种情形合并为
\[
 |K_\lambda(x,z)|\le C_N(1+\lambda|x-z|)^{-N}.
\]
在 \(X\) 外的输出点上核为零。取整数 \(N>d\)，由\cref{b2:lem:21-schur}及上一小节的换元公式得到 \(TT^*\) 的 \(C\lambda^{-d}\) 界，再开平方即得局部结论。

对于一般紧支撑，在其上选有限个上述乘积邻域，用从属光滑分割写成
\(a=\sum_{\nu=1}^L a_\nu\)。对应算子为
\(T_\lambda=\sum_{\nu=1}^L T_{\lambda,\nu}\)，因而
\[
 \|T_\lambda\|_{2\to2}
 \le\sum_{\nu=1}^L\|T_{\lambda,\nu}\|_{2\to2}
 \le\left(\sum_{\nu=1}^L C_\nu\right)\lambda^{-d/2}.
\]
有限求和完成整体结论，不需要估计不同小片之间的交叉核。
\end{proof}

小片条件在证明中有明确用途：平均矩阵仍接近同一个可逆矩阵。仅知每一点的混合矩阵可逆，不能推出它沿长线段的平均可逆，也不能推出 \(x\mapsto\nabla_y\Phi(x,y)\) 在大区域上单射。因此先分别估计局部算子，再求和。每个小片内的核集中在对应的对角线附近，整体核则可能具有更复杂的集中结构。

\subsection{参数依赖性}
\label{b2:subsec:210404}

在参数族中，要统一上述证明，需要同时控制支撑、小片半径、逆矩阵和分部积分系数。一个充分条件是：所有振幅支撑在同一紧集 \(K\Subset U\)，相位在 \(K\) 的同一邻域上定义，混合矩阵最小奇异值一致不小于 \(c>0\)，并且所需的相位、振幅导数一致有界。相位三阶导数界使混合矩阵有统一的连续模，因而可选择统一的小片半径；\(K\) 的紧性再给出片数及分割函数导数的统一控制。

\ProseParagraphBreak
对于一个给定的 \(L^2\) 估计，只需有限阶光滑性。固定整数 \(N>d\) 后，在上述固定邻域上控制相位到 \(N+2\) 阶、振幅到 \(N\) 阶，是保证这套证明中常数统一的一个充分条件。相位的这一阶数还包含均值公式所多用的一次 \(x\) 微分。若要求任意阶渐近展开，则要对每一阶分别提出相应导数界，不能从一次有限阶估计推出全部展开。

\ProseParagraphBreak
仅记录行列式不为零，尚不足以核对这套证明的统一常数。例如
\(M_\epsilon=\operatorname{diag}(\epsilon,\epsilon^{-1})\)
的行列式恒为一，但最小奇异值趋零，矩阵范数趋于无穷。这个例子表明上述充分条件失效，却不单独证明对应算子范数变坏：无截断的线性 Fourier 模型，其准确范数由行列式决定。若要证明某族的统一衰减确实失败，应直接检验算子。

\ProseParagraphBreak
一个直接检验是 \(\Phi_\epsilon(x,y)=\epsilon x\cdot y\)，其中 \(0<\epsilon\le1\)。取固定分离振幅 \(\chi(x)\eta(y)\)，则
\[
 T_{\lambda,\epsilon}f(x)
 =\chi(x)\widehat{\eta f}(-\lambda\epsilon x).
\]
当 \(\lambda\epsilon\ge1\) 时，Plancherel 给出的上界为
\(C(\lambda\epsilon)^{-d/2}\|f\|_2\)；当 \(\lambda\epsilon<1\) 时，紧支撑给出 \(C\|f\|_2\)。这两个界说明真正控制振荡的是 \(\lambda\epsilon\)。若选择固定 \(g\) 使 \(\int\eta g\ne0\)，并令 \(\epsilon=\lambda^{-2}\)，则
\[
 T_{\lambda,\epsilon}g\longrightarrow
 \chi\int\eta(y)g(y)\,\mathrm dy
 \quad\text{于 }L^2.
\]
只要 \(\chi\not\equiv0\)，极限非零，故不可能有独立于 \(\epsilon\) 的 \(C\lambda^{-d/2}\) 界。

\ProseParagraphBreak
在固定非退化情形，\(\lambda^{-d/2}\) 本身也不能改善。仍取双线性相位和分离振幅，设 \(\chi(0)\ne0\)，固定 \(g\in C_c^\infty\) 满足 \(\eta g\not\equiv0\)。则
\[
 \|T_\lambda g\|_2^2
 =\lambda^{-d}\int |\chi(v/\lambda)|^2\cdot
                     |\widehat{\eta g}(-v)|^2\,\mathrm dv.
\]
支配收敛及 Plancherel 给出
\(\lambda^d\|T_\lambda g\|_2^2
 \to |\chi(0)|^2\cdot\|\eta g\|_2^2>0\)，所以算子范数至少是 \(c\lambda^{-d/2}\)。这个测试无需让输入依赖于 \(\lambda\)。

\ProseParagraphBreak
若输入与输出变量的维数不同，混合矩阵不再是方阵，应改为考虑其秩及无振荡变量的作用。多加一个仅作普通积分的变量，并不会多产生半个幂次的衰减；衰减来自真正配对的振荡方向。后续限制问题中出现的秩与曲率条件，需要在相应几何下重新分析。

% !TeX root = ../../Book2.tex
\OptionalSection{Fourier 积分算子初步}{b2:sec:2105}

前几节从相位的导数得到积分或算子范数估计。本节进一步问：输入在某个位置、某个频率方向的振荡，会传到输出的哪里？Fourier 积分算子的相位可以同时记录这两类信息。从驻点关系中可以读出输入和输出的位置、频率，再考察改变相位参数时这些信息怎样保持。相关背景可参阅\appref{b2:app:G}及专著 \cite{Hrmander2009AnalysisIV}。

\subsection{非退化相位}
\label{b2:subsec:210501}

先考虑不会产生收敛困难的表示。设 \(U\subset\mathbb R^d\) 开，实值 \(\phi\) 光滑，\(a\in C_c^\infty(U\times\mathbb R^d)\)，\(f\in\mathcal S(\mathbb R^d)\)。由 Fourier 变换的定义及 Fubini 定理，
\[
 \begin{aligned}
 Af(x)
 &=\iint e^{2\pi\mathrm i(\phi(x,\theta)-y\cdot\theta)}
          a(x,\theta)f(y)\,\mathrm dy\,\mathrm d\theta\\
 &=\int e^{2\pi\mathrm i\phi(x,\theta)}
          a(x,\theta)\widehat f(\theta)\,\mathrm d\theta.
 \end{aligned}
\]
频率支撑紧时，对 \(x,y\) 作任意次微分都可以移入核积分，所以核光滑。这个情形适合先检查相位和换元，但一般 Fourier 积分算子的奇性来自无界频率，不能由紧频率核本身看出。

\ProseParagraphBreak
例如 \(\phi(x,\theta)=x\cdot\theta\) 给出左量子化的伪微分表示。取 \(a\equiv1\) 后，积分对 Schwartz 输入仍绝对收敛，由 Fourier 反演得到恒等算子。若用 \(\chi(\theta/R)\) 截断，其中 \(\chi\in C_c^\infty\) 在零点附近等于一，则各个截断核都光滑，而它们只在分布意义下趋于 \(\delta(x-y)\)。这是光滑核序列产生奇异极限的一个基本例子。

\ProseParagraphBreak
另一个模型是
\[
 \phi_t(x,\theta)=x\cdot\theta+t|\theta|,
 \quad
 A_tf=\mathcal F^{-1}\bigl(e^{2\pi\mathrm i t|\theta|}\widehat f(\theta)\bigr).
\]
这个 Fourier 乘子定义了半波算子 \(e^{2\pi\mathrm i t|D|}\)。它的乘子模为一，故由 Plancherel 延拓为 \(L^2\) 上的酉算子。其相位只在 \(\theta\ne0\) 光滑；低频部分将另行处理。以下的齐次相位描述均针对非零频率。

\Needspace{6\baselineskip}
\begin{definition}[齐次非退化相位]\label{b2:def:21-nondegenerate-phase}
设整数 \(d\ge1\)，\(U,V\subset\mathbb R^d\) 开，\(\Gamma\subset\mathbb R^k\setminus\{0\}\) 为开锥，\(k\ge1\)。实值光滑函数 \(\Psi(x,y,\theta)\) 若满足
\(\Psi(x,y,r\theta)=r\Psi(x,y,\theta)\)（\(r>0\)），且在临界集合
\[
 C_\Psi\coloneqq\{(x,y,\theta):\partial_\theta\Psi(x,y,\theta)=0\}
\]
上，\(k\) 个微分
\(\mathrm d_{x,y,\theta}(\partial_{\theta_j}\Psi)\)
线性无关，则称为\term[qi-ci-fei-tui-hua-xiang-wei]{齐次非退化相位}[homogeneous nondegenerate phase]。本节只在
\(\partial_x\Psi\ne0\)、\(\partial_y\Psi\ne0\)
的部分讨论输入、输出均为非零的余切方向。
\end{definition}

这里的非退化是驻点方程作为约束的正则性。隐函数定理说明 \(C_\Psi\) 的维数为 \(2d\)，因为总变量数为 \(2d+k\)，有 \(k\) 个独立约束。它并不要求频率 Hesse 矩阵\(\partial_\theta^2\Psi\)可逆。事实上，Euler 齐次恒等式及其微分给出
\[
 \Psi=\theta\cdot\partial_\theta\Psi,
 \quad (\partial_\theta^2\Psi)\theta=0;
\]
因而频率 Hesse 矩阵总有径向零方向。

\ProseParagraphBreak
对于混合表示 \(\Psi=\phi(x,\theta)-y\cdot\theta\)，约束微分中的 \(y\) 分块是 \(-I\)，所以这组约束自动独立。若再要求 \(\partial_x\partial_\theta\phi\) 可逆，才可从
\(y=\partial_\theta\phi(x,\theta)\) 局部解出 \(x\)，将所得关系写成输入到输出的局部图。前几节还有另一种条件：标量驻相要求积分变量上的二阶 Hesse 矩阵可逆。三种条件分别服务于驻相渐近、\(L^2\) 算子估计和驻点集合的几何，使用时必须说明求导变量。

\subsection{典范关系}
\label{b2:subsec:210502}

\Needspace{6\baselineskip}
\begin{definition}[相位所生成的典范关系]\label{b2:def:21-canonical-relation}
在\cref{b2:def:21-nondegenerate-phase}的条件下，相位 \(\Psi\) 所生成的\term[dianfan-guanxi]{典范关系}[canonical relation]是集合
\[
 \mathcal C_\Psi\coloneqq
 \bigl\{(x,\partial_x\Psi;\ y,-\partial_y\Psi):
       (x,y,\theta)\in C_\Psi\bigr\},
\]
并只保留两端余切向量均非零的部分。这里 \((x,\xi)\) 表示输出位置及频率，\((y,\eta)\) 表示输入位置及频率。
\end{definition}

输入端的负号来自核的约定。若输入局部含有正频率振荡
\(e^{2\pi\mathrm i y\cdot\eta}\)，则总相位对 \(y\) 的驻点条件为
\(\partial_y\Psi+\eta=0\)，而输出相位梯度是 \(\partial_x\Psi\)。因此上式记录的正是可能相互耦合的输入和输出频率。

\ProseParagraphBreak
非退化条件也保证上述参数化在局部没有额外的切向冗余。若 \(C_\Psi\) 上一个切向量在 \((x,\partial_x\Psi;y,-\partial_y\Psi)\) 下的像为零，则其 \(\delta x,\delta y\) 均为零，并且
\[
 \Psi_{x\theta}\delta\theta=0,
 \quad\Psi_{y\theta}\delta\theta=0,
 \quad\Psi_{\theta\theta}\delta\theta=0.
\]
最后一个等式来自临界约束的微分。把三个式子合并，就是约束微分矩阵的转置作用于 \(\delta\theta\) 为零；其行满秩迫使 \(\delta\theta=0\)。这说明关系具有正确的局部维数，但不排除不同参数点给出同一像。Euler 恒等式又使 \(\Psi\) 在 \(C_\Psi\) 上等于零，从而沿该集合有
\(\partial_x\Psi\cdot\mathrm dx+\partial_y\Psi\cdot\mathrm dy=0\)。这些恒等式将相位的微分与输入、输出频率联系起来，是进一步研究其几何结构的出发点。

\ProseParagraphBreak
恒等相位 \(\Psi=(x-y)\cdot\theta\) 的约束为 \(x=y\)，所以
\[
 \mathcal C_\Psi=\{(x,\theta;x,\theta):\theta\ne0\}.
\]
半波相位则给出
\[
 \Psi_t=(x-y)\cdot\theta+t|\theta|,
 \quad x-y+t\frac{\theta}{|\theta|}=0,
\]
于是
\[
 \mathcal C_{\Psi_t}
 =\left\{\left(y-t\frac{\theta}{|\theta|},\theta;\ y,\theta\right):
          \theta\ne0\right\}.
\]
对本节采用的正指数 \(e^{2\pi\mathrm i t|D|}\)，输出位置为
\(x=y-t\theta/|\theta|\)；若改成负指数，方向反转。位置与频率必须同时记录，单凭空间中的球面无法恢复这个符号信息。

\subsection{相位等价}
\label{b2:subsec:210503}

首先看真正的频率坐标变换。设
\(\theta=\Theta(x,y,\eta)\) 对每个 \((x,y)\) 都是两个开锥之间的光滑微分同胚，并且对 \(\eta\) 正齐次一次。换元后的相位及振幅为
\[
 \widetilde\Psi(x,y,\eta)=\Psi(x,y,\Theta(x,y,\eta)),
 \quad
 \widetilde a=a(x,y,\Theta(x,y,\eta))
              |\det D_\eta\Theta|.
\]
对于紧频率支撑，这是普通换元公式。其驻点方程为
\[
 \partial_\eta\widetilde\Psi
 =(D_\eta\Theta)^{\mathsf T}\partial_\theta\Psi=0,
\]
与原约束等价。在临界集合上，对 \(x,y\) 微分时由 \(\Theta\) 产生的额外项均乘着 \(\partial_\theta\Psi=0\)，故两端余切向量相同。约束的微分也只多乘可逆矩阵及作坐标变换，因此非退化性与典范关系都保持不变。Jacobi 因子则必须保留，否则积分算子会改变。

\ProseParagraphBreak
增加辅助二次变量也能产生等价表示，但需要顾及齐次性。直接加 \(|\eta|^2/2\) 会破坏一次齐次；在 \(\theta\ne0\) 上应改用
\[
 \widetilde\Psi(x,y,\theta,\eta)
 =\Psi(x,y,\theta)+\frac{|\eta|^2}{2|\theta|},
 \quad\eta\in\mathbb R^m.
\]
它对 \((\theta,\eta)\) 同时伸缩为一次齐次。新的约束先给出 \(\eta=0\)，再给出旧约束 \(\partial_\theta\Psi=0\)，所以典范关系相同；新约束的 \(\eta\) 分块为 \(|\theta|^{-1}I\)，保证其独立性。

\ProseParagraphBreak
为核对算子而不只核对驻点集合，可先把 \(\theta\) 限制在远离零的紧集。\cref{b2:prop:21-complex-gaussian}逐坐标给出 Abel 正则化积分
\[
 \lim_{\epsilon\downarrow0}\int_{\mathbb R^m}
 e^{\pi\mathrm i|\eta|^2/|\theta|}
 e^{-\pi\epsilon|\eta|^2}\,\mathrm d\eta
 =e^{\mathrm i\pi m/4}|\theta|^{m/2}.
\]
因此把振幅换成
\(e^{-\mathrm i\pi m/4}|\theta|^{-m/2}a(x,y,\theta)\)，再加入这组辅助变量，正则化后恰恢复原核。该等式在固定的紧频率集上一致，允许把极限移入 \(\theta\) 积分。对于无界频率，因子 \(|\theta|^{-m/2}\) 还会改变振幅的阶；必须连同符号条件及分布极限一起处理。

\subsection{微局部观点}
\label{b2:subsec:210504}

微局部观察先用空间截断选定位置，再用锥形频率截断选定方向。相位在该方向没有驻点时，可以通过反复分部积分排除相应核奇性。为了说明这句话在无界频率中为什么成立，需要把每个有限频率壳层的估计求和，而不是直接把紧支撑非驻相公式套在整个频率空间上。

\ProseParagraphBreak
具体地，在一个固定的 \((x,y)\) 紧邻域及频率开锥内，设 \(\Psi\) 是一次齐次光滑相位，并且在紧包含于该开锥的单位球面截面内的角向支撑上
\( |\partial_\theta\Psi|\ge c>0\)。设高频振幅 \(a\) 满足：对所有多重指标 \(\alpha,\beta\)，存在常数使
\[
 |\partial_{x,y}^\alpha\partial_\theta^\beta a(x,y,\theta)|
 \le C_{\alpha\beta}(1+|\theta|)^{m-|\beta|},
 \quad |\theta|\ge1.
\]
把频率分为 \(|\theta|\asymp2^j\) 的光滑壳层，令 \(\theta=2^j\omega\)。壳层核变成固定环域上的振荡积分，其参数为 \(2^j\)，Jacobi 因子为 \(2^{jk}\)，振幅规模至多为 \(2^{jm}\)。对 \((x,y)\) 作 \(|\alpha|\) 次微分，最坏再增加 \(2^{j|\alpha|}\)。利用非驻相估计任意 \(N\) 次，便得到
\[
 |\partial_{x,y}^\alpha K_j(x,y)|
 \le C_{\alpha,N}\,2^{j(k+m+|\alpha|-N)}.
\]
取 \(N>k+m+|\alpha|\)，右端可求和。对每个微分阶数分别如此选择 \(N\)，壳层级数及其全部导数一致收敛，所以这个位置、方向片上的核光滑。

\ProseParagraphBreak
在恒等算子模型，\(\partial_\theta\Psi=x-y\)，上述论证解释了核在对角线外光滑。在半波模型，条件变为
\(x-y+t\theta/|\theta|\ne0\)，因此高频贡献只能在满足传播关系的方向上产生核奇性。存在驻点只是可能产生奇性的条件；振幅若在相关位置消失，相应贡献仍可能消失。一般的波前集映射需要进一步的局部化论证，不能只由这一单向判别替代。

\ProseParagraphBreak
半波的低频核可直接处理。若 \(\chi\) 紧支撑于零点附近，则
\[
 K_{t,\mathrm{low}}(x,y)
 =\int e^{2\pi\mathrm i((x-y)\cdot\theta+t|\theta|)}
          \chi(\theta)\,\mathrm d\theta
\]
对 \(x,y\) 的任何微分只引入频率多项式，仍是绝对可积函数，故核光滑。因频率支撑有界，所有频率矩都可积，空间变量的各阶导数因而都有绝对收敛的积分表示。分离低频之后，零点的不光滑与高频传播的几何便不会混淆。完整波前集定义、伪微分椭圆正则性与传播定理见\appref{b2:app:G}，其中分别说明各结论的证明范围。


\begin{chaptersummary}
非驻相法把相位导数下界转为反复分部积分；端点不消失时，留下的边界项本身就可能是主项。van der Corput 引理通过小导数区域的测度控制处理有限阶退化。非退化驻相则给出更细的展开：每个有效二次方向贡献半个衰减幂次，行列式给出大小，惯性指数给出复相位，坐标换元的体积因子进入振幅。

\ProseParagraphBreak

估计应连同所控制的对象一起使用。固定振幅的导数界、随参数缩小的支撑，以及允许输入自由变化的算子单位球，不能彼此替代。振荡算子的混合 Hesse 矩阵控制相位梯度在两个输出点之间的变化；\(TT^*\) 在各个小片上先得到一个集聚在对角线附近的核，再由 Schur 估计和开平方给出 \(L^2\) 增益。有限个小片的估计相加，便得到紧支集振幅上的整体界。

\ProseParagraphBreak

遇到参数退化，先把参数与空间尺度共同规范化，再判断哪一种估计仍然一致。振幅的消失、部分方向的非退化和不同驻点之间的干涉，都可能改变实际衰减。选读部分进一步用位置与频率方向描述传播；它提供几何信息，而具体算子的大小与误差仍要由相应估计验证。
\end{chaptersummary}

\BookExercises

% 自拟；基础到中等：由端点消失阶求非线性非驻相边界首项，并验证单端贡献的尖锐阶。
\Needspace{6\baselineskip}
\begin{exercise}\label{b2:ex:21-new-1}
设整数 \(m\geq0\)，\(a\in C^\infty([0,1];\mathbb C)\)、\(\phi\in C^\infty([0,1];\mathbb R)\)，且 \(|\phi'|\geq c>0\)。假定
\(a^{(j)}(0)=a^{(j)}(1)=0\) 对 \(0\leq j<m\) 成立；\(m=0\) 时这一条件为空。证明当 \(\lambda\to+\infty\) 时，
\[
 \begin{aligned}
 I_\lambda
 &=\int_0^1 e^{2\pi\mathrm i\lambda\phi(t)}\cdot a(t)\,\mathrm dt\\
 &=\frac{(-1)^m}{(2\pi\mathrm i\lambda)^{m+1}}
 \left[\frac{e^{2\pi\mathrm i\lambda\phi(t)}\cdot a^{(m)}(t)}
                  {\phi'(t)^{m+1}}\right]_{t=0}^{t=1}
 +O(\lambda^{-m-2}).
 \end{aligned}
\]
取 \(\phi(t)=t+t^2\)、\(a(t)=t^m(1-t)^{m+1}\)，求出非零首项，并说明由此为何不能把一般结论改为 \(o(\lambda^{-m-1})\)。
\end{exercise}
\begin{solution}
记 \(\kappa=2\pi\mathrm i\lambda\)，并递归令
\(b_0=a\)、\(b_{j+1}=(b_j/\phi')'\)。每次分部积分给出
\[
 I_\lambda
 =\sum_{j=0}^{M-1}(-1)^j\kappa^{-j-1}
       \left[\frac{e^{\kappa\phi}\cdot b_j}{\phi'}\right]_0^1
   +(-1)^M\kappa^{-M}
       \int_0^1e^{\kappa\phi}\cdot b_M\,\mathrm dt.
\]
归纳展开递推式，\(b_j\) 的最高阶振幅导数项为
\((\phi')^{-j}\cdot a^{(j)}\)，其余项只含低于 \(j\) 阶的振幅导数，系数由 \(\phi\) 的导数及 \((\phi')^{-1}\) 构成。因此在两个端点，\(j<m\) 时有 \(b_j=0\)，而
\(b_m=(\phi')^{-m}\cdot a^{(m)}\)。取 \(M=m+2\)，第 \(m+1\) 个边界项和最后的积分都是 \(O(\lambda^{-m-2})\)；它们的系数有限，因为区间紧且 \(|\phi'|\geq c\)。剩余第 \(m\) 个边界项就是所求公式。

具体例中，\(a^{(m)}(0)=m!\)、\(a^{(m)}(1)=0\)、\(\phi(0)=0\)、\(\phi'(0)=1\)，故
\[
 I_\lambda
 =\frac{(-1)^{m+1}m!}{(2\pi\mathrm i\lambda)^{m+1}}
       +O(\lambda^{-m-2}),
 \quad
 \lambda^{m+1}|I_\lambda|
       \longrightarrow\frac{m!}{(2\pi)^{m+1}}>0.
\]
这个例子只有左端贡献非零，所以没有两端首项相互抵消的问题。一般情况下，两个端点即使都贡献非零系数，它们的和也可能沿某些参数抵消；因此总和在某些大参数处可能小于各首项的量级。
\end{solution}

% 自拟；中等：退化单侧相位配消失振幅，证明振荡常数非零并控制无穷远尾项，不仅作形式缩放。
\Needspace{6\baselineskip}
\begin{exercise}\label{b2:ex:21-new-2}
设整数 \(k\geq2\)、\(0\leq m\leq k-2\)，\(\chi\in C_c^\infty([0,\infty);\mathbb R)\) 在零点附近恒为一；这里光滑指能够延拓到零点的一个邻域。令
\[
 I_{k,m}(\lambda)=\int_0^\infty
       e^{2\pi\mathrm i\lambda t^k}\cdot t^m\cdot\chi(t)\,\mathrm dt,
 \quad \alpha=\frac{m+1}{k}.
\]
证明反常积分
\[
 C_{k,m}=\int_0^\infty e^{2\pi\mathrm i s^k}\cdot s^m\,\mathrm ds
\]
收敛且虚部严格为正，并且对每个整数 \(N\geq1\)，
\[
 I_{k,m}(\lambda)
 =C_{k,m}\lambda^{-(m+1)/k}+O_N(\lambda^{-N}).
\]
可在替换 \(v=2\pi s^k\) 后，将正弦积分的相邻正、负半周期配对。另求 \(m=k-1\) 时的首项，并解释这一端点为什么不再由上面普通收敛的反常积分给出。
\end{exercise}
\begin{solution}
换元后，
\[
 C_{k,m}=\frac1{k(2\pi)^\alpha}
          \int_0^\infty e^{\mathrm iv}\cdot v^{\alpha-1}\,\mathrm dv,
 \quad 0<\alpha<1.
\]
零点附近绝对可积。在无穷远对 \(e^{\mathrm iv}\) 积分分部，边界因子 \(v^{\alpha-1}\) 趋零，导数 \((\alpha-1)v^{\alpha-2}\) 绝对可积，故反常积分收敛。其正弦部分在第 \(j\) 对半周期上的积分为
\[
 \int_0^\pi\sin v\cdot
 \left[(2j\pi+v)^{\alpha-1}
       -\bigl((2j+1)\pi+v\bigr)^{\alpha-1}\right]\,\mathrm dv>0
 \quad(j\geq0).
\]
方括号为正，因为幂函数严格递减；第一项在 \(j=0\) 时仍可积。按整周期取部分积分再取极限，每一项为正且第一项非零，所以总虚部严格为正。

在完整的反常积分中作 \(s=\lambda^{1/k}t\)，得到
\[
 \int_0^\infty e^{2\pi\mathrm i\lambda t^k}\cdot t^m\,\mathrm dt
       =\lambda^{-\alpha}C_{k,m}.
\]
因此两者之差是振幅 \(q_0(t)=t^m\cdot\bigl(\chi(t)-1\bigr)\) 的振荡积分。这个振幅在零点附近消失，在足够大的 \(t\) 上等于 \(-t^m\)。令
\[
 q_{j+1}(t)=\left(\frac{q_j(t)}{kt^{k-1}}\right)'.
\]
所有 \(q_j\) 在零点附近都为零；在无穷远有
\(q_j(t)=O(t^{m-jk})\)。分部积分的边界因子
\(q_j(t)/(kt^{k-1})=O(t^{m-(j+1)k+1})\)
趋零，而 \(q_N\in L^1\) 对 \(N\geq1\) 成立。先在有限区间积分分部，再令上端趋于无穷，因而得到
\[
 I_{k,m}(\lambda)-\lambda^{-\alpha}C_{k,m}
 =(-2\pi\mathrm i\lambda)^{-N}
       \int_0^\infty e^{2\pi\mathrm i\lambda t^k}\cdot q_N(t)\,\mathrm dt.
\]
取 \(L^1\) 上界即得任意阶余项。由于 \(C_{k,m}\ne0\)，还得到
\(\lambda^{(m+1)/k}|I_{k,m}(\lambda)|\to|C_{k,m}|\)，所以这里的幂次是实际大小。

当 \(m=k-1\) 时，令 \(u=t^k\)，有
\[
 I_{k,k-1}(\lambda)
 =\frac1k\int_0^\infty e^{2\pi\mathrm i\lambda u}
                         \cdot\chi(u^{1/k})\,\mathrm du.
\]
新振幅在零点附近为一，其余部分光滑且紧支撑；它的导数远离零点。一次分部积分留下左端项，余项可任意次积分分部，故
\[
 I_{k,k-1}(\lambda)
 =-\frac1{2\pi\mathrm i k\lambda}+O_N(\lambda^{-N}).
\]
此时未经截断的积分变成常数倍 \(\int_0^\infty e^{\mathrm iv}\,\mathrm dv\)，普通反常极限不存在。非零首项来自保留下来的端点，而非上一个收敛论证。
\end{solution}

% 自拟；中等：由凸相位的参数临界值导出逆Hesse矩阵，再将驻相首项与梯度换元、Plancherel的能量配平。
\Needspace{6\baselineskip}
\begin{exercise}\label{b2:ex:21-new-3}
设整数 \(d\geq1\)，\(\psi\in C^\infty(\mathbb R^d;\mathbb R)\)，且
\(D^2\psi(x)\geq cI\) 对所有 \(x\) 成立，其中 \(c>0\)。设 \(a\in C_c^\infty(\mathbb R^d;\mathbb C)\)，对 \(\lambda\geq1\) 令
\[
 I_\lambda(p)=\int_{\mathbb R^d}
   e^{2\pi\mathrm i\lambda(\psi(x)-p\cdot x)}\cdot a(x)\,\mathrm dx.
\]
证明每个 \(p\in\mathbb R^d\) 都对应唯一的 \(x(p)\)，使 \(\nabla\psi\bigl(x(p)\bigr)=p\)。令
\(S(p)=\psi\bigl(x(p)\bigr)-p\cdot x(p)\)，求 \(\nabla S\) 与 \(D^2S\)，并求 \(I_\lambda(p)\) 的驻相首项；要求余项在每个紧参数集上一致。

\ProseParagraphBreak

若将首项写为
\(e^{2\pi\mathrm i\lambda S(p)}\lambda^{-d/2}e^{\mathrm i\pi d/4}\cdot b(p)\)，证明
\[
 \int_{\mathbb R^d}|b(p)|^2\,\mathrm dp
       =\int_{\mathbb R^d}|a(x)|^2\,\mathrm dx.
\]
再直接用 Plancherel 求 \(\int|I_\lambda(p)|^2\,\mathrm dp\)，解释它与上述恒等式为何相符。
\end{exercise}
\begin{solution}
二阶 Taylor 积分公式给出
\[
 \psi(x)-p\cdot x
 \geq\psi(0)+(\nabla\psi(0)-p)\cdot x+\frac c2|x|^2.
\]
因此它在无穷远趋于正无穷，连续性使最小值在某个紧球内达到。最小点满足 \(\nabla\psi(x)=p\)。对任意 \(x,z\)，沿线段积分又有
\[
 (\nabla\psi(x)-\nabla\psi(z))\cdot(x-z)
       \geq c|x-z|^2,
\]
故梯度映射单射。它处处导数可逆，逆函数定理说明全局逆映射 \(x(p)\) 光滑，并满足
\[
 D_px(p)=\bigl[D^2\psi\bigl(x(p)\bigr)\bigr]^{-1}.
\]
对 \(S\) 求导时，来自 \(\psi\bigl(x(p)\bigr)\) 的链式项与 \(p\cdot x(p)\) 中的对应项相消，所以
\[
 \nabla S(p)=-x(p),
 \quad D^2S(p)=-\bigl[D^2\psi\bigl(x(p)\bigr)\bigr]^{-1}.
\]
这个符号取决于题中选取 \(\psi(x)-p\cdot x\) 的次序。

相位有唯一驻点 \(x(p)\)，其 Hesse 矩阵正定，故\cref{b2:thm:21-stationary-phase}给出
\[
 I_\lambda(p)=e^{2\pi\mathrm i\lambda S(p)}\lambda^{-d/2}
 e^{\mathrm i\pi d/4}\cdot
 \frac{a\bigl(x(p)\bigr)}{\det\Bigl(D^2\psi\bigl(x(p)\bigr)\Bigr)^{1/2}}
 +O(\lambda^{-d/2-1}).
\]
给定紧参数集 \(P\)，先取包含它且闭包紧的参数开邻域 \(V\)，再取一个紧集 \(K\)，包含 \(\operatorname{supp}a\) 和所有 \(p\in\overline V\) 对应的驻点。这样的 \(K\) 存在，因为 \(x(p)\) 连续。每个 \(p\in V\) 的驻点都在同一个 \(K\) 内，Hesse 矩阵有同一个正下界，所需导数在相应紧邻域内也一致有界。因此\cref{b2:prop:21-parametric-stationary}可一致应用。若驻点不在振幅支撑内，首项为零；围绕该驻点局部化后，余下积分仍由统一非驻相估计处理。

梯度映射是全局微分同胚，Jacobi 行列式为正数 \(\det D^2\psi(x)\)。上述首项中的 \(b\) 因而满足
\[
 \int|b(p)|^2\,\mathrm dp
 =\int\frac{|a(x)|^2}{\det D^2\psi(x)}
                    \cdot\det D^2\psi(x)\,\mathrm dx
 =\|a\|_2^2.
\]
另一方面，令 \(g_\lambda=e^{2\pi\mathrm i\lambda\psi}\cdot a\)，则
\(I_\lambda(p)=\widehat g_\lambda(\lambda p)\)。Plancherel 与换元直接给出准确等式
\[
 \int|I_\lambda(p)|^2\,\mathrm dp
       =\lambda^{-d}\cdot\|g_\lambda\|_2^2
       =\lambda^{-d}\cdot\|a\|_2^2.
\]
驻相首项的平方积分恰有同一个值。梯度换元给出首项的平方积分，Plancherel 则给出原振荡积分的准确范数，两者由此相互对应。
\end{solution}

% 自拟；中等综合：固定有限维的分离双线性模型，逐坐标证明含退化参数的尖锐两侧范数。
\Needspace{6\baselineskip}
\begin{exercise}\label{b2:ex:21-new-4}
固定整数 \(d\geq1\) 及实函数 \(\chi\in C_c^\infty(\mathbb R)\)，满足 \(\chi\geq0\)，且在 \([-1,1]\) 上恒为一。设 \(\lambda\geq1\)、\(\varepsilon_1,\ldots,\varepsilon_d\geq0\)，对 \(f\in L^2(\mathbb R^d)\) 定义
\[
 T_{\lambda,\varepsilon}f(x)
 =\left(\prod_{j=1}^d\chi(x_j)\right)
   \int_{\mathbb R^d}
     e^{2\pi\mathrm i\lambda\sum_{j=1}^d\varepsilon_jx_jy_j}
     \cdot\left(\prod_{j=1}^d\chi(y_j)\right)\cdot f(y)\,\mathrm dy.
\]
令 \(w(\mu)=1\) 当 \(0\leq\mu\leq1\)，\(w(\mu)=\mu^{-1/2}\) 当 \(\mu>1\)。证明存在只依赖 \(d,\chi\) 的正数 \(c_d,C_d\)，使
\[
 c_d\prod_{j=1}^dw(\lambda\varepsilon_j)
 \leq\|T_{\lambda,\varepsilon}\|_{2\to2}
 \leq C_d\prod_{j=1}^dw(\lambda\varepsilon_j).
\]
先证明一维结论：上界分别用核的平方积分和缩放 Plancherel；下界可固定一个非负光滑输入，再利用其 Fourier 变换在原点附近不为零。最后用 Fubini 与乘积输入处理 \(d\) 维。
\end{exercise}
\begin{solution}
在一维记
\[
 S_\mu f(x)=\chi(x)\cdot
       \int e^{2\pi\mathrm i\mu xy}\cdot\chi(y)\cdot f(y)\,\mathrm dy.
\]
核的平方积分范数为 \(\|\chi\|_2^2\)，所以对所有 \(\mu\geq0\)，
\(\|S_\mu\|\leq\|\chi\|_2^2\)。当 \(\mu>0\) 时，
\(S_\mu f(x)=\chi(x)\cdot\widehat{\chi\cdot f}(-\mu x)\)，故
\[
 \|S_\mu f\|_2
 \leq\|\chi\|_\infty\cdot\mu^{-1/2}\cdot\|\chi\cdot f\|_2
 \leq\|\chi\|_\infty^2\cdot\mu^{-1/2}\cdot\|f\|_2.
\]
合并两种界得到 \(\|S_\mu\|\leq C\cdot w(\mu)\)。零参数时使用第一个界，不对零作负幂运算。

取固定的非零非负 \(g\in C_c^\infty\bigl((-1,1)\bigr)\)，归一化使 \(\|g\|_2=1\)。因为 \(\widehat g(0)=\int g>0\)，连续性给出 \(\delta\in(0,1]\)、\(c_0>0\)，使 \(|\widehat g(\xi)|\geq c_0\) 于 \(|\xi|\leq\delta\)。在区间
\[
 |x|\leq\delta\cdot w(\mu)^2
\]
上，\(\chi(x)=1\)，且 \(|S_\mu g(x)|=|\widehat g(-\mu x)|\geq c_0\)。于是
\[
 \|S_\mu g\|_2^2
       \geq2\delta c_0^2\cdot w(\mu)^2.
\]
这给出与参数无关的一维下界。

多维算子等于按各坐标先后应用 \(S_{\lambda\varepsilon_j}\)。对其他变量固定后使用一维上界，再用 Fubini 积分，可逐个得到乘积上界；先对紧支撑简单函数证明，再由密度延拓到 \(L^2\) 即可。取乘积单位输入 \(f(y)=\prod_jg(y_j)\)，输出也分解成 \(\prod_jS_{\lambda\varepsilon_j}g(x_j)\)，其平方范数由 Tonelli 分解为一维平方范数的乘积，得到下界。

因此每个坐标都在 \(\lambda\varepsilon_j\) 越过一时才开始贡献半个衰减幂次。零参数方向没有振荡增益；若所有有效参数都至少为一，才可将乘积改写为 \(\lambda^{-d/2}\prod_j\varepsilon_j^{-1/2}\)。
\end{solution}

% 自拟；反例与计算：混合矩阵处处可逆不保证全局单射，两个重复片仍满足局部算子衰减。
\Needspace{6\baselineskip}
\begin{exercise}\label{b2:ex:21-new-5}
在 \(\mathbb R_x^2\times\mathbb R_y^2\) 上令
\[
 F(s,t)=(e^s\cos t,e^s\sin t),
 \quad x=(s,t),\quad \Phi(x,y)=F(x)\cdot y.
\]
固定 \(0<\delta<\pi/4\)，取实函数 \(\chi_0\in C_c^\infty\bigl((-\delta,\delta)^2\bigr)\)，在原点附近恒为一；令 \(v=(0,2\pi)\)、\(\chi_1(x)=\chi_0(x-v)\)、\(\chi=\chi_0+\chi_1\)。再取非零 \(\eta\in C_c^\infty(\mathbb R^2)\)，对 \(\lambda\geq1\)、\(f\in L^2(\mathbb R^2)\) 定义
\[
 T_\lambda f(x)=\chi(x)\cdot
       \int e^{2\pi\mathrm i\lambda\Phi(x,y)}\cdot\eta(y)\cdot f(y)\,\mathrm dy.
\]
证明混合 Hesse 矩阵处处可逆，但 \(T_\lambda T_\lambda^*\) 的核不能在整个支撑上满足任何固定 \(N>0\) 的统一估计
\(C_N(1+\lambda|x-z|)^{-N}\)。另一方面，直接利用两个坐标片中的变量替换证明
\(\|T_\lambda\|_{2\to2}\leq C\lambda^{-1}\)。若 \(T_{\lambda,0}\) 表示仅用 \(\chi_0\) 的算子，进一步证明
\(\|T_\lambda f\|_2^2=2\|T_{\lambda,0}f\|_2^2\)。
\end{exercise}
\begin{solution}
直接求导给出
\[
 DF(s,t)=e^s
 \begin{pmatrix}\cos t&-\sin t\\\sin t&\cos t\end{pmatrix},
 \quad\det DF(s,t)=e^{2s}>0.
\]
按正文的约定，混合矩阵第 \((j,k)\) 项为 \(\partial_{y_j}\partial_{x_k}\Phi\)，此处它就是 \(DF(x)\)，所以处处可逆。但 \(F(x+v)=F(x)\)，映射并不全局单射。

相应的 \(TT^*\) 核是
\[
 K_\lambda(x,z)=\chi(x)\cdot\overline{\chi(z)}\cdot
       \int e^{2\pi\mathrm i\lambda(F(x)-F(z))\cdot y}
                                      \cdot|\eta(y)|^2\,\mathrm dy.
\]
取 \(x=0\)、\(z=v\)，两个截断值都为一，且相位差恒零，因此
\(K_\lambda(0,v)=\|\eta\|_2^2>0\)。而所提上界在这两点为 \(C_N(1+2\pi\lambda)^{-N}\to0\)，矛盾。

在第一个矩形中，极角限制于 \((-\delta,\delta)\)，所以 \(F\) 单射并为其像上的微分同胚。写 \(q=\eta\cdot f\)，则
\[
 T_{\lambda,0}f(x)=\chi_0(x)\cdot\widehat q(-\lambda F(x)).
\]
令 \(u=F(x)\)。该矩形中 \(e^{2s}\) 有正下界，且 \(\chi_0\) 有界，故
\[
 \begin{aligned}
 \|T_{\lambda,0}f\|_2^2
 &\leq C\int_{F\bigl((-\delta,\delta)^2\bigr)}
                      |\widehat q(-\lambda u)|^2\,\mathrm du\\
 &\leq C\lambda^{-2}\cdot\|q\|_2^2
 \leq C\lambda^{-2}\cdot\|\eta\|_\infty^2\cdot\|f\|_2^2.
 \end{aligned}
\]
输入 \(q\) 因紧支撑和 Cauchy--Schwarz 属于 \(L^1\cap L^2\)，所用 Fourier 表达合法。

两个截断支撑互不相交。在第二片作平移 \(x\mapsto x+v\)，有
\(\chi_1(x+v)=\chi_0(x)\) 和 \(F(x+v)=F(x)\)，所以第二片输出是第一片输出的平移副本。平方积分相加恰得到题中的二倍恒等式，继而得到整体 \(C\lambda^{-1}\) 界。失效的是跨越两片的单一梯度差下界；各片中的估计及其有限求和仍然有效。
\end{solution}

% 自拟；中等：不同输入输出维数中的秩模型，分离无振荡变量，计算准确范数并证明达到。
\Needspace{6\baselineskip}
\begin{exercise}\label{b2:ex:21-new-6}
设整数 \(r\geq1\)、\(m>r\)、\(n>r\)，实矩阵 \(B\in\mathbb R^{r\times r}\) 可逆。写
\(x=(x',x'')\in\mathbb R^r\times\mathbb R^{m-r}\)、
\(y=(y',y'')\in\mathbb R^r\times\mathbb R^{n-r}\)。给定非零
\(u\in L^2(\mathbb R^{m-r})\)、\(v\in L^2(\mathbb R^{n-r})\)，对 \(\lambda>0\) 及 \(f\in C_c^\infty(\mathbb R^n)\) 定义
\[
 T_\lambda f(x)
 =u(x'')\cdot\iint
       e^{2\pi\mathrm i\lambda x'\cdot By'}
       \cdot\overline{v(y'')}\cdot f(y',y'')\,\mathrm dy''\,\mathrm dy'.
\]
证明此算子唯一延拓为 \(L^2(\mathbb R^n)\to L^2(\mathbb R^m)\) 的有界算子，并求其准确范数。构造达到范数的非零 \(L^2\) 输入，说明衰减幂次为何由 \(r\) 决定，而不由 \(m\)、\(n\) 中任一个单独决定。
\end{exercise}
\begin{solution}
先对 \(y''\) 作配对，令
\[
 h(y')=\int\overline{v(y'')}\cdot f(y',y'')\,\mathrm dy''.
\]
对一般 \(L^2\) 输入，Fubini 与 Cauchy--Schwarz 保证此积分对几乎每个 \(y'\) 绝对收敛，并有
\[
 \|h\|_2\leq\|v\|_2\cdot\|f\|_2.
\]
在原来的光滑紧支撑输入上，\(h\) 还属于 \(L^1\)，原积分因而可识别为
\[
 T_\lambda f(x)=u(x'')\cdot\widehat h(-\lambda B^{\mathsf T}x').
\]
在 \(x'\) 上换元并使用 Plancherel，得到
\[
 \|T_\lambda f\|_2^2
 =\|u\|_2^2\cdot\lambda^{-r}\cdot|\det B|^{-1}\cdot\|h\|_2^2
 \leq\lambda^{-r}\cdot|\det B|^{-1}
             \cdot\|u\|_2^2\cdot\|v\|_2^2\cdot\|f\|_2^2.
\]
因此由 \(C_c^\infty\) 在 \(L^2\) 中的稠密性可唯一延拓；也可以直接用一般 \(L^2\) 的部分配对与 Plancherel 变换定义同一延拓。

取任意单位向量 \(h_0\in L^2(\mathbb R^r)\)，令
\[
 f(y',y'')=h_0(y')\cdot\frac{v(y'')}{\|v\|_2}.
\]
它是单位输入，部分配对恰为 \(h=\|v\|_2\cdot h_0\)，所以前面的不等式取等，得到
\[
 \boxed{\ \|T_\lambda\|_{2\to2}
 =\lambda^{-r/2}\cdot|\det B|^{-1/2}
                   \cdot\|u\|_2\cdot\|v\|_2\ }.
\]
相位 \(x'\cdot By'\) 的混合矩阵按输入变量作行指标，是一个左上块为 \(B^{\mathsf T}\)、其余为零的 \(n\times m\) 矩阵，秩为 \(r\)。变量 \(y''\) 只参与有界线性配对，变量 \(x''\) 只提供固定的 \(L^2\) 因子，两者都不增加振荡衰减。这个模型把混合相位的振荡方向与其余变量的配对作用明确分开。
\end{solution}

% 自拟；中等：坐标换元的酉归一化，推导两侧换元后的振幅及混合矩阵，检验局部估计的坐标一致性。
\Needspace{6\baselineskip}
\begin{exercise}\label{b2:ex:21-new-7}
设整数 \(d\geq1\)，\(U,V,X,Y\subset\mathbb R^d\) 为开集，\(F:U\to V\)、\(G:X\to Y\) 为光滑微分同胚，记
\(J_F=|\det DF|\)、\(J_G=|\det DG|\)。证明
\[
 (U_Ff)(x)=J_F(x)^{1/2}\cdot f\bigl(F(x)\bigr)
\]
给出 \(L^2(V)\to L^2(U)\) 的酉算子，并求其逆；\(U_G\) 同样定义。对 \(\lambda\geq1\) 及 \(f\in L^2(Y)\)，设
\[
 T_\lambda f(v)=\int_Y e^{2\pi\mathrm i\lambda\Phi(v,y)}
                                  \cdot a(v,y)\cdot f(y)\,\mathrm dy,
\]
其中实相位 \(\Phi\) 在 \(V\times Y\) 上光滑，\(a\in C_c^\infty(V\times Y)\)。求
\(\widetilde T_\lambda=U_FT_\lambda U_G^{-1}\) 的相位、振幅与混合 Hesse 矩阵，证明两算子的 \(L^2\) 范数完全相同，并说明局部非退化条件如何在换元下保持。不能只改变相位而省去振幅中的体积因子。
\end{exercise}
\begin{solution}
变量替换给出
\[
 \int_U J_F(x)\cdot|f\bigl(F(x)\bigr)|^2\,\mathrm dx
       =\int_V|f(v)|^2\,\mathrm dv.
\]
其逆公式为
\[
 (U_F^{-1}h)(v)
 =J_F\bigl(F^{-1}(v)\bigr)^{-1/2}\cdot h\bigl(F^{-1}(v)\bigr).
\]
代入并换元可验证逆映射也保范数，且两次复合为恒等，所以 \(U_F\) 为满射等距映射，即酉算子。\(G\) 的结论相同。

把逆公式代入原算子，再令 \(y=G(z)\)。输入端的逆平方根与 \(\mathrm dy=J_G(z)\,\mathrm dz\) 相乘，得到
\[
 \begin{aligned}
 \widetilde\Phi(x,z)&=\Phi\bigl(F(x),G(z)\bigr),\\
 \widetilde a(x,z)&=J_F(x)^{1/2}\cdot J_G(z)^{1/2}
                                \cdot a\bigl(F(x),G(z)\bigr),\\
 \widetilde T_\lambda h(x)
 &=\int_Xe^{2\pi\mathrm i\lambda\widetilde\Phi(x,z)}
                            \cdot\widetilde a(x,z)\cdot h(z)\,\mathrm dz.
 \end{aligned}
\]
原振幅的紧支撑经乘积微分同胚的逆像仍是紧集，故新振幅也是光滑紧支撑的。

令 \(M(v,y)\) 的第 \((j,k)\) 项为 \(\partial_{y_j}\partial_{v_k}\Phi\)，与正文采用同一约定。因两套换元分别只依赖自己的变量，链式法则没有额外的二阶交叉项，给出
\[
 \widetilde M(x,z)
 =DG(z)^{\mathsf T}\,M\bigl(F(x),G(z)\bigr)\,DF(x).
\]
两端矩阵可逆，因此 \(\widetilde M\) 可逆当且仅当对应的 \(M\) 可逆。固定紧支撑附近的局部估计仍可应用；变换后的具体导数界应按新坐标计算。

另一方面，酉映射把整个单位球双射到单位球，故直接有准确等式
\[
 \|\widetilde T_\lambda\|_{L^2(X)\to L^2(U)}
       =\|T_\lambda\|_{L^2(Y)\to L^2(V)}.
\]
这比仅分别应用两次估计更直接。平方根权重正是平方范数与变量替换相容所需的因子；例如在一维取实数 \(c\ne0\)，单独拉回 \(f\mapsto f(cx)\) 的范数为 \(|c|^{-1/2}\)，一般并不保范数。遗漏权重就改变了当前比较的算子。
\end{solution}

% 自拟；提高且依赖选读内容：对频率局部波包作Taylor误差估计，以实际能量集中检验半波传播正负号。
\Needspace{6\baselineskip}
\begin{optionalexercise}\label{b2:ex:21-new-8}
设整数 \(d\geq1\)，\(\omega\in\mathbb R^d\)、\(|\omega|=1\)，\(x_0\in\mathbb R^d\)，\(0<\alpha<1/2\)。取 \(h\in\mathcal S(\mathbb R^d)\)、\(\|h\|_2=1\)，并假定 \(\operatorname{supp}\widehat h\subset B(0,R)\)，其中 \(R>0\)。对 \(N\to\infty\)，定义单位波包
\[
 f_N(x)=N^{\alpha d/2}e^{2\pi\mathrm iN\omega\cdot x}
                       \cdot h\bigl(N^\alpha(x-x_0)\bigr).
\]
令 \(\sigma\in\{1,-1\}\)，半波算子 \(U_t^\sigma\) 的 Fourier 乘子为
\(e^{\sigma2\pi\mathrm it|\xi|}\)，其中 \(t\in\mathbb R\)。证明当 \(N^{1-\alpha}\geq2R\) 时，
\[
 \begin{aligned}
 v_N^\sigma(t,x)
 &=N^{\alpha d/2}e^{2\pi\mathrm i(N\omega\cdot x+\sigma tN)}
                    \cdot h\bigl(N^\alpha(x-x_0+\sigma t\omega)\bigr),\\
 \|U_t^\sigma f_N-v_N^\sigma(t,\cdot)\|_2
 &\leq C|t|N^{2\alpha-1}
                \bigl\|\,|\eta|^2\widehat h(\eta)\,\bigr\|_2,
 \end{aligned}
\]
其中 \(C\) 与 \(N,t\) 无关。由此证明，对固定 \(t\) 和任意有界连续函数 \(b\)，
\[
 \int b(x)\cdot|U_t^\sigma f_N(x)|^2\,\mathrm dx
       \longrightarrow b(x_0-\sigma t\omega).
\]
用所得位置说明正、负两个演化指数的传播方向；并指出允许时间依赖 \(N\) 时，上述近似保证收敛的时间尺度。
\end{optionalexercise}
\begin{solution}
调制、平移与缩放公式给出
\[
 \widehat f_N(\xi)
 =N^{-\alpha d/2}e^{-2\pi\mathrm ix_0\cdot(\xi-N\omega)}
                \cdot\widehat h\bigl(N^{-\alpha}(\xi-N\omega)\bigr).
\]
它只支撑于 \(\xi=N\omega+N^\alpha\eta\)、\(|\eta|\leq R\) 的区域。在题设范围内，连接 \(N\omega\) 与这些频率的线段距原点至少为 \(N/2\)。函数 \(\xi\mapsto|\xi|\) 的 Hesse 矩阵在该线段上的算子范数至多为 \(2/N\)，Taylor 积分余项因而给出
\[
 \bigl||N\omega+N^\alpha\eta|-N-N^\alpha\omega\cdot\eta\bigr|
       \leq C N^{2\alpha-1}|\eta|^2.
\]
将半波乘子中的长度换成线性项
\(N+\omega\cdot(\xi-N\omega)\)，Fourier 反演恰好得到题中的 \(v_N^\sigma\)：常数项产生 \(e^{\sigma2\pi\mathrm itN}\)，线性项使包络平移为 \(x-x_0+\sigma t\omega\)。

利用 \(|e^{\mathrm ia}-e^{\mathrm ib}|\leq|a-b|\)，再用 Plancherel 和变量替换 \(\xi=N\omega+N^\alpha\eta\)，得到
\[
 \|U_t^\sigma f_N-v_N^\sigma(t,\cdot)\|_2^2
 \leq C^2t^2N^{4\alpha-2}
       \int|\eta|^4\cdot|\widehat h(\eta)|^2\,\mathrm d\eta.
\]
开平方即为误差界。两种函数的 \(L^2\) 范数都是一，所以 Cauchy--Schwarz 还给出
\[
 \bigl\|\,|U_t^\sigma f_N|^2-|v_N^\sigma(t,\cdot)|^2\,\bigr\|_1
       \leq2\|U_t^\sigma f_N-v_N^\sigma(t,\cdot)\|_2.
\]
对固定 \(t\)，右侧因 \(2\alpha-1<0\) 而趋零。

包络的概率密度则可以直接换元：
\[
 \int b(x)\cdot|v_N^\sigma(t,x)|^2\,\mathrm dx
 =\int b\bigl(x_0-\sigma t\omega+N^{-\alpha}z\bigr)
                                      \cdot|h(z)|^2\,\mathrm dz.
\]
被积函数由 \(\|b\|_\infty\cdot|h|^2\) 控制，连续性与支配收敛给出极限 \(b(x_0-\sigma t\omega)\)。再用上面的密度 \(L^1\) 误差，即得真实演化的同一极限。

因此正指数 \(e^{2\pi\mathrm it|\xi|}\) 的包络中心为 \(x_0-t\omega\)，负指数的中心为 \(x_0+t\omega\)。若时间取 \(t=t_N\)，误差估计保证近似成立的充分条件是
\(|t_N|N^{2\alpha-1}\to0\)，即 \(|t_N|=o(N^{1-2\alpha})\)。这只是当前线性近似保证的尺度；在更长时间上，估计不能直接断言近似成立或失效。
\end{solution}

% !TeX root = ../../Book2.tex
\chapter{Fourier 乘子}
\label{b2:ch:22}
\BookChapterNumber{b2:ch:22}
\begin{chapterabstract}
有界符号在 Fourier 侧定义了一个 \(L^2\) 算子，但同一个算子是否保持 \(L^p\)，还取决于频率之间如何共同作用。本章先建立乘子的代数、对偶、极限与对称性质，再把环带上的导数或 Sobolev 控制转为奇异积分估计。应用部分区分零阶变换、分数阶正则性、同指数椭圆估计和低频逆算子，说明符号大小、光滑性与核的可积性分别能够提供什么。
\end{chapterabstract}

\begin{learninggoals}\begin{itemize}
\item\label{b2:item:22-goal-1}从有界符号的 \(L^2\) 算子出发定义 \(L^p\) 乘子，准确使用复合、对偶、极限及尺度对称，并区分强收敛与算子范数收敛。
\item\label{b2:item:22-goal-2}把 Mihlin 导数条件转为归一化环带核估计，说明低、高频求和与 \(L^2\) 起点各自的作用。
\item\label{b2:item:22-goal-3}理解 Hörmander 环带 Sobolev 条件怎样给出积分差分控制，识别严格指数条件实际用于哪一步。
\item\label{b2:item:22-goal-4}在 Riesz 变换、Bessel 位势及椭圆方程中选择匹配的空间，区分同指数有界性、正则性提升和低频存在性问题。
\end{itemize}\end{learninggoals}

对一个固定频率作乘法很简单，困难在于不同频率重新叠加后，空间中的 \(L^p\) 大小会如何变化。Plancherel 使 \(p=2\) 的答案完全由符号的本性上确界决定；其他指数没有同样的等距性质。即使符号处处具有模一，或者支撑在一个有界集合中，也不能只凭这些信息得出全部 \(L^p\) 界。

\ProseParagraphBreak

\secref{b2:sec:2201}先说明所讨论的是哪一个算子。对于有界可测符号，\(L^2\) 定义不需要预先寻找普通函数核；向有限指数 \(L^p\) 延拓时，再核对共同定义域上的一致性。无穷指数不能通过 Schwartz 函数的范数稠密性处理，而应使用有界 \(L^1\) 算子的伴随。恒等算子、平移以及区间投影也提醒我们：乘子理论不局限于具有可积卷积核的算子。

\ProseParagraphBreak

主定理将\chapref{b2:ch:18}的奇异积分方法与\chapref{b2:ch:19}的频率分解接起来。在每个环带上缩放到固定大小，再估计逆 Fourier 变换，便可把符号条件转为核条件。Mihlin 版本使用有限阶逐点导数，Hörmander 版本使用环带上的 Sobolev 范数；后者保留证明需要的积分差分控制，而不要求先得到极限核的逐点梯度界。单个环带的可积核估计不能不经求和论证就用于无限多个尺度。

\ProseParagraphBreak

阅读应用时，应先分开观察零频和无穷频率。两个符号可以在高频具有相同阶数，却因零频行为不同而作用于不同的空间；对最高阶导数的先验界，也不自动提供一个同指数的无质量逆算子。本章所需预备知识是 Fourier 反演、Plancherel、Young 不等式、插值与对偶，以及已建立的 Calderón--Zygmund 和平方函数估计。章末题沿这些工具安排，不要求预先使用后章的球面乘子结论。

% !TeX root = ../../Book2.tex
\section{乘子算子}
\label{b2:sec:2201}

常系数微分算子经过 Fourier 变换成为多项式乘法，卷积算子则成为其核的 Fourier 变换所给出的乘法。求解方程时，还会遇到多项式的倒数、平方根或更一般的函数。因此，一个共同的问题是：对每个频率分别乘上 \(m(\xi)\)，所得操作能否在原来的函数空间中保持有界？

\ProseParagraphBreak
在 \(L^2\) 上，Plancherel 等式直接回答这个问题；在其他 \(L^p\) 上，频率之间的叠加不能由平方和完全描述，符号的振荡和不连续性便会产生影响。本节先把算子的定义和基本运算说清，再说明平移不变性为何引出乘子。下一节的尺度不变光滑性条件，将给出容易应用的一类充分判据。

\subsection{Fourier 乘子}
\label{b2:subsec:220101}

本章沿用
\[
 \widehat f(\xi)=\int_{\mathbb R^d}
 f(x)e^{-2\pi\mathrm i x\cdot\xi}\,dx
\]
的归一化。若 \(m\) 只是可测函数，一般不能把 \(m\widehat f\) 解释为任意两个分布的乘积。先取 \(m\in L^\infty(\mathbb R^d)\)，便可对每个 \(f\in L^2\) 定义
\[
 T_m f\coloneqq\mathcal F^{-1}(m\widehat f)\in L^2.
\]
这一步不要求 \(m\) 连续，也不要求它有普通函数意义下的逆 Fourier 变换。

\Needspace{6\baselineskip}
\begin{definition}[Fourier 乘子]\label{b2:def:22-multiplier}
设 \(d\geq1\)、\(1\leq p<\infty\)，且
\(m\in L^\infty(\mathbb R^d)\)。若存在常数 \(C\)，使以上在
\(L^2\) 中定义的算子对全部 \(f\in\mathcal S(\mathbb R^d)\) 满足
\[
 \|T_m f\|_p\leq C\|f\|_p,
\]
称 \(m\) 为 \(L^p\) \term[fourier-chengzi]{Fourier 乘子}[Fourier multiplier]，记 \(m\in M_p\)。该算子唯一延拓到
\(L^p\)，其范数记为 \(\|m\|_{M_p}\)。
对于 \(p=\infty\)，本章约定
\[
 m\in M_\infty\quad\Longleftrightarrow\quad
 \overline m\in M_1,
\]
并把 \(T_m\) 定义为 \(T_{\overline m}:L^1\to L^1\) 在
\(\langle f,g\rangle=\int f\overline g\) 配对下的伴随；
记 \(\|m\|_{M_\infty}=\|\overline m\|_{M_1}\)。
\end{definition}

当 \(p<\infty\) 时，唯一性来自 Schwartz 函数在 \(L^p\) 中的稠密性。
若 \(f\in L^2\cap L^p\)，先截断再作光滑化，可以找到在这两个范数中同时趋近 \(f\) 的 Schwartz 函数。因此，延拓后的 \(T_m f\) 与原来的 \(L^2\) 定义相同：两种范数收敛各自给出依测度收敛的子列，极限必一致。以后在共同定义域内使用 Fourier 计算，依据的就是这一相容性。

\ProseParagraphBreak
\(L^\infty\) 的情形必须另作说明。Schwartz 函数并不在其中范数稠密，例如常数 \(1\) 到任意趋于零的连续函数的 \(L^\infty\) 距离都不小于 \(1\)。上述伴随约定选出与 \(L^1\) 配对相容的 \WeakStar{}连续算子；它在 \(L^\infty\cap L^2\) 上也与 \(L^2\) 定义一致。事实上，对 \(g\in L^1\cap L^2\)，Plancherel 等式给出
\[
 \langle T_m f,g\rangle
 =\langle f,T_{\overline m}g\rangle,
\]
这个配对确定了同一个函数。以后涉及端点 \(\infty\) 的乘子，均采用此约定，而不从 Schwartz 类上的范数估计声称得到了任意意义下的唯一延拓。

\ProseParagraphBreak
最直接的例子来自卷积。若 \(k\in L^1\)、\(m=\widehat k\)，则
\[
 T_m f=k*f,\quad
 \|m\|_{M_p}\leq\|k\|_1,\quad 1\leq p\leq\infty.
\]
这里先在 Schwartz 类上用 Fubini 验证 Fourier 恒等式，再应用 Young 不等式。对于有限复 Borel 测度 \(\mu\)，同样有
\[
 m(\xi)=\int e^{-2\pi\mathrm i y\cdot\xi}\,d\mu(y),
 \quad T_m f(x)=\int f(x-y)\,d\mu(y),
 \quad \|m\|_{M_p}\leq|\mu|(\mathbb R^d).
\]
因此，常数符号 \(1\) 的核为 \(\delta_0\)，指数符号
\(e^{-2\pi\mathrm i y\cdot\xi}\) 对应平移。可积核只是乘子的一部分；Hilbert 变换的核则要用主值分布解释。下一节从局部化符号重构核时，首先得到的是原点外的普通函数，原点处可能存在的分布项仍由 \(L^2\) 定义确定。

\subsection{\texorpdfstring{\(L^2\)}{L2} 理论}
\label{b2:subsec:220102}

微分算子的符号通常无界，仍可在 \(L^2\) 的稠密子空间上定义。为准确地区别无界算子与有界乘子，对任意几乎处处有限的可测 \(m\)，令
\[
 D_m\coloneqq\{f\in L^2:m\widehat f\in L^2\},
 \quad T_m f=\mathcal F^{-1}(m\widehat f),\quad f\in D_m.
\]
该定义域稠密：把任意 \(\widehat f\in L^2\) 乘上
\(\mathbf1_{\{|\xi|\leq n,\ |m(\xi)|\leq n\}}\)，再令 \(n\to\infty\)，即可由支配收敛得到 \(L^2\) 逼近。

\Needspace{6\baselineskip}
\begin{proposition}\label{b2:prop:22-l2}
设 \(m:\mathbb R^d\to\mathbb C\) 可测且几乎处处有限。在自然定义域
\(D_m\) 上定义的 \(T_m\) 能延拓为 \(L^2\) 有界算子，当且仅当
\(m\in L^\infty\)。此时
\[
 \|m\|_{M_2}=\|m\|_\infty.
\]
\end{proposition}
\begin{proof}
若 \(m\in L^\infty\)，Plancherel 等式直接给出
\[
 \|T_m f\|_2^2=\int|m(\xi)|^2\cdot|\widehat f(\xi)|^2\,d\xi
 \leq\|m\|_\infty^2\|f\|_2^2.
\]
反过来，设在 \(D_m\) 上有界，范数上界为 \(C\)。对任意有限测度集合 \(E\)，令
\[
 E_n=E\cap\{|\xi|\leq n,\ |m(\xi)|\leq n\}.
\]
取 \(\widehat f=\mathbf1_{E_n}\) 时确有 \(f\in D_m\)，故
\(\int_{E_n}|m|^2\leq C^2|E_n|\)。由单调收敛，
\(\int_E|m|^2\leq C^2|E|\)。
如果 \(\{|m|>C+\epsilon\}\) 有正测度，取其中有限正测度子集代入，就产生矛盾。
因此 \(\|m\|_\infty\leq C\)，与充分性合起来得到范数等式。
\end{proof}

例如 \(|m|=1\) 时，\(T_m\) 在 \(L^2\) 上等距；若 \(m=\mathbf1_E\)，则 \(T_m\) 是频率集合 \(E\) 上的正交投影。无论 \(E\) 的边界多么复杂，这个 \(L^2\) 结论都成立。其他 \(p\) 的困难不在逐点乘法本身，而在 Fourier 变换不能再把原空间范数等同于频率空间的 \(L^p\) 范数。

\ProseParagraphBreak
有界符号的最大值还是任何乘子范数的下界。证明的办法是把 Fourier 质量集中在一个频率附近：在空间侧用越来越缓慢变化的函数，再乘上一个固定频率的指数。此处不能直接代入 \(e^{2\pi\mathrm i x\cdot\xi_0}\)，因为它不属于有限指数的 \(L^p(\mathbb R^d)\)。

\Needspace{6\baselineskip}
\begin{proposition}[符号的必要界]\label{b2:prop:22-symbol-bound}
设 \(d\geq1\)、\(1\leq p\leq\infty\)，且 \(m\in M_p\)，其中
\(M_\infty\) 采用\cref{b2:def:22-multiplier}的伴随约定。则
\[
 \|m\|_\infty\leq\|m\|_{M_p}.
\]
\end{proposition}
\begin{proof}
先设 \(1<p<\infty\)，\(p'=p/(p-1)\)。取
\(f(x)=e^{-\pi|x|^2}\)、\(g(x)=f(x)^{p-1}\)，并在 \(m\) 的 Lebesgue 点
\(\xi_0\) 处定义
\[
 f_\epsilon(x)=\epsilon^{d/p}e^{2\pi\mathrm i x\cdot\xi_0}f(\epsilon x),
 \quad
 g_\epsilon(x)=\epsilon^{d/p'}e^{2\pi\mathrm i x\cdot\xi_0}g(\epsilon x).
\]
两者的相应范数不随 \(\epsilon\) 改变。变换变量和 Plancherel 配对给出
\[
 \langle T_m f_\epsilon,g_\epsilon\rangle
 =\int m(\xi_0+\epsilon\eta)\widehat f(\eta)
       \overline{\widehat g(\eta)}\,d\eta.
\]
两个 Fourier 变换都是正 Gaussian 函数；其乘积除以积分后是一个
Gaussian 近似恒等核。Lebesgue 点性质因此使右端趋于
\(m(\xi_0)\int f\overline g\)。另一方面，左端的绝对值不超过
\(\|m\|_{M_p}\cdot\|f\|_p\cdot\|g\|_{p'}\)。
所选 \(g\) 恰使
\(\int f\overline g=\|f\|_p\cdot\|g\|_{p'}\)，故
\(|m(\xi_0)|\leq\|m\|_{M_p}\)。

当 \(p=1\) 时，仍取同一个 \(f\)，改用
\(g_\delta(x)=e^{-\pi\delta|x|^2}\)，其中 \(\delta>0\)。
对每个固定 \(\delta\)，以上计算以及 \(\|g_\delta\|_\infty=1\) 给出
\[
 |m(\xi_0)|\int f g_\delta\leq\|m\|_{M_1}\cdot\|f\|_1.
\]
再令 \(\delta\downarrow0\)，由支配收敛得到所需不等式。
\(p=\infty\) 的结论应用于 \(\overline m\in M_1\) 即得。
\end{proof}

这里从 Lebesgue 点得到的是几乎处处的结论，恰好对应本质上确界。
某个频率点上的任意赋值不影响乘子。这个下界通常不是充分条件：
\(L^2\) 的特殊性正是只凭它就能完全刻画有界性。

\subsection{\texorpdfstring{\(L^p\)}{Lp} 乘子问题}
\label{b2:subsec:220103}

在判别具体符号以前，先整理三种常用操作。它们使有限频率截断的计算能够进入无穷尺度的结论，也说明何时可以把两个已经构造的算子复合。

\Needspace{6\baselineskip}
\begin{proposition}[乘子的运算、对偶与极限]\label{b2:prop:22-multiplier-calculus}
设 \(d\geq1\)，并采用\cref{b2:def:22-multiplier}中的乘子定义。
\begin{enumerate}
 \item 对 \(1\leq p\leq\infty\)，若 \(m_1,m_2\in M_p\)，则
 \(m_1+m_2,m_1m_2\in M_p\)，且
 \[
 T_{m_1m_2}=T_{m_1}T_{m_2},\quad
 \|m_1m_2\|_{M_p}\leq\|m_1\|_{M_p}\cdot\|m_2\|_{M_p}.
 \]
 加法满足相应的三角不等式。
 \item 若 \(1<p<\infty\)，则
 \(m\in M_p\) 当且仅当 \(\overline m\in M_{p'}\)，并且
 \(\|m\|_{M_p}=\|\overline m\|_{M_{p'}}\)。
 两个算子在复 \(L^p\) 对偶下互为伴随。
 \item 对固定 \(1\leq p\leq\infty\)，若
 \(m_n\in M_p\)、\(\sup_n\|m_n\|_{M_p}\leq C\)，且
 \(m_n\to m\) 几乎处处，则 \(m\in M_p\)，
 \(\|m\|_{M_p}\leq C\)。
\end{enumerate}
\end{proposition}
\begin{proof}
先取 \(p<\infty\)。对 \(f\in L^2\cap L^p\)，
\(T_{m_2}f\) 仍在这两个空间内，故 Fourier 乘法可以连续使用两次：
\[
 \widehat{T_{m_1}T_{m_2}f}=m_1m_2\widehat f.
\]
两种延拓的相容性及 \(L^2\cap L^p\) 的稠密性给出算子恒等式与范数界。
加法的论证相同。

对第二条，在 \(f,g\in\mathcal S\) 上先用 \(L^2\) 配对写成
\[
 \langle T_m f,g\rangle
 =\int m\widehat f\,\overline{\widehat g}
 =\langle f,T_{\overline m}g\rangle.
\]
若 \(T_m\) 在 \(L^p\) 上有界，则右端作为 \(f\) 的泛函范数不超过
\(\|m\|_{M_p}\cdot\|g\|_{p'}\)。由 \(L^p\) 对偶性，
\(T_{\overline m}g\) 属于 \(L^{p'}\) 并满足相应界；
交换 \(p,p'\) 后得到范数等式。端点 \(p=\infty\) 的第一条则由
\(L^1\) 算子的伴随得到，注意标量乘子可交换，所以取伴随时复合次序的反转不改变结果。

对第三条，\cref{b2:prop:22-symbol-bound}先给出
\(\|m_n\|_\infty\leq C\)，因而 \(|m|\leq C\) 几乎处处。
当 \(f\in\mathcal S\) 时，支配收敛和 Plancherel 等式给出
\[
 \|T_{m_n}f-T_mf\|_2^2
 =\int|m_n-m|^2\cdot|\widehat f|^2\longrightarrow0.
\]
选取几乎处处收敛的子列。对于 \(p<\infty\)，Fatou 引理给出
\(\|T_mf\|_p\leq C\|f\|_p\)，正好是乘子定义所需的估计。
对于 \(p=\infty\)，把已经证明的 \(p=1\) 情形应用于
\(\overline{m_n}\)，再取伴随即可。
\end{proof}

极限结论只保证有界性，不保证算子范数收敛；即使符号逐点收敛，也可能在不断远移的频率处保留很大的差异。若 \(m\) 实值，对偶性给出 \(p\) 与 \(p'\) 的相同乘子范数。再与 \(L^2\) 界插值，可以得到已知指数与 \(2\) 之间的全部指数；这不能自动推出 \(1\) 或 \(\infty\) 端点。

\ProseParagraphBreak
乘子范数还具有若干精确的对称性。令
\[
 D_rf(x)=f(rx),\quad M_\eta f(x)=e^{2\pi\mathrm i\eta\cdot x}f(x),
 \quad \tau_yf(x)=f(x-y),\quad r>0.
\]
直接计算 Fourier 变换可得
\[
 \begin{aligned}
 T_{m(r\,\cdot)}&=D_{1/r}T_mD_r,\\
 T_{m(\,\cdot-\eta)}&=M_\eta T_mM_{-\eta},\\
 T_{e^{-2\pi\mathrm i y\cdot\xi}m(\xi)}&=\tau_yT_m.
 \end{aligned}
\]
\(D_r\) 的范数为 \(r^{-d/p}\)，其中 \(d/\infty=0\)，所以两个缩放范数相消。
调制和平移都是等距同构；把恒等式反向使用，得到这些操作都保持
\(\|m\|_{M_p}\)。正交坐标变换也有相同性质。对于端点 \(\infty\)，上述算子均与 \(L^1\) 配对相容，恒等式可由伴随验证。
这解释了下一节为什么估计
\(|\xi|^{|\alpha|}\partial^\alpha m(\xi)\)：缩放频率时，导数产生的尺度因子恰被前面的幂抵消。

\ProseParagraphBreak
一个不光滑但确实有界的例子是一维区间投影。Hilbert 变换 \(H\) 的符号为
\(-\mathrm i\,\operatorname{sgn}\xi\)，所以
\[
 P_+\coloneqq\tfrac12(I+\mathrm i H)
\]
在零点以外的符号为 \(\mathbf1_{(0,\infty)}\)。由\chapref{b2:ch:18}的
\(L^p\) 有界性，\(P_+\) 对 \(1<p<\infty\) 有界。若 \(a<b\)，则
\[
 T_{\mathbf1_{(a,b)}}=
 M_aP_+M_{-a}-M_bP_+M_{-b},
 \quad
 \|\mathbf1_{(a,b)}\|_{M_p}\leq2\|P_+\|_{p\to p},
\]
与区间的位置和长度无关。符号光滑性因此是充分条件而非必要条件；
也不能把这个一维结论直接推广到高维球的特征函数。边界曲率带来的困难将在\chapref{b2:ch:24}说明。

\ProseParagraphBreak
最后说明乘子与平移不变性的准确关系。卷积与平移交换很直观；
下面的反向结论说明，在 \(L^2\) 上，除此以外没有另一类有界平移不变算子。

\Needspace{6\baselineskip}
\begin{proposition}\label{b2:prop:22-translation-multiplier}
设 \(d\geq1\)。若复线性算子
\(T:L^2(\mathbb R^d)\to L^2(\mathbb R^d)\) 有界，且
\(T\tau_y=\tau_yT\) 对全部 \(y\in\mathbb R^d\) 成立，
其中 \(\tau_yf(x)=f(x-y)\)，则存在在几乎处处意义下唯一的
\(m\in L^\infty(\mathbb R^d)\)，使
\[
 T=T_m,\quad \|T\|_{2\to2}=\|m\|_\infty.
\]
反过来，每个有界 \(L^2\) 乘子都与所有平移交换。
\end{proposition}
\begin{proof}
令 \(A=\mathcal FT\mathcal F^{-1}\)。平移变成乘法
\(e^{-2\pi\mathrm i y\cdot\xi}\)，故 \(A\) 与所有这些乘法算子交换。
对任意 \(\varphi\in C_c^\infty\)，Fourier 反演把 \(\varphi\) 写成这些指数函数的
\(L^1\) 加权积分。更具体地，对每个 \(h\in L^2\)，
\[
 \varphi h=\int_{\mathbb R^d}
  (\mathcal F^{-1}\varphi)(y)e^{-2\pi\mathrm i y\cdot\xi}h(\xi)\,dy
\]
是 \(L^2\) 中的强积分，被积函数的范数不超过
\(|(\mathcal F^{-1}\varphi)(y)|\cdot\|h\|_2\)。
有界算子可移入这个积分，因而 \(A\) 与乘法 \(M_\varphi\) 交换。

接着把光滑截断换成可测集投影。若 \(E\) 有界，将 \(\mathbf1_E\) 与非负光滑近似恒等核卷积，可得
\(0\leq\varphi_n\leq1\)、\(\varphi_n\in C_c^\infty\)，且
\(\varphi_n\to\mathbf1_E\) 几乎处处。对每个 \(h\in L^2\)，
支配收敛说明 \(M_{\varphi_n}h\to M_{\mathbf1_E}h\)。
故 \(A\) 也与 \(M_{\mathbf1_E}\) 交换。再令有界截断集合递增到任意可测集，就得到一般情形。

取递增的有限测度立方体 \(Q_n\) 覆盖全空间。
由交换投影，\(A\mathbf1_{Q_n}\) 支持在 \(Q_n\) 内；记它为 \(m_n\)。
若 \(n<k\)，则
\[
 \mathbf1_{Q_n}m_k
 =A(\mathbf1_{Q_n}\mathbf1_{Q_k})=m_n.
\]
删去可数个零测集后，这些函数可以拼成可测函数 \(m\)。
对任意可测 \(E\subset Q_n\)，又有
\[
 A\mathbf1_E=\mathbf1_E m,\quad
 \int_E|m|^2\leq\|A\|^2|E|.
\]
取 \(E=\{|m|>\|A\|+\epsilon\}\cap Q_n\) 可知
\(\|m\|_\infty\leq\|A\|\)。由有限测度支集简单函数在 \(L^2\) 中的稠密性，
\(A=M_m\)。范数等式来自\cref{b2:prop:22-l2}；若两个符号给出同一算子，
将其差乘以任意有限测度集的特征函数，即得二者几乎处处相等。
反向结论只需注意两个逐点乘法算子可交换。
\end{proof}

该证明使用了 \(T\) 的平方界，尤其使用 Fourier 变换在 \(L^2\) 上的酉性。
若某算子最初只在另一个 \(L^p\) 空间给定，不能未经说明便把它搬到
\(L^2\) 上应用此命题。实际应用通常从一个有界符号出发，先建立
\(L^2\) 算子，再通过核估计证明它在其他 \(L^p\) 上也有界；这正是下面采用的顺序。

% !TeX root = ../../Book2.tex
\section{Mihlin–Hörmander 定理}
\label{b2:sec:2202}

有界符号在 \(L^2\) 上的作用已经由 Plancherel 完全描述；要得到其他指数上的界，还须控制不同频率的输出如何在空间中叠加。本节用两种条件完成这一步。Mihlin 条件逐点控制符号导数，先得到核的大小与梯度估计；Hörmander 条件只控制每个环带上的 Sobolev 范数，得到的是核差分的积分估计。两条路线都以 \(L^2\) 有界性为起点，再使用\chapref{b2:ch:18}的分解与插值方法。

\subsection{符号导数}
\label{b2:subsec:220201}

设 \(d\geq1\)，\(N\) 为非负整数，\(m\) 是定义在 \(\mathbb R^d\setminus\{0\}\) 上的复值 \(C^N\) 函数。尺度不变的导数条件写成
\[
 |\xi|^{|\alpha|}\cdot|\partial^\alpha m(\xi)|\leq A,
 \quad 0\leq|\alpha|\leq N,\quad \xi\ne0.
\]
这里 \(\alpha=(\alpha_1,\ldots,\alpha_d)\) 是多重指标，\(|\alpha|=\alpha_1+\cdots+\alpha_d\)，\(\partial^\alpha=\partial_1^{\alpha_1}\cdots\partial_d^{\alpha_d}\)。零阶条件就是 \(|m|\leq A\)。若观察半径约为 \(2^j\) 的频率环带，把它缩放到单位尺度后，一阶导数带来 \(2^j\)，而假设中的 \(|\xi|^{-1}\) 恰好抵消这个因子；更高阶导数遵循同一规律。

\ProseParagraphBreak
例如 \(m(\xi)=\xi_\ell/|\xi|\) 是零次齐次的：\(m(t\xi)=m(\xi)\)，\(t>0\)。将此式对 \(\xi\) 求导可得
\[
 (\partial^\alpha m)(t\xi)
 =t^{-|\alpha|}\cdot(\partial^\alpha m)(\xi).
\]
每个导数在单位球面上连续，因此任意固定有限阶的尺度导数界都成立。这里常数可以随阶数增长，并不要求一个常数同时控制所有阶导数。

\ProseParagraphBreak
另一个例子是 \(m(\xi)=b(\log|\xi|)\)，其中 \(b\in C^N(\mathbb R;\mathbb C)\)，并且 \(b,b',\ldots,b^{(N)}\) 全部有界。链式法则把 \(\partial^\alpha m\) 写成若干项之和，每一项都是 \(|\xi|^{-|\alpha|}\)、单位方向的有界光滑函数和某个 \(b^{(r)}(\log|\xi|)\) 的乘积，所以也满足尺度导数界。这样的符号可以沿尺度持续振荡，在原点不必有极限。乘子由几乎处处的符号决定，原点的取值不影响算子；定理要求的是离开原点后的光滑性。

\ProseParagraphBreak
尺度因子也划定了本节的范围。对 \(\alpha>0\)，\(|\xi|^{-\alpha}\) 在零频无界，因而不能是 \(L^2\) 上的有界零阶乘子；它适合放到改变可积性或正则性的分数次积分理论中。另一方面，普通的“各阶导数有界”未必包含高频处所需的衰减，不能直接代替尺度导数条件。

\subsection{二进分解}
\label{b2:subsec:220202}

先恢复\chapref{b2:ch:19}所用的分割。取实值径向函数 \(\chi\in C_c^\infty(\mathbb R^d)\)，使 \(0\leq\chi\leq1\)，在 \(|\xi|\leq1\) 时等于 \(1\)，在 \(|\xi|\geq2\) 时等于 \(0\)，并沿半径不增。令
\[
 \varphi(\xi)=\chi(\xi)-\chi(2\xi).
\]
于是 \(\varphi\geq0\)，其支撑包含于 \(\{\,1/2\leq|\xi|\leq2\,\}\)。有限和的望远镜恒等式为
\[
 \sum_{j=-M}^{N}\varphi(2^{-j}\xi)
 =\chi(2^{-N}\xi)-\chi(2^{M+1}\xi).
\]
对每个 \(\xi\ne0\)，令 \(M,N\) 趋于无穷便得到 \(\sum_{j\in\mathbb Z}\varphi(2^{-j}\xi)=1\)。每个非零频率只落在有限个相邻环带中，而且任何有限指标集 \(J\subset\mathbb Z\) 都满足
\[
 0\leq\sum_{j\in J}\varphi(2^{-j}\xi)\leq1.
\]
最后这个不等式将在取算子极限时提供共同的 \(L^2\) 控制。

\ProseParagraphBreak
把符号局部化后，令
\[
 \begin{aligned}
  m_j(\xi)&=m(\xi)\cdot\varphi(2^{-j}\xi),\\
  b_j(\zeta)&=m(2^j\zeta)\cdot\varphi(\zeta),\\
  h_j&=\mathcal F^{-1}b_j,\quad
  k_j(x)=2^{jd}\cdot h_j(2^jx).
 \end{aligned}
\]
这里第二行把所有频率环带送到同一个固定环带；Fourier 变量换元给出 \(k_j=\mathcal F^{-1}m_j\)。这样就可以先证明一个不依赖尺度的估计，再把尺度因子放回去。

\Needspace{6\baselineskip}
\begin{lemma}[固定环带上的有限导数核估计]\label{b2:lem:22-annular-kernel}
设 \(d\geq1\)，\(N\) 为非负整数，\(a\in C_c^N(\mathbb R^d;\mathbb C)\) 的支撑包含于 \(\{\,1/2\leq|\xi|\leq2\,\}\)，且
\[
 \max_{|\alpha|\leq N}\|\partial^\alpha a\|_\infty\leq M.
\]
则 \(h=\mathcal F^{-1}a\) 是光滑函数，并满足
\[
 |h(x)|+|\nabla h(x)|
 \leq C_{d,N}\cdot M(1+|x|)^{-N},
 \quad x\in\mathbb R^d.
\]
这里空间变量的一阶导数不要求符号具有 \(N+1\) 阶导数。
\end{lemma}
\begin{proof}
由于 \(a\) 支撑有界，任何多项式乘以 \(a\) 都可积。因此可以在逆 Fourier 积分内对 \(x\) 任意多次求导；这说明 \(h\) 光滑。在 \(|x|\leq1\) 时，直接估计积分给出 \(|h(x)|\leq C_dM\)。对 \(|x|>1\)，选取满足 \(|x_\ell|\geq|x|/\sqrt d\) 的坐标，分部积分 \(N\) 次：
\[
 h(x)=(-2\pi\mathrm i x_\ell)^{-N}
 \int_{\mathbb R^d}e^{2\pi\mathrm i x\cdot\xi}
 \cdot\partial_\ell^Na(\xi)\,\mathrm d\xi.
\]
边界项为零，因为 \(a\) 及所需低阶导数都紧支撑。积分区域的体积固定，于是
\[
 |h(x)|\leq C_{d,N}\cdot M|x|^{-N}.
\]
结合 \(|x|\leq1\) 的估计得到第一项的结论。

对于 \(\partial_{x_r}h\)，其频率符号是 \(2\pi\mathrm i\xi_r\cdot a(\xi)\)。在固定环带上，Leibniz 公式表明这个符号截至 \(N\) 阶的导数仍由 \(C_{d,N}M\) 控制：落在坐标 \(\xi_r\) 上的导数最多一次，落在 \(a\) 上的导数不超过 \(N\) 次。将刚才的估计用于该符号，再对 \(r=1,\ldots,d\) 求和，即得梯度界。
\end{proof}

现在对归一化符号 \(b_j\) 使用这个引理。Leibniz 公式明确给出
\[
 \partial^\alpha b_j(\zeta)
 =\sum_{\beta\leq\alpha}\binom{\alpha}{\beta}
 2^{j|\beta|}\cdot(\partial^\beta m)(2^j\zeta)
 \cdot\partial^{\alpha-\beta}\varphi(\zeta).
\]
在 \(\varphi\) 的支撑上，\(|\zeta|\) 介于 \(1/2\) 与 \(2\) 之间，因而
\[
 2^{j|\beta|}\cdot|(\partial^\beta m)(2^j\zeta)|
 \leq A|\zeta|^{-|\beta|}
 \leq 2^{|\beta|}A.
\]
所以 \(b_j\) 的 \(C^N\) 界与 \(j\) 无关。引理给出
\[
 \begin{aligned}
 |k_j(x)|&\leq C A\,2^{jd}(1+2^j|x|)^{-N},\\
 |\nabla k_j(x)|&\leq C A\,2^{j(d+1)}(1+2^j|x|)^{-N}.
 \end{aligned}
\]
本节中分割函数 \(\chi,\varphi\) 一经选定便保持不变，与它们有关的常数并入 \(C\)。

\ProseParagraphBreak
在 \(x\ne0\) 处，选整数 \(j_0\)，使 \(2^{j_0}|x|\leq1<2^{j_0+1}|x|\)。低频部分不需要衰减因子，高频部分则使用它。若 \(N>d\)，两段大小估计分别为
\[
 \begin{aligned}
 \sum_{j\leq j_0}2^{jd}
 &\leq C_d\,2^{j_0d}\leq C_d|x|^{-d},\\
 |x|^{-N}\sum_{j>j_0}2^{j(d-N)}
 &\leq C_{d,N}|x|^{-N}\cdot2^{j_0(d-N)}
 \leq C_{d,N}|x|^{-d}.
 \end{aligned}
\]
梯度估计的计算完全由同样两个级数给出，只需把其中的 \(d\) 改成 \(d+1\)，故要求 \(N>d+1\)。于是
\[
 \sum_{j\in\mathbb Z}|k_j(x)|\leq C A|x|^{-d},
 \quad
 \sum_{j\in\mathbb Z}|\nabla k_j(x)|
 \leq C A|x|^{-d-1}.
\]
在任意离开原点的紧集上，低频尾可由 \(\sum2^{jd}\) 与 \(\sum2^{j(d+1)}\) 一致控制，高频尾可由相应的衰减级数一致控制。因此核及其一阶导数的级数都局部一致收敛，
\[
 K(x)=\sum_{j\in\mathbb Z}k_j(x),\quad x\ne0,
\]
是 \(\mathbb R^d\setminus\{0\}\) 上的 \(C^1\) 函数，并满足上述大小与梯度界。这个核在离开原点处有普通函数意义，将用于给定 \(L^2\) 算子的支集外表示；其原点附近的奇性仍需由算子估计处理。

\subsection{Mihlin 乘子定理}
\label{b2:subsec:220203}

核的梯度界要求本证明选取 \(N>d+1\)。为避免在陈述中反复保留整数参数，以下采用 \(N=d+2\)。这得到\term[mihlin-chengzi-dingli]{Mihlin 乘子定理}[Mihlin multiplier theorem]的一种方便版本。

\Needspace{6\baselineskip}
\begin{theorem}[Mihlin 乘子定理的有限导数版本]\label{b2:thm:22-mihlin}
设 \(d\geq1\)，\(m\in C^{d+2}(\mathbb R^d\setminus\{0\};\mathbb C)\)，且
\[
 \sup_{\xi\ne0}|\xi|^{|\alpha|}\cdot|\partial^\alpha m(\xi)|
 \leq A<\infty,\quad |\alpha|\leq d+2.
\]
在原点任意指定 \(m\) 的值，并由
\(\widehat{T_mf}=m\cdot\widehat f\) 定义 \(L^2(\mathbb R^d;\mathbb C)\) 上的乘子。则对每个 \(1<p<\infty\)，\(T_m\) 唯一延拓为有界算子 \(L^p\to L^p\)，并且
\[
 \|T_mf\|_p\leq C_{d,p}\cdot A\|f\|_p.
\]
不同指数的延拓在相应空间的交集上一致。它还具有弱型 \((1,1)\) 延拓，弱型常数不超过 \(C_dA\)。
\end{theorem}
\begin{proof}
若 \(A=0\)，则 \(m=0\) 几乎处处，结论直接成立。以下设 \(A>0\)。零阶条件和\cref{b2:prop:22-l2}给出 \(\|T_m\|_{2\to2}\leq A\)。按前一小节构造 \(k_j\) 与 \(K\)，得到
\[
 |K(z)|\leq C A|z|^{-d},
 \quad |\nabla K(z)|\leq C A|z|^{-d-1},
 \quad z\ne0.
\]
还须证明这个 \(K\) 确实表示已经由 Fourier 变换定义的算子。令
\[
 m^{(L)}(\xi)=m(\xi)\cdot
 \sum_{|j|\leq L}\varphi(2^{-j}\xi),
 \quad K_L=\sum_{|j|\leq L}k_j.
\]
每个 \(k_j\) 都属于 \(L^1\)，故 \(T_{m^{(L)}}f=K_L*f\) 在 \(L^2\) 中成立。有限分割和介于 \(0\) 与 \(1\) 之间，且在非零频率趋于 \(1\)，于是
\[
 \|T_{m^{(L)}}f-T_mf\|_2^2
 =\int_{\mathbb R^d}|m^{(L)}-m|^2\cdot|\widehat f|^2
 \,\mathrm d\xi\longrightarrow0.
\]
这里共同控制函数是 \(4A^2\cdot|\widehat f|^2\)，可由支配收敛处理。

取紧支撑的 \(f\in L^2\)，并取紧集 \(E\subset\mathbb R^d\setminus\operatorname{supp}f\)。此时 \(f\in L^1\)，而所有差 \(x-y\)，其中 \(x\in E\)、\(y\in\operatorname{supp}f\)，组成离开原点的紧集 \(D\)。因此
\[
 \sup_{x\in E}\left|
 (K_L*f)(x)-\int K(x-y)\cdot f(y)\,\mathrm dy
 \right|
 \leq\|f\|_1\cdot\sup_{z\in D}|K_L(z)-K(z)|
 \longrightarrow0.
\]
另一方面，\(L^2\) 收敛有一个几乎处处收敛的子列。比较两个极限，得到 \(T_mf(x)=\int K(x-y)\cdot f(y)\,\mathrm dy\) 对 \(E\) 中几乎处处的 \(x\) 成立。以可数个这样的紧集穷尽支撑补集，便确认支撑外表示；该积分在支撑外绝对收敛。

把两变量核记为 \(\mathcal K(x,y)=K(x-y)\)。当 \(|x-x'|\leq|x-y|/2\) 时，连接 \(x-y\) 与 \(x'-y\) 的线段离原点至少 \(|x-y|/2\)，故微积分基本定理给出
\[
 |\mathcal K(x,y)-\mathcal K(x',y)|
 \leq C A\,\frac{|x-x'|}{|x-y|^{d+1}}.
\]
改变 \(y\) 得到另一变量的同样估计。于是 \(\mathcal K\) 满足标准 Calderón--Zygmund 核的大小与双变量光滑性条件，\(T_m\) 又已有 \(L^2\) 界。将\cref{b2:thm:18-cz-weak}和\cref{b2:thm:18-cz-lp}应用于 \(A^{-1}T_m\)，其核常数与 \(L^2\) 常数只依赖维数，得到所述弱型界与强型界。后者先在 \(L^2\cap L^p\) 上成立，再由该子空间在 \(L^p\) 中的稠密性唯一延拓；不同 \(p\) 的延拓用同一截断逼近识别。
\end{proof}

支撑外核的识别不能省略。例如常数符号 \(m=1\) 给出恒等算子，其分布核集中在原点，支撑外核却是零。一般地，支撑外核无法独自决定一个乘法项。本证明从 \(L^2\) 乘子出发，再核对它的支撑外表示，所以这种集中在对角线上的部分仍由 \(L^2\) 算子保留。

\ProseParagraphBreak
定理通常不能补上强型 \((1,1)\) 或 \(L^\infty\to L^\infty\) 两个端点。\chapref{b2:ch:18}的 Hilbert 变换给出了同时满足尺度导数条件而端点失败的例子。另一方面，若符号是支撑于一个固定环带的 \(C_c^{d+1}\) 函数，\cref{b2:lem:22-annular-kernel}已经使其核属于 \(L^1\)，Young 不等式便给出全部 \(1\leq p\leq\infty\) 的界。无限尺度之所以不同，是因为每个环带核的 \(L^1\) 范数虽一致有界，却未必对环带指标可求和。

\subsection{Hörmander 条件}
\label{b2:subsec:220204}

上面的论证把符号导数转成核的逐点衰减。若只知道导数的平方积分，仍能通过 Plancherel 得到核的带权平方积分；只要权足够强，就能控制核尾部的 \(L^1\) 积分。这条路线允许用非整数阶的光滑性，并避免要求极限核存在逐点梯度界。

\ProseParagraphBreak
本小节的 \(H^s\) 范数作用在符号的频率变量上。为免混淆，把该变量写成 \(\zeta\)，其对偶变量写成 \(x\)。对 \(s\geq0\)，采用
\[
 \|q\|_{H^s_\zeta}^2
 =\int_{\mathbb R^d}(1+4\pi^2|x|^2)^s
 \cdot|\mathcal F_\zeta q(x)|^2\,\mathrm dx.
\]
这是 \(p=2\) 的 Sobolev 范数。权关于 \(x\) 对称，因此用逆 Fourier 变换替代 \(\mathcal F_\zeta q\) 得到同一范数。

\Needspace{6\baselineskip}
\begin{lemma}[Sobolev 范数中的光滑截断]\label{b2:lem:22-sobolev-cutoff}
设 \(d\geq1\)，\(s\geq0\)，\(\psi\in C_c^\infty(\mathbb R^d;\mathbb C)\)。则乘法 \(q\mapsto\psi\cdot q\) 是 \(H^s(\mathbb R^d;\mathbb C)\) 上的有界算子，且
\[
 \|\psi\cdot q\|_{H^s}\leq C_{d,s,\psi}\cdot\|q\|_{H^s}.
\]
\end{lemma}
\begin{proof}
令 \(w_s(x)=(1+|x|)^s\)，它与 \((1+4\pi^2|x|^2)^{s/2}\) 等价。三角不等式蕴含
\[
 w_s(x)\leq w_s(x-y)\cdot w_s(y).
\]
先取 \(q\) 为 Schwartz 函数。乘积的 Fourier 变换是卷积，所以
\[
 w_s(x)\cdot|\widehat{\psi\cdot q}(x)|
 \leq
 \bigl((w_s\cdot|\widehat\psi|)*(w_s\cdot|\widehat q|)\bigr)(x).
\]
由于 \(\widehat\psi\) 是 Schwartz 函数，\(\|w_s\cdot\widehat\psi\|_1<\infty\)。Young 不等式给出
\[
 \|w_s\cdot\widehat{\psi\cdot q}\|_2
 \leq\|w_s\cdot\widehat\psi\|_1
 \cdot\|w_s\cdot\widehat q\|_2,
\]
即所需估计。

一般 \(q\in H^s\) 可在 \(H^s\) 中用 Schwartz 函数 \(q_n\) 逼近：在 Fourier 一侧先截断到有界区域，再以紧支撑光滑函数逼近，权在每个有界区域上与常数等价，逆变换即得所需 \(q_n\)。上式使 \(\psi\cdot q_n\) 在 \(H^s\) 中为 Cauchy 列；同时 \(q_n\to q\) 于 \(L^2\)，而 \(\psi\) 有界，所以 \(\psi\cdot q_n\to\psi\cdot q\) 于 \(L^2\)。两个极限一致，从而估计延续到 \(q\)。
\end{proof}

还需要把核的积分差分条件接到弱型估计上。为便于以后复用，下面单独给出卷积核版本。这个结果只要求支撑外表示和核差分积分有界，不要求核本身满足逐点大小条件。

\Needspace{6\baselineskip}
\begin{lemma}[积分 Hörmander 条件与弱型估计]\label{b2:lem:22-integral-hormander-weak}
设 \(d\geq1\)，\(T\) 是复线性有界算子 \(L^2(\mathbb R^d)\to L^2(\mathbb R^d)\)，\(\|T\|_{2\to2}\leq M\)。设 Borel 函数 \(K\in L^1_{\mathrm{loc}}(\mathbb R^d\setminus\{0\})\) 具有下列支撑外表示：对每个紧支撑的 \(f\in L^2\)，
\[
 Tf(x)=\int_{\mathbb R^d}K(x-y)\cdot f(y)\,\mathrm dy
\]
在 \(\operatorname{supp}f\) 外几乎处处成立，且右侧在这些点绝对收敛。若
\[
 H\coloneqq\sup_{v\ne0}
 \int_{|x|>2|v|}|K(x-v)-K(x)|\,\mathrm dx<\infty,
\]
则对 \(f\in L^1\cap L^2\) 和 \(t>0\)，
\[
 \bigl|\{\,x:|Tf(x)|>t\,\}\bigr|
 \leq \frac{C_d(M+H)}{t}\cdot\|f\|_1.
\]
因此 \(T\) 唯一延拓为与原 \(L^2\) 算子相容的弱型 \((1,1)\) 算子。
\end{lemma}
\begin{proof}
记 \(B=M+H\)。若 \(B=0\)，则 \(T=0\) 于 \(L^2\)，结论成立。设 \(B>0\)，并对 \(f\in L^1\cap L^2\) 在高度 \(\alpha=t/B\) 使用\cref{b2:thm:18-cz-decomposition}。该结果适用于复值函数，给出
\[
 f=g+\sum_Qb_Q,\quad
 \|g\|_2^2\leq 2^d\alpha\cdot\|f\|_1,\quad
 \sum_Q|Q|\leq\alpha^{-1}\cdot\|f\|_1,
\]
其中所选二进立方体两两不交，\(b_Q\) 支撑于 \(Q\)，\(\int b_Q=0\)，并且
\[
 \sum_Q\|b_Q\|_1\leq2\|f\|_1.
\]
每个坏块属于 \(L^2\)，坏块级数在 \(L^2\) 中收敛。这个收敛性质随后使有限坏块估计可以传到算子作用下的极限。

把每个 \(Q\) 关于其中心 \(c_Q\) 放大 \(8\sqrt d\) 倍，所得立方体记为 \(Q^*\)，并令
\[
 \Omega^*=\bigcup_QQ^*,\quad G=\mathbb R^d\setminus\Omega^*.
\]
由体积估计，
\[
 |\Omega^*|\leq C_d\cdot\frac{B}{t}\cdot\|f\|_1.
\]
好部用 \(L^2\) 有界性和 Chebyshev 不等式处理：
\[
 \bigl|\{\,|Tg|>t/2\,\}\bigr|
 \leq\frac{4M^2}{t^2}\cdot\|g\|_2^2
 \leq C_d\cdot\frac{M^2}{Bt}\cdot\|f\|_1
 \leq C_d\cdot\frac{B}{t}\cdot\|f\|_1.
\]

坏部所需的抵消正是积分条件发挥作用的地方。对 \(x\notin Q^*\)，支撑外表示和 \(\int b_Q=0\) 给出
\[
 Tb_Q(x)=\int_Q
 [K(x-y)-K(x-c_Q)]\cdot b_Q(y)\,\mathrm dy
\]
几乎处处成立。如果 \(y\in Q\) 且 \(x\notin Q^*\)，则 \(|x-c_Q|>2|y-c_Q|\)。令 \(z=x-c_Q\)、\(v=y-c_Q\)，Tonelli 定理和积分假设因此给出
\[
 \begin{aligned}
 \int_{\mathbb R^d\setminus Q^*}|Tb_Q(x)|\,\mathrm dx
 &\leq\int_Q|b_Q(y)|\cdot
   \left(\int_{\mathbb R^d\setminus Q^*}
    |K(x-y)-K(x-c_Q)|\,\mathrm dx\right)\,\mathrm dy\\
 &\leq H\|b_Q\|_1.
 \end{aligned}
\]
当 \(y=c_Q\) 时差分为零，其余点由 \(|z|>2|v|\) 的积分控制；不需为该单点另设条件。

枚举坏块并令 \(b^{(n)}\) 为前 \(n\) 项之和。在共同区域 \(G\) 上，有限和满足
\[
 \int_G|Tb^{(n)}(x)|\,\mathrm dx
 \leq H\sum_Q\|b_Q\|_1\leq2H\|f\|_1.
\]
由 \(b^{(n)}\to b=f-g\) 于 \(L^2\)，有 \(Tb^{(n)}\to Tb\) 于 \(L^2\)。选取几乎处处收敛的子列，再用 Fatou 引理，得到 \(\int_G|Tb|\leq2H\|f\|_1\)。于是
\[
 \bigl|\{\,x\in G:|Tb(x)|>t/2\,\}\bigr|
 \leq\frac{4H}{t}\cdot\|f\|_1.
\]
把这一界、好部界和 \(|\Omega^*|\) 相加，得到 \(L^1\cap L^2\) 上的弱型估计。

最后说明延拓的含义。对 \(f\in L^1\)，选有界紧支撑 \(f_n\)，使 \(\|f_n-f\|_1\leq2^{-3n}\)。对差 \(f_{n+1}-f_n\) 使用刚证的估计，可得
\[
 \sum_{n=1}^{\infty}
 \bigl|\{\,|Tf_{n+1}-Tf_n|>2^{-n}\,\}\bigr|<\infty.
\]
几乎每点最终都避开这些集合，故 \(Tf_n\) 几乎处处收敛。以其极限定义 \(Tf\)，再对 \(|Tf_n-Tf_m|\) 的分布估计令 \(m\to\infty\)，得到
\[
 \bigl|\{\,|Tf_n-Tf|>t\,\}\bigr|
 \leq\frac{C_dB}{t}\cdot\|f_n-f\|_1.
\]
同样的极限过程把弱型界传给 \(f\)，并说明不同逼近所得极限相同。用两组逼近的和与数乘可验证线性。若 \(f\in L^1\cap L^2\)，可选同时在 \(L^1,L^2\) 中逼近的截断，所得极限又与原算子的 \(L^2\) 极限相同。这给出相容性与弱型延拓的唯一性。
\end{proof}

固定 \(\eta\in C_c^\infty(\mathbb R^d\setminus\{0\})\)，使它在闭环带 \(\{\,1/2\leq|\zeta|\leq2\,\}\) 的一个邻域等于 \(1\)。对可测有界符号 \(m\)，令
\[
 B_s(m)\coloneqq\sup_{j\in\mathbb Z}
 \|m(2^j\,\cdot)\cdot\eta\|_{H^s}.
\]
固定截断 \(\eta\) 留出了一点环带边界余量，使 \(\varphi\) 和坐标乘以 \(\varphi\) 都能作为光滑乘法作用于被测量的对象。以下是\term[hormander-chengzi-dingli]{Hörmander 乘子定理}[Hörmander multiplier theorem]的环带版本。

\Needspace{6\baselineskip}
\begin{theorem}[Hörmander 环带条件]\label{b2:thm:22-hormander}
设 \(d\geq1\)，\(s>d/2\)，\(m\in L^\infty(\mathbb R^d;\mathbb C)\)。固定 \(\eta\in C_c^\infty(\mathbb R^d\setminus\{0\})\)，使它在 \(\{\,1/2\leq|\xi|\leq2\,\}\) 的邻域等于 \(1\)，并假设
\[
 B_s(m)=\sup_{j\in\mathbb Z}
 \|m(2^j\,\cdot)\cdot\eta\|_{H^s(\mathbb R^d)}<\infty.
\]
则 \(L^2\) 乘子 \(T_m\) 对每个 \(1<p<\infty\) 唯一延拓为有界算子 \(L^p\to L^p\)，且
\[
 \|T_mf\|_p\leq C_{d,p,s,\eta}
 \bigl(\|m\|_\infty+B_s(m)\bigr)\cdot\|f\|_p.
\]
它还具有与 \(L^2\) 定义相容的弱型 \((1,1)\) 延拓，弱型常数不超过
\(C_{d,s,\eta}\bigl(\|m\|_\infty+B_s(m)\bigr)\)。
\end{theorem}
\begin{proof}
仍用本节固定的 \(\varphi\)，并记
\[
 q_j(\zeta)=m(2^j\zeta)\cdot\eta(\zeta),\quad
 b_j=\varphi\cdot q_j,\quad h_j=\mathcal F^{-1}b_j,\quad
 k_j(x)=2^{jd}\cdot h_j(2^jx).
\]
因为 \(\eta=1\) 于 \(\operatorname{supp}\varphi\) 的邻域，\(b_j=m(2^j\,\cdot)\cdot\varphi\)，与二进分解中的符号相同。\cref{b2:lem:22-sobolev-cutoff}分别用于截断 \(\varphi\) 和 \(2\pi\mathrm i\zeta_r\cdot\varphi(\zeta)\)，给出
\[
 \|b_j\|_{H^s}+\sum_{r=1}^{d}
 \|2\pi\mathrm i\zeta_r\cdot b_j\|_{H^s}
 \leq C_{d,s,\eta}\cdot B_s(m).
\]
Plancherel 与 Sobolev 范数的定义因此说明
\[
 \|(1+|x|)^s\cdot h_j\|_2
 +\|(1+|x|)^s\cdot\nabla h_j\|_2
 \leq C\cdot B_s(m).
\]
这些 \(h_j\) 是光滑函数：\(b_j\) 有界且紧支撑，任何多项式乘以 \(b_j\) 都属于 \(L^1\)，可以在逆变换中对空间变量求导。光滑性来自有限频带上的可积性与空间变量求导。

令 \(\gamma=s-d/2>0\)。由 Cauchy--Schwarz 不等式，
\[
 \|h_j\|_1+\|\nabla h_j\|_1
 \leq C\cdot B_s(m)\cdot
 \left(\int_{\mathbb R^d}(1+|x|)^{-2s}\,\mathrm dx\right)^{1/2}
 \leq C_{d,s,\eta}\cdot B_s(m).
\]
更精确地，对 \(R\geq1\)，极坐标给出
\[
 \int_{|x|>R}(1+|x|)^{-2s}\,\mathrm dx
 \leq C_d\int_R^\infty r^{d-1-2s}\,\mathrm dr
 =\frac{C_d}{2s-d}\cdot R^{d-2s},
\]
所以
\[
 \int_{|x|>R}|h_j(x)|\,\mathrm dx
 \leq C_{d,s,\eta}\cdot B_s(m)R^{-\gamma}.
\]
这一步把带权 \(L^2\) 信息变成了两个不同用途的 \(L^1\) 信息：梯度的全空间积分控制小位移，核的尾积分控制大位移。

取有限指标集 \(J\subset\mathbb Z\)，令
\[
 K_J=\sum_{j\in J}k_j,\quad
 m_J(\xi)=m(\xi)\cdot\sum_{j\in J}\varphi(2^{-j}\xi),
 \quad T_Jf=K_J*f.
\]
每个 \(k_j\in L^1\)，故 \(K_J\in L^1\)，卷积与 \(L^2\) 乘子 \(T_{m_J}\) 一致。分割的非负性保证 \(|m_J|\leq\|m\|_\infty\)，因而
\[
 \|T_J\|_{2\to2}\leq\|m\|_\infty
\]
与有限集 \(J\) 无关。下一步证明积分 Hörmander 常数也与 \(J\) 无关。

固定 \(v\ne0\)。若 \(2^j|v|\leq1\)，由
\[
 k_j(x-v)-k_j(x)
 =-\int_0^1v\cdot\nabla k_j(x-tv)\,\mathrm dt
\]
及平移不变性可得
\[
 \begin{aligned}
 \int_{|x|>2|v|}|k_j(x-v)-k_j(x)|\,\mathrm dx
 &\leq|v|\cdot\|\nabla k_j\|_1\\
 &=2^j|v|\cdot\|\nabla h_j\|_1
 \leq C\cdot B_s(m)\cdot2^j|v|.
 \end{aligned}
\]
若 \(2^j|v|>1\)，则在 \(|x|>2|v|\) 上同时有 \(|x|>|v|\) 与 \(|x-v|>|v|\)。分别在两个积分中换元，使用上面的尾积分界，得到
\[
 \begin{aligned}
 &\int_{|x|>2|v|}|k_j(x-v)-k_j(x)|\,\mathrm dx\\
 &\quad\leq
 \int_{|z|>2^j|v|}|h_j(z)|\,\mathrm dz+
 \int_{|z|>2^j|v|}|h_j(z)|\,\mathrm dz
 \leq C\cdot B_s(m)\cdot(2^j|v|)^{-\gamma}.
 \end{aligned}
\]
选择 \(j_0\) 满足 \(2^{j_0}|v|\leq1<2^{j_0+1}|v|\)。两部分级数满足
\[
 \sum_{j\leq j_0}2^j|v|\leq2,\quad
 \sum_{j>j_0}(2^j|v|)^{-\gamma}
 \leq\frac{1}{1-2^{-\gamma}}.
\]
对任意有限 \(J\) 取子和不会增大这些上界，所以
\[
 \sup_{v\ne0}\int_{|x|>2|v|}
 |K_J(x-v)-K_J(x)|\,\mathrm dx
 \leq C_{d,s,\eta}B_s(m).
\]

有限和核是光滑的 \(L^1\) 函数，满足所需支撑外表示。于是\cref{b2:lem:22-integral-hormander-weak}给出所有 \(T_J\) 的一致弱型界。其 \(L^2\) 伴随的核为
\[
 K_J^*(x)=\overline{K_J(-x)}.
\]
变量替换 \(x\mapsto-x\)、\(v\mapsto-v\) 表明 \(K_J^*\) 具有同一个积分条件，伴随的 \(L^2\) 范数也相同，所以 \(T_J^*\) 具有同样的弱型界。记 \(D=\|m\|_\infty+B_s(m)\)；若 \(D=0\) 则结论平凡，否则先对 \(D^{-1}T_J\) 使用\cref{b2:thm:17-marcinkiewicz}，由弱型 \((1,1)\) 和强型 \((2,2)\) 得到 \(1<p<2\) 的界。对 \(D^{-1}T_J^*\) 作相同插值，再由 \(p'\in(1,2)\) 的对偶性得到 \(2<p<\infty\) 的界。因此
\[
 \|T_Jf\|_p\leq C_{d,p,s,\eta}\cdot D\|f\|_p,
 \quad1<p<\infty,
\]
常数与 \(J\) 的大小和位置无关。

最后令 \(J=\{\,-L,\ldots,L\,\}\)。分割恒等式和 \(|m_J|\leq\|m\|_\infty\) 使 \(T_Jf\to T_mf\) 于 \(L^2\)。对 \(f\in L^p\cap L^2\)，选取几乎处处收敛的子列，再用 Fatou 引理，将统一 \(L^p\) 界传到 \(T_mf\)。在 \(L^1\cap L^2\) 上也可把弱型分布估计传到极限：先把阈值 \(t\) 降为任意 \(t'<t\)，使用几乎处处收敛和 Fatou，再令 \(t'\uparrow t\)。强型延拓由 \(L^p\cap L^2\) 的稠密性得到，弱型延拓由前一引理证明末尾的逼近步骤得到；所有这些延拓与原 \(L^2\) 算子相容。
\end{proof}

严格不等式 \(s>d/2\) 在两处承担同一任务：使逆权 \((1+|x|)^{-s}\) 属于 \(L^2\)，并使尾积分指数 \(\gamma=s-d/2\) 为正。在 \(s=d/2\) 时，当前尾积分计算出现对数发散，高频环带级数也不再衰减。这指出了本证明不能到达端点的原因；证明失效本身并不构成某个附加条件下端点结论的反例。

\ProseParagraphBreak
有限环带和的估计通过 \(L^2\) 极限传到原算子。与 Mihlin 证明相比，保留下来的是足以估计坏块的核差分积分，而不是逐点梯度。这正是 Sobolev 环带条件可以降低光滑性要求的原因。

\ProseParagraphBreak
还可用平移乘子看清“充分条件”的含义。若 \(a\in\mathbb R^d\)，则 \(m(\xi)=e^{2\pi\mathrm i a\cdot\xi}\) 给出 \(T_mf(x)=f(x+a)\)，它在每个 \(L^p\) 上等距。但当 \(a\ne0\) 时，至少一个一阶导数的模恒为非零常数，乘上 \(|\xi|\) 后无界，所以不满足 Mihlin 条件。Mihlin 条件的作用在于直接检验符号，并使所得控制在尺度分解下保持一致。

\ProseParagraphBreak
最后给出下一节所需的一族乘子版本。标量算子范数一致有界并不能单独推出先在位置上求平方和的估计；此处统一的 Mihlin 条件还提供统一核结构，因而可以调用\cref{b2:prop:18-hilbert-valued-cz}。

\Needspace{6\baselineskip}
\begin{proposition}[统一 Mihlin 条件下的平方和估计]\label{b2:prop:22-uniform-mihlin-square}
设 \(d\geq1\)，\(1<p<\infty\)，\(I\) 是有限或可数指标集。对每个 \(\nu\in I\)，设 \(m_\nu\in C^{d+2}(\mathbb R^d\setminus\{0\};\mathbb C)\)，并假设
\[
 \sup_{\nu\in I}\sup_{\xi\ne0}
 |\xi|^{|\alpha|}\cdot|\partial^\alpha m_\nu(\xi)|
 \leq A<\infty,\quad|\alpha|\leq d+2.
\]
以各自 \(L^2\) 定义的相容 \(L^p\) 延拓记 \(T_{m_\nu}\)。则对任意复值函数族 \((f_\nu)_{\nu\in I}\)，只要右侧有限，就有
\[
 \left\|\left(\sum_{\nu\in I}|T_{m_\nu}f_\nu|^2\right)^{1/2}\right\|_p
 \leq C_{d,p}\cdot A
 \left\|\left(\sum_{\nu\in I}|f_\nu|^2\right)^{1/2}\right\|_p.
\]
常数与指标集的大小无关。
\end{proposition}
\begin{proof}
先设 \(I=\{1,\ldots,n\}\)，并令 \(H=\mathbb C^n\)，赋予 Euclidean 范数。在 \(L^2(\mathbb R^d;H)\) 上定义对角算子
\[
 \mathbb T(f_1,\ldots,f_n)
 =(T_{m_1}f_1,\ldots,T_{m_n}f_n).
\]
逐坐标 Plancherel 给出
\[
 \|\mathbb Tf\|_{L^2(H)}^2
 =\sum_{\nu=1}^{n}\int|m_\nu(\xi)|^2
 \cdot|\widehat f_\nu(\xi)|^2\,\mathrm d\xi
 \leq A^2\cdot\|f\|_{L^2(H)}^2.
\]
对每个符号，用本节二进分解构造支撑外核 \(K_\nu\)。全部符号共享常数 \(A\)，因此核的大小、梯度和支撑外表示也共享同一组常数。令
\[
 \mathbb K(z)=\operatorname{diag}
 \bigl(K_1(z),\ldots,K_n(z)\bigr),\quad z\ne0.
\]
对角矩阵在 \(H\) 上的算子范数等于对角元模的最大值，故
\[
 \|\mathbb K(z)\|_{H\to H}
 =\max_{1\leq\nu\leq n}|K_\nu(z)|
 \leq C A|z|^{-d}.
\]
当 \(|h|\leq|z|/2\) 时，各坐标的梯度界还给出
\[
 \|\mathbb K(z+h)-\mathbb K(z)\|_{H\to H}
 \leq C A\,\frac{|h|}{|z|^{d+1}}.
\]
对两变量核 \(\mathbb K(x-y)\)，改变任一变量都满足相应差分估计。支撑外表示由有限个坐标的表示逐项组合；此处有限维 Bochner 积分就是坐标积分。

这样，\(\mathbb T\) 满足\cref{b2:prop:18-hilbert-valued-cz}的全部条件：有 \(L^2(H)\) 起点、算子范数意义的标准核估计，以及紧支撑输入的支撑外表示。该命题的常数不依赖 Hilbert 空间的维数。应用它并对 \(A\) 归一化，即得
\[
 \|\mathbb Tf\|_{L^p(H)}
 \leq C_{d,p}\cdot A\|f\|_{L^p(H)},
\]
这正是有限族的平方和估计。先对 \(L^2(H)\cap L^p(H)\) 输入应用，再由稠密性推广，得到陈述中的相容 \(L^p\) 算子。

若 \(I\) 可数，枚举其元素并对前 \(n\) 项使用已证估计。右侧的有限平方和由完整平方和控制；左侧随着 \(n\) 增大逐点单调增加。对其 \(p\) 次方使用单调收敛定理，得到可数族结论。
\end{proof}

这个向量估计同时控制一组频带上的乘子。将它与\chapref{b2:ch:19}的平方函数定理结合，就能把逐频带的估计重新合成为原函数的范数；下一节将据此识别 Bessel 势范数。

% !TeX root = ../../Book2.tex
\section{应用}
\label{b2:sec:2203}

应用乘子定理时，先判断符号阶数，再选取匹配的函数空间。零阶符号可保持 \(L^p\)，负阶符号在不同频率区间可能改善正则性或改变可积指数，正阶符号则需输入具有相应导数。

\subsection{Riesz 变换}
\label{b2:subsec:220301}

Riesz 变换把一个方向上的微分除以总微分尺度。沿用\chapref{b2:ch:18}的归一化，令
\[
 r_j(\xi)\coloneqq-\mathrm i\frac{\xi_j}{|\xi|},
 \quad R_j\coloneqq T_{r_j},
 \quad 1\leq j\leq d,\quad \xi\ne0.
\]
先在 \(L^2(\mathbb R^d;\mathbb C)\) 中定义这个算子，零频处的取值不影响它。符号 \(r_j\) 为零次齐次函数，每次微分使齐次次数降低一。因此
\[
 (\partial^\alpha r_j)(t\omega)
 =t^{-|\alpha|}(\partial^\alpha r_j)(\omega),
 \quad t>0,\quad |\omega|=1.
\]
单位球面紧，右侧的角向因子有界，故满足\cref{b2:thm:22-mihlin}要求的 \(d+2\) 阶条件。这给出全部 \(1<p<\infty\) 上的有界延拓，并与\chapref{b2:ch:18}由奇异核得到的算子一致：两者在共同的 \(L^2\cap L^p\) 输入上相同，再由稠密性得到一致性。奇异积分证明直接显示核的零平均与尾部，乘子证明则便于把 Riesz 变换同微分算子组合。

\Needspace{6\baselineskip}
\begin{proposition}[Riesz 变换与导数]\label{b2:prop:22-riesz-gradient}
设整数 \(d\geq1\)、\(1<p<\infty\)，\(R_j\) 是符号 \(-\mathrm i\xi_j/|\xi|\) 的 \(L^p(\mathbb R^d;\mathbb C)\) 乘子。则
\[
 -\sum_{j=1}^dR_j^2f=f\quad(f\in L^p).
\]
对 \(f\in\mathcal S(\mathbb R^d)\)，令 \((-\Delta)^{1/2}f\) 为符号 \(2\pi|\xi|\) 的 Fourier 乘子，并把 \(\|\nabla f\|_p\) 理解为欧氏向量长度的 \(L^p\) 范数。有
\[
 \partial_jf=-R_j(-\Delta)^{1/2}f,\quad
 (-\Delta)^{1/2}f=\sum_{j=1}^dR_j\partial_jf,\quad
 \|(-\Delta)^{1/2}f\|_p\asymp\|\nabla f\|_p.
\]
此外，对每个 \(1\leq j,k\leq d\)，
\[
 \partial_j\partial_k f=R_jR_k(-\Delta f),
 \quad
 \sum_{j,k=1}^d\|\partial_j\partial_k f\|_p
 \leq C_{d,p}\|\Delta f\|_p.
\]
\end{proposition}
\begin{proof}
符号恒等式
\[
 -\sum_j\left(-\mathrm i\frac{\xi_j}{|\xi|}\right)^2=1
 \quad(\xi\ne0)
\]
先由 Plancherel 给出 \(L^2\) 上的算子恒等式。有限次复合在 \(L^p\) 有界，因此先限制到 \(L^2\cap L^p\)，再逼近一般 \(L^p\) 输入，得到第一式。

对 Schwartz 输入，\(\sum_jR_j\partial_jf\) 属于 \(L^p\cap L^2\)，其 Fourier 符号为
\[
 \sum_j\left(-\mathrm i\frac{\xi_j}{|\xi|}\right)
                (2\pi\mathrm i\xi_j)=2\pi|\xi|.
\]
这同时证明 \((-\Delta)^{1/2}f\in L^p\)，并给出
\(\|(-\Delta)^{1/2}f\|_p\leq C\sum_j\|\partial_jf\|_p\leq C'\|\nabla f\|_p\)。
再将 \(r_j\) 与 \(2\pi|\xi|\) 相乘，得到
\(R_j(-\Delta)^{1/2}f=-\partial_jf\)，所以反向估计也由 Riesz 有界性得到。这里两处负号由本书的负号 Fourier 变换约定决定。

最后，
\[
 \left(-\mathrm i\frac{\xi_j}{|\xi|}\right)
 \left(-\mathrm i\frac{\xi_k}{|\xi|}\right)
       4\pi^2|\xi|^2=-4\pi^2\xi_j\xi_k.
\]
因此二阶导数恒等式成立，对有限个复合算子应用 \(L^p\) 界即可。
\end{proof}
在 \(p=2\) 时，Plancherel 进一步给出准确等式
\(\|(-\Delta)^{1/2}f\|_2^2=\sum_j\|\partial_jf\|_2^2\)；一般 \(p\) 上则是范数可比。这个结果允许从 \(\Delta f\) 恢复每个二阶导数，而不仅是它们的和。对于 \(d=1\)，\(R_1\) 就是 Hilbert 变换，恒等式成为 \(H^2=-I\) 与 \((-\Delta)^{1/2}f=Hf'\)。区间示性函数的 Hilbert 变换有边界对数奇性及不可积的远场尾部，故这里的强型范围不能顺手补上 \(p=1,\infty\)。

\ProseParagraphBreak
还可处理任意光滑零次齐次符号 \(m\)，因为单位球面上的有限阶导数有界。零次齐次性并不要求球面平均为零；常数符号的核有一个原点质量。只有在提取原点外的奇异核时，才需要辨认其球面抵消，不能把抵消条件错误地加到所有零阶符号上。

\subsection{分数幂}
\label{b2:subsec:220302}

分数阶导数要同时指定高频的阶数与零频的处理。先对 Schwartz 函数记
\[
 J^sf\coloneqq\mathcal F^{-1}
       \bigl[(1+4\pi^2|\xi|^2)^{s/2}\cdot\widehat f\bigr]
 \quad(s\in\mathbb R),
\]
\[
 (-\Delta)^{s/2}f\coloneqq
 \mathcal F^{-1}\bigl[(2\pi|\xi|)^s\cdot\widehat f\bigr]
 \quad(s>0).
\]
前者是不齐次分数阶算子，后者是齐次分数阶算子。前者的符号在全空间光滑，各阶导数至多多项式增长，因而乘法保持 \(\mathcal S\)，并由对偶作用定义在全部 \(\mathcal S'\) 上；\(J^sJ^t=J^{s+t}\)。后者在零频不一定光滑，不能未经说明就与任意温和分布相乘。当前定义的 Schwartz 输入使频率乘积局部可积，已经足以进行下面的计算。

\ProseParagraphBreak

当 \(s>0\)，\(J^{-s}\) 的符号满足 Mihlin 条件，故在 \(1<p<\infty\) 上保持 \(L^p\)。齐次逆 \((-\Delta)^{-s/2}\) 在零频发散，应在适当齐次空间中定义，或在 \(0<s<d\) 时以\chapref{b2:ch:19}的 Riesz 势解释。后者在 \(1<p<d/s\) 时把 \(L^p\) 送入 \(L^q\)、\(1/q=1/p-s/d\)，改变了指数；不能据此声称它在同一个 \(L^p\) 上有界。

\Needspace{6\baselineskip}
\begin{definition}[Bessel 位势空间]\label{b2:def:22-bessel-space}
设整数 \(d\geq1\)、\(s\in\mathbb R\)、\(1<p<\infty\)。令 \(J^s\) 为符号
\((1+4\pi^2|\xi|^2)^{s/2}\) 的温和分布乘子。满足 \(J^sf\) 可由一个 \(L^p(\mathbb R^d;\mathbb C)\) 函数表示的温和分布组成\term[bessel-weishi-kongjian]{Bessel 位势空间}[Bessel potential space]，记作
\[
 H^{s,p}(\mathbb R^d)\coloneqq
 \{\,f\in\mathcal S'(\mathbb R^d):J^sf\in L^p(\mathbb R^d)\,\},
 \quad
 \|f\|_{H^{s,p}}\coloneqq\|J^sf\|_p.
\]
\end{definition}
这里把函数与它所定义的温和分布视为同一对象。这个空间完备，而且 \(J^s:H^{s,p}\to L^p\) 是等距双射，逆为 \(J^{-s}\)。事实上，若 \(g_n\to g\) 于 \(L^p\)，则对任意测试函数 \(\psi\in\mathcal S\)，
\[
 |\langle J^{-s}(g_n-g),\psi\rangle|
 \leq\|g_n-g\|_p\cdot\|J^{-s}\psi\|_{p'}
 \longrightarrow0.
\]
测试函数右侧仍为 Schwartz 函数。于是 \(J^{-s}g_n\) 在 \(\mathcal S'\) 中趋于 \(J^{-s}g\)，同时在所定义范数中收敛。取 \(g_n\in\mathcal S\) 逼近 \(g\)，还说明 Schwartz 类在 \(H^{s,p}\) 中稠密。相同的符号恒等式给出
\[
 J^t:H^{s,p}\longrightarrow H^{s-t,p},
 \quad \|J^tf\|_{H^{s-t,p}}=\|f\|_{H^{s,p}}.
\]
负阶空间允许分布输入；并不是每个元素都预先具有普通函数的代表。

\Needspace{6\baselineskip}
\begin{proposition}[Bessel 范数的平方函数描述]\label{b2:prop:22-bessel-square}
设整数 \(d\geq1\)、\(s\in\mathbb R\)、\(1<p<\infty\)。取实径向
\(\chi\in C_c^\infty(\mathbb R^d)\)，满足 \(0\leq\chi\leq1\)，在
\(|\xi|\leq1\) 为一、\(|\xi|\geq2\) 为零，并沿半径单调不增。令
\[
 \varphi=\chi-\chi(2\cdot),\quad
 P_0=\chi(D),\quad P_j=\varphi(2^{-j}D)\quad(j\geq1).
\]
这里 \(m(D)f=\mathcal F^{-1}(m\cdot\widehat f)\)，高频块 \(P_j\) 就是齐次块 \(\Delta_j\)，但 \(P_0\neq\Delta_0\)。对每个 \(f\in\mathcal S(\mathbb R^d)\)，有
\[
 \|J^sf\|_p\asymp
 \left\|\left(|P_0f|^2+
              \sum_{j\geq1}2^{2js}|P_jf|^2\right)^{1/2}\right\|_p.
\]
常数只依赖 \(d,p,s,\chi\)。更一般地，温和分布 \(f\) 的右侧有限当且仅当 \(f\in H^{s,p}(\mathbb R^d)\)。因此与\cref{b2:def:19-inhomogeneous-square}采用同一分割时，
\(F^s_{p,2}(\mathbb R^d)=H^{s,p}(\mathbb R^d)\)，且范数等价。
\end{proposition}
\begin{proof}
记 \(w_s(\xi)=(1+4\pi^2|\xi|^2)^{s/2}\)。先证明分析方向。令 \(g=J^sf\)，则
\(2^{js}P_jf=U_jg\)，其中
\[
 u_j(\xi)=2^{js}\varphi(2^{-j}\xi)w_{-s}(\xi).
\]
缩放到固定环带，得到
\[
 u_j(2^j\eta)=
 \varphi(\eta)(2^{-2j}+4\pi^2|\eta|^2)^{-s/2}.
\]
在 \(1/2\leq|\eta|\leq2\) 上，括号的值有与 \(j\geq1\) 无关的正下界和有限上界。反复求导只产生括号的实数幂与 \(\eta\) 的多项式，所以所有预先指定阶数的导数一致有界。

令 \(k_j=\check u_j\)。固定环带上的积分分部使归一化逆变换及其一阶导数具有任意阶一致衰减。因此对任意整数 \(M>d+1\)，
\[
 |k_j(x)|\leq C_M2^{jd}(1+2^j|x|)^{-M},\quad
 |\nabla k_j(x)|\leq C_M2^{j(d+1)}(1+2^j|x|)^{-M}.
\]
这正是\cref{b2:lem:19-uniform-dyadic-kernels}求和时使用的两种界；该求和只用一致常数，不要求归一化核逐个完全相同。将 \(2^j|x|\leq1\) 与 \(>1\) 分开求几何级数，得到
\[
 \sum_{j\geq1}|k_j(x)|\leq C|x|^{-d},\quad
 \sum_{j\geq1}|\nabla k_j(x)|\leq C|x|^{-d-1}.
\]

对有限指标集 \(E\subset\{1,2,\ldots\}\)，把 \(U_Eg=(U_jg)_{j\in E}\) 看作从标量到 \(\mathbb C^E\) 的算子。每个频率只落入固定数目的环带，故
\[
 \|U_Eg\|_{L^2(\ell^2)}^2
 =\int\sum_{j\in E}|u_j(\xi)|^2\cdot|\widehat g(\xi)|^2\,\mathrm d\xi
 \leq C\|g\|_2^2.
\]
列向量核的算子范数为 \((\sum_{j\in E}|k_j|^2)^{1/2}\)，不超过上面的绝对和；梯度同理。沿线段积分梯度给出两个变量的标准核差分界。有限个 Schwartz 核又保证支集外的普通卷积表示。因此\cref{b2:prop:18-hilbert-valued-cz}给出与 \(|E|\) 无关的 \(L^p\to L^p(\ell^2)\) 界。让 \(E\) 增加并用单调收敛，得到全部高频分析界。低频 \(P_0J^{-s}\) 的符号光滑紧支撑，单独由 Young 不等式控制；两者合起来给出平方范数不超过 \(C\|J^sf\|_p\)。

反向需要合成，而不能把各块的标量估计相加。取实偶的
\(\widetilde\varphi\in C_c^\infty(\{\,1/4<|\xi|<4\,\})\)，在
\(\operatorname{supp}\varphi\) 的邻域等于一；另取实偶
\(\widetilde\chi\in C_c^\infty\)，在 \(\operatorname{supp}\chi\) 的邻域等于一。令 \(V_0,V_j\) 的符号分别为
\[
 v_0=w_s\widetilde\chi,\quad
 v_j(\xi)=2^{-js}w_s(\xi)\widetilde\varphi(2^{-j}\xi).
\]
对高频族，归一化符号为
\[
 v_j(2^j\eta)=
 (2^{-2j}+4\pi^2|\eta|^2)^{s/2}\widetilde\varphi(\eta).
\]
与前面完全相同的固定环带导数界、有限重叠与列向量核证明给出
\[
 \left\|\left(\sum_{j\in E}|V_jh|^2\right)^{1/2}\right\|_{p'}
 \leq C\|h\|_{p'}.
\]
所有 \(v_j\) 均为实值，故 \(V_j\) 在积分配对下自伴随。对任意有限族 \(b_j\in L^p\)，逐点 Cauchy--Schwarz 与 Hölder 因而给出
\[
 \begin{aligned}
 \left|\int\left(\sum_{j\in E}V_jb_j\right)\cdot\overline h\right|
 &\leq\int
       \left(\sum_{j\in E}|b_j|^2\right)^{1/2}
       \cdot\left(\sum_{j\in E}|V_jh|^2\right)^{1/2}\\
 &\leq C\left\|\left(\sum_{j\in E}|b_j|^2\right)^{1/2}\right\|_p
          \cdot\|h\|_{p'}.
 \end{aligned}
\]
先在 Schwartz 输入上核对配对，固定有限族的 \(L^p\) 连续性再推广到一般 \(b_j\)。对偶性得合成界；低频项 \(V_0b_0\) 再用其 Schwartz 核单独估计。于是
\[
 \left\|V_0b_0+\sum_{j\in E}V_jb_j\right\|_p
 \leq C\left\|\left(|b_0|^2+\sum_{j\in E}|b_j|^2\right)^{1/2}\right\|_p.
\]
常数仍独立于有限指标集。

由于辅助截断在相应频率支撑上等于一，有限重构 \(f_N=P_0f+\sum_{j=1}^NP_jf\) 满足
\[
 J^sf_N=V_0P_0f+\sum_{j=1}^NV_j(2^{js}P_jf).
\]
对 Schwartz \(f\)，非齐次重构在 Schwartz 拓扑收敛。用刚才的合成界并取极限，得到所求反向不等式。

最后处理任意温和分布。若 \(g=J^sf\in L^p\)，取 \(g_n\in\mathcal S\) 在 \(L^p\) 中逼近它，令 \(f_n=J^{-s}g_n\)。分析界说明加权块序列在 \(L^p(\ell^2)\) 中收敛；逐个块在 \(\mathcal S'\) 中的极限又是 \(P_0f,2^{js}P_jf\)，故这个向量极限就是所需的频率块，得到右侧有限及上界。

反之，若右侧有限，每个块都是 \(L^p\) 函数。合成界可用于这些块，且对 \(N>M\geq0\)，
\[
 \|J^s(f_N-f_M)\|_p
 \leq C\left\|\left(
       \sum_{j=M+1}^N2^{2js}|P_jf|^2\right)^{1/2}\right\|_p
 \longrightarrow0.
\]
平方尾逐点趋零，并由题设的 \(L^p\) 平方函数控制。支配收敛说明 \(J^sf_N\) 在 \(L^p\) 中有极限。另一方面，\cref{b2:eq:19-inhomogeneous-reconstruction}给出 \(f_N\to f\) 于 \(\mathcal S'\)，而 \(J^s\) 在该空间连续，所以同一序列的分布极限是 \(J^sf\)。两个极限一致，证明 \(J^sf\in L^p\) 及反向界。
\end{proof}
估计常数独立于有限频带的个数，因而可以让频带数趋于无穷。平方和必须位于空间 \(L^p\) 范数内部；\(\sum_j2^{2js}\|P_jf\|_p^2\) 一般是另一种量。

\ProseParagraphBreak

整数阶时，频率权重应当恢复通常的弱导数。下面给出一个明确的符号恒等式，既处理混合导数，也避免把“椭圆性”留作没有写出的分解。

\Needspace{6\baselineskip}
\begin{proposition}[整数阶 Bessel 空间与 Sobolev 空间]\label{b2:prop:22-integer-sobolev}
设整数 \(d,k\geq1\)、\(1<p<\infty\)。令 \(W^{k,p}(\mathbb R^d;\mathbb C)\) 为所有弱导数 \(\partial^\alpha f\)、\(|\alpha|\leq k\)，均属于 \(L^p(\mathbb R^d)\) 的函数空间。则
\[
 H^{k,p}(\mathbb R^d)=W^{k,p}(\mathbb R^d),
\]
并且
\[
 \|J^kf\|_p
 \asymp\sum_{|\alpha|\leq k}\|\partial^\alpha f\|_p
 \asymp\|f\|_p+\sum_{j=1}^d\|\partial_j^kf\|_p.
\]
最后一项有限的 \(L^p\) 函数，已经自动具有全部至多 \(k\) 阶的 \(L^p\) 弱导数。
\end{proposition}
\begin{proof}
先从 \(H^{k,p}\) 出发。符号
\[
 \frac{(2\pi\mathrm i\xi)^\alpha}
      {(1+4\pi^2|\xi|^2)^{k/2}},
 \quad |\alpha|\leq k,
\]
在零点附近光滑有界；在高频处，\(\beta\) 阶导数为
\(O(|\xi|^{|\alpha|-k-|\beta|})\)。因而它满足 Mihlin 条件。把它作用于 \(J^kf\) 就是分布导数 \(\partial^\alpha f\)，得到第一方向及所有弱导数的存在。

反向只假定 \(f,\partial_1^kf,\ldots,\partial_d^kf\in L^p\)，其中导数按分布定义。令
\[
 q_j(\xi)=(2\pi\mathrm i\xi_j)^k,\quad
 Q(\xi)=1+\sum_{j=1}^d|q_j(\xi)|^2,\quad
 w_k(\xi)=(1+4\pi^2|\xi|^2)^{k/2}.
\]
有限维范数比较给出 \(Q(\xi)\asymp(1+|\xi|)^{2k}\)。多项式求导及恒等式 \(Q\cdot Q^{-1}=1\) 的反复微分说明
\[
 |\partial^\beta Q^{-1}(\xi)|
 \leq C_\beta(1+|\xi|)^{-2k-|\beta|}.
\]
例如一次微分为 \(-Q^{-2}\partial_\ell Q\)；更高阶每项都是 \(Q^{-1}\) 的若干因子乘上若干个 \(Q\) 的导数，总导数阶数仍为 \(|\beta|\)，给出所列幂次。

定义
\[
 a_0=\frac{w_k}{Q},\quad
 a_j=\frac{\overline{q_j}\,w_k}{Q}\quad(1\leq j\leq d).
\]
由上述界和 Leibniz 法则，\(a_0\) 的 \(\beta\) 阶导数为
\(O((1+|\xi|)^{-k-|\beta|})\)，\(a_j\) 的相应导数为
\(O((1+|\xi|)^{-|\beta|})\)，故全是 Mihlin 符号。精确恒等式
\[
 w_k=a_0+\sum_{j=1}^d a_jq_j
\]
于是给出分布等式
\[
 J^kf=T_{a_0}f+\sum_{j=1}^dT_{a_j}\partial_j^kf.
\]
这些光滑符号及其各阶导数至多多项式增长，分布乘法合法；其 \(L^p\) 延拓与该分布作用一致，可对 Schwartz 逼近使用对偶配对验证。右侧属于 \(L^p\)，且范数至多为题中最后一项的常数倍。故 \(f\in H^{k,p}\)，前半段再恢复全部混合弱导数。
\end{proof}
例如 \(k=2\) 时，只知道 \(f,\partial_1^2f,\ldots,\partial_d^2f\in L^p\)，就能控制所有 \(\partial_j\partial_kf\)，也能控制一阶导数。低阶项 \(\|f\|_p\) 保留了零频信息。该结论在整个 \(\mathbb R^d\) 上利用 Fourier 乘子建立，不自动给出任意区域或任意边界条件下的同一结论。

\ProseParagraphBreak

分数阶算子还具有一个反映非局部性的实空间公式。对 \(0<s<2\)，设
\[
 A_{d,s}=\int_{\mathbb R^d}\frac{1-\cos(2\pi z_1)}{|z|^{d+s}}\,\mathrm dz,
 \quad c_{d,s}=\frac{(2\pi)^s}{A_{d,s}}.
\]
原点处 \(1-\cos(2\pi z_1)\leq C|z|^2\)，无穷处它有界，故 \(0<A_{d,s}<\infty\)；这同时解释范围 \(0<s<2\)。
\Needspace{6\baselineskip}
\begin{proposition}[分数阶 Laplace 算子的二阶差公式]\label{b2:prop:22-fractional-kernel}
设整数 \(d\geq1\)、\(0<s<2\)，
\(c_{d,s}=(2\pi)^s/\int_{\mathbb R^d}(1-\cos(2\pi z_1))|z|^{-d-s}\,\mathrm dz\)，并以符号 \((2\pi|\xi|)^s\) 定义 \((-\Delta)^{s/2}\)。对复 Schwartz 函数 \(f\in\mathcal S(\mathbb R^d)\)，下面的第一式逐点绝对收敛，且
\[
 (-\Delta)^{s/2}f(x)=\frac{c_{d,s}}2
 \int_{\mathbb R^d}\frac{2f(x)-f(x+y)-f(x-y)}{|y|^{d+s}}\,\mathrm dy.
\]
以 \(\langle g,f\rangle=\int g\cdot\overline f\) 记复积分配对，还有
\[
 \langle(-\Delta)^{s/2}f,f\rangle
 =\frac{c_{d,s}}2
 \iint\frac{|f(x)-f(y)|^2}{|x-y|^{d+s}}\,\mathrm dx\,\mathrm dy.
\]
\end{proposition}
\begin{proof}
Taylor 积分公式使对称二阶差在原点为 \(O(|y|^2)\)，远场则由 \(f\) 的有界性控制。因为 \(0<s<2\)，这两端都可积，先得到逐点绝对收敛。还需使 Fourier 变换中的换序合法。对二阶差在 \(x\) 上取 \(L^1\) 范数，平移不变性与同一 Taylor 公式给出
\[
 \|2f-f(\cdot+y)-f(\cdot-y)\|_1
 \leq\min\left\{4\|f\|_1,\,
             |y|^2\sum_{j,k}\|\partial_j\partial_k f\|_1\right\}.
\]
该上界乘以 \(|y|^{-d-s}\) 后可积，Fubini 因而允许在全积分中作 Fourier 变换。所得乘子为
\[
 c_{d,s}\int_{\mathbb R^d}
       (1-\cos(2\pi y\cdot\xi))|y|^{-d-s}\,\mathrm dy.
\]
在 \(\xi\ne0\) 时先旋转，再作 \(z=|\xi|y\)，它等于
\(c_{d,s}|\xi|^s A_{d,s}=(2\pi|\xi|)^s\)；在 \(\xi=0\) 两边均为零。这证明第一式与 Fourier 定义相同。
两边都连续：Fourier 侧的符号乘 \(\widehat f\) 可积，实空间侧用有界二阶导数及有界函数值给出与 \(x\) 无关的可积控制。因此由 Fourier 变换得到的几乎处处等式实际逐点成立。

能量公式可直接用非负积分计算。Plancherel 给出
\[
 \|f(\cdot+y)-f\|_2^2
 =2\int_{\mathbb R^d}
        (1-\cos(2\pi y\cdot\xi))\cdot|\widehat f(\xi)|^2\,\mathrm d\xi.
\]
两边乘以 \(c_{d,s}|y|^{-d-s}/2\) 再积分。Tonelli 和刚算出的乘子把右边变为
\(\int(2\pi|\xi|)^s\cdot|\widehat f(\xi)|^2\,\mathrm d\xi\)；
左边令第二个位置变量等于 \(x+y\)，便是题中的双重积分。积分有限性也由
\(\|f(\cdot+y)-f\|_2^2\leq
\min\{4\|f\|_2^2,|y|^2\cdot\|\nabla f\|_2^2\}\)
直接保证。
\end{proof}
若 \(f\geq0\) 且紧支集，在支集外有
\((-\Delta)^{s/2}f(x)=-c_{d,s}\int f(y)|x-y|^{-d-s}\,\mathrm dy<0\)，除非 \(f=0\)。
因此这个导数一般不保持支集，也不保持非负性；它的二次型却非负。这三个性质并不矛盾：非负的是总能量，而不是每一点的算子输出。与整数 Laplace 算子的局部性相比，这个尾部正是分数阶模型的本质区别。

频率定义中的 Bessel 势也有一个不振荡的实空间核。取整数 \(d\ge1\)、\(s>0\)、\(x\in\mathbb R^d\setminus\{0\}\)，定义（原点值按可积函数的等价类解释）
\[
 G_s(x)=\frac1{\Gamma(s/2)}
 \int_0^\infty t^{s/2-1}e^{-t}
 (4\pi t)^{-d/2}e^{-|x|^2/(4t)}\,\mathrm dt.
\]
这是把热核作正权平均，权的总质量为一。它与齐次 Riesz 势的区别恰在额外因子 \(e^{-t}\)：大时间被截去，因而不会保留齐次核的长距离幂尾。
\Needspace{6\baselineskip}
\begin{proposition}[Bessel 位势的正核]\label{b2:prop:22-bessel-kernel}
设整数 \(d\geq1\)、\(s>0\)，定义
\[
 G_s(x)=\frac1{\Gamma(s/2)}
 \int_0^\infty t^{s/2-1}e^{-t}
 (4\pi t)^{-d/2}e^{-|x|^2/(4t)}\,\mathrm dt,
 \quad x\ne0.
\]
原点值按可积函数的等价类解释。令 \(J^{-s}\) 为符号 \((1+4\pi^2|\xi|^2)^{-s/2}\) 的乘子。有 \(G_s\geq0\)、\(\int G_s=1\)、\(\widehat G_s(\xi)=(1+4\pi^2|\xi|^2)^{-s/2}\)。故 \(J^{-s}f=G_s*f\)，并对全部 \(1\leq p\leq\infty\)、\(f\in L^p(\mathbb R^d;\mathbb C)\) 满足 \(\|J^{-s}f\|_p\leq\|f\|_p\)。
当 \(0<s<d\) 时，\(G_s(x)\asymp |x|^{s-d}\) 于 \(0<|x|\leq1\)；当 \(s=d\) 时量级为 \(1+\log(1/|x|)\)；当 \(s>d\) 时在原点附近有界。无穷远处三种情形均有指数衰减。
\end{proposition}
\begin{proof}
非负性与 Tonelli 给出总质量一。热核 Fourier 变换为 \(e^{-4\pi^2t|\xi|^2}\)，对 \(t\) 的 gamma 积分得到所述符号，Young 即给出压缩界。卷积定义对包括 \(p=\infty\) 的所有输入都有意义；其分布 Fourier 变换也满足该乘子公式。若 \(x\) 限制在远离零点的紧集，热核积分在 \(t\to0\) 的指数衰减控制任意次 \(x\) 导数，\(t\to\infty\) 则由 \(e^{-t}\) 控制，所以 \(G_s\) 在原点之外光滑。

记 \(r=|x|\)。在 \(0<r\leq1\) 时，积分的 \(t\geq1\) 部分有统一有限上界。对 \(0<t<1\)，用 \(e^{-1}\leq e^{-t}\leq1\)，再换元 \(u=r^2/(4t)\)，得到与
\[
 r^{s-d}\int_{r^2/4}^\infty
 u^{(d-s)/2-1}e^{-u}\,\mathrm du
\]
可比的表达。若 \(s<d\)，积分趋于一个正的 gamma 常数；若 \(s=d\)，零点附近是 \(\int_{r^2/4}^1\mathrm du/u\)；若 \(s>d\)，它的下限奇性恰好抵消外面的 \(r^{s-d}\)，结果有界。下界可以分别在固定 \(u\) 区间或 \(r^2/4<u<1\) 中取出，故这些不仅是上界。

当 \(r\geq1\)，使用 \(t+r^2/(4t)\geq r\)，有
\[
 e^{-t-r^2/(4t)}
 \leq e^{-r/2}e^{-t/2-r^2/(8t)}
 \leq e^{-r/2}e^{-t/2-1/(8t)}.
\]
最后一个因子乘上任意固定实数幂的 \(t\) 后，在零点及无穷远都可积。因此
\(G_s(x)\leq C_{d,s}e^{-|x|/2}\) 于 \(|x|\geq1\)，给出明确的远场上界。
\end{proof}
例如在 \(d=1,s=2\) 时，
\[
 G_2(x)=\frac12e^{-|x|}.
\]
事实上，右侧的 Fourier 变换为
\(\int_0^\infty e^{-x}\cos(2\pi x\xi)\,\mathrm dx
=(1+4\pi^2\xi^2)^{-1}\)，所以由 \(L^1\) Fourier 变换的单射性得到公式。它在原点连续，但一阶导数的右极限减左极限等于 \(-1\)，因而分布恒等式
\((1-\partial_x^2)G_2=\delta_0\) 成立。可积核与导数中的点质量同时出现，正好解释卷积解为什么满足二阶方程。

\ProseParagraphBreak

对任意 \(s,t>0\)，两个正核的卷积也属于 \(L^1\)，Fourier 变换相乘得到
\[
 G_s*G_t=G_{s+t}.
\]
因此连续施加两个 Bessel 位势把阶数相加。这与 \(J^{-s}J^{-t}=J^{-(s+t)}\) 一致，但这里还给出实空间的正卷积表示及端点压缩。

\ProseParagraphBreak

这个核解释了非齐次空间的低频项：常数频率附近的符号有界，所以 \(J^{-s}\) 保持同一 \(L^p\)；齐次符号 \(|\xi|^{-s}\) 在零频发散，改善可积性时必须改变目标指数或规定低频归一化。两者在高频上具有相同阶数，却不能在全空间上互相替代。

\ProseParagraphBreak

核估计还给出 Sobolev 嵌入的另一种表述。若 \(1<p<\infty\)、\(0<s<d/p\)，置 \(1/q=1/p-s/d\)。由 \(G_s(x)\leq C|x|^{s-d}\) 和 HLS，
\[
 \|f\|_q=\|G_s*(J^sf)\|_q\leq C\|J^sf\|_p.
\]
若 \(s>d/p\)，则 \(G_s\in L^{p'}\)：在 \(s<d\) 时，原点幂次条件正是
\((d-s)p'<d\)；\(s=d\) 时的对数在有限指数下可积，\(s>d\) 时局部有界，远场均由指数衰减处理。因此
\(\|f\|_\infty\leq\|G_s\|_{p'}\cdot\|J^sf\|_p\)，并得到连续代表。具体地，写 \(g=J^sf\)，则
\[
 |f(x+h)-f(x)|
 \leq\|G_s(\cdot+h)-G_s\|_{p'}\cdot\|g\|_p.
\]
有限指数下的核平移连续性使右端趋零，且与 \(x\) 无关。把 \(g\) 用有界紧支撑函数在 \(L^p\) 中逼近，卷积误差就在一致范数中趋零；每个近似卷积由核的远场指数界在无穷远趋零。因此该代表属于 \(C_0(\mathbb R^d)\)。

\ProseParagraphBreak

在临界 \(s=d/p\) 时，核不属于 \(L^{p'}\)，不能原样保留无穷端点。取
\[
 h(x)=|x|^{-s}\bigl(\log(e/|x|)\bigr)^{-\beta}
          \cdot\mathbf1_{\{\,0<|x|<1/2\,\}},
 \quad \frac1p<\beta\leq1.
\]
径向积分给出
\(\int|h|^p=C_d\int_0^{1/2}r^{-1}(\log(e/r))^{-\beta p}\,\mathrm dr<\infty\)。
于是 \(f=G_s*h\in H^{s,p}\)，并且 \(f\) 几乎处处有限。另一方面，正核的近原点下界给出
\[
 f(0)\geq c\int_0^{1/2}
         r^{-1}\bigl(\log(e/r)\bigr)^{-\beta}\,\mathrm dr=\infty.
\]
核在原点之外连续，所以对任意 \(x_n\to0\)，被积函数
\(G_s(x_n-y)\cdot h(y)\) 在几乎每个 \(y\) 趋于 \(G_s(-y)\cdot h(y)\)。
Fatou 因而给出 \(\liminf_n f(x_n)=\infty\)。这说明输出在原点的整个小邻域内可任意大，排除了通过修改一个零测集恢复本性有界性的可能。故临界 \(H^{d/p,p}\) 一般不嵌入 \(L^\infty\)。

\ProseParagraphBreak

最后，椭圆导数估计与这些嵌入处理不同问题。前者从方程右端恢复同指数的最高阶导数，后者把已经拥有的导数换成更高的可积指数或连续性。应用偏微分方程时，常见路线是先证明解处于 \(H^{s,p}\)，再按 \(s\) 与 \(d/p\) 的相对位置使用嵌入。若只看方程的阶数就声称解连续，会遗漏右端的可积指数及维数这两个决定性参数。

\ProseParagraphBreak

齐次导数与非齐次范数之间还有一个常用的精确比较。对 \(s>0\)、\(1<p<\infty\)，在 Schwartz 类上有
\[
 \|J^sf\|_p\asymp\|f\|_p+\|(-\Delta)^{s/2}f\|_p.
\]
正向用两个乘子：\(J^{-s}\) 以及
\((2\pi|\xi|)^s(1+4\pi^2|\xi|^2)^{-s/2}\)。
后者在原点不一定普通光滑，但定义域只要求 \(\xi\ne0\)，其尺度导数依然受控，故满足 Mihlin 条件。反向取低频截断 \(\chi\)，低频 \(J^s\chi\) 的核可积；高频符号
\[
 (1-\chi(\xi))\frac{(1+4\pi^2|\xi|^2)^{s/2}}{(2\pi|\xi|)^s}
\]
是零阶 Mihlin 符号，把它作用于齐次导数即可恢复高频 \(J^sf\)。两部分合并证明比较；Schwartz 类在 \(H^{s,p}\) 中稠密，因此按
\(\|f\|_p+\|(-\Delta)^{s/2}f\|_p\) 完备化得到同一个空间。这里的齐次导数是由这个图范数闭包定义的闭算子：若 \(f_n\to0\) 于 \(L^p\) 且齐次导数在 \(L^p\) 中收敛，范数比较使 \(J^sf_n\) 收敛；分布连续性迫使其极限为零，进而齐次导数极限也为零。它没有借此被定义到任意 \(\mathcal S'\) 上。

\ProseParagraphBreak

因此非齐次空间不是拒绝齐次导数，而是在其之外补上低阶 \(L^p\) 大小。高频上的两个定义等价，真正差别留在零频附近。使用能量方法时，可以在已说明的空间范围内自由改写这个图范数；使用全空间齐次逆、缩放不变估计或商掉多项式时，却必须保留齐次定义自己的归一化。一个范数比较不能自动解除所有齐次分布的低频歧义。

\subsection{椭圆型估计}
\label{b2:subsec:220303}

从方程右端恢复二阶导数时，要除以二阶主符号。若它沿某个方向为零，除法可能失去控制；正定性则给出与方向无关的分母下界。先处理常系数全空间问题，随后说明低频项、边界和系数变化分别增加什么困难。本小节记
\[
 \|D^2u\|_p\coloneqq\sum_{j,k=1}^d\|\partial_j\partial_ku\|_p.
\]

\Needspace{6\baselineskip}
\begin{proposition}[常系数椭圆导数估计]\label{b2:prop:22-elliptic-estimate}
设整数 \(d\geq1\)、\(1<p<\infty\)，\(A=(a_{jk})\) 是实对称常矩阵，并且存在 \(0<\nu\leq\Lambda<\infty\)，使
\[
 \nu|\xi|^2\leq\xi^{\mathsf T}A\xi\leq\Lambda|\xi|^2
 \quad(\xi\in\mathbb R^d).
\]
令 \(Lu\coloneqq-\sum_{j,k}a_{jk}\partial_j\partial_ku\)。则每个
\(u\in W^{2,p}(\mathbb R^d;\mathbb C)\) 满足
\[
 \|D^2u\|_p\leq C_{d,p,\nu,\Lambda}\|Lu\|_p.
\]
所有导数及方程均按分布理解；对光滑函数，它们与通常导数相同。
\end{proposition}
\begin{proof}
先设 \(u\in C_c^\infty\)。记 \(q(\xi)=\xi^{\mathsf T}A\xi\)，则
\[
 \widehat{Lu}(\xi)=4\pi^2q(\xi)\cdot\widehat u(\xi),\quad
 \widehat{\partial_j\partial_ku}(\xi)
 =-\frac{\xi_j\xi_k}{q(\xi)}\widehat{Lu}(\xi)
 \quad(\xi\ne0).
\]
商 \(m_{jk}=-\xi_j\xi_k/q\) 零次齐次。在单位球面上，\(q\geq\nu\)，而 \(q\) 的至多二阶导数由 \(\Lambda\) 控制，更高阶导数为零。商的反复求导只产生这些导数的有限乘积及 \(q\) 的负幂，故单位球面上的至多 \(d+2\) 阶导数由 \(d,\nu,\Lambda\) 一致控制。齐次性再给出
\[
 |\partial^\alpha m_{jk}(\xi)|
 \leq C_{d,\alpha,\nu,\Lambda}|\xi|^{-|\alpha|}.
\]
由\cref{b2:thm:22-mihlin}得到每个二阶导数的界，有限相加即可。

对一般 \(W^{2,p}\) 函数，先以光滑单位积分核磨光，其至多二阶弱导数均在 \(L^p\) 中收敛。再乘以在 \(B_R\) 为一、在 \(B_{2R}\) 外为零的缩放截断。乘积求导产生的附加项由
\(CR^{-1}\|\nabla u\|_p\)、\(CR^{-2}\|u\|_p\) 控制，原有函数及导数在 \(|x|>R\) 上的尾范数趋零。因此可以取得 \(u_n\in C_c^\infty\)，使 \(u_n\to u\) 于 \(W^{2,p}\)。常系数保证 \(Lu_n\to Lu\) 于 \(L^p\)，将已证估计传给极限即得结论。
\end{proof}
这个命题是先验估计：它从已经存在的 \(u\) 及 \(Lu\) 比较范数，没有构造对每个右端 \(f\) 的逆解。即使 \(A=I\)，低频仍阻止同指数的无质量逆。取非零 \(\psi\in C_c^\infty\)、\(u_R(x)=\psi(x/R)\)，则
\[
 \|u_R\|_p=R^{d/p}\|\psi\|_p,\quad
 \|Lu_R\|_p=R^{d/p-2}\|L\psi\|_p.
\]
让 \(R\to\infty\)，便排除独立于 \(u\) 的估计 \(\|u\|_p\leq C\|Lu\|_p\)。二阶导数在相同缩放下也带有 \(R^{-2}\)，所以刚证明的最高阶估计没有矛盾。

\ProseParagraphBreak

添加正的零阶项以后，频率原点的除法也成为合法操作。这个变化既给出解的存在唯一性，也产生一个尺度一致的参数族。

\Needspace{6\baselineskip}
\begin{proposition}[带正参数的椭圆预解算子]\label{b2:prop:22-elliptic-resolvent}
设整数 \(d\geq1\)、\(1<p<\infty\)，实对称常矩阵 \(A\) 满足
\(\nu|\xi|^2\leq\xi^{\mathsf T}A\xi\leq\Lambda|\xi|^2\)，其中
\(0<\nu\leq\Lambda<\infty\)。令
\(L=-\sum_{j,k}a_{jk}\partial_j\partial_k\)。对每个 \(\lambda>0\) 和
\(f\in L^p(\mathbb R^d;\mathbb C)\)，分布方程
\[
 (\lambda+L)u=f
\]
有唯一的温和分布解。该解属于 \(W^{2,p}(\mathbb R^d)\)，记为
\(u=(\lambda+L)^{-1}f\)，并满足
\[
 \lambda\|u\|_p+\sqrt\lambda\|\nabla u\|_p+\|D^2u\|_p
 \leq C_{d,p,\nu,\Lambda}\|f\|_p.
\]
常数与 \(\lambda,f\) 无关。特别地，\((1+L)^{-1}:L^p\to W^{2,p}\) 是有界解算子。
\end{proposition}
\begin{proof}
函数 \(b_\lambda(\xi)=(\lambda+4\pi^2\xi^{\mathsf T}A\xi)^{-1}\) 在全空间光滑，且对固定 \(\lambda>0\)，它及其各阶导数至多多项式增长。因此可在 \(\mathcal S'\) 中定义
\(\widehat u=b_\lambda\cdot\widehat f\)，并立即得到分布方程。若齐次方程的温和分布解为 \(v\)，把
\((\lambda+4\pi^2\xi^{\mathsf T}A\xi)\widehat v=0\)
乘以同一个光滑逆，即得 \(\widehat v=0\)。这证明的是整个温和分布类中的唯一性。

为证明正则性与统一界，考察三个加权符号
\[
 \frac{\lambda}{\lambda+4\pi^2\xi^{\mathsf T}A\xi},\quad
 \frac{\sqrt\lambda\,2\pi\mathrm i\xi_j}
      {\lambda+4\pi^2\xi^{\mathsf T}A\xi},\quad
 \frac{-4\pi^2\xi_j\xi_k}
      {\lambda+4\pi^2\xi^{\mathsf T}A\xi}.
\]
令 \(\xi=\sqrt\lambda\,\zeta\)，它们分别变为
\[
 \frac1{1+4\pi^2\zeta^{\mathsf T}A\zeta},\quad
 \frac{2\pi\mathrm i\zeta_j}{1+4\pi^2\zeta^{\mathsf T}A\zeta},\quad
 \frac{-4\pi^2\zeta_j\zeta_k}{1+4\pi^2\zeta^{\mathsf T}A\zeta}.
\]
在 \(|\zeta|\leq1\) 时，分母至少为一，有限阶导数有界。在
\(|\zeta|\geq1\) 时，分母至少为 \(4\pi^2\nu|\zeta|^2\)；反复求导分别给出阶数 \(-2,-1,0\) 的符号界。因此三者都有只依赖
\(d,\nu,\Lambda\) 的 Mihlin 常数。尺度导数中的
\(|\xi|^{|\alpha|}\) 正好抵消缩放微分产生的
\(\lambda^{-|\alpha|/2}\)，所以原来的参数族有相同常数。

对 Schwartz 右端应用 Mihlin 定理，就得到题中估计。对一般
\(f\in L^p\)，取 Schwartz 逼近。该估计使对应 \(u_n\) 及一、二阶导数分别在 \(L^p\) 中收敛；分布求导与极限相容，所以极限属于 \(W^{2,p}\) 并满足方程。前面已证的分布唯一性将它与 Fourier 定义的 \(u\) 识别，统一界也传给极限。
\end{proof}
取 \(A=I,\lambda=1\)，解就是前面的 \(J^{-2}f=G_2*f\)。在一维还可直接写为
\[
 u(x)=\frac12\int_{\mathbb R}e^{-|x-y|}\cdot f(y)\,\mathrm dy.
\]
该公式把正核、分布中的导数跳跃与 \(W^{2,p}\) 解连接起来。对于一般对角矩阵
\(A=\operatorname{diag}(a_1,\ldots,a_d)\)，换元
\(x_j=\sqrt{a_j}\,y_j\) 将 \(L\) 化为 \(-\Delta_y\)，而
\(\partial_{x_j}=a_j^{-1/2}\partial_{y_j}\)。在一维，
\(\|u''\|_p=a_1^{-1}\|Lu\|_p\) 恰为等式，表明椭圆常数趋零时最高阶估计确会退化。一般正定矩阵先作固定正交对角化即可理解同一现象。

\ProseParagraphBreak

当 \(\lambda\downarrow0\)，最高阶导数界仍然统一，但零阶估计只给出
\(\|u\|_p\leq C\lambda^{-1}\|f\|_p\)。这没有产生同指数的 \(L^{-1}\)。温和分布中的无质量齐次方程还允许满足 \(LP=0\) 的多项式；其中仿射函数的二阶导数为零，而在 \(d\geq2\) 时，更高次的这类多项式可以具有非零二阶导数。这些非零多项式不属于本节有限指数的 \(W^{2,p}(\mathbb R^d)\)，不能代入先验命题。自由度来自所选解空间及低频归一化，不能统称为二阶导数都看不见的项。

\ProseParagraphBreak

全空间与有边界区域也要区别。在 \(\mathbb R^2\) 的单位球上，
\(u(x)=x_1x_2\) 满足 \(\Delta u=0\)，却有
\(\partial_1\partial_2u=1\)。所以不加边界条件，不能把全空间结论改写为
\(\|D^2u\|_{L^p(B_1)}\leq C\|\Delta u\|_{L^p(B_1)}\)。
若只求内部估计，截断后的低阶项会记录边界之外的信息。

\ProseParagraphBreak

变系数时，空间乘法不与 Fourier 变换交换，不能把 \(a_{jk}(x)\) 当成常数塞进上述商。下面的冻结步骤说明，在一个小球内接近固定矩阵时，哪一部分误差可以吸收。这里选择非散度形式，因而不要求对系数求导。

\Needspace{6\baselineskip}
\begin{proposition}[冻结系数的先验估计]\label{b2:prop:22-freezing-estimate}
设整数 \(d\geq1\)、\(1<p<\infty\)、\(r>0\)，以
\(\|D^2v\|_p=\sum_{j,k=1}^d\|\partial_j\partial_kv\|_p\)
记二阶导数和范数。令 \(A_0\) 为实对称正定常矩阵，并取常数 \(C_0\) 使
\[
 \|D^2v\|_p\leq C_0\|L_{A_0}v\|_p
 \quad(v\in W^{2,p}(\mathbb R^d)),\quad
 L_{A_0}v=-\sum_{j,k}a_{jk}^0\partial_j\partial_kv.
\]
设实矩阵场 \(A=(a_{jk})\) 在 \(B_r\subset\mathbb R^d\) 上有界可测，满足
\[
 \max_{j,k}\|a_{jk}-a_{jk}^0\|_{L^\infty(B_r)}
 \leq\varepsilon,\quad C_0\varepsilon\leq\frac12.
\]
对 \(u\in W^{2,p}(B_r)\)，以可测函数
\(L_Au=-\sum_{j,k}a_{jk}\partial_j\partial_ku\) 定义非散度算子。
若 \(u\) 的支撑紧包含于 \(B_r\)，则
\[
 \|D^2u\|_{L^p(B_r)}\leq2C_0\|L_Au\|_{L^p(B_r)}.
\]
一般 \(u\in W^{2,p}(B_r)\) 则满足
\[
 \begin{aligned}
 \|D^2u\|_{L^p(B_{r/2})}
 \leq C\bigl(&\|L_Au\|_{L^p(B_r)}
       +r^{-1}\|\nabla u\|_{L^p(B_r)}\\
       &+r^{-2}\|u\|_{L^p(B_r)}\bigr),
 \end{aligned}
\]
其中 \(C\) 可依赖 \(d,p,C_0\) 及各系数的 \(L^\infty\) 上界，但与 \(u,r\) 无关。
\end{proposition}
\begin{proof}
先设 \(u\) 紧支于球内。它在球外延零后属于全空间 \(W^{2,p}\)：在边界的一个邻域内原函数已经为零，所以分布求导没有边界项。由系数差的上界及本小节选定的二阶导数和范数，
\[
 \|(L_{A_0}-L_A)u\|_p
 \leq\varepsilon\sum_{j,k}\|\partial_j\partial_ku\|_p
 =\varepsilon\|D^2u\|_p.
\]
常系数命题于是给出
\[
 \|D^2u\|_p
 \leq C_0\|L_Au\|_p+C_0\varepsilon\|D^2u\|_p.
\]
移项即得第一式。

对一般 \(u\)，取 \(\eta\in C_c^\infty(B_r)\)，在 \(B_{r/2}\) 上为一，并使
\(\|\nabla\eta\|_\infty\leq C_dr^{-1}\)、
\(\|D^2\eta\|_\infty\leq C_dr^{-2}\)。弱导数的乘积公式给出
\[
 \begin{aligned}
 L_A(\eta\cdot u)
 ={}&\eta\cdot L_Au
 -\sum_{j,k}a_{jk}
       \bigl((\partial_j\eta)\cdot\partial_ku
              +(\partial_k\eta)\cdot\partial_ju\bigr)\\
 &-\sum_{j,k}a_{jk}(\partial_j\partial_k\eta)\cdot u.
 \end{aligned}
\]
对 \(\eta\cdot u\) 使用第一式，再以系数上界及截断导数界控制右端。在
\(B_{r/2}\) 内，\(D^2(\eta\cdot u)=D^2u\)，因而得到内部估计。
\end{proof}
连续系数在足够小的球上可以满足该小振荡条件；结论的输入仍是已经属于 \(W^{2,p}\) 的函数。若起始解只有 \(H^1\)，尚不能先把
\(a_{jk}\partial_j\partial_ku\) 当作 \(L^p\) 函数，再以这个估计证明二阶导数存在。存在性或由弱解提升正则性还需要合法的逼近、解的构造及收敛论证。

\ProseParagraphBreak

散度算子是另一个对象。只有系数有足够的弱导数时，才可把
\(-\sum_{j,k}\partial_j(a_{jk}\partial_ku)\) 展开成非散度部分及
\(-\sum_{j,k}(\partial_ja_{jk})\partial_ku\)。任意有界可测系数不允许这样把额外项当作普通函数移到右侧。\chapref{b2:ch:16}和\appref{b2:app:H}分别处理相应椭圆正则问题及低正则系数；乘子估计则为全空间常系数方程及小振荡扰动提供先验控制。

\subsection{球面乘子 预告}
\label{b2:subsec:220304}

径向符号只依赖 \(|\xi|\)，仍可能在球面上发生跳跃。
光滑环带截断 \(\chi(|\xi|)\) 的核属于 \(L^1\)；球特征函数 \(\mathbf1_{\{|\xi|\le1\}}\) 则不满足任何跨边界的 Mihlin 光滑条件。
把跳跃平滑成宽度 \(\epsilon\) 的过渡层时，第 \(N\) 阶导数通常增长如 \(\epsilon^{-N}\)。对每个固定 \(\epsilon\) 应用定理可得有界性，却得不到 \(\epsilon\downarrow0\) 的统一界。

\ProseParagraphBreak
Bochner–Riesz 平均以 \((1-|\xi|^2)_+^\delta\) 逐渐消减边界跳跃，其所需光滑阶数与 \(p\)、维数有关。Mihlin 定理给出方便的充分条件，但常不是锐利范围。\chapref{b2:ch:24}用径向核的振荡和曲率研究这一问题；\chapref{b2:ch:23}的限制估计控制频率薄壳上的集中，则说明为什么乘子问题会与曲面延拓相遇。

\ProseParagraphBreak
因此“径向”“有界”“紧支撑”三项都不能单独替代乘子条件。对核使用 Young 不等式、对符号使用尺度导数、对薄壳使用曲面几何，是三条不同证明路线。应用时应明确正在使用哪一条，以及它在端点或退化极限中会失去哪一个条件。

平滑乘子还有一个时间演化模型，可与频率硬截断比较。对整数 \(d\geq1\)、\(t>0\)，热算子 \(e^{t\Delta}\) 的符号是 \(e^{-4\pi^2t|\xi|^2}\)，核为
\[
 H_t(x)=(4\pi t)^{-d/2}e^{-|x|^2/(4t)}.
\]
其空间核非负，积分为一，因而在全部 \(L^p\) 上压缩。写
\(H_t(x)=t^{-d/2}H_1(x/\sqrt t)\)，对多重指标 \(\alpha\) 求导得到
\[
 \partial^\alpha H_t(x)
 =t^{-(d+|\alpha|)/2}
       (\partial^\alpha H_1)(x/\sqrt t),\quad
 \|\partial^\alpha H_t\|_1
 =t^{-|\alpha|/2}\|\partial^\alpha H_1\|_1.
\]
固定时间核的每个导数都是多项式乘 Gaussian，故右端常数有限。Young 不等式于是对每个整数 \(k\geq0\) 给出
\[
 \|D^ke^{t\Delta}f\|_p\leq C_k t^{-k/2}\|f\|_p,
 \quad 1\leq p\leq\infty,\quad t>0.
\]
这里 \(\|D^k u\|_p\) 可取全部 \(k\) 阶导数范数的有限和。这个证明连两个端点也包括。每个正时间的任意阶导数还可由核的平移连续性得到连续代表，因而产生光滑函数。单纯的开区间 Mihlin 定理已经足以处理 \(1<p<\infty\)，但核方法在这个模型中提供了额外端点。

\ProseParagraphBreak

任意实 \(s\geq0\) 时同样有
\[
 \|e^{t\Delta}f\|_{H^{s,p}}\leq C_{s,p}(1+t^{-s/2})\|f\|_p,
 \quad1<p<\infty.
\]
在 \(0<t\leq1\) 中，把 \(t^{s/2}(1+4\pi^2|\xi|^2)^{s/2}e^{-4\pi^2t|\xi|^2}\) 以 \(\zeta=\sqrt t\,\xi\) 重写，成为
\((t+4\pi^2|\zeta|^2)^{s/2}e^{-4\pi^2|\zeta|^2}\)。
为看清常数为何统一，先只对幂函数求 \(\alpha\) 阶导数。每一项都由
\[
 C_{\alpha,\beta,\ell}\,
 \zeta^\beta(t+4\pi^2|\zeta|^2)^{s/2-\ell},
 \quad |\beta|=2\ell-|\alpha|
\]
组成。乘上 \(|\zeta|^{|\alpha|}\)，再用
\(|\zeta|^{2\ell}\leq C_\ell(t+4\pi^2|\zeta|^2)^\ell\)，就控制为
\(C(t+4\pi^2|\zeta|^2)^{s/2}\)。乘积求导中落到 Gaussian 的部分只增加固定多项式，在无穷处由 Gaussian 衰减吸收；在原点处没有剩下负幂。因此 \(0<t\leq1\) 的 Mihlin 常数统一。若 \(t\geq1\)，则把一次固定时间的平滑与剩余时间的压缩复合；\(t=1\) 时后者是恒等算子。这说明分数阶平滑不必先把 \(s\) 取整。

\ProseParagraphBreak

当 \(t\downarrow0\)，正阶范数常数发散，与一般 \(L^p\) 数据可能没有任何正阶导数完全相容。对每个固定正时间得到光滑函数，不意味着这些函数在更强空间中一致趋向原数据。对 \(1\leq p<\infty\)，热核的近似恒等性保证在原来的 \(L^p\) 中趋于初值，提升正则性则付出时间尺度。一般 \(L^\infty\) 输入没有这种范数强连续性：取半空间示性函数 \(f=\mathbf1_{\{x_1>0\}}\)，热卷积在边界附近连续趋于 \(1/2\)，原函数在两侧分别为零和一，因此每个 \(t>0\) 都有 \(\|e^{t\Delta}f-f\|_\infty\geq1/2\)。

\ProseParagraphBreak

把时间倒过来，固定 \(t>0\) 后符号成为 \(e^{4\pi^2t|\xi|^2}\)，甚至失去 \(L^2\) 乘子所需的本性有界性。取非零 \(g\in\mathcal S\)，使
\(\widehat g\in C_c^\infty(B(0,1))\)，并令
\(f_N(x)=e^{2\pi\mathrm iNx_1}\cdot g(x)\)、\(N>1\)。
其 Fourier 支撑落在 \(B(Ne_1,1)\)，所以对这些频率紧支撑输入，逆热像仍可定义，且 Plancherel 给出
\[
 \|f_N\|_2=\|g\|_2,\quad
 \|e^{-t\Delta}f_N\|_2
 \geq e^{4\pi^2t(N-1)^2}\|g\|_2.
\]
再取 \(h_N=e^{-2\pi^2t(N-1)^2}f_N\)，便有 \(\|h_N\|_2\to0\)，而
\(\|e^{-t\Delta}h_N\|_2\to\infty\)。正向平滑因此不能通过形式上的指数取反得到稳定逆过程。

\ProseParagraphBreak

与球面乘子相比，热符号没有一条频率边界上的跳跃，所有尺度的导数也具有统一控制；与齐次分数积分相比，它在零频不发散。良好的端点压缩、正时间平滑和低频稳定是这些特殊性质共同产生的，不应当归给任意一个写成 \(m(\xi)\) 的符号。

\ProseParagraphBreak

频率过渡层还可以用一个明确的算子范数检验。取固定 \(\chi\in C_c^\infty\)，令 \(m_R(\xi)=\chi(\xi/R)\)。其核为 \(R^d\check\chi(Rx)\)，所以全部 \(R>0\) 的 \(L^1\) 核范数都等于 \(\|\check\chi\|_1\)，在每个 \(L^p\) 上有同一个界。若 \(\chi(0)=1\)，在 \(R\to\infty\) 时，有限指数下这族算子强趋向恒等：先对 Schwartz 输入用核的近似恒等性，再用一致范数和稠密性推广。将频率支撑推向无穷本身并不造成常数发散；关键是过渡轮廓经过尺度归一化后仍固定。

\ProseParagraphBreak

若改用
\[
 m_\epsilon(\xi)=\theta\bigl((|\xi|-1)/\epsilon\bigr),
 \quad0<\epsilon<\frac12,
\]
其中 \(\theta\in C^\infty(\mathbb R)\)、\(0\leq\theta\leq1\)，在
\((-\infty,-1]\) 为一、在 \([1,\infty)\) 为零，则过渡层紧贴单位球面。符号在原点附近恒为一，故原点处也光滑；但归一化整个频率空间不能同时固定球的半径和层的厚度。在某个固定的过渡参数处 \(\theta'\neq0\)，径向一阶导数就达到
\(c\epsilon^{-1}\)，而 \(|\xi|\) 仍约为一，尺度导数条件因此真实增长。前一个模型只有整体半径，后一个还带有厚度与半径之比；两者的参数一致性不同。

\ProseParagraphBreak

在 \(L^2\) 中只需符号的本性上确界，故这族越来越薄的过渡层始终满足
\(\|T_{m_\epsilon}\|_{2\to2}\leq1\)。对 Schwartz 输入，支配收敛还给出
\[
 T_{m_\epsilon}f\longrightarrow
 T_{\mathbf1_{\{|\xi|\leq1\}}}f\quad\text{于 }L^2.
\]
若某个 \(1<p<\infty\) 上另外已有与 \(\epsilon\) 无关的算子界，就可用\cref{b2:prop:22-multiplier-calculus}的极限结论，把该界传给球乘子。仅有每个固定 \(\epsilon\) 的 Mihlin 界不能完成这一步；导数常数增长本身也不证明极限算子一定无界。后两章将利用频率边界的几何进一步分析实际有界范围。


\begin{chaptersummary}
乘子问题先有一个明确的 \(L^2\) 算子，再问它能否保持其他空间。复合、对偶、坐标变化与受控极限提供了基本运算，但每一步都要保留指数、定义域和常数。有限指数下的稠密性与无穷指数下的伴随定义承担不同职责；逐个输入的强收敛，也不能替代整个单位球上的范数收敛。

\ProseParagraphBreak

Mihlin 与 Hörmander 条件都是可核验的充分条件。它们通过固定环带上的统一估计控制各尺度之间的叠加，再由奇异积分理论得到 \(1<p<\infty\) 的结论。端点若要成立，常常还须使用正核、有限测度或可积核等额外结构。某个导数判据失效，只说明该判据不能直接使用，不能据此断言乘子本身无界。

\ProseParagraphBreak

在应用中，Riesz 变换组织齐次导数，Bessel 位势补上低频大小，常系数椭圆符号恢复同指数的高阶导数。热核的正性与可积性则提供额外的端点和平滑性质。符号的阶数、原点行为、边界几何和参数尺度应分别核对；将这些信息合在一个没有说明的“光滑符号”条件中，会掩盖真正使用的假设。
\end{chaptersummary}

\BookExercises

% 自拟；基础核验与反例：有限原子卷积在不同指数上的准确范数，区分总变差、符号大小和可积核。
\Needspace{6\baselineskip}
\begin{exercise}\label{b2:ex:22-new-1}
在 \(\mathbb R\) 上令 \(\tau_af(x)=f(x-a)\)，并定义
\[
 Tf=f+\tau_1f-\tau_2f.
\]
求其 Fourier 符号，证明
\[
 \|T\|_{L^1\to L^1}=3,
 \quad \|T\|_{L^2\to L^2}=\sqrt5,
 \quad \|T\|_{L^\infty\to L^\infty}=3.
\]
说明为什么卷积核的总变差只给出一般上界，以及该算子为何不能由一个普通 \(L^1\) 卷积核表示。
\end{exercise}
\begin{solution}
平移公式给出
\[
 m(\xi)=1+e^{-2\pi\mathrm i\xi}-e^{-4\pi\mathrm i\xi}.
\]
三个平移都是等距映射，三角不等式给出所有 \(1\leq p\leq\infty\) 上的范数上界三。为求 \(L^1\) 下界，取非负 \(g\in C_c^\infty\bigl((-1/4,1/4)\bigr)\)，使 \(\|g\|_1=1\)。三个平移后的支撑互不相交，因此 \(\|Tg\|_1=3\)。

令 \(I=(-1/4,1/4)\)，取
\(f=\mathbf1_I+\mathbf1_{I-1}-\mathbf1_{I-2}\)。三个区间互不相交，\(\|f\|_\infty=1\)，而对 \(x\in I\) 有 \(Tf(x)=3\)。这在一个正测度集合上达到三，证明了本性上确界意义下的 \(L^\infty\) 下界。

令 \(z=e^{-2\pi\mathrm i\xi}\)，直接展开有
\[
 |1+z-z^2|^2=3-z^2-z^{-2}=3-2\cos(4\pi\xi).
\]
该连续周期函数的本性上确界为五。由\cref{b2:prop:22-l2}，\(L^2\) 范数恰为 \(\sqrt5\)。分布核为 \(\delta_0+\delta_1-\delta_2\)，其总变差为三，却大于这个平方空间范数。

若存在 \(k\in L^1\) 给出同一算子，则 \(\widehat k=m\) 几乎处处。两边连续，故处处相等；Riemann--Lebesgue 引理要求 \(\widehat k\) 在无穷远趋零，但 \(m(n)=1\) 对所有整数 \(n\) 成立，矛盾。有限原子核可以定义有界卷积算子，并不需要成为普通可积函数。
\end{solution}

% 自拟；中等：虚幂乘子群的参数导数界、强连续性与不范数连续，超出单一Mihlin验证。
\Needspace{6\baselineskip}
\begin{exercise}\label{b2:ex:22-new-2}
设整数 \(d\geq1\)、\(1<p<\infty\)，对 \(t\in\mathbb R\) 和 \(\xi\ne0\) 令
\[
 m_t(\xi)=|\xi|^{\mathrm it}
       =e^{\mathrm it\log|\xi|},
 \quad U_t=T_{m_t}.
\]
证明它们在 \(L^p(\mathbb R^d)\) 上形成一个由有界算子组成的一参数群，并有
\[
 \|U_t\|_{p\to p}\leq C_{d,p}(1+|t|)^{d+2}.
\]
再证明 \(t\mapsto U_tf\) 对每个 \(f\in L^p\) 都范数连续，但当 \(t\ne0\) 时
\[
 \|U_t-I\|_{p\to p}\geq2.
\]
说明这两个结论为何不矛盾，并求 \(p=2\) 时最后一个范数的准确值。原点处的符号值任取，不影响问题。
\end{exercise}
\begin{solution}
求一次导数得到
\(\partial_jm_t=\mathrm it\xi_j|\xi|^{-2}\cdot m_t\)。继续求导，每个 \(r\) 阶导数都可写成 \(|\xi|^{\mathrm it-r}\) 乘以关于 \(t\) 次数不超过 \(r\) 的多项式，其系数是单位球面上的光滑函数。因此
\[
 |\partial^\alpha m_t(\xi)|
 \leq C_\alpha(1+|t|)^{|\alpha|}|\xi|^{-|\alpha|}.
\]
取到 \(d+2\) 阶并应用\cref{b2:thm:22-mihlin}，得到题中的界。点态恒等式 \(m_{s+t}=m_sm_t\)、\(m_0=1\) 与乘子复合规则给出
\(U_{s+t}=U_sU_t\)、\(U_0=I\)、\(U_t^{-1}=U_{-t}\)。这里每个参数对应一个有界算子，并未断言全体参数的范数具有同一个上界。

先证明在零点强连续。对 \(f\in\mathcal S\)，有 \(m_t\to1\) 几乎处处，且 \(|m_t-1|\leq2\)，所以 Plancherel 与支配收敛给出
\(\|(U_t-I)f\|_2\to0\)。若 \(p\ne2\)，选 \(q\in(1,\infty)\)，使 \(p\) 严格介于 \(2\) 和 \(q\) 之间。对 \(|t|\leq1\)，已证乘子界使 \(\|(U_t-I)f\|_q\) 一致有界。空间范数的 Hölder 插值便使 \(\|(U_t-I)f\|_p\to0\)。再用 Schwartz 类在 \(L^p\) 中的稠密性及 \(|t|\leq1\) 的一致算子界，把结论推广到所有 \(f\in L^p\)。群性质将零点连续性移到任意参数点。

当 \(t\ne0\) 时，\(t\log r\) 随 \(r>0\) 走遍实轴。因此 \(|e^{\mathrm it\log r}-1|\) 的上确界为二，而且接近二的值出现在正测度环带上，故
\(\|m_t-1\|_\infty=2\)。由\cref{b2:prop:22-symbol-bound}，
\[
 \|U_t-I\|_{p\to p}\geq\|m_t-1\|_\infty=2.
\]
当 \(p=2\) 时等号成立。强连续固定一个输入再令参数趋零；算子范数则对整个单位球取上确界，所测试的频率范围可以随参数改变，两种收敛要求不同。
\end{solution}

% 自拟；基础到中等：带非零下界的Mihlin符号求逆，展开倒数导数而不假定算子范数小扰动。
\Needspace{6\baselineskip}
\begin{exercise}\label{b2:ex:22-new-3}
设整数 \(d\geq1\)，\(N=d+2\)，\(m\in C^N(\mathbb R^d\setminus\{0\};\mathbb C)\)，并满足
\[
 |m(\xi)|\geq\delta>0,
 \quad
 \sup_{\xi\ne0}|\xi|^{|\alpha|}\cdot
           |\partial^\alpha m(\xi)|\leq A
 \quad(|\alpha|\leq N).
\]
证明对每个 \(1<p<\infty\)，\(T_m\) 是 \(L^p\) 上的有界可逆算子，逆为 \(T_{1/m}\)，并给出仅依赖 \(d,p,A,\delta\) 的逆范数上界。可证明倒数的尺度导数统一不超过
\[
 C_N\delta^{-1}\sum_{j=0}^N(A/\delta)^j.
\]
将结果用于 \(m(\xi)=2+e^{\mathrm it\log|\xi|}\)、\(t\in\mathbb R\)，说明为何无需先证明 \(\|T_m-2I\|_{p\to p}<2\)。
\end{exercise}
\begin{solution}
记 \(n=1/m\)。零阶界为 \(|n|\leq\delta^{-1}\)。若 \(|\alpha|=r\geq1\)，反复微分倒数，得到有限项之和；每一项具有形式
\[
 c\,m^{-j-1}\prod_{\ell=1}^j\partial^{\beta_\ell}m,
 \quad
 1\leq j\leq r,
 \quad |\beta_\ell|\geq1,
 \quad\sum_{\ell=1}^j\beta_\ell=\alpha.
\]
该形式可由一次求导归纳：导数落到 \(m^{-j-1}\) 上时增加一个一阶导数因子，落到已有因子上时增加其多重指标。乘以 \(|\xi|^r\) 后，每一项至多为 \(C\delta^{-j-1}A^j\)。对有限项求和，并包括零阶界，即得到所给统一控制。

Mihlin 定理于是给出
\[
 \|T_{1/m}\|_{p\to p}
 \leq C_{d,p}\delta^{-1}\sum_{j=0}^{d+2}(A/\delta)^j.
\]
符号乘积 \(m\cdot(1/m)=1\) 几乎处处成立，乘子复合规则便说明
\(T_mT_{1/m}=T_{1/m}T_m=I\)。这同时给出每个右端的存在性和解的唯一性，而不只是一条先验不等式。

在给定例中，三角不等式给出 \(|m|\geq1\)。虚幂符号的导数计算又给出有限的 Mihlin 常数，故每个实 \(t\) 都有上述逆。算子级的几何级数需要一个小范数条件，本题直接控制倒数符号的导数，因而不需要验证该条件。结论使用了完整的 Mihlin 假设；这里只凭一个一般乘子符号远离零，并没有证明其倒数必属于同一个乘子空间。
\end{solution}

% 自拟；中等：张量乘子的准确范数乘积，推广一维投影并构造端点失败的Schwartz测试输入。
\Needspace{6\baselineskip}
\begin{exercise}\label{b2:ex:22-new-4}
设整数 \(d\geq1\)、\(1\leq p<\infty\)，且 \(m_j\in M_p(\mathbb R)\)、\(1\leq j\leq d\)。证明
\[
 m(\xi)=\prod_{j=1}^dm_j(\xi_j)
 \quad\Longrightarrow\quad
 \|m\|_{M_p(\mathbb R^d)}
       =\prod_{j=1}^d\|m_j\|_{M_p(\mathbb R)}.
\]
对上界按坐标使用 Fubini，对下界使用乘积输入。由此证明任意有限边长的矩形
\(Q=\prod_{j=1}^d[a_j,b_j]\)、\(a_j<b_j\)，其特征函数都是 \(1<p<\infty\) 上的乘子，且范数与各边长、中心无关。最后证明它不是 \(L^1\) 乘子，也没有规范的\WeakStar{}连续 \(L^\infty\) 乘子延拓。
\end{exercise}
\begin{solution}
先在有限和的乘积输入上，按各坐标先后应用一维算子。每一步由 Fubini 和一维 \(L^p\) 界控制，得到乘积上界；这类输入在有限指数 \(L^p(\mathbb R^d)\) 中稠密，所以算子唯一延拓。在同时属于 \(L^2\) 的乘积输入上，逐坐标 Fourier 变换说明该算子的符号正是 \(m\)，故所得延拓与乘子定义一致。

若各个范数均非零，对任意 \(\varepsilon>0\)，可以取单位输入 \(f_j\in\mathcal S(\mathbb R)\)，使
\(\|T_{m_j}f_j\|_p\geq\|m_j\|_{M_p}-\varepsilon\)。这是由 Schwartz 稠密性和有界性得到的。令 \(f(x)=\prod_jf_j(x_j)\)，则 \(\|f\|_p=1\)，而输出的范数恰为各个一维输出范数的乘积。令 \(\varepsilon\downarrow0\) 得下界。若某个范数为零，则对应符号为零几乎处处，等式仍成立。

一维区间投影已在正文由 Hilbert 变换建立。设 \(C_p=\|\mathbf1_{[0,1]}\|_{M_p(\mathbb R)}\)，频率平移和缩放保持范数，故
\[
 \|\mathbf1_Q\|_{M_p(\mathbb R^d)}=C_p^d,
 \quad1<p<\infty.
\]
这给出与矩形位置和各边长无关的界，但保留维数依赖。

为检查端点，取 \(\vartheta_j\in C_c^\infty(\mathbb R)\)，在 \([a_j,b_j]\) 的邻域恒为一，并令 \(\widehat f=\prod_j\vartheta_j\)。则 \(f\in\mathcal S\)，而其投影为各个区间逆变换的乘积。一维因子为
\[
 \mathcal F^{-1}\mathbf1_{[a,b]}(x)
 =e^{\pi\mathrm i(a+b)x}\cdot
       \frac{\sin\bigl(\pi(b-a)x\bigr)}{\pi x},
\]
原点按连续值补充。其绝对值不可积：每个充分远的周期中，正弦绝对值至少为 \(1/2\) 的部分占固定比例，相应积分下界形成发散的调和级数。Tonelli 因而使多维乘积的 \(L^1\) 范数也为无穷。这给出一个 \(L^1\cap L^2\) 输入，其 \(L^2\) 投影不属于 \(L^1\)，排除了 \(L^1\) 延拓。符号实值，规范 \(L^\infty\) 延拓若存在，其 \(L^1\) 预伴随正是同一符号的乘子，也将与上述结论矛盾。
\end{solution}

% 自拟；反例与尺度：全部消失矩也不消除同指数低频障碍，求带低频截断的逆算子范数增长。
\Needspace{6\baselineskip}
\begin{exercise}\label{b2:ex:22-new-5}
设整数 \(d\geq1\)、\(1<p<\infty\)、\(\alpha>0\)。对 Fourier 变换光滑紧支撑且远离原点的函数，定义
\[
 I_\alpha f=\mathcal F^{-1}\bigl[(2\pi|\xi|)^{-\alpha}\widehat f(\xi)\bigr].
\]
构造一族 \(f_r\in\mathcal S\)、\(r\downarrow0\)，使 \(\|f_r\|_p=1\)，每个多重指标 \(\beta\) 都满足
\(\int x^\beta f_r(x)\,\mathrm dx=0\)，但 \(\|I_\alpha f_r\|_p\to\infty\)。

\ProseParagraphBreak

再固定 \(\chi\in C_c^\infty(\mathbb R^d)\)，满足 \(0\leq\chi\leq1\)，在 \(|\xi|\leq1\) 为一、\(|\xi|\geq2\) 为零。对 \(\varepsilon>0\) 令
\[
 q_\varepsilon(\xi)=(2\pi|\xi|)^{-\alpha}
                    \cdot\bigl(1-\chi(\xi/\varepsilon)\bigr),
\]
在原点定义为零。证明其乘子范数恰为
\(\varepsilon^{-\alpha}\|q_1\|_{M_p}\)，其中 \(0<\|q_1\|_{M_p}<\infty\)。解释为何“有消失矩”与“具有固定正半径的频率缺口”是不同的控制。
\end{exercise}
\begin{solution}
取非零 \(\widehat f\in C_c^\infty(\{\,\xi:1<|\xi|<2\,\})\)，并归一化使 \(\|f\|_p=1\)。因为 \(\widehat f\) 在原点邻域恒零，其原点处的全部导数为零；Fourier 求导公式给出
\(\int x^\beta f(x)\,\mathrm dx=0\)。令
\[
 f_r(x)=r^{d/p}\cdot f(rx),\quad 0<r<1.
\]
缩放保持其 \(L^p\) 范数为一。每个 \(f_r\) 的 Fourier 变换仍在原点附近恒零，因而全部矩仍然消失。

令 \(g=I_\alpha f\)。频率乘积光滑紧支撑，所以 \(g\in\mathcal S\)，且非零。直接缩放给出
\[
 I_\alpha f_r(x)=r^{-\alpha}r^{d/p}\cdot g(rx),
 \quad\|I_\alpha f_r\|_p=r^{-\alpha}\|g\|_p\longrightarrow\infty.
\]
每一个输入都有频率缺口，但该缺口的半径与 \(r\) 一同趋零；消失矩本身没有给出统一的正半径。

符号 \(q_1\) 在原点附近为零，在固定过渡环带上光滑，在 \(|\xi|\geq2\) 上为 \((2\pi|\xi|)^{-\alpha}\)。其任意有限阶尺度导数都有界，故由 Mihlin 定理属于 \(M_p\)。它不是零符号，\cref{b2:prop:22-symbol-bound}保证其乘子范数为正。恒等式
\[
 q_\varepsilon(\xi)=\varepsilon^{-\alpha}
                                  \cdot q_1(\xi/\varepsilon)
\]
和符号缩放的等范数性质给出准确关系
\(\|q_\varepsilon\|_{M_p}=\varepsilon^{-\alpha}\|q_1\|_{M_p}\)。所以移除一个固定零频邻域确实给出同指数有界逆，但该邻域缩小时，常数具有明确且不可消去的增长。
\end{solution}

% 自拟；综合应用：Riesz变换组成向量场互补投影，核对符号、散度、旋度及平方空间的准确范数。
\Needspace{6\baselineskip}
\begin{exercise}\label{b2:ex:22-new-6}
设整数 \(d\geq2\)、\(1<p<\infty\)，向量空间 \(L^p(\mathbb R^d;\mathbb C^d)\) 使用逐点欧氏范数。令 \(R_j\) 为符号 \(-\mathrm i\xi_j/|\xi|\) 的 Riesz 变换，并对向量场 \(F=(F_1,\ldots,F_d)\) 定义
\[
 (QF)_j=-\sum_{k=1}^dR_jR_kF_k,
 \quad P=I-Q.
\]
证明 \(P,Q\) 有界，且
\[
 P^2=P,\quad Q^2=Q,\quad PQ=QP=0,\quad P+Q=I.
\]
在分布意义下验证
\[
 \sum_j\partial_j(PF)_j=0,
 \quad
 \partial_j(QF)_k-\partial_k(QF)_j=0.
\]
最后证明在 \(L^2\) 中两部分正交，两个投影的范数都等于一。可用紧支撑光滑标量函数的梯度和向量场 \((\partial_2\phi,-\partial_1\phi,0,\ldots,0)\) 检验两个下界。
\end{exercise}
\begin{solution}
Riesz 变换在 \(L^p\) 上有界，有限个分量的求和使 \(P,Q\) 在所给向量空间上有界，常数允许依赖 \(d,p\)。它们的矩阵符号分别为
\[
 q(\xi)=\frac{\xi\xi^{\mathsf T}}{|\xi|^2},
 \quad p(\xi)=I-\frac{\xi\xi^{\mathsf T}}{|\xi|^2},
 \quad\xi\ne0.
\]
这里 \((-\mathrm i)^2=-1\)，故定义中的负号使 \(q\) 具有正号。矩阵 \(q\) 是沿 \(\xi\) 方向的正交投影，\(p\) 是其正交补投影。因此上述四个代数恒等式先在 \(L^2\cap L^p\) 上由 Fourier 变换成立，再由稠密性推广到 \(L^p\)。

对 Schwartz 向量场，符号满足
\[
 \xi^{\mathsf T}p(\xi)=0,
 \quad \xi_jq_{k\ell}(\xi)-\xi_kq_{j\ell}(\xi)=0.
\]
乘上微分对应的 \(2\pi\mathrm i\) 后，这就是所需分布恒等式。对于一般 \(F\in L^p\)，取 Schwartz 逼近。投影的有界性给出输出的 \(L^p\) 收敛，而 \(L^p\) 收敛蕴含温和分布收敛；分布求导连续，所以两条恒等式仍成立。

在 \(L^2\) 中，逐频率正交分解与 Plancherel 给出
\[
 \|F\|_2^2=\|PF\|_2^2+\|QF\|_2^2,
 \quad\langle PF,QF\rangle_{L^2}=0.
\]
两投影因而均为压缩。取一个非恒定 \(\phi\in C_c^\infty\)，则非零梯度场 \(F=\nabla\phi\) 满足 \(QF=F\)，给出 \(\|Q\|=1\)。再选择 \(\phi\) 使 \((\partial_2\phi,-\partial_1\phi)\not\equiv0\)；该场的 Fourier 向量与 \(\xi\) 垂直，所以 \(PF=F\)，给出 \(\|P\|=1\)。在一般 \(L^p\) 中仍有有界的互补分解，但这里的正交恒等式专属于平方范数。
\end{solution}

% 自拟；中等：光滑频率截断与空间Lipschitz系数的交换子，明确初始定义并证明尖锐尺度。
\Needspace{6\baselineskip}
\begin{exercise}\label{b2:ex:22-new-7}
设整数 \(d\geq1\)，\(\chi\in C_c^\infty(\mathbb R^d)\)，\(k=\mathcal F^{-1}\chi\)，对 \(R>0\) 令
\(T_R=T_{\chi(\,\cdot\,/R)}\)。设 \(b:\mathbb R^d\to\mathbb C\) 为 Lipschitz 函数，常数为 \(L\)，不要求 \(b\) 有界。先对 \(f\in C_c^\infty\) 定义交换子
\[
 [b,T_R]f=b\cdot T_Rf-T_R(bf).
\]
证明它可由一个绝对积分定义于全部 \(L^p\)，并满足
\[
 \|[b,T_R]f\|_p
 \leq\frac L R\left(\int|y|\cdot|k(y)|\,\mathrm dy\right)\|f\|_p,
 \quad1\leq p\leq\infty.
\]
当 \(b(x)=x_j\) 且 \(\partial_j\chi\not\equiv0\) 时，求准确的交换子符号，证明对 \(1<p<\infty\) 其范数恰为一个与 \(R\) 无关的正数乘 \(R^{-1}\)。
\end{exercise}
\begin{solution}
核为 \(k_R(y)=R^d\cdot k(Ry)\)，它是 Schwartz 函数。在初始输入上，两项的差可写为
\[
 [b,T_R]f(x)
 =\int\bigl(b(x)-b(x-y)\bigr)\cdot k_R(y)\cdot f(x-y)\,\mathrm dy.
\]
因为 \(b\) 至多线性增长、\(f\) 紧支撑，初始两项有意义。差分满足
\(|b(x)-b(x-y)|\leq L|y|\)，从而绝对值由
\(L\bigl((|\,\cdot\,|\cdot|k_R|)*|f|\bigr)(x)\) 控制。这个非负卷积对有限指数输入几乎处处有限，对有界输入处处有限，因此上式也直接定义了一般输入上的交换子。

Young 不等式和换元给出
\[
 \|[b,T_R]f\|_p
 \leq L\left(\int|y|\cdot|k_R(y)|\,\mathrm dy\right)\|f\|_p
 =\frac L R\left(\int|z|\cdot|k(z)|\,\mathrm dz\right)\|f\|_p.
\]
有限指数时它与稠密延拓一致；无穷指数时使用上述直接积分，不使用 Schwartz 范数稠密性。若 \(b\) 无界，两个单独的算子项未必对全部 \(L^p\) 输入都有意义，真正延拓的是这个有抵消的差。

当 \(b(x)=x_j\) 时，核差分恰为 \(y_jk_R(y)\)。本卷 Fourier 约定下，
\[
 \mathcal F(y_jk_R)(\xi)
 =\frac{\mathrm i}{2\pi}\partial_{\xi_j}\bigl(\chi(\xi/R)\bigr)
 =\frac{\mathrm i}{2\pi R}\cdot(\partial_j\chi)(\xi/R).
\]
符号缩放保持乘子范数，所以
\[
 \|[x_j,T_R]\|_{p\to p}
 =\frac1{2\pi R}\|\partial_j\chi\|_{M_p}.
\]
导数符号光滑紧支撑，故范数有限；它非零，\cref{b2:prop:22-symbol-bound}使范数严格为正。因此一般的 \(R^{-1}\) 阶不能改善。
\end{solution}

% 自拟；中等：Bessel正核收缩的最优常数与不取到，区分每个输入严格和算子范数严格。
\Needspace{6\baselineskip}
\begin{exercise}\label{b2:ex:22-new-8}
设整数 \(d\geq1\)、\(s>0\)、\(1<p<\infty\)，\(J^{-s}\) 为符号 \((1+4\pi^2|\xi|^2)^{-s/2}\) 的 Bessel 位势算子。使用正文已建立的核 \(G_s\)，即
\(G_s>0\) 几乎处处、\(\int G_s=1\)、\(J^{-s}f=G_s*f\)，证明
\[
 \|J^{-s}\|_{p\to p}=1,
 \quad
 \|J^{-s}f\|_p<\|f\|_p
 \quad\text{若 }f\in L^p\setminus\{0\}.
\]
给出趋近算子范数的单位输入序列，并说明在 \(p=1\) 和 \(p=\infty\) 时为何不能保留上述逐个非零输入的严格不等式。
\end{exercise}
\begin{solution}
由正核的单位质量与 Jensen 不等式，几乎处处有
\[
 |G_s*f(x)|^p
 \leq\int G_s(y)\cdot|f(x-y)|^p\,\mathrm dy.
\]
Tonelli 给出右侧的总积分为 \(\|f\|_p^p\)，所以算子范数至多为一。若范数取等，两个非负可积函数的差积分为零，故上述 Jensen 不等式对几乎每个 \(x\) 都取等。选取一个等号成立且右侧有限的点 \(x\)。函数 \(z\mapsto|z|^p\) 在复数视作实二维向量时严格凸，因此取等迫使 \(f(x-y)\) 关于概率测度 \(G_s(y)\,\mathrm dy\) 几乎处处为同一个常数。因 \(G_s>0\) 几乎处处，\(f\) 也就关于 Lebesgue 测度几乎处处为常数。有限指数 \(L^p(\mathbb R^d)\) 中只有零常数符合要求，故每个非零输入都严格缩小。

为求算子范数下界，取 \(\psi\in C_c^\infty\)、\(\|\psi\|_p=1\)，并令
\[
 f_R(x)=R^{-d/p}\cdot\psi(x/R),\quad R\longrightarrow\infty.
\]
这些都是单位输入。缩放后用 Minkowski 不等式，得到
\[
 \|G_s*f_R-f_R\|_p
 \leq\int G_s(y)\cdot
                 \|\psi(\,\cdot-y/R)-\psi\|_p\,\mathrm dy.
\]
对每个固定 \(y\)，平移连续性使范数趋零；它又不超过 \(2\|\psi\|_p=2\)，而 \(G_s\) 可积。支配收敛给出误差趋零，所以
\(\|J^{-s}f_R\|_p\to1\)，算子范数恰为一。这同时说明范数并未由某个非零 \(L^p\) 输入达到。

当 \(p=1\) 时，对任意非零非负 \(f\in L^1\)，正性与 Tonelli 直接给出
\(\|G_s*f\|_1=\|f\|_1\)。当 \(p=\infty\) 时，常数输入满足 \(G_s*1=1\)。两个端点都有非零取等输入，不能把严格凸性与有限指数空间的论证直接搬过去。
\end{solution}

% !TeX root = ../../Book2.tex
\chapter{Fourier 限制问题}
\label{b2:ch:23}
\BookChapterNumber{b2:ch:23}
\begin{chapterabstract}
Fourier 变换的环境空间等价类不能直接在零测度曲面上取值，限制估计为这种取值建立连续意义。本章先辨认帽集中、振荡尾部和曲率退化造成的必要条件，再借助延拓对偶、正二次型与解析族证明球面及非零曲率紧片的 Stein--Tomas 定理。选读部分继续讨论极值函数、紧性损失和现代几何方法，区分已证范围、条件性结论与仍需深层输入的目标。
\end{chapterabstract}

\begin{learninggoals}\begin{itemize}
\item\label{b2:item:23-goal-1}从稠密测试类定义曲面限制，分清普通 Lebesgue 范数控制、带矩迹估计以及不同曲面测度的作用。
\item\label{b2:item:23-goal-2}分别核对帽面积、相干管体积和远场振荡尾部，在曲率退化或测度改变时重新判断必要条件。
\item\label{b2:item:23-goal-3}把限制估计、延拓估计与卷积二次型联系起来，准确保留插值区域、半径因子及平方范数所在的空间。
\item\label{b2:item:23-goal-4}理解解析族的分布延拓、两条边界估计和目标参数，并辨认经典端点之后的紧性问题与几何问题。
\end{itemize}\end{learninggoals}

Hausdorff--Young 控制的是整个频率空间中的大小，不能直接约束一个零测度曲面上的取值。即使从 Schwartz 函数开始，这种取值也可能因频率集中在越来越薄的邻域而失控。本章首先寻找能够排除这种失控的输入空间；曲面上的限制随后由连续延拓定义，而不依赖某个任意选择的 Fourier 代表。

\ProseParagraphBreak

曲率带来两种看似相反的现象。分布在整个曲面上的光滑数据产生振荡抵消，给出远场衰减；集中在很小一顶帽内的数据，却可在一根很大的细长管中保持同相。\secref{b2:sec:2302}把这些现象转成必要条件。它们检验一个猜测是否可能成立，却不替代对不同帽相互作用的估计。

\ProseParagraphBreak

主线按限制、延拓、解析族的次序推进。对偶性把曲面取值变为一个绝对收敛的振荡积分，\(TT^*\) 再把平方限制界变为具有曲面 Fourier 支集的卷积估计。Stein--Tomas 证明把 Fourier 侧的有界性与空间侧的核控制放在不同参数边界上，插值到球面测度所在的位置。这里既要计算指数，也要解释分布族的定义及其增长，不能把一个形式参数代入发散积分。

\ProseParagraphBreak

所需基础包括 Fourier 反演与 Plancherel、Hölder 不等式、算子对偶、\chapref{b2:ch:17}的插值和\chapref{b2:ch:21}的驻相方法。选读的\secref{b2:sec:2305}与\secref{b2:sec:2306}分别讨论最佳常数之后的紧性，以及经典范围之外的几何组织；略过它们不影响本章端点定理的证明。章末前六题沿主线训练，后两题通过可显式计算的极值与 Gaussian 波包补充选读内容。

% !TeX root = ../../Book2.tex
\section{限制算子}
\label{b2:sec:2301}

如果 \(f\in L^1(\mathbb R^d)\)，其 Fourier 变换是一个确定的连续函数，可以在任何曲面上取值。若只知道 \(f\in L^2\)，Plancherel 给出的却是相对于 \(d\) 维 Lebesgue 测度的等价类。一个光滑超曲面的 \(d\) 维测度为零，在这个集合上任意改值都不改变等价类。因此，限制问题首先是一个定义问题：是否存在一个连续算子，与光滑输入的曲面取值一致？

\ProseParagraphBreak
这个问题与通常的迹问题相似，但所控制的范数不同。频率变量的 Sobolev 正则性可以产生迹；限制定理则希望只用原函数的普通 \(L^p\) 范数，并由曲面的振荡几何取得额外控制。两种办法的条件和作用不能混用。

\subsection{\texorpdfstring{\(\widehat f|_S\)}{Fourier 变换在 S 上的限制} 的意义}
\label{b2:subsec:230101}

先设 \(d\geq2\)，\(S\subset\mathbb R^d\) 是非空紧致光滑嵌入超曲面，取面积测度 \(\sigma\)。仍用
\(\widehat f(\xi)=\int f(x)e^{-2\pi\mathrm i x\cdot\xi}\,dx\)。
对 \(f\in\mathcal S(\mathbb R^d)\)，定义
\[
 Rf\coloneqq\widehat f|_S.
\]
\symidx{\(Rf\)：曲面 Fourier 限制}
一个 \(L^p\) 到 \(L^r\) 的\term[fourier-xianzhi]{Fourier 限制估计}[Fourier restriction estimate]是指
\[
 \|Rf\|_{L^r(S,d\sigma)}\leq C\|f\|_{L^p(\mathbb R^d)},
 \quad f\in\mathcal S(\mathbb R^d).
\]
这里 \(1\leq p<\infty\)、\(1\leq r\leq\infty\)，常数与 \(f\) 无关。对于带紧支撑光滑密度的局部曲面片，也使用同一定义，只需把面积测度换成指定的有限测度。

\ProseParagraphBreak
为什么这个估计足以解决定义问题？若 \(f_n\in\mathcal S\) 在 \(L^p\) 中趋于 \(f\)，则
\[
 \|Rf_n-Rf_m\|_r\leq C\|f_n-f_m\|_p\longrightarrow0.
\]
目标 \(L^r(S)\) 是 Banach 空间，包括 \(r=\infty\) 的情形，故有极限。用同一不等式比较两组逼近，极限不依赖选择，从而定义了唯一的有界延拓 \(R:L^p\to L^r(S)\)。这个 \(Rf\) 是曲面测度意义下的等价类；它并不要求先给环境 Fourier 等价类选一个逐点代表。

\ProseParagraphBreak
最简单的估计来自绝对积分：
\[
 |\widehat f(\xi)|\leq\|f\|_1,\quad
 \|Rf\|_r\leq\sigma(S)^{1/r}\|f\|_1,
\]
其中 \(1/\infty=0\)。所以 \(p=1\) 对所有目标指数都成立。一般 \(p\) 的延拓若已建立，则对 \(f\in L^1\cap L^p\) 仍与连续 Fourier 变换的取值一致：选同时在两种范数中逼近的 Schwartz 函数，\(L^1\) 收敛使 Fourier 变换在整个频率空间一致收敛，而限制估计给出曲面范数收敛，两种极限必相同。

\ProseParagraphBreak
相反，\(p=2\) 的困难不能仅靠选一个较弱的曲面范数消除。下面把薄壳反例写成实际的 Schwartz 函数，避免仅从零测集问题作形式判断。

\Needspace{6\baselineskip}
\begin{proposition}[薄壳输入]\label{b2:prop:23-thin-shell}
设 \(d\geq2\)，\(S\subset\mathbb R^d\) 是非空光滑嵌入超曲面，取面积测度。对任意 \(1\leq r\leq\infty\)，不存在常数 \(C\)，使
\[
 \|\widehat f|_S\|_{L^r(S)}\leq C\|f\|_2
 \quad\text{对全部 }f\in\mathcal S(\mathbb R^d)\text{ 成立}.
\]
即使只要求控制 \(S\) 上一个非空的小坐标片，也没有这样的界。
\end{proposition}
\begin{proof}
经平移和正交变换，可把一小片写成
\(S=\{(\eta,h(\eta)):\eta\in U\}\)，其中 \(U\subset\mathbb R^{d-1}\) 开，\(h\) 光滑。在这一局部表示中选非零
\(b\in C_c^\infty(U)\)，再取 \(\chi\in C_c^\infty(\mathbb R)\)，使 \(\chi(0)=1\)。对充分小的 \(\epsilon>0\)，令
\[
 \widehat f_\epsilon(\eta,\xi_d)
 =b(\eta)\chi\bigl((\xi_d-h(\eta))/\epsilon\bigr),
\]
并在该坐标邻域外延零。支撑可选在邻域内部，故这确实是一个全空间的光滑紧支撑函数，\(f_\epsilon\in\mathcal S\)。
在曲面片上的值恒为 \(b(\eta)\)，因此其 \(L^r\) 范数是一个不随 \(\epsilon\) 变化的正数。

另一方面，作 \(\xi_d=h(\eta)+\epsilon t\) 换元，Jacobi 因子为 \(\epsilon\)，Plancherel 给出
\[
 \|f_\epsilon\|_2^2
 =\|\widehat f_\epsilon\|_2^2
 =\epsilon\|b\|_{L^2(\mathbb R^{d-1})}^2
       \cdot\|\chi\|_{L^2(\mathbb R)}^2.
\]
假定的限制界会把固定正数控制为 \(O(\epsilon^{1/2})\)，令 \(\epsilon\downarrow0\) 即矛盾。
\end{proof}

薄壳越来越窄，环境平方质量趋零，曲面上的值却没有变小。若增加空间矩的控制，频率函数便获得法向正则性，情况就不同了。下面的球面模型把两种理论的差别显示出来。

\Needspace{6\baselineskip}
\begin{proposition}[带空间矩的平方迹]\label{b2:prop:23-moment-trace}
设 \(d\geq2\)，单位球面 \(S^{d-1}\) 取面积测度。对每个
\(f\in\mathcal S(\mathbb R^d)\)，有
\[
 \|\widehat f|_{S^{d-1}}\|_2^2
 \leq C_d\|f\|_2\cdot
       \bigl(\|f\|_2+\||x|f\|_2\bigr).
\]
因此球面限制唯一连续延拓到
\[
 X\coloneqq\{f\in L^2(\mathbb R^d):|x|f\in L^2(\mathbb R^d)\},
 \quad \|f\|_X=\|f\|_2+\||x|f\|_2,
\]
并保持上述估计。
\end{proposition}
\begin{proof}
对固定 \(\omega\in S^{d-1}\)，令 \(F(r)=\widehat f(r\omega)\)。
对任意 \(s\in(1/2,3/2)\)，微积分基本定理给出
\[
 |F(1)|^2\leq |F(s)|^2+
      2\int_{1/2}^{3/2}|F(r)|\cdot|F'(r)|\,dr.
\]
把 \(s\) 在长度为一的区间上积分，再对 \(\omega\) 积分。因为
\(r^{d-1}\) 在这个区间有正的统一上下界，并且
\(|F'(r)|\leq|\nabla\widehat f(r\omega)|\)，在壳
\(A=\{\xi:1/2<|\xi|<3/2\}\) 上得到
\[
 \|\widehat f|_{S^{d-1}}\|_2^2
 \leq C_d\|\widehat f\|_{L^2(A)}^2
       +C_d\|\widehat f\|_{L^2(A)}
                \cdot\|\nabla\widehat f\|_{L^2(A)}.
\]
Plancherel 和
\(\partial_{\xi_j}\widehat f=-2\pi\mathrm i\,\widehat{x_jf}\)
给出
\(\|\nabla\widehat f\|_2=2\pi\||x|f\|_2\)，即得估计。

最后，\(C_c^\infty\) 在 \(X\) 中稠密。先截断 \(f\) 的空间支撑，由两项平方可积性控制尾部；对紧支撑输入再卷积光滑近似恒等核。此时
\[
 x_j(\rho_\epsilon*f)
 =\rho_\epsilon*(x_jf)+(x_j\rho_\epsilon)*f,
\]
后一个核的 \(L^1\) 范数为 \(O(\epsilon)\)，故两项加权范数都收敛。估计应用于逼近差，使球面迹为 Cauchy 列，完成延拓。
\end{proof}

带矩范数提供的是 \(\widehat f\) 在频率变量中的一阶平方正则性。
Stein–Tomas 定理将改用一个小于 \(2\) 的普通 Lebesgue 指数，不要求空间矩。
这两种输入条件并不相互等同：在无限测度空间上，单纯改变 \(p\) 不能直接替代位置权重。

\subsection{曲面测度}
\label{b2:subsec:230102}

曲面限制必须指定相对于哪一个测度计算范数。在图形坐标
\(\xi=(\eta,h(\eta))\) 中，面积测度为
\[
 d\sigma(\xi)=\sqrt{1+|\nabla h(\eta)|^2}\,d\eta.
\]
这个因子是切向面积的 Jacobi 因子。若改用
\(d\mu=a\,d\sigma\)，其中 \(a\geq0\) 光滑且紧支撑，则在 \(a\) 有正上下界的紧片上，两个有限指数范数等价。密度若消失，反例的集中质量可能随之变小，不能仍按恒定正密度计算必要条件。以后局部化曲面时，都把密度和面积因子写进振幅中。

\ProseParagraphBreak
球面的极坐标公式具体说明了面测度怎样由薄壳得到：
\[
 \int_{\mathbb R^d}F(\xi)\,d\xi
 =\int_0^\infty r^{d-1}
       \int_{S^{d-1}}F(r\omega)\,d\sigma(\omega)\,dr.
\]
例如对连续紧支撑的 \(F\)，有
\[
 \lim_{\epsilon\downarrow0}
 \frac1{2\epsilon}\int_{1-\epsilon<|\xi|<1+\epsilon}F(\xi)\,d\xi
 =\int_{S^{d-1}}F(\omega)\,d\sigma(\omega).
\]
证明只需在固定紧壳上用一致连续性，并注意
\(r^{d-1}\to1\) 一致成立。这里已经除以厚度 \(2\epsilon\)；
没有这个归一化，壳积分本身趋于零。

\ProseParagraphBreak
对于一般定义函数 \(\rho\)，假设 \(\nabla\rho\neq0\) 于
\(S=\{\rho=0\}\) 附近，沿法向换元得到分布记号
\[
 \int F(\xi)\delta(\rho(\xi))\,d\xi
 =\int_S\frac{F(\xi)}{|\nabla\rho(\xi)|}\,d\sigma(\xi).
\]
它可先在坐标片内将 \(\rho\) 作为一个坐标验证，再用有限单位分解拼接。
特别地，
\(\delta(1-|\xi|^2)=\sigma/2\) 在单位球面上成立。
后面解析族出现的二分之一来自此处的换元，不来自 Fourier 归一化。
上述公式对测试函数成立，并不自动控制任意 \(L^p\) 输入的薄壳极限；取得这种控制仍需要限制估计。

\subsection{曲率}
\label{b2:subsec:230103}

对图形曲面 \((\eta,h(\eta))\)，选定单位法向后，其 Gaussian 曲率为
\[
 K(\eta)=\frac{\det D^2h(\eta)}
                 {(1+|\nabla h(\eta)|^2)^{(d+1)/2}},
\]
符号可能随法向约定改变，但非零性不变。因此本章使用的条件是
\(\det D^2h\neq0\)，并不要求所有主曲率同号。
在一个紧支撑片上，行列式绝对值有正下界；这使\chapref{b2:ch:21}的非退化驻相估计具有统一常数。
球面在每个点都满足这个条件。

\ProseParagraphBreak
曲率的作用有两个方面。对于光滑密度，\(d-1\) 个切向变量的二次振荡带来
\((1+|x|)^{-(d-1)/2}\) 衰减。对于小帽数据，曲面离开切平面的高度却只有帽半径的平方，这允许输出在法向保持更长距离的相干。前者用于证明充分估计，后者产生必要条件；两者是同一二次几何的不同后果。

\ProseParagraphBreak
平面上的局部问题已经能显示曲率消失的影响。取
\[
 S_0=\{(\eta,0):|\eta|<1\}\subset\mathbb R^d
\]
及其面积测度。选 \(g\in\mathcal S(\mathbb R^{d-1})\)，使
\(\widehat g\) 在该片上不恒为零，另选 \(h\in\mathcal S(\mathbb R)\)，使
\(\int h\neq0\)。令
\[
 f_L(x',x_d)=g(x')h(x_d/L),\quad L\geq1.
\]
对 \(1<p<\infty\)，
\[
 \|f_L\|_p=L^{1/p}\|g\|_p\cdot\|h\|_p,\quad
 \widehat f_L(\eta,0)=L\widehat g(\eta)\int h.
\]
任何固定 \(1\leq r\leq\infty\) 的限制范数都按 \(L\) 增长，比输入的
\(L^{1/p}\) 快。故这个有界平面片也没有 \(p>1\) 的限制界。
无穷大的平面面积不是反例的原因；法向上没有振荡损失才是关键。

\subsection{缩放}
\label{b2:subsec:230104}

单位球面不是尺度不变对象。把它放大时，既改变面测度，又改变 Fourier 变换中的空间尺度，两个因子都要保留。

\Needspace{6\baselineskip}
\begin{proposition}[曲面半径的缩放]\label{b2:prop:23-surface-scaling}
设 \(d\geq2\)，\(S\) 是紧致光滑超曲面，\(1\leq p<\infty\)、
\(1\leq r\leq\infty\)。假设相对于面积测度的限制算子
\(L^p(\mathbb R^d)\to L^r(S)\) 的最佳常数为 \(C\)。对 \(R>0\)，
\(S_R=RS\) 也取面积测度，则其最佳常数恰为
\[
 C_R=C R^{(d-1)/r-d/p'},
 \quad \frac1p+\frac1{p'}=1.
\]
\end{proposition}
\begin{proof}
取 \(g(x)=f(x/R)\)。直接换元给出
\[
 \widehat g(\xi)=R^d\widehat f(R\xi),
 \quad \|g\|_p=R^{d/p}\|f\|_p.
\]
同时 \(d\sigma_R(R\xi)=R^{d-1}d\sigma(\xi)\)，所以
\[
 \|\widehat f\|_{L^r(S_R)}
 =R^{(d-1)/r-d}\|\widehat g\|_{L^r(S)}
 \leq C R^{(d-1)/r-d/p'}\|f\|_p.
\]
\(r=\infty\) 时第一因子按 \(R^0\) 理解。缩放是 Schwartz 类的双射，
反向代入近乎达到最佳常数的输入，便知这个因子不只是一个上界。
\end{proof}

例如三维球面 Stein–Tomas 端点的 \(p=4/3,r=2\) 将产生
\(R^{1/4}\)，而不是与半径无关的常数。单位球面估计一旦成立，缩放便确定了其他半径所需的因子。若把面测度改为总质量一的概率测度，还须再乘
\(R^{-(d-1)/r}\)；使用不同测度归一化的公式时，应先对齐这一约定。

% !TeX root = ../../Book2.tex
\section{必要条件}
\label{b2:sec:2302}

限制估计不是单凭 Fourier 变换的有界性就能得到的。曲面上的质量可以集中在小帽，也可以遍布整个球面；前一种输入产生很长的相干管，后一种输入留下缓慢衰减的尾部。本节分别计算这些障碍。计算所得是不等式成立的必要条件，后面的充分估计还需要独立证明。

\subsection{缩放}
\label{b2:subsec:230201}

先直接研究曲面上的积分
\[
 Eg(x)\coloneqq\int_S e^{2\pi\mathrm i x\cdot\xi}
                     g(\xi)\,d\sigma(\xi),
 \quad \|Eg\|_{L^q(\mathbb R^d)}
                 \leq C\|g\|_{L^s(S)}.
\]
\symidx{\(Eg\)：曲面延拓算子}
其中 \(S\) 暂取单位球面，\(1\leq s,q\leq\infty\)。球面测度有限，故
\(g\in L^s(S)\) 总属于 \(L^1(S)\)，积分处处绝对收敛。
这里先把 \(E\) 当作一个明确的积分算子；它与限制估计的对偶等价将在
\cref{b2:prop:23-duality}中证明。对应的限制指数是
\(p=q'\)、\(r=s'\)。

\ProseParagraphBreak
设半径 \(R\) 的球面取原生面积测度，\(g_R\) 是其上的函数，并令
\(\widetilde g(\omega)=g_R(R\omega)\)。直接换元得
\[
 E_Rg_R(x)=R^{d-1}E\widetilde g(Rx),\quad
 \|g_R\|_{L^s(RS)}
       =R^{(d-1)/s}\|\widetilde g\|_{L^s(S)}.
\]
若单位球面的最佳常数为 \(C\)，半径 \(R\) 的最佳常数便为
\[
 C_R=C R^{\,d-1-d/q-(d-1)/s}.
\]
因为输入缩放可逆，这里也是准确的最佳常数关系。
它与\cref{b2:prop:23-surface-scaling}在共轭指数下完全一致。
固定单位球面并不要求这个幂等于零；只有要求一族半径上的常数统一时，才能由 \(R\to0\) 和 \(R\to\infty\) 进一步限制该幂。
本节真正限制固定球面指数的，是球面内部任意小帽的尺度。

\ProseParagraphBreak
还须区分两个空间的测度。曲面有限测度给出
\[
 \|g\|_{s_1}\leq \sigma(S)^{1/s_1-1/s_2}\|g\|_{s_2},
 \quad 1\leq s_1\leq s_2\leq\infty.
\]
环境空间的测度却是无穷大，不能据 \(q\) 的大小直接写出一般
\(L^q(\mathbb R^d)\) 之间的包含。不过延拓函数另有
\(\|Eg\|_\infty\leq\|g\|_1\)。如果已经证明某个输出指数 \(q_0<\infty\)
的界，那么
\[
 \|Eg\|_q\leq
 \|Eg\|_{q_0}^{q_0/q}
       \cdot\|Eg\|_\infty^{\,1-q_0/q},
 \quad q_0\leq q<\infty,
\]
才给出更大输出指数的界。这一步使用的是输出的额外有界性。

\subsection{Knapp 例子}
\label{b2:subsec:230202}

球面小帽在切向上的半径为 \(\delta\)，离开切平面的高度只有
\(O(\delta^2)\)。为了使不同频率的相位保持接近，环境变量的切向分量可以达到
\(\delta^{-1}\)，法向分量则可以达到 \(\delta^{-2}\)。
这就是下面的\term[knapp-lizi]{Knapp 例子}[Knapp example]所用的几何。

\Needspace{6\baselineskip}
\begin{proposition}[帽与管的必要条件]\label{b2:prop:23-knapp}
设 \(d\geq2\)，单位球面 \(S^{d-1}\) 取面积测度，
\(1\leq s,q\leq\infty\)。若其延拓积分对全部
\(g\in L^s(S^{d-1})\) 满足
\(\|Eg\|_q\leq C\|g\|_s\)，则必须有
\[
 \frac{d+1}{q}\leq\frac{d-1}{s'},
 \quad \frac1s+\frac1{s'}=1.
\]
\end{proposition}
\begin{proof}
在北极附近写
\(\xi=(\eta,h(\eta))\)，其中
\(h(\eta)=\sqrt{1-|\eta|^2}\)。取 \(0<\delta<1/4\)，令
\[
 \Theta_\delta=\{(\eta,h(\eta)):|\eta|\leq\delta\},
 \quad g_\delta=\mathbf1_{\Theta_\delta}.
\]
面积因子 \((1-|\eta|^2)^{-1/2}\) 在这些帽上有统一正上下界，所以
\[
 \sigma(\Theta_\delta)\asymp\delta^{d-1},\quad
 \|g_\delta\|_s\asymp\delta^{(d-1)/s}.
\]
当 \(s=\infty\) 时后一范数等于一，仍与所写幂一致。

由 \(h(\eta)-1=-|\eta|^2/2+O(|\eta|^4)\)，存在与 \(\delta\)
无关的 \(C_0\)，使 \(|h(\eta)-1|\leq C_0\delta^2\)。
取一个固定的小 \(c>0\)，在集合
\[
 \mathcal T_\delta=
 \{x:|x'|\leq c\delta^{-1},\ |x_d|\leq c\delta^{-2}\}
\]
上，有
\[
 2\pi|x'\cdot\eta+x_d(h(\eta)-1)|
       \leq 2\pi c(1+C_0).
\]
选 \(c\) 使右侧不超过 \(\pi/3\)。从积分中提出公共因子
\(e^{2\pi\mathrm i x_d}\)，余下被积函数的实部至少为 \(1/2\)，故
\[
 |Eg_\delta(x)|\geq\tfrac12\sigma(\Theta_\delta)
       \gtrsim\delta^{d-1}\quad(x\in\mathcal T_\delta).
\]
管的体积为常数乘 \(\delta^{-(d+1)}\)。因此，包括 \(q=\infty\)
的通常解释在内，
\[
 \|Eg_\delta\|_q\gtrsim
       \delta^{d-1-(d+1)/q}.
\]
假定的估计会要求
\(\delta^{d-1-(d+1)/q}\lesssim\delta^{(d-1)/s}\)。
令 \(\delta\downarrow0\)，左侧幂不能小于右侧幂，整理即得结论。
\end{proof}

同样的构造可以用非负光滑帽代替特征函数：令帽内部保留一个固定比例的核心，振幅在核心上有正下界即可。因此，必要条件并非由输入的不光滑造成。对任意光滑曲面的局部片，减去切平面后只需使用高度的 \(O(\delta^2)\) 上界，也可得到同一个必要条件；若曲率退化，还可能有更强的条件。

\ProseParagraphBreak
当 \(s=2\) 时，上式成为
\[
 q\geq q_*\coloneqq\frac{2(d+1)}{d-1}.
\]
本章将证明等号成立，因此平方曲面输入的最小输出指数已经由一个帽确定。
例如 \(d=3\) 时帽面积约为 \(\delta^2\)，管体积约为 \(\delta^{-4}\)，
比较 \(\delta^{2-4/q}\) 与 \(\delta\) 就得到 \(q\geq4\)。
若 \(s=1\)，必要条件要求 \(q=\infty\)，正好对应绝对积分界。

\cref{b2:fig:23-cap-tube}把两个小相位条件画成空间宽度与长度。两幅图分别属于频率空间和位置空间；帽的法向高度与管的长度不是同一空间中的两个长度。

\begin{figure}[htbp]
\centering
\begin{tikzpicture}[x=1cm,y=1cm]
\begin{scope}[shift={(-2.8,0)}]
\draw[->] (-1.5,0)--(1.5,0) node[right]{\(\eta\)};
\draw[->] (0,-.3)--(0,1.8) node[above]{法向};
\draw[AnalysisBlue,thick,domain=-1.3:1.3,samples=35] plot(\x,{1.3-.35*\x*\x});
\draw[AnalysisBlue,very thick,domain=-.5:.5,samples=15] plot(\x,{1.3-.35*\x*\x});
\draw[<->] (-.5,.7)--(.5,.7) node[midway,below]{\(\delta\)};
\draw[<->] (.7,1.2125)--(.7,1.3) node[right]{\(\delta^2\)};
\node at (0,-.6){频率帽：切向与法向高度};
\end{scope}
\draw[->] (-.7,.8)--(.6,.8);
\begin{scope}[shift={(2.8,0)}]
\fill[AnalysisBlue!15] (-.4,0) rectangle(.4,1.8);
\draw[AnalysisBlue,thick] (-.4,0) rectangle(.4,1.8);
\draw[<->] (-.4,-.2)--(.4,-.2) node[midway,below]{\(\delta^{-1}\)};
\draw[<->] (.65,0)--(.65,1.8) node[midway,right]{\(\delta^{-2}\)};
\node at (0,2.05){\(x_d\) 法向};
\node at (0,-.85){物理空间相干管};
\end{scope}
\end{tikzpicture}
\caption[频率帽与相干管]{Knapp 尺度的对应，不按比例。移去北极的常相位后，要让 \(x'\cdot\eta\) 与 \(x_dO(|\eta|^2)\) 都小，分别需要 \(|x'|\lesssim\delta^{-1}\) 与 \(|x_d|\lesssim\delta^{-2}\)。细长管的长轴沿帽中心的法向；帽的面积和管的体积仍由正文计算。}
\label{b2:fig:23-cap-tube}
\end{figure}

\subsection{球面尾部的可积性条件}
\label{b2:subsec:230203}

小帽检验只利用一个大而有限的管。要检查无穷远处的积分，取一个不会随尺度收缩的输入更合适。球面上的常数函数便提供了这样的检验。

\Needspace{6\baselineskip}
\begin{proposition}[常密度的严格尾部条件]\label{b2:prop:23-sphere-tail}
设 \(d\geq2\)，\(E\) 是单位球面面积测度的延拓算子。
对于 \(1\leq q<\infty\)，有
\[
 E1\in L^q(\mathbb R^d)
 \quad\Longleftrightarrow\quad
 q>\frac{2d}{d-1}.
\]
因而任何 \(L^s(S^{d-1})\to L^q(\mathbb R^d)\) 的有界性，
若 \(q<\infty\)，都必须满足这个严格不等式。
\end{proposition}
\begin{proof}
球面的旋转不变性使 \(E1(r\omega)\) 只依赖 \(r\)。
相位 \(\omega\cdot\xi\) 在球面上的驻点恰为 \(\omega\) 与 \(-\omega\)，
两处的切向 Hessian 都非退化。按\cref{b2:thm:21-stationary-phase}
分别取两个小驻点邻域，其余部分用非驻相分部积分，可得
\[
 E1(r\omega)=r^{-(d-1)/2}
       \bigl(c_+e^{2\pi\mathrm i r}+c_-e^{-2\pi\mathrm i r}\bigr)
       +O(r^{-(d+1)/2}),\quad r\geq1.
\]
系数 \(c_\pm\) 非零；由球面对称性，它们互为共轭，因而模相等。
旋转不变性也使余项常数不依赖 \(\omega\)。
这给出 \(|E1(r\omega)|\leq Cr^{-(d-1)/2}\)。
原积分在有界区域有界，所以当
\(d-1-q(d-1)/2<-1\) 时，极坐标积分收敛。

反过来，令
\(H(r)=c_+e^{2\pi\mathrm i r}+c_-e^{-2\pi\mathrm i r}\)。
这是非恒零的一周期连续函数。因此存在 \(\kappa>0\) 和
\(A\subset[0,1]\)，使 \(|A|>0\)，且 \(|H(t)|\geq2\kappa\)
对 \(t\in A\) 成立。对充分大的整数 \(n\)，余项在
\(n+A\) 上至多为 \(\kappa r^{-(d-1)/2}\)，从而
\[
 |E1(r\omega)|\geq\kappa r^{-(d-1)/2}
       \quad(r\in n+A).
\]
于是
\[
 \|E1\|_q^q\gtrsim
 \sum_{n\geq n_0} n^{\,d-1-q(d-1)/2}.
\]
这个级数只有在指数小于 \(-1\) 时才收敛。等号时出现调和级数，
也发散。常数函数属于球面上的每个 \(L^s\)，故得到最后的必要条件。
\end{proof}

这里不能把“驻相衰减上界”直接当作发散证明；真正排除端点的是每段径向区间内都有固定比例的下界集合。余弦式领先项确有零点，但这些零点无法消去其余径向区域的质量。

\ProseParagraphBreak
两种必要条件的作用不同：常密度给出严格的水平门槛
\(q>2d/(d-1)\)，小帽给出可含等号的关系
\((d+1)/q\leq(d-1)/s'\)。
在 \(s=2\) 时，小帽条件更强；若允许更大的 \(s\)，两个边界便可能分别起作用。
三维中它们是 \(q>3\) 和 \(4/q\leq2/s'\)。

\subsection{几何障碍}
\label{b2:subsec:230204}

曲率退化时，圆形帽未必反映最强的集中机制。考虑三维中的局部曲面
\[
 S=\{(u,v,u^2+v^4):|u|,|v|<a\},
\]
其中 \(a>0\) 固定，取原生面积测度。对小 \(\delta>0\)，取帽
\[
 |u|\leq\delta^{1/2},\quad |v|\leq\delta^{1/4}.
\]
其高度至多为 \(2\delta\)，面积约为 \(\delta^{3/4}\)。
令 \(g_\delta\) 为该帽的特征函数。若
\[
 |x_1|\leq c\delta^{-1/2},\quad
 |x_2|\leq c\delta^{-1/4},\quad
 |x_3|\leq c\delta^{-1},
\]
则相位 \(2\pi(x_1u+x_2v+x_3(u^2+v^4))\) 的绝对值至多为
\(8\pi c\)。取 \(c\) 足够小，仍可用实部正下界得到
\(|Eg_\delta|\gtrsim\delta^{3/4}\)。
这个长方管的体积约为 \(\delta^{-7/4}\)，所以有界性必须要求
\[
 \delta^{3/4-7/(4q)}
       \lesssim\delta^{3/(4s)},\quad
 \frac7q\leq\frac3{s'}.
\]
特别地，平方输入必须有 \(q\geq14/3\)，比非零曲率曲面的三维指数
\(4\) 更严格。这里曲率在 \(v=0\) 处消失，第四次方方向让帽比二次帽更宽，才产生了更大的相干体积。

\ProseParagraphBreak
这些检验也说明了必要条件的限度。它们都选取了一类特别的输入；通过这些检验，只能说明该类输入没有排除某个指数。任意输入还可能包含许多尺度、位置与方向不同的帽，叠加后的管束相互作用没有包含在单帽计算中。下一节建立对偶与 \(TT^*\) 表述，随后用解析插值证明平方输入的临界估计；更广的目标放在选读部分讨论。

% !TeX root = ../../Book2.tex
\section{延拓算子}
\label{b2:sec:2303}

限制算子的难点是在零测集上解释 Fourier 变换，而延拓算子从一开始就是绝对收敛的曲面积分。对偶性把两者联系起来；当曲面目标是平方范数时，再作一次复合，就把问题转为环境空间的卷积估计。本节逐步建立这两座桥，并说明插值实际能给出哪些指数。

\subsection{对偶限制}
\label{b2:subsec:230301}

以下配对统一采用
\(\langle u,v\rangle=\int u\overline v\)，对第一变量线性。
设 \(S\subset\mathbb R^d\) 紧致，取面积测度。对于
\(g\in L^1(S)\)，定义
\[
 Eg(x)=\int_S e^{2\pi\mathrm i x\cdot\xi}g(\xi)\,d\sigma(\xi).
\]
因为指数因子的模为一，
\(\|Eg\|_\infty\leq\|g\|_1\)；支配收敛还给出 \(Eg\) 连续。
这些事实不需要曲率条件。

\Needspace{6\baselineskip}
\begin{proposition}[限制与延拓的对偶]\label{b2:prop:23-duality}
设 \(d\geq2\)，\(S\subset\mathbb R^d\) 为非空紧致光滑超曲面，
取面积测度 \(\sigma\)。令
\(Rf=\widehat f|_S\) 先定义在 \(\mathcal S(\mathbb R^d)\) 上，
\(Eg(x)=\int_S e^{2\pi\mathrm i x\cdot\xi}g(\xi)\,d\sigma(\xi)\)。
对于 \(1\leq p,r<\infty\)，
下列两个估计等价，且最佳常数相同：
\[
 \|Rf\|_{L^r(S)}\leq C\|f\|_{L^p(\mathbb R^d)}
       \quad(f\in\mathcal S),
\]
\[
 \|Eg\|_{L^{p'}(\mathbb R^d)}
       \leq C\|g\|_{L^{r'}(S)}
       \quad(g\in L^{r'}(S)).
\]
其中 \(p'\)、\(r'\) 为共轭指数；第一式也等价于 \(R\) 的唯一有界
延拓 \(L^p(\mathbb R^d)\to L^r(S)\)。
\end{proposition}
\begin{proof}
由于 \(\sigma(S)<\infty\)，每个 \(g\in L^{r'}(S)\) 都属于
\(L^1(S)\)，包括 \(r'= \infty\)。对 \(f\in\mathcal S\)，
二重积分的绝对值由 \(\|f\|_1\cdot\|g\|_1\) 控制，
所以 Fubini 给出
\[
 \langle Rf,g\rangle_{S}
       =\langle f,Eg\rangle_{\mathbb R^d}.
\]
若第一式成立，则
\[
 |\langle f,Eg\rangle|
       \leq C\|f\|_p\cdot\|g\|_{r'}.
\]
先把这个配对看成 Schwartz 类上的有界泛函，再由其在有限指数
\(L^p\) 中的稠密性延拓。\(L^p\) 对偶表示说明表示函数属于
\(L^{p'}\)，范数至多为 \(C\|g\|_{r'}\)；\(p=1\) 时使用
\((L^1)^*=L^\infty\)。表示函数与原来的有界连续积分 \(Eg\)
在全部测试函数上配对相同，因而几乎处处相等。这便证明第二式。

反过来，由第二式和 Hölder 不等式，
\[
 |\langle Rf,g\rangle|
       \leq C\|f\|_p\cdot\|g\|_{r'}.
\]
对 \(L^{r'}(S)\) 单位球取上确界，利用 \(L^r\) 范数的对偶公式，
得到第一式。对 \(r=1\)，这个上确界取遍 \(L^\infty(S)\) 单位球，
仍是通常的 \(L^1\) 范数公式。两次推导没有增大常数，所以最佳常数相同。
最后由\subsecref{b2:subsec:230101}的稠密延拓论证，得到 \(R\)
在全部 \(L^p\) 上的定义。
\end{proof}

光滑函数在 \(L^\infty(S)\) 中一般并不按范数稠密。对于这个无穷对偶指数，有限曲面测度保证延拓积分对全部输入有意义，因而可以直接使用配对与对偶范数。后面若把 \(E\) 写成 \(R^*\)，指的正是这个复配对下的伴随关系。

\subsection{Fourier 衰减}
\label{b2:subsec:230302}

若 \(S\) 的 Gaussian 曲率在紧支撑片上处处非零，
\cref{b2:prop:21-curved-surface-decay}给出光滑曲面密度的衰减：
\[
 \left|\int_S e^{2\pi\mathrm i x\cdot\xi}a(\xi)\,d\sigma(\xi)\right|
       \leq C_a(1+|x|)^{-(d-1)/2},
 \quad a\in C_c^\infty(S).
\]
所用条件包括非退化切向 Hessian、紧支撑和有限阶振幅导数的控制。
常数 \(C_a\) 因而包含 \(a\) 的导数，不能只换成 \(C\|a\|_2\)。
若作这种替换，会把对每个固定光滑输入的衰减，误当成整个
\(L^2(S)\) 单位球上的统一算子界。

\ProseParagraphBreak
平移对称性提供了一个直接反证。设 \(a\in\mathbb R^d\)，则
\[
 E(e^{-2\pi\mathrm i a\cdot\xi}g)(x)=Eg(x-a),
 \quad
 \|e^{-2\pi\mathrm i a\cdot\xi}g\|_2=\|g\|_2.
\]
取 \(g=1\)，在 \(x=a\) 处输出恒为 \(\sigma(S)>0\)。
如果有只由输入平方范数控制的界
\[
 |Eh(x)|\leq C\|h\|_2(1+|x|)^{-\beta},
 \quad \beta>0,
\]
将 \(h=e^{-2\pi\mathrm i a\cdot\xi}\)、\(x=a\) 代入并令
\(|a|\to\infty\)，便矛盾。问题不是这些输入缺少光滑性；它们全都光滑，
只是导数随调制参数增大。

\ProseParagraphBreak
这不妨碍证明平移不变的 \(L^q\) 估计：平移峰值不会改变环境范数。
同时，对每个固定 \(g\in L^1(S)\)，在上述曲率条件下仍有
\(Eg(x)\to0\) 当 \(|x|\to\infty\)。事实上，取光滑 \(g_n\) 在
\(L^1(S)\) 中逼近 \(g\)，则
\(\|Eg-Eg_n\|_\infty\leq\|g-g_n\|_1\)；
每个 \(Eg_n\) 都趋于零，故其一致极限也趋于零。
这个论证没有给出仅依赖 \(\|g\|_1\) 的统一衰减速度。

\subsection{\texorpdfstring{\(L^p\to L^q\)}{Lp→Lq} 估计}
\label{b2:subsec:230303}

本小节说明一个端点估计一旦成立，会自动产生怎样的范围。
下一节将用解析族方法建立这里所需的端点估计。
设 \(S\) 是有限面积的紧曲面，并已知
\[
 \|Eg\|_{q_*}\leq C_*\|g\|_2,\quad
 q_*=\frac{2(d+1)}{d-1}.
\]
另一个始终可用的端点是 \(L^1(S)\to L^\infty(\mathbb R^d)\)，常数为一。
两个算子都是同一个曲面积分，因而在输入空间交集上一致。
Riesz--Thorin 插值给出
\[
 \frac1s=1-\frac\theta2,\quad
 \frac1q=\frac\theta{q_*},\quad
 \|Eg\|_q\leq C_*^\theta\|g\|_s,
 \quad 0\leq\theta\leq1.
\]
消去 \(\theta\)，得到
\[
 \frac1q=\frac{d-1}{d+1}\left(1-\frac1s\right),
 \quad 1\leq s\leq2.
\]
这恰是这一段输入指数对应的 Knapp 边界。对固定 \(s\)，再用
\(\|Eg\|_\infty\leq\sigma(S)^{1-1/s}\|g\|_s\)
对输出范数插值，可把 \(q\) 增大。若 \(s>2\)，使用有限测度包含
\[
 \|g\|_2\leq\sigma(S)^{1/2-1/s}\|g\|_s
\]
则得到所有 \(q\geq q_*\) 的估计。这里不能把原来的公式延长到
\(\theta>1\)：插值定理只控制两端之间的参数。

\ProseParagraphBreak
例如三维球面有 \(q_*=4\)。由平方端点和绝对积分界，可以推出
\(L^{4/3}(S^2)\to L^8(\mathbb R^3)\)，因为取 \(\theta=1/2\)。
对 \(L^4(S^2)\) 输入，直接包含只能给出 \(L^4(\mathbb R^3)\)
及更大的输出指数。必要条件却仅要求 \(q>3\) 且 \(q\geq8/3\)，
所以 \(3<q<4\) 仍没有由这次插值解决。
这个例子显示了已证范围与必要范围之间的实际空缺。

\ProseParagraphBreak
在限制形式中，平方端点是
\[
 p_*=q_*'=\frac{2(d+1)}{d+3},\quad
 R:L^{p_*}(\mathbb R^d)\longrightarrow L^2(S).
\]
与 \(R:L^1(\mathbb R^d)\to L^\infty(S)\) 插值，得
\[
 \frac1p=1-\frac\theta{q_*},\quad
 \frac1r=\frac\theta2,\quad 0\leq\theta\leq1.
\]
这些恰是前述延拓指数的共轭关系。曲面目标可以借有限测度降低指数
\(r\)，但环境输入不能仅凭指数大小随意改变。
例如把某个 \(L^{p_*}\) 输入换成 \(L^p\)、\(p>p_*\)，并不能由
无限测度空间上的包含关系完成。

\subsection{\texorpdfstring{\(TT^*\)}{TT*} 表述}
\label{b2:subsec:230304}

当曲面范数取 \(L^2\) 时，可以把限制算子与延拓算子相接。
令
\[
 K(x)=\check\sigma(x)=\int_S e^{2\pi\mathrm i x\cdot\xi}\,d\sigma(\xi),
 \quad Af=f*K.
\]
核 \(K\) 有界，故 \(Af\) 对 Schwartz 输入绝对收敛。
下面的等价关系还适用于任意有限正紧支撑测度；测度的正性正是平方范数能够出现的原因。

\Needspace{6\baselineskip}
\begin{proposition}[平方限制与卷积]\label{b2:prop:23-ttstar}
设 \(\mu\) 是 \(\mathbb R^d\) 上的有限正紧支撑 Borel 测度，
\(1\leq p<\infty\)，并置
\[
 R_\mu f=\widehat f\quad\text{作为 }L^2(\mu)\text{ 中的函数},\quad
 E_\mu g(x)=\int e^{2\pi\mathrm i x\cdot\xi}g(\xi)\,d\mu(\xi).
\]
令 \(K=\check\mu\)，\(Af=f*K\) 先定义于 \(\mathcal S(\mathbb R^d)\)。
以下三个有界性等价：
\[
 R_\mu:L^p\to L^2(\mu),\quad
 E_\mu:L^2(\mu)\to L^{p'},\quad
 A:L^p\to L^{p'}.
\]
其中第一、第三项指从 Schwartz 类唯一连续延拓。
若这些界成立，则有准确的范数关系
\[
 \|A\|_{L^p\to L^{p'}}
       =\|R_\mu\|_{L^p\to L^2(\mu)}^2
       =\|E_\mu\|_{L^2(\mu)\to L^{p'}}^2.
\]
对 Schwartz 函数还恒有
\[
 Af=E_\mu R_\mu f,\quad
 \langle Af,f\rangle=\|R_\mu f\|_{L^2(\mu)}^2\geq0.
\]
\end{proposition}
\begin{proof}
对偶等价的证明与\cref{b2:prop:23-duality}完全相同：
有限性使 \(L^2(\mu)\subset L^1(\mu)\)，Fubini 给出
\(\langle R_\mu f,g\rangle_\mu=\langle f,E_\mu g\rangle\)，
两侧取对偶范数即可，且两算子范数相等。

对 \(f\in\mathcal S\)，再用 Fubini，
\[
 \begin{aligned}
 E_\mu R_\mu f(x)
 &=\int e^{2\pi\mathrm i x\cdot\xi}
       \left(\int f(y)e^{-2\pi\mathrm i y\cdot\xi}\,dy\right)d\mu(\xi)\\
 &=\int f(y)K(x-y)\,dy=Af(x).
 \end{aligned}
\]
绝对可积性由 \(\mu(\mathbb R^d)\|f\|_1\) 保证。将它再与 \(f\)
配对，二重空间积分由 \(\mu(\mathbb R^d)\|f\|_1^2\) 控制，得到
\[
 \langle Af,f\rangle
 =\int|\widehat f(\xi)|^2\,d\mu(\xi).
\]
若 \(A\) 有界，Hölder 立即给出
\[
 \|R_\mu f\|_2^2
       \leq\|Af\|_{p'}\cdot\|f\|_p
       \leq\|A\|\|f\|_p^2,
\]
所以 \(R_\mu\) 有界且 \(\|R_\mu\|^2\leq\|A\|\)。
反之，若 \(R_\mu\) 有界，其伴随 \(E_\mu\) 也有界，复合算子
\(E_\mu R_\mu\) 在测试函数上等于 \(A\)，故
\(\|A\|\leq\|E_\mu\|\cdot\|R_\mu\|=\|R_\mu\|^2\)。
由稠密性完成延拓，并得到范数等式。
所有步骤在 \(p=1,p'=\infty\) 时仍成立；这里只要求输入端的稠密性。
\end{proof}

通常把这称为 \(TT^*\) 方法：若把延拓算子记作 \(T\)，则
\(T^*=R_\mu\)、\(TT^*=A\)。本节另用 \(A\) 记卷积算子，是为了分清三个映射的定义域。二次型非负并不说 \(K(x)\geq0\)，也不说 \(A\) 把非负函数映成非负函数；球面核本身就会振荡变号。

\ProseParagraphBreak
先看只保留核绝对值能走多远。在非零曲率的紧光滑片上，
\[
 |K(x)|\leq C(1+|x|)^{-\beta}\leq C|x|^{-\beta},
 \quad \beta=\frac{d-1}{2}.
\]
右侧是阶
\(\alpha=d-\beta=(d+1)/2\) 的分数积分核，满足 \(0<\alpha<d\)。
由\cref{b2:thm:19-hls}，要取得 \(A:L^p\to L^{p'}\)，须令
\[
 \frac1p-\frac1{p'}=\frac{\alpha}{d}
       =\frac{d+1}{2d}.
\]
解得
\[
 p=\frac{4d}{3d+1},\quad p'=\frac{4d}{d-1}.
\]
此处 \(1<p<d/\alpha\)，符合 HLS 的强型条件。因此
\(TT^*\) 已给出 \(E:L^2\to L^{4d/(d-1)}\)。
但
\[
 \frac{4d}{d-1}>
       \frac{2(d+1)}{d-1}\quad(d\geq2),
\]
所以它比所求临界输出指数大。三维中，这个办法只能到 \(q=6\)，
而帽例子允许的端点是 \(q=4\)。
下一节把核的振荡与 Fourier 侧的曲面结构放在解析族的两条边界上，
正是为了补回取绝对值时丢掉的信息。

\ProseParagraphBreak
最后说明平方范数的位置。对球面上的 \(g\in L^2(S^{d-1})\)，
有限测度使 \(g\in L^1\)。频率支集紧致，允许逐阶对积分求导，且
\[
 \partial^\gamma Eg(x)
       =\int_S(2\pi\mathrm i\xi)^\gamma
                  e^{2\pi\mathrm i x\cdot\xi}g(\xi)\,d\sigma(\xi).
\]
各阶导数都有只依赖该阶数和 \(\|g\|_1\) 的上界，特别地
\[
 (-\Delta-4\pi^2)Eg=0.
\]
这是光滑的 Helmholtz 解，其频率集中在单位球面。
光滑性来自紧频率支集，全空间的可积性则是另一件事。

\ProseParagraphBreak
事实上，非零 \(g\) 的这个输出不属于 \(L^2(\mathbb R^d)\)。
若 \(Eg\in L^2\)，Plancherel 说明其 Fourier 变换由一个
\(L^2\) 函数表示。另一方面，对测试函数使用 Fubini 可知该变换就是测度
\(g\,d\sigma\)，支于 Lebesgue 零测的球面。一个绝对连续的 \(L^2\)
分布若支于零测集只能为零，所以 \(g=0\)，矛盾。
若 \(1\leq q<2\) 且 \(Eg\in L^q\)，再利用有界性，
\[
 \int|Eg|^2\leq\|Eg\|_\infty^{\,2-q}\int|Eg|^q<\infty,
\]
也得同一矛盾。因而 \(TT^*\) 中的平方空间位于曲面数据端；
环境输出使用较大的指数。将这两种测度下的平方空间混在一起，
就会失去限制理论的核心区别。

% !TeX root = ../../Book2.tex
\section{Stein–Tomas 定理}
\label{b2:sec:2304}

前一节已把限制估计化为卷积算子的正二次型。本节把球面测度放进一个复参数族：一条边界线上，频率符号有界，Plancherel 控制平方范数；另一条边界线上，空间核有界，得到 \(L^1\to L^\infty\) 界。插值点恰好恢复球面测度。要完成这条路线，必须同时说明分布怎样随参数解析变化、两条竖边上的常数怎样增长，以及变化测试函数的配对为何合法。最后把径向构造改为局部法向构造，得到一般非零曲率紧片的同一端点。

\subsection{球面测度的 Fourier 衰减}
\label{b2:subsec:230401}

由\cref{b2:thm:A2-sphere-fourier}，采用\cref{b2:def:A2-bessel-j}的 Bessel 函数归一化，有
\[
 \check\sigma(x)=2\pi |x|^{1-d/2}J_{d/2-1}(2\pi|x|).
\]
原点处按连续极限解释。由\cref{b2:thm:A2-bessel-asymptotic}的实阶渐近或\chapref{b2:ch:21}直接驻相，得到
\(|\check\sigma(x)|\lesssim(1+|x|)^{-(d-1)/2}\)。
这个指数取决于曲面维数，而不是环境维数的一半；球面有 \(d-1\) 个切向振荡变量。

\ProseParagraphBreak
在小 \(r\) 时，\(J_\nu(r)\sim(r/2)^\nu/\Gamma(\nu+1)\)，故上式有有限极限，等于球面面积。大 \(r\) 时则有两个相反法向驻点的贡献。将小尺度与大尺度分开，是后面的解析族核界能够在整个空间成立的原因。

\subsection{Gamma 归一化与球面解析族}
\label{b2:subsec:230402}

我们的目标是把球面测度放进一个可插值的参数族。球内的普通幂 \((1-|\xi|^2)_+^z\) 在 \(\operatorname{Re}z>-1\) 时可积，却不能直接取目标参数 \(z=-1\)；除以适当的 Gamma 因子，将使这个边界值成为球面 delta，而不是发散的函数积分。下面先确定这个因子，再说明分布延拓怎样得到目标值。

\ProseParagraphBreak
先恢复本处所需的复分析用语。函数在一个复开域中全纯，是指处处存在复导数；整函数是在整个复平面全纯的函数。这里使用基础复分析的局部 Cauchy 公式：若闭圆盘及其邻域都在全纯域中，则盘内有
\[
 f(w)=\frac1{2\pi\mathrm i}\int_{|\zeta-w_0|=r}
          \frac{f(\zeta)}{\zeta-w}\,\mathrm d\zeta,
 \quad |w-w_0|<r,
\]
圆周取正向。它的两个局部推论供下文使用：全纯函数的局部一致极限仍全纯，因为圆周上一致收敛可通过该积分，而极限式可对 \(w\) 求导；全纯函数有局部幂级数，故非恒零函数的零点孤立，从而连通域中在一个开子域相等的两个全纯函数处处相等。后面将在紧含于参数域的小圆盘上应用 Cauchy 公式，用这些局部性质处理无限乘积和分布的解析延拓。

\ProseParagraphBreak
前文的正实 Gamma 积分在右半平面有同一个复参数定义
\[
 \Gamma(w)=\int_0^\infty e^{-t}t^{w-1}\,\mathrm dt,
 \quad \operatorname{Re}w>0,
 \quad t^{w-1}=e^{(w-1)\log t},
\]
其中 \(\log t\) 为实对数。若参数紧集满足 \(0<a\leq\operatorname{Re}w\leq b\)，积分及任意固定阶参数导数可由 \(e^{-t}(t^{a-1}+t^{b-1})(1+|\log t|^k)\) 控制，故积分全纯，正实轴上的值与前文一致。分部积分仍给出 \(\Gamma(w+1)=w\Gamma(w)\)。为跨过这个半平面的边界，使用倒数而不是在奇点处逐个给 Gamma 赋值。

\ProseParagraphBreak
记 \(\gamma=\lim_{N\to\infty}(\sum_{n=1}^Nn^{-1}-\log N)\)；积分比较表明该序列递减且有下界，所以极限存在。对任意 \(w\in\mathbb C\)，我们需要一个在整个复平面全纯的函数，即整函数，其零点将抵消负参数的极点。定义
\[
 P(w)=w e^{\gamma w}
      \prod_{n=1}^\infty(1+w/n)e^{-w/n}.
\]
单独的 \(1+w/n\) 因子含有不可和的线性项；乘上 \(e^{-w/n}\) 正是为了消掉它。在每个紧参数集上，对充分大的 \(n\) 有 \(|w/n|<1/2\)，可用 \(\log(1+\zeta)\) 在零点的幂级数，并且
\(\log(1+w/n)-w/n=O(n^{-2})\) 一致成立。尾部对数绝对一致可和，前面有限因子保留原式，故乘积局部一致收敛为整函数。避开有限头部的零点后，尾部是一个全纯函数的指数，不会为零；所以 \(P\) 的零点恰为 \(0,-1,-2,\ldots\)，每一个都是单零点。

\ProseParagraphBreak
这个乘积确实对应原来的 Gamma 积分。对 \(\operatorname{Re}w>0\)，有限次分部积分给出
\[
 \int_0^1t^{w-1}(1-t)^N\,\mathrm dt
 =\frac{N!}{w(w+1)\cdots(w+N)}.
\]
乘以 \(N^w\) 再令 \(u=Nt\)，左侧趋于 \(\Gamma(w)\)：在 \(0<u<N\) 上有 \((1-u/N)^N\leq e^{-u}\)，紧参数集可用上述不含对数的共同可积控制，先截断积分再控制尾部，得到局部一致收敛。同时，右侧的倒数可写成
\[
 \frac{w(w+1)\cdots(w+N)}{N!N^w}
 =w e^{w(H_N-\log N)}
       \prod_{n=1}^N(1+w/n)e^{-w/n},
 \quad H_N=\sum_{n=1}^Nn^{-1},
\]
并局部一致趋于 \(P(w)\)。有限两式相乘为一，取极限得到 \(P(w)\Gamma(w)=1\)。因此 \(1/P\) 给出 Gamma 的亚纯延拓：它在零点以外全纯，在每个零点附近则为一个全纯非零因子除以 \(w-w_0\)，即有简单极点。它没有零点，极点恰在非正整数处；其倒数 \(1/\Gamma=P\) 在这些位置仍是良定义的整函数。递推先在右半平面成立；两侧 \(P(w)\) 与 \(wP(w+1)\) 都整，由刚说明的唯一性推论，在全平面相等，包括零点处的等式。

\ProseParagraphBreak
后面还要对虚部增长作统一估计。若 \(u\) 落在固定正区间 \([\alpha,\beta]\subset(0,\infty)\)，取无限乘积的模得到
\[
 \frac{|P(u+\mathrm iy)|}{P(u)}
 =\left(1+\frac{y^2}{u^2}\right)^{1/2}
 \prod_{n=1}^{\infty}
 \left(1+\frac{y^2}{(n+u)^2}\right)^{1/2}.
\]
使用 \(\log(1+t)\leq t\)、\(t\geq0\)，以及 \(\sum n^{-2}<\infty\)，得到 \(|P(u+\mathrm iy)|\leq C e^{C(1+y^2)}\)，其中常数对 \(u\in[\alpha,\beta]\) 一致。若实部落在任意固定有限区间，先用
\[
 P(w)=w(w+1)\cdots(w+N-1)\cdot P(w+N)
\]
把它移到正半平面；有限个多项式因子仍可由 \(e^{C(1+y^2)}\) 控制。因此对任意实数 \(c_0<c_1\)，均有
\[
 |P(u+\mathrm iy)|\leq C_{c_0,c_1}e^{C_{c_0,c_1}(1+y^2)},
 \quad c_0\leq u\leq c_1.
\]
我们只需要这种较粗的 Gaussian 型增长，不需要复 Gamma 函数的精确渐近式。

\ProseParagraphBreak
现在把归一化用于分布。未归一化的 \(t_+^z\) 在负整数参数处出现极点，Gamma 倒数的单零点恰好将它们消掉。下面用分布导数给出实际定义；随后再用 Taylor 减法解释这些极点的来源。

\Needspace{6\baselineskip}
\begin{lemma}[法向幂的解析分布族]\label{b2:lem:23-normal-analytic-family}
对 \(z\in\mathbb C\)、\(\operatorname{Re}z>-1\)，在实轴上令
\[
 U_z(t)=\frac{t_+^z}{\Gamma(z+1)},
 \quad
 t_+^z=
 \begin{cases}
  e^{z\log t},&t>0,\\
  0,&t<0,
 \end{cases}
\]
其中 \(\log t\) 是实对数，\(1/\Gamma=P\) 是上述整函数。这个局部可积函数族唯一延拓为支撑于 \([0,\infty)\) 的整分布族，并满足
\[
 \partial_tU_{z+1}=U_z,\quad
 U_0=\mathbf1_{(0,\infty)},\quad
 U_{-1}=\delta_0,\quad U_{-2}=\delta_0'.
\]
对任意固定实部区间 \([c_0,c_1]\) 和紧区间 \(I\subset\mathbb R\)，存在整数 \(N\geq0\) 与常数 \(C\)，使所有支撑于 \(I\) 的 \(\phi\in C_c^\infty(\mathbb R;\mathbb C)\) 满足
\[
 |\langle U_{u+\mathrm iy},\phi\rangle|
 \leq C e^{C(1+y^2)}
 \max_{0\leq k\leq N}\|\phi^{(k)}\|_\infty,
 \quad c_0\leq u\leq c_1.
\]
整分布族的含义是：对每个测试函数 \(\phi\)，配对 \(z\mapsto\langle U_z,\phi\rangle\) 在整个复平面全纯。
\end{lemma}
\begin{proof}
在 \(\operatorname{Re}z>-1\) 中，Gamma 递推和实变量分部积分给出
\(\partial_tU_{z+1}=U_z\)。在零点没有边界项，因为 \(t^{z+1}\to0\)；在无穷远没有边界项，因为测试函数紧支撑。对任意 \(z\)，选整数 \(N\geq0\)，使 \(\operatorname{Re}z+N>-1\)，定义
\[
 \langle U_z,\phi\rangle
 =(-1)^N P(z+N+1)
 \int_0^\infty t^{z+N}\cdot\phi^{(N)}(t)\,\mathrm dt.
\]
积分绝对收敛。再多取一次导数，利用已经建立的递推与分部积分，所得配对不变；因此定义不依赖 \(N\)。在任意参数紧集上可固定同一个 \(N\)，积分及其参数导数由可积的幂乘有限次 \(|\log t|\) 控制，所以每个配对全纯。不同局部定义在重叠区域一致，给出整个平面的延拓。另一整延拓若在原半平面相同，其每个测试配对由全纯唯一性处处相同，因此延拓唯一。

对给定实部区间，取 \(N\) 使 \(c_0+N>-1\)。由于 \(\phi\) 支撑于固定有界区间，积分绝对值由 \(C\|\phi^{(N)}\|_\infty\) 控制，常数不依赖虚部。前面的 \(P\) 增长估计于是给出所述有限阶分布界。测试函数若支撑在负半轴，积分为零，故支撑性质也保留。

最后 \(P(1)=1\)，所以 \(U_0=\mathbf1_{(0,\infty)}\)。对测试函数，
\[
 \langle\partial_tU_0,\phi\rangle
 =-\int_0^\infty\phi'(t)\,\mathrm dt=\phi(0).
\]
因此 \(U_{-1}=\delta_0\)，再求一次导数得到 \(U_{-2}=\delta_0'\)。
\end{proof}

Taylor 减法给出同一延拓的另一种读法。对 \(F\in C_c^\infty(\mathbb R)\)，原来收敛的配对可写成
\[
 \begin{aligned}
 \langle U_z,F\rangle
 =P(z+1)\biggl[
 &\int_0^1 t^z\cdot
 \left(F(t)-\sum_{k=0}^{N-1}\frac{F^{(k)}(0)}{k!}t^k\right)\,\mathrm dt\\
 &+\sum_{k=0}^{N-1}\frac{F^{(k)}(0)}{k!(z+k+1)}
 +\int_1^\infty t^z\cdot F(t)\,\mathrm dt
 \biggr].
 \end{aligned}
\]
减去 \(N-1\) 阶 Taylor 多项式后，余项为 \(O(t^N)\)，故第一项在 \(\operatorname{Re}z>-N-1\) 时可积。各个显式分式的简单极点被 \(P(z+1)\) 的单零点消去；最后一项在全平面解析。这与引理的定义在原半平面一致，因而给出同一族。在 \(z=-1\) 时只有 \(k=0\) 项留下 \(F(0)\)，正是球面测度将出现的原因。

\ProseParagraphBreak
取整数 \(d\geq2\)，在频率空间中置
\[
 G_z(\xi)=\frac{(1-|\xi|^2)_+^z}{\Gamma(z+1)},
 \quad T_zf=\mathcal F^{-1}(G_z\cdot\widehat f),
 \quad f\in\mathcal S(\mathbb R^d).
\]
当 \(\operatorname{Re}z>-1\) 时，\(G_z\) 是支撑于闭单位球的可积函数。跨过这条线时，只在球面附近使用法向坐标 \(t=1-|\xi|^2\)，将 \(U_z\) 与光滑测度因子配对；在球心附近，\(1-|\xi|^2\) 严格为正，原来的普通幂已经全纯，不需要在梯度为零处作法向换元。用固定分割把这两部分相加，便得到支撑于同一闭球的整分布族，且在任意固定实部区间上具有引理所给的有限阶、Gaussian 型增长。

\ProseParagraphBreak
具体地，在球面附近，极坐标换成 \(t=1-r^2\) 后的测度为
\[
 r^{d-1}\,\mathrm dr\,\mathrm d\sigma(\omega)
 =\frac12(1-t)^{(d-2)/2}\,\mathrm dt\,\mathrm d\sigma(\omega),
\]
其中积分方向已调正。因为 \(U_{-1}=\delta_0\)，得到
\[
 G_{-1}=\frac12\,\mathrm d\sigma,\quad
 T_{-1}f=\frac12(f*\check\sigma).
\]
这个二分之一来自法向函数的梯度长度 \(|\nabla(1-|\xi|^2)|=2\)，不能在讨论归一化时遗漏。右竖边 \(z=\mathrm iy\) 上，符号是球内模为 \(|P(1+\mathrm iy)|\) 的普通函数，所以 Plancherel 给出
\[
 \|T_{\mathrm iy}\|_{2\to2}
 \leq |P(1+\mathrm iy)|
 \leq C e^{C(1+y^2)}.
\]

\ProseParagraphBreak
为了找到另一条竖边，先把\appref{b2:app:A}中的 Bessel 级数扩展到复阶。对 \(\nu\in\mathbb C\)、\(r>0\)，仍令
\[
 J_\nu(r)=\sum_{k=0}^\infty
 \frac{(-1)^kP(k+\nu+1)}{k!}
 \left(\frac r2\right)^{2k+\nu},
\]
复幂使用正实数 \(r/2\) 的实对数。对紧参数集中的 \(\nu\)，Gamma 递推从某个固定大 \(k\) 起给出阶乘型分母界，因此级数在紧 \(r\) 区间和紧 \(\nu\) 集上局部一致收敛，固定次数的 \(r\) 导数也一样。逐项求导仍验证 Bessel 方程。这里扩展的是级数与方程；\cref{b2:thm:A2-bessel-asymptotic}的固定实阶渐近并未因此自动变成复阶统一渐近。

\ProseParagraphBreak
由\cref{b2:thm:A2-bochner-riesz}，在实参数 \(z>-1\) 时有
\[
 \check G_z(x)
 =\pi^{-z}|x|^{-d/2-z}J_{d/2+z}(2\pi|x|).
\]
对每个固定 \(x\)，左侧可将指数函数乘一个在闭球附近等于一的光滑截断，再与紧支撑分布 \(G_z\) 配对，所以全纯依赖 \(z\)。右侧在 \(r=|x|\) 处可写成
\[
 \pi^{d/2}\sum_{k=0}^\infty
 \frac{(-\pi^2r^2)^k}{k!}\cdot P(k+d/2+z+1),
\]
同样为整函数，并且在 \(r=0\) 时也有意义。两侧在实区间 \(z>-1\) 一致，全纯唯一性遂把恒等式延拓到所有 \(z\in\mathbb C\)。这种论证同时说明原点处应取级数值
\[
 \check G_z(0)=\pi^{d/2}P(d/2+z+1).
\]

\ProseParagraphBreak
令 \(b=(d+1)/2\)。在左竖边 \(z=-b+\mathrm iy\) 上，Bessel 阶数成为
\(\nu=-1/2+\mathrm iy\)。我们要证明 \(\check G_z\) 在全部空间有界，且其界至多以 \(e^{C(1+y^2)}\) 增长。先对 \(0<r\leq1\) 使用归一化级数
\[
 r^{-\nu}J_\nu(2\pi r)
 =\pi^\nu\sum_{k=0}^\infty
 \frac{(-\pi^2r^2)^k}{k!\Gamma(k+\nu+1)}.
\]
Gamma 递推及 \(|\ell+1/2+\mathrm iy|\geq\ell+1/2\) 给出
\[
 \left|\frac1{\Gamma(k+1/2+\mathrm iy)}\right|
 \leq
 \frac{|P(1/2+\mathrm iy)|}
      {\prod_{\ell=0}^{k-1}(\ell+1/2)}.
\]
空乘积取一。右侧代入级数后可一致求和，固定次数的 \(r\) 导数也可求和。因此 \(r^{-\nu}J_\nu(2\pi r)\) 在 \(0\leq r\leq1\) 有界，\(J_\nu(2\pi)\) 及对 \(r\) 的导数在 \(r=1\) 的值也受 \(C e^{C(1+y^2)}\) 控制；微分中出现的 \(1+|\nu|\) 因子可并入同样的增长界。

\ProseParagraphBreak
大半径不再对级数取绝对值，而使用方程。令
\[
 w(r)=r^{1/2}J_\nu(2\pi r),\quad
 E(r)=|w'(r)|^2+4\pi^2|w(r)|^2.
\]
Bessel 方程化为
\[
 w''+4\pi^2w=(\nu^2-1/4)r^{-2}w.
\]
即使 \(\nu\) 为复数，\(E\) 仍是非负的实能量；微分后自由振荡项抵消，故
\[
 \begin{aligned}
 |E'(r)|
 &=\left|2\operatorname{Re}
   \bigl((\nu^2-1/4)r^{-2}w\overline{w'}\bigr)\right|\\
 &\leq C|\nu^2-1/4|r^{-2}\cdot E(r)
 \leq C(1+y^2)r^{-2}\cdot E(r).
 \end{aligned}
\]
把这一微分不等式积分，或直接使用 Gronwall 不等式，并注意 \(\int_1^\infty r^{-2}\,\mathrm dr=1\)，得到
\[
 E(r)\leq E(1)e^{C(1+y^2)}
 \leq C e^{C(1+y^2)},\quad r\geq1.
\]
所以 \(|J_\nu(2\pi r)|\leq C e^{C(1+y^2)}r^{-1/2}\)。在 \(\check G_{-b+\mathrm iy}\) 的公式中，另一幂因子的模恰为 \(\pi^b r^{1/2}\)，两者抵消。连同小半径估计，最终得到
\[
 \|\check G_{-b+\mathrm iy}\|_\infty
 \leq C_d e^{C_d(1+y^2)},\quad
 \|T_{-b+\mathrm iy}\|_{1\to\infty}
 \leq C_d e^{C_d(1+y^2)}.
\]
左竖边的位置由此确定：归一化核的增长和 Bessel 振荡的半阶衰减在这里恰好平衡。

\subsection{Stein–Tomas 定理}
\label{b2:subsec:230403}

两条边界估计已经齐备，但还需要一个定义域上的桥梁。\chapref{b2:ch:17}的插值证明使用随参数变化的简单函数；若算子原来只定义于 Schwartz 类，不能直接把这些函数代入分布符号。紧支撑频率分布恰好使这一问题可以明确解决。

\Needspace{6\baselineskip}
\begin{lemma}[紧支撑分布族的解析插值]\label{b2:lem:23-compact-analytic-interpolation}
设 \(d\geq1\)、\(b>0\)，\((G_z)_{z\in\mathbb C}\) 是 \(\mathbb R^d\) 上的整分布族，全部支撑于同一个紧集 \(K\)。假设存在整数 \(N\)、常数 \(C,c\) 及固定紧邻域 \(K_1\) 包含 \(K\)，使
\[
 |\langle G_{u+\mathrm iy},\phi\rangle|
 \leq C e^{cy^2}\!
 \max_{\substack{|\alpha|\leq N\\ \xi\in K_1}}
 |\partial^\alpha\phi(\xi)|,
 \quad -b\leq u\leq0,
\]
对全部光滑测试函数 \(\phi\) 成立。再假设每个 \(y\in\mathbb R\) 都满足：\(G_{\mathrm iy}\) 是有界函数，且 \(\check G_{-b+\mathrm iy}\) 是有界函数，并有
\[
 \|G_{\mathrm iy}\|_\infty\leq A_0e^{cy^2},
 \quad
 \|\check G_{-b+\mathrm iy}\|_\infty\leq A_1e^{cy^2}.
\]
对 \(0<\theta<1\)，令
\[
 \frac1{p_\theta}=\frac{1+\theta}{2},
 \quad \frac1{q_\theta}=\frac{1-\theta}{2}.
\]
则初始定义于 Schwartz 类的
\(T_{-b\theta}f=\mathcal F^{-1}(G_{-b\theta}\cdot\widehat f)\)
具有唯一有界延拓 \(L^{p_\theta}\to L^{q_\theta}\)，其范数不超过
\(C_{b,c,\theta}A_0^{1-\theta}A_1^\theta\)。
\end{lemma}
\begin{proof}
紧支撑分布的逆变换是光滑的至多多项式增长函数。具体地，取 \(\rho\in C_c^\infty\) 在 \(K\) 附近等于一，则
\[
 \check G_z(x)=\langle G_z,\rho(\xi)e^{2\pi\mathrm i x\cdot\xi}\rangle.
\]
有限阶分布估计允许对 \(x\) 求导，并给出多项式增长。因而对有界紧支撑简单函数 \(f,g\)，可以定义
\[
 \begin{aligned}
 B_z(f,g)
 &=\iint\check G_z(x-y)\cdot f(y)\cdot g(x)\,\mathrm dy\,\mathrm dx\\
 &=\langle G_z,\rho\cdot\widehat f(\xi)\cdot\widehat g(-\xi)\rangle.
 \end{aligned}
\]
第一行只在固定有界的 \(x,y\) 集上积分，绝对收敛；第二行通过把有限阶分布配对移过积分得到，因为紧支撑输入的 Fourier 变换光滑，所需频率导数由输入的有限阶矩控制。故此处的简单函数确实在配对的共同定义域内。

记 \(p=p_\theta\)，取 \(\|f\|_p=\|g\|_p=1\) 的上述简单函数。将非零值分开写成 \(f=\sum_j a_j\mathbf1_{E_j}\)，各 \(E_j\) 两两不交且有界。令
\[
 f_w=\sum_j\frac{a_j}{|a_j|}|a_j|^{p(1+w)/2}
 \cdot\mathbf1_{E_j},\quad0\leq\operatorname{Re}w\leq1,
\]
并同样构造 \(g_w\)。复幂采用正数 \(|a_j|\) 的实对数。它们是有限个固定简单函数的解析线性组合，在 \(w=\theta\) 时分别等于 \(f,g\)，而
\[
 \|f_{\mathrm it}\|_2=\|g_{\mathrm it}\|_2=1,
 \quad
 \|f_{1+\mathrm it}\|_1=\|g_{1+\mathrm it}\|_1=1.
\]
采用双线性配对而不对 \(g_w\) 取共轭，函数
\[
 F(w)=B_{-bw}(f_w,g_w)
\]
在条带内部全纯并连续到边界。其全纯性来自有限解析组合和 \(G_z\) 的分布解析性，不是对任意非光滑函数作未经定义的乘法。

固定 \(f,g\) 后，\(f_w,g_w\) 的支撑不变，各系数的模只随 \(\operatorname{Re}w\) 改变。因此上面出现的 Fourier 变换及其有限阶导数在 \(K_1\) 上有与 \(\operatorname{Im}w\) 无关的界。分布增长假设给出
\[
 |F(s+\mathrm it)|\leq C_{f,g}e^{cb^2t^2},
 \quad0\leq s\leq1.
\]
在右、左两条竖边上，Plancherel 和有界核卷积估计分别给出
\[
 |F(\mathrm it)|\leq A_0e^{cb^2t^2},
 \quad
 |F(1+\mathrm it)|\leq A_1e^{cb^2t^2}.
\]
取 \(\kappa>cb^2\)，乘上 \(e^{\kappa(w-\theta)^2}\)，便使内部增长有界；边界范数只多出分别为 \(e^{\kappa\theta^2}\) 和 \(e^{\kappa(1-\theta)^2}\) 的固定因子。\cref{b2:thm:17-three-lines}遂给出
\[
 |B_{-b\theta}(f,g)|
 \leq e^{\kappa\theta(1-\theta)}
 \cdot A_0^{1-\theta}A_1^\theta.
\]
常数中的 \(\kappa\) 对所有测试输入相同。若某个 \(A_i=0\)，先增加一个正数再取极限，结论仍成立。

由齐次性和 \(L^p\) 对偶性，这给出紧支撑简单函数上的 \(L^p\to L^{p'}\) 界；这里 \(p'=q_\theta\)。这些简单函数在 \(L^p\) 中稠密，所以得到唯一有界延拓。它与原来 Schwartz 类上的分布乘子一致：先对光滑紧支撑输入用有界简单函数一致逼近，所有固定频率矩导数也收敛，故分布配对极限相同；再将任意 Schwartz 函数作空间光滑截断，同时在 Schwartz 拓扑和 \(L^p\) 范数中趋近，便识别两个极限。
\end{proof}

引理中，指数关系使 \(q_\theta=p_\theta'\)。这一对称性正好适合上一节的正二次型，而内部增长条件保证三线引理真正适用。现在采用\term[stein-tomas-dingli]{Stein--Tomas 定理}[Stein--Tomas theorem]的球面版本。

\Needspace{6\baselineskip}
\begin{theorem}[Stein--Tomas 定理：球面]\label{b2:thm:23-stein-tomas}
设整数 \(d\geq2\)，单位球面 \(S=S^{d-1}\subset\mathbb R^d\) 取面积测度 \(\sigma\)。令
\[
 p_*=\frac{2(d+1)}{d+3},\quad
 q_*=p_*'=\frac{2(d+1)}{d-1}.
\]
对每个 \(1\leq p\leq p_*\) 和复值 Schwartz 函数 \(f\)，有
\[
 \|\widehat f|_S\|_{L^2(S,\sigma)}
 \leq C_{d,p}\cdot\|f\|_{L^p(\mathbb R^d)}.
\]
限制算子由此唯一连续延拓到 \(L^p(\mathbb R^d)\)。对偶地，
\[
 Eg(x)=\int_Se^{2\pi\mathrm i x\cdot\xi}
 \cdot g(\xi)\,\mathrm d\sigma(\xi)
\]
对每个 \(q_*\leq q\leq\infty\) 满足
\(\|Eg\|_q\leq C_{d,q}\cdot\|g\|_{L^2(S)}\)。
在曲面输入指数固定为 \(2\) 时，输出端点 \(q_*\) 不能减小。
\end{theorem}
\begin{proof}
球面族 \(G_z\) 已具有共同紧支撑、分布解析性及条带上的统一有限阶增长。上面又证明了两条边界
\[
 \operatorname{Re}z=0:\ L^2\to L^2,
 \quad
 \operatorname{Re}z=-b:\ L^1\to L^\infty,
 \quad b=\frac{d+1}{2}.
\]
目标参数为 \(z=-1\)，因此在前一引理中取
\[
 \theta=\frac1b=\frac2{d+1}.
\]
因为 \(d\geq2\)，有 \(0<\theta<1\)，插值点确实位于开条带。直接计算
\[
 \frac1{p_\theta}=\frac{1-\theta}{2}+\theta
 =\frac{d+3}{2(d+1)},\quad
 \frac1{q_\theta}=\frac{1-\theta}{2}
 =\frac{d-1}{2(d+1)}.
\]
于是 \(T_{-1}=\tfrac12(f\mapsto f*\check\sigma)\) 有界
\(L^{p_*}\to L^{p_*'}\)。\cref{b2:prop:23-ttstar}的正二次型恒等式给出
\[
 \|\widehat f|_S\|_2^2
 =2\langle T_{-1}f,f\rangle
 \leq2\|T_{-1}f\|_{p_*'}\cdot\|f\|_{p_*}
 \leq C_d\|f\|_{p_*}^2.
\]
开平方得到限制端点。

另一方面，\(|\widehat f(\xi)|\leq\|f\|_1\)，所以
\(\|\widehat f|_S\|_2\leq\sigma(S)^{1/2}\cdot\|f\|_1\)。
将这个强型界与刚得的端点界作 Riesz--Thorin 插值，得到全部 \(1\leq p\leq p_*\)。Schwartz 类在这些有限指数空间中稠密，故有唯一连续限制延拓。\cref{b2:prop:23-duality}给出对偶延拓界，\(p'\) 恰好遍历 \([q_*,\infty]\)。最后，\subsecref{b2:subsec:230202}的 Knapp 例子在曲面输入 \(L^2\) 时要求 \(q\geq q_*\)，证明此端点在线上的锐利性。
\end{proof}

解析插值在这里同时使用了两种信息：一条边界保留频率符号的大小，另一条边界保留振荡核的抵消。只把球面核换成 \((1+|x|)^{-(d-1)/2}\) 的绝对值，就会失去前一种结构，无法由相同方法得到这个端点。Gaussian 标量因子只负责压住虚部方向的增长，在实插值点是固定非零数，因此影响常数而不改变指数；本证明并不计算最佳限制常数。

\ProseParagraphBreak
球面族依赖旋转对称性，但端点的几何来源并不限于球面。对一般曲面，下面用局部法向幂代替球内径向幂。Bessel 核的半径计算将被拆成两个部分：法向分布产生增长，非零切向曲率产生等量衰减。

\subsection{非零曲率超曲面}
\label{b2:subsec:230404}

先单独估计法向因子。这个估计不涉及曲面，只有一维分布与频率尺度；将它与\chapref{b2:ch:21}的统一曲面衰减配合，就能看出临界竖边为何仍是 \(\operatorname{Re}z=-(d+1)/2\)。

\Needspace{6\baselineskip}
\begin{lemma}[截断法向幂的核增长]\label{b2:lem:23-normal-kernel-growth}
设 \(b>0\)，\(\tau\in C_c^\infty(\mathbb R;\mathbb C)\) 在零点邻域等于 \(1\)，\(U_z\) 为\cref{b2:lem:23-normal-analytic-family}中的整分布族。令
\[
 F_z(v)=\mathcal F^{-1}(\tau\cdot U_z)(v).
\]
则 \(F_z\) 是 \(v\) 的光滑函数，且存在只依赖 \(b,\tau\) 的常数 \(C\)，使
\[
 |F_{-b+\mathrm iy}(v)|
 \leq C e^{C(1+y^2)}(1+|v|)^{b-1},
 \quad v,y\in\mathbb R.
\]
\end{lemma}
\begin{proof}
选整数 \(N>b\)，并令 \(\beta=N+1-b>0\)。分布恒等式
\(U_z=\partial_t^NU_{z+N}\) 与乘积求导公式给出
\[
 \partial_t^N(\tau\cdot U_{z+N})
 =\tau\cdot U_z+R_z.
\]
这里 \(R_z\) 是若干项
\(\binom N\ell\tau^{(\ell)}\cdot U_{z+\ell}\)、\(1\leq\ell\leq N\)
之和。所有 \(\tau^{(\ell)}\) 都在零点附近为零，所以这些项是紧支撑光滑函数。对 \(z=-b+\mathrm iy\)，它们的任意固定阶导数范数都不超过 \(C_M e^{C_M(1+y^2)}\)：求导只产生 \(z\) 的多项式因子，Gamma 倒数由已证的竖线增长控制。因此
\[
 |\check R_{-b+\mathrm iy}(v)|
 \leq C_M e^{C_M(1+y^2)}(1+|v|)^{-M}
\]
对任意固定整数 \(M\geq0\) 成立。

剩下要估计可积函数
\[
 a_z(t)=\tau(t)\cdot U_{z+N}(t)
 =P(z+N+1)\cdot\tau(t)t_+^{z+N}.
\]
它在零点的实幂次为 \(N-b>-1\)。远离零点的光滑部分可任意多次积分分部；在零点附近，取光滑二进分割，将正半轴分成 \(t\asymp2^{-j}\)、\(j\geq0\) 的环带。固定小尺度的常数由 \(\tau\) 吸收，必要时先作一次线性缩放。第 \(j\) 块的 \(L^1\) 范数满足
\[
 \|a_{z,j}\|_1\leq C e^{C(1+y^2)}2^{-j\beta}.
\]
若在其 Fourier 积分中积分分部 \(M\) 次，则截断导数与幂函数导数各带来至多 \(2^j\)，从而
\[
 |\check a_{z,j}(v)|
 \leq C_M e^{C_M(1+y^2)}2^{-j\beta}
 \min\{\,1,(|v|2^{-j})^{-M}\,\}.
\]
每个环带的截断光滑紧支撑，故没有积分边界项。

对 \(|v|\geq1\)，取 \(j_0\geq0\) 使
\(2^{j_0}\leq|v|<2^{j_0+1}\)，并固定整数 \(M>\beta\)。小环带 \(j>j_0\) 用直接积分界：
\[
 \sum_{j>j_0}2^{-j\beta}\leq C_\beta|v|^{-\beta}.
\]
其余环带 \(0\leq j\leq j_0\) 用分部积分界：
\[
 |v|^{-M}\sum_{j=0}^{j_0}2^{j(M-\beta)}
 \leq C_{\beta,M}|v|^{-\beta}.
\]
因此 \(|\check a_z(v)|\leq C e^{C(1+y^2)}|v|^{-\beta}\)。
逆 Fourier 变换把 \(\partial_t^N\) 变成乘以 \((-2\pi\mathrm iv)^N\)，故
\[
 F_z(v)=(-2\pi\mathrm iv)^N\cdot\check a_z(v)-\check R_z(v).
\]
选择余项衰减阶足够大，并使用 \(N-\beta=b-1\)，得到大 \(|v|\) 的目标界。对 \(|v|\leq1\)，\(\check a_z\) 由 \(a_z\) 的 \(L^1\) 范数控制，\(\check R_z\) 也一致有界，便得到同样结论。光滑性则来自 \(\tau\cdot U_z\) 的紧支撑分布性质。
\end{proof}

\Needspace{6\baselineskip}
\begin{theorem}[非零曲率超曲面的 Stein--Tomas 估计]\label{b2:thm:23-curved-tomas}
设整数 \(d\geq2\)，\(S\subset\mathbb R^d\) 是光滑嵌入超曲面，
\(a\in C_c^\infty(S;\mathbb R)\)、\(a\geq0\)。假设 \(S\) 在
\(\operatorname{supp}a\) 附近没有边界，且 Gaussian 曲率在
\(\operatorname{supp}a\) 上处处非零。令
\[
 \mathrm d\mu=a\,\mathrm d\sigma,\quad
 p_*=\frac{2(d+1)}{d+3},\quad q_*=\frac{2(d+1)}{d-1},
\]
其中 \(\sigma\) 是诱导曲面测度。则对每个复值 \(f\in\mathcal S(\mathbb R^d)\)，有
\[
 \|\widehat f\|_{L^2(S,\mu)}
 \leq C_{S,a}\cdot\|f\|_{p_*}.
\]
该限制唯一延拓为 \(L^{p_*}\to L^2(S,\mu)\) 的有界算子。对偶延拓
\[
 E_\mu g(x)=\int_Se^{2\pi\mathrm i x\cdot\xi}
 \cdot g(\xi)\,\mathrm d\mu(\xi)
\]
有界 \(L^2(S,\mu)\to L^{q_*}(\mathbb R^d)\)。同样得到所有
\(1\leq p\leq p_*\) 的限制界及 \(q_*\leq q\leq\infty\) 的延拓界。常数只涉及紧支撑附近的曲面几何及给定密度，未要求不同曲面具有统一常数。
\end{theorem}
\begin{proof}
若 \(a=0\)，结论平凡。否则由紧性选取有限个图形坐标片，并在 \(\operatorname{supp}a\) 附近取非负光滑单位分解。每个片都可经刚性运动写成
\[
 \xi=(\xi',h(\xi')),\quad \xi'\in U\subset\mathbb R^{d-1},
\]
其密度为
\[
 A(\xi')\,\mathrm d\xi',\quad
 A(\xi')=\rho(\xi',h(\xi'))\cdot
 a(\xi',h(\xi'))\cdot\sqrt{1+|\nabla h(\xi')|^2},
\]
其中 \(A\in C_c^\infty(U)\)、\(A\geq0\)。平移和正交变换分别对应环境输入的模为一调制和正交换元，保持 \(L^p\) 范数，所以只需在这些坐标中证明估计。曲率公式
\[
 \mathcal K(\xi')
 =\frac{\det D^2h(\xi')}
 {\bigl(1+|\nabla h(\xi')|^2\bigr)^{(d+1)/2}}
\]
保证 \(D^2h\) 在 \(\operatorname{supp}A\) 上可逆；缩小坐标邻域后，所需非退化下界在支撑附近统一成立。

取 \(\tau\in C_c^\infty(\mathbb R)\) 在零点附近等于一。使用本片的法向坐标
\(t=\xi_d-h(\xi')\)，定义
\[
 G_z(\xi)=A(\xi')\cdot\tau(t)U_z(t),
 \quad t=\xi_d-h(\xi').
\]
严格地说，它对测试函数 \(\phi\) 的配对是
\[
 \langle G_z,\phi\rangle
 =\int_U A(\xi')\cdot
 \left\langle U_z(t),
 \tau(t)\cdot\phi(\xi',h(\xi')+t)\right\rangle\,\mathrm d\xi'.
\]
换元的环境 Jacobi 行列式为一；图形的面积因子已放入 \(A\)。引理中的统一有限阶估计允许在 \(\xi'\) 上积分，并表明 \(G_z\) 是共同紧支撑的整分布族，具有条带上的统一 Gaussian 型增长。这里 \(\partial_t\) 只对测试函数的最后一个变量求导，没有额外的法向尺度因子。

在 \(z=-1\) 时，\(U_{-1}=\delta_0\)，所以
\[
 G_{-1}=A(\xi')\delta(\xi_d-h(\xi')),
\]
恰好是该片的曲面测度。在 \(\operatorname{Re}z=0\) 上，
\[
 \|G_{\mathrm iy}\|_\infty
 \leq\|A\|_\infty\cdot\|\tau\|_\infty
 \cdot|P(1+\mathrm iy)|
 \leq C e^{C(1+y^2)}.
\]
因此这条竖边仍提供 \(L^2\) 起点。

令 \(b=(d+1)/2\)。法向与切向变量分离后，逆变换核为
\[
 \check G_z(x)=F_z(x_d)\cdot I(x),
 \quad
 I(x)=\int_U e^{2\pi\mathrm i(x'\cdot\xi'+x_dh(\xi'))}
 \cdot A(\xi')\,\mathrm d\xi'.
\]
先在 \(\operatorname{Re}z>-1\) 中由绝对可积性和 Fubini 得到此式，再由两侧的分布解析性延拓；也可直接从上一配对公式代入指数测试函数得到。法向引理给出
\[
 |F_{-b+\mathrm iy}(x_d)|
 \leq C e^{C(1+y^2)}(1+|x_d|)^{(d-1)/2}.
\]

切向积分正是图形曲面上光滑密度的逆 Fourier 变换。将其写为曲面测度时，密度为 \(A/\sqrt{1+|\nabla h|^2}\)，仍光滑紧支撑。因此\cref{b2:prop:21-curved-surface-decay}适用，并给出对全部方向一致的
\[
 |I(x)|\leq C(1+|x|)^{-(d-1)/2}.
\]
恢复该结论所用的几何条件：当 \(x_d\ne0\) 且 \(x'/x_d\) 在紧参数集内时，归一化切向相位的 Hesse 矩阵为 \(D^2h\)，可用\cref{b2:prop:21-parametric-stationary}；若 \(|x'|>C_0|x_d|\)，其中 \(C_0>2\sup_{\operatorname{supp}A}|\nabla h|\)，则切向梯度长度至少为 \(|x'|/2\)，可用非驻相积分分部。前章已经通过位置与方向的有限覆盖，使这些下界、分割振幅的导数和全部常数统一。这里正好满足该命题的紧支撑与非零曲率假设。

因为 \(1+|x_d|\leq1+|x|\)，两种因子相乘后便有
\[
 \|\check G_{-b+\mathrm iy}\|_\infty
 \leq C e^{C(1+y^2)}.
\]
法向增长指数 \(b-1\) 与切向衰减指数 \((d-1)/2\) 相等，这就是选择该竖边的理由。对小 \(x_d\) 或接近切向的方向，统一曲面估计仍然成立，不需要把非退化驻相公式用于不存在驻点的区域。

于是\cref{b2:lem:23-compact-analytic-interpolation}适用于本片的 \(G_z\)。取 \(\theta=2/(d+1)\)，得到 \(T_{-1}\) 的
\(L^{p_*}\to L^{p_*'}\) 界。本片没有球面径向换元的 \(1/2\) 因子，故其正二次型直接为
\[
 \langle T_{-1}f,f\rangle
 =\int_U|\widehat f(\xi',h(\xi'))|^2
 \cdot A(\xi')\,\mathrm d\xi'.
\]
Hölder 不等式使它不超过 \(C\|f\|_{p_*}^2\)。对有限个非负密度片相加，左侧恰为 \(\|\widehat f\|_{L^2(\mu)}^2\)，右侧常数也只作有限相加；开平方得到所述端点。

由 Schwartz 稠密性和\cref{b2:prop:23-ttstar}中的对偶性得到限制与延拓的相容定义。有限测度 \(\mu\) 还给出
\(\|\widehat f\|_{L^2(\mu)}\leq\mu(S)^{1/2}\cdot\|f\|_1\)；
与端点作插值，再取对偶，得到陈述中的其余指数。
\end{proof}

一般曲面证明没有要求曲率同号。切向 Hesse 矩阵可以是不定二次型；驻相上界使用的是行列式不为零及所需的统一下界，而不是凸性。因此本定理也适用于具有非零负 Gaussian 曲率的紧片。若曲率趋于零，统一常数可能发散；若密度碰到曲面边界，则还要处理边界振荡，不能直接删除紧支撑于无边界邻域这一条件。

\ProseParagraphBreak
球面的 Bessel 族和一般图形的法向族得到同一个经典端点：前者展示精确径向归一化，后者把法向奇性与切向曲率的平衡分开。它们都只完成了曲面输入为 \(L^2\) 时的端点及其基本插值后果，尚未处理改变输入指数后全部帽之间的相互作用；那是后面现代限制问题的内容。


% !TeX root = ../../Book2.tex
\OptionalSection{极值问题}{b2:sec:2305}

算子有界以后，可以进一步研究最佳常数及达到它的输入。极值问题比有界性多出一个紧性要求：极值化序列可能因调制而在输出空间中平移，也可能在输入曲面上集中。本节先证明非零弱极限如何恢复强紧性，再利用这一机制取得严格高于 Stein--Tomas 指数时的极值函数。临界指数处的尺度集中需要额外分析，不能由同一个论证自动解决。

\subsection{最佳常数}
\label{b2:subsec:230501}

本节固定紧致光滑超曲面 \(S\subset\mathbb R^d\)，取非零面积测度 \(\sigma\)；同样的论证适用于紧支撑的非零非负光滑曲面密度。记
\[
 Eg(x)=\int_S e^{2\pi\mathrm i x\cdot\xi}g(\xi)\,\mathrm d\sigma(\xi).
\]
除非另有说明，假设已经建立 \(E:L^2(S)\to L^q(\mathbb R^d)\) 的有界性，其中 \(2<q<\infty\)。曲率不用于下面的条件紧性论证；它在前面的延拓估计以及后面的具体应用中出现。

\Needspace{6\baselineskip}
\begin{definition}[最佳常数与极值函数]\label{b2:def:23-extremizer}
设 \(d\geq2\)，\(S\subset\mathbb R^d\) 为非空紧致光滑超曲面，取面积测度 \(\sigma\)，并设 \(2<q<\infty\)。假设延拓算子
\[
 Eg(x)=\int_S e^{2\pi\mathrm i x\cdot\xi}g(\xi)\,\mathrm d\sigma(\xi)
\]
从 \(L^2(S,\sigma)\) 到 \(L^q(\mathbb R^d)\) 有界。其最佳常数为
\[
 \mathcal C_q=\sup_{\|g\|_{L^2(S)}=1}\|Eg\|_{L^q(\mathbb R^d)}.
\]
若单位向量 \(g\) 满足 \(\|Eg\|_q=\mathcal C_q\)，称它为归一化极值函数。若 \(\|g_n\|_2=1\) 且 \(\|Eg_n\|_q\to\mathcal C_q\)，则称 \((g_n)\) 为极值化序列。
\end{definition}

有界性给出 \(\mathcal C_q<\infty\)。取 \(g=\sigma(S)^{-1/2}\)，则 \(Eg(0)=\sigma(S)^{1/2}>0\)；输出连续，所以在原点附近也不为零，故 \(\mathcal C_q>0\)。最佳常数的数值、是否达到及所有极值函数的分类是三个不同问题。证明存在性并不同时计算常数，更不自动给出分类。

\ProseParagraphBreak
即使算子范数有限，也未必有极值向量。例如 \(\ell^2(\{2,3,\ldots\})\) 上的对角算子
\[
 T(x_n)_{n\geq2}=\bigl((1-1/n)x_n\bigr)_{n\geq2}
\]
范数为 \(1\)，因为 \(\|Te_n\|_2=1-1/n\to1\)。但是对每个非零 \(x\)，至少一个非零坐标的平方被严格缩小，因而 \(\|Tx\|_2<\|x\|_2\)。单位坐标向量构成极值化序列，弱收敛到零，却没有强收敛子列。延拓问题必须找到排除这类损失的具体机制。

\subsection{极值化序列}
\label{b2:subsec:230502}

上确界定义保证极值化序列存在。Hilbert 空间的弱紧性给出子列 \(g_n\rightharpoonup g\)，但可能 \(g=0\)，也可能 \(\|g\|_2<1\)。弱收敛在延拓端提供了比一般有界算子更多的局部信息。

\Needspace{6\baselineskip}
\begin{lemma}[弱输入收敛与局部延拓紧性]\label{b2:lem:23-local-extension-compactness}
设 \(d\geq2\)，\(S\subset\mathbb R^d\) 为非空紧致光滑超曲面，取面积测度 \(\sigma\)，并定义
\[
 Ef(x)=\int_S e^{2\pi\mathrm i x\cdot\xi}f(\xi)\,\mathrm d\sigma(\xi)
 \quad(f\in L^2(S,\sigma)).
\]
若 \(g_n\rightharpoonup g\) 于 \(L^2(S,\sigma)\)，则对每个非负整数 \(k\)，有
\[
 Eg_n\longrightarrow Eg\quad\text{于 }C^k_{\mathrm{loc}}(\mathbb R^d).
\]
这里的结论只使用曲面测度有限且支集紧致，不要求全空间 \(L^q\) 估计。
\end{lemma}
\begin{proof}
弱收敛序列有界。对每个多重指标 \(\alpha\)，积分可以逐阶求导，并且
\[
 \partial_x^\alpha Eg_n(x)
 =\int_S(2\pi\mathrm i\xi)^\alpha
       e^{2\pi\mathrm i x\cdot\xi}g_n(\xi)\,\mathrm d\sigma(\xi),
\]
\[
 |\partial_x^\alpha Eg_n(x)|
 \leq(2\pi)^{|\alpha|}
       \left(\int_S|\xi^\alpha|^2\,\mathrm d\sigma(\xi)\right)^{1/2}
       \cdot\|g_n\|_2.
\]
紧支集使右侧有限且对 \(n,x\) 一致。固定 \(x\) 时，积分中的测试函数属于 \(L^2(S)\)，所以弱收敛给出全部导数的逐点收敛。再用高一阶导数的一致界，得到各阶导数的一致 Lipschitz 界。

固定紧球及一个导数阶数，取该球的有限 \(\eta\)-网。网格各点的差在 \(n\) 充分大时同时趋于零；任意点到最近网格点的差由统一 Lipschitz 常数乘 \(\eta\) 控制。先取小 \(\eta\)，再取大 \(n\)，便得到这个导数在球上的一致收敛。对有限个 \(|\alpha|\leq k\) 重复即可。
\end{proof}

局部一致收敛仍然允许质量在越来越远的区域出现，因此不能直接推出全空间 \(L^q\) 强收敛。另一方面，若已经得到输入的强收敛 \(g_n\to g\)，则
\[
 \|g\|_2=1,\quad
 \|Eg_n-Eg\|_q\leq\mathcal C_q\|g_n-g\|_2\longrightarrow0,
\]
所以极限就是极值函数。真正需要解决的是如何从弱紧性恢复全部输入范数。

\subsection{对称性}
\label{b2:subsec:230503}

球面上的精确对称可以直接代入积分验证。若 \(U\) 为正交矩阵、\(|c|=1\)，定义
\[
 R_Ug(\xi)=g(U^{-1}\xi),\quad
 M_ag(\xi)=e^{-2\pi\mathrm i a\cdot\xi}g(\xi).
\]
则
\[
 E(cR_Ug)(x)=c\,Eg(U^{\mathsf T}x),\quad
 E(M_ag)(x)=Eg(x-a).
\]
这些变换都保持输入 \(L^2\) 范数和输出 \(L^q\) 范数。调制对应于环境空间的平移；它对一般曲面密度仍是精确对称，而旋转只有在测度本身保持不变时才是对称。

\ProseParagraphBreak
在球面上，若 \(g\in L^2(S)\) 固定且 \(|a_n|\to\infty\)，则 \(M_{a_n}g\rightharpoonup0\)。事实上，对任意 \(h\in L^2(S)\)，乘积 \(g\overline h\in L^1(S)\)。先以光滑函数在 \(L^1(S)\) 中逼近这个乘积；光滑曲面密度的 Fourier 变换由驻相估计趋于零，逼近误差则由其 \(L^1\) 范数一致控制。因此
\[
 \int_S e^{-2\pi\mathrm i a_n\cdot\xi}
      g(\xi)\overline{h(\xi)}\,\mathrm d\sigma(\xi)\longrightarrow0.
\]
若 \(g\) 已经是极值函数，便得到一个弱极限为零的极值化序列。极值函数的存在与所有极值化序列原样强紧，显然是不同的结论。

\ProseParagraphBreak
另一种损失发生在输入曲面。设 \(\theta_\delta\) 为球面上以固定点为中心、直径趋于零的帽，令
\[
 g_\delta=\sigma(\theta_\delta)^{-1/2}\mathbf1_{\theta_\delta}.
\]
它的 \(L^2\) 范数为一，而对每个 \(h\in L^2(S)\)，
\[
 |\langle g_\delta,h\rangle|
 \leq\left(\int_{\theta_\delta}|h|^2\,\mathrm d\sigma\right)^{1/2}
 \longrightarrow0.
\]
所以它也弱收敛到零。帽测度约为 \(\delta^{d-1}\) 时，Knapp 管给出的输出下界为
\[
 \|Eg_\delta\|_q
 \gtrsim \delta^{(d-1)/2-(d+1)/q}.
\]
在 \(q_*=2(d+1)/(d-1)\) 处，右侧不趋于零。这组帽的面积虽趋于零，临界输出却保留正的质量，显示了集中可能造成的紧性损失。

帽的局部放大趋于抛物面几何，但这不是单位球面的精确缩放对称。归一化必须区别两种操作：调制、相位与旋转在原问题内保持范数；帽重标度可能带来另一个极限曲面问题。消除精确对称并不自动排除尺度集中，也不自动排除多个远离的输出峰。

\subsection{极值函数}
\label{b2:subsec:230504}

先假设极值函数存在。它满足一个可以在现有算子范围内严格解释的变分方程。

\Needspace{6\baselineskip}
\begin{proposition}[极值函数的 Euler--Lagrange 方程]\label{b2:prop:23-extremizer-euler}
设 \(d\geq2\)，\(S\subset\mathbb R^d\) 为非空紧致光滑超曲面，取面积测度；\(E\) 为\cref{b2:def:23-extremizer}中定义的延拓算子。设 \(2<q<\infty\)，\(E:L^2(S)\to L^q(\mathbb R^d)\) 有界，并记其最佳常数为
\(\mathcal C_q=\sup_{\|f\|_2=1}\|Ef\|_q\)。
若 \(\|g\|_2=1\)、\(\|Eg\|_q=\mathcal C_q\)，则
\[
 E^*(|Eg|^{q-2}Eg)=\mathcal C_q^q g
 \quad\text{于 }L^2(S).
\]
这里 \(E^*:L^{q'}(\mathbb R^d)\to L^2(S)\) 为有界伴随算子。
\end{proposition}
\begin{proof}
首先
\[
 \bigl\||Eg|^{q-2}Eg\bigr\|_{q'}
 =\|Eg\|_q^{q-1}<\infty,
\]
所以左侧有定义。固定 \(h\in L^2(S)\)，对实参数 \(t\) 考察
\[
 F(t)=\frac{\|E(g+th)\|_q^q}{\|g+th\|_2^q}.
\]
分母在零附近不为零，且极值性给出 \(F(t)\leq F(0)=\mathcal C_q^q\)。对 \(u,v\in\mathbb C\) 及 \(|t|\leq1\)，差商的模由
\[
 C_q\bigl(|u|^{q-1}\cdot|v|+|v|^q\bigr)
\]
控制。代入 \(u=Eg(x)\)、\(v=Eh(x)\)，由 Hölder 不等式可积，因而允许在积分号下求导。分子与分母在零点的导数分别为
\[
 q\,\operatorname{Re}\int_{\mathbb R^d}
       |Eg|^{q-2}\overline{Eg}\,Eh\,\mathrm dx,\quad
 q\,\operatorname{Re}\int_S\overline g\,h\,\mathrm d\sigma.
\]
由 \(F'(0)=0\) 及 \(\|g\|_2=1\)，得到
\[
 \operatorname{Re}\int_{\mathbb R^d}
       |Eg|^{q-2}\overline{Eg}\,Eh\,\mathrm dx
 =\mathcal C_q^q\operatorname{Re}
       \int_S\overline g\,h\,\mathrm d\sigma.
\]
再把 \(h\) 换成 \(\mathrm i h\)，便得到对应的虚部等式。有界伴随的定义遂给出所述 \(L^2(S)\) 等式。
\end{proof}

这里的 \(E^*\) 按有界算子对偶定义，不是对任意 \(L^{q'}\) 代表逐点取 Fourier 限制。这个方程只是必要条件：它的解还可能是其他临界点，不能仅凭方程就断言达到全局最佳常数。特殊曲面和特殊指数有时允许用卷积几何计算最佳常数与极值分类，但这需要进一步证明。

\subsection{集中紧性}
\label{b2:subsec:230505}

极值化序列的有利之处在于：一旦找到非零弱极限，严格的幂次不等式便有机会排除剩余质量。下面的范数分裂公式把这一思路变成可计算的不等式。

\Needspace{6\baselineskip}
\begin{lemma}[幂次范数分裂]\label{b2:lem:23-power-splitting}
设 \((X,\mu)\) 为测度空间，\(1<q<\infty\)，\(u_n\in L^q(X,\mu)\)，\(u\) 为可测函数。若 \(u_n\to u\) \(\mu\)-几乎处处，且 \(\sup_n\|u_n\|_{L^q(X,\mu)}<\infty\)，则 \(u\in L^q(X,\mu)\)，并且
\[
 \|u_n\|_q^q-\|u_n-u\|_q^q\longrightarrow\|u\|_q^q.
\]
\end{lemma}
\begin{proof}
Fatou 引理给出 \(u\in L^q\)。先注意，对每个 \(\varepsilon>0\)，存在 \(C_{\varepsilon,q}\)，使所有复数 \(a,b\) 满足
\[
 \bigl||a+b|^q-|a|^q\bigr|
 \leq\varepsilon|a|^q+C_{\varepsilon,q}|b|^q.
\]
当 \(|b|\leq\eta|a|\) 时，沿线段用实变量微分估计，左侧至多为
\(q(1+\eta)^{q-1}\eta|a|^q\)，取小 \(\eta\) 即可；在其余情形，\(|a|\leq\eta^{-1}|b|\)，两项都由常数乘 \(|b|^q\) 控制。

记 \(v_n=u_n-u\) 及
\[
 D_n=|u_n|^q-|v_n|^q-|u|^q.
\]
则 \(D_n\to0\) 几乎处处，并有
\[
 \bigl(|D_n|-\varepsilon|v_n|^q\bigr)_+
 \leq(C_{\varepsilon,q}+1)|u|^q.
\]
左侧也几乎处处趋于零。支配收敛与 \(\sup_n\|v_n\|_q<\infty\) 给出
\[
 \limsup_n\int|D_n|
 \leq\varepsilon\sup_n\|v_n\|_q^q.
\]
令 \(\varepsilon\downarrow0\)，便得到 \(D_n\to0\) 于 \(L^1\)，特别得到所述积分等式。
\end{proof}

\Needspace{6\baselineskip}
\begin{proposition}[非零弱极限恢复全部质量]\label{b2:prop:23-nonzero-profile}
设 \(d\geq2\)，\(S\subset\mathbb R^d\) 为非空紧致光滑超曲面，取面积测度；\(E\) 为\cref{b2:def:23-extremizer}中的延拓算子。设 \(2<q<\infty\)，\(E:L^2(S)\to L^q(\mathbb R^d)\) 有界，并记
\(\mathcal C_q=\sup_{\|f\|_2=1}\|Ef\|_q\)。
若 \(\|g_n\|_2=1\)、\(\|Eg_n\|_q\to\mathcal C_q\)，且
\[
 g_n\rightharpoonup g\neq0\quad\text{于 }L^2(S),
\]
则 \(g_n\to g\) 强收敛于 \(L^2(S)\)，且 \(g\) 为归一化极值函数。
\end{proposition}
\begin{proof}
令 \(r_n=g_n-g\)、\(a=\|g\|_2^2\)。弱收敛与 Hilbert 空间范数展开给出
\[
 0<a\leq1,\quad \|r_n\|_2^2\longrightarrow1-a.
\]
由\cref{b2:lem:23-local-extension-compactness}，有 \(Eg_n(x)\to Eg(x)\)；有界性使输出在 \(L^q\) 中有界。于是\cref{b2:lem:23-power-splitting}给出
\[
 \|Eg_n\|_q^q=\|Eg\|_q^q+\|Er_n\|_q^q+o(1).
\]
对两项分别使用同一个最佳常数，
\[
 \mathcal C_q^q
 \leq\mathcal C_q^q\bigl[a^{q/2}+(1-a)^{q/2}\bigr].
\]
若 \(0<a<1\)，则因 \(q/2>1\)，有
\[
 a^{q/2}+(1-a)^{q/2}<a+(1-a)=1,
\]
与 \(\mathcal C_q>0\) 矛盾。因此 \(a=1\)、\(\|r_n\|_2\to0\)。算子连续性再给出 \(\|Eg\|_q=\mathcal C_q\)。
\end{proof}

这个命题证明了所需的输入平方范数分裂和输出幂次范数分裂，数值上的严格不等式才因此可以使用。非零弱极限对于任意有界序列并不足以推出强收敛；这里还使用了极值化性质及 \(q>2\)。它把极值存在问题归结为一个具体目标：通过允许的归一化找到非零弱极限。

\Needspace{6\baselineskip}
\begin{theorem}[较大输出指数下的极值存在]\label{b2:thm:23-supercritical-extremizer}
设 \(d\geq2\)，\(S\subset\mathbb R^d\) 为非空紧致光滑超曲面，取面积测度；\(E\) 为\cref{b2:def:23-extremizer}中的延拓算子。假设对某个 \(2<q_0<\infty\) 及有限常数 \(A\)，有
\[
 \|Eg\|_{q_0}\leq A\|g\|_2.
\]
则对每个有限 \(q>q_0\)，\(E:L^2(S)\to L^q(\mathbb R^d)\) 有界，且其最佳常数
\(\mathcal C_q=\sup_{\|f\|_2=1}\|Ef\|_q\)
都能达到。更确切地，每个满足 \(\|g_n\|_2=1\)、\(\|Eg_n\|_q\to\mathcal C_q\) 的序列，在适当调制和常数相位变换后，具有强 \(L^2(S)\) 收敛子列。
\end{theorem}
\begin{proof}
有限测度给出 \(\|Eg\|_\infty\leq\sigma(S)^{1/2}\|g\|_2\)，与假设插值得到目标 \(L^q\) 有界性。对单位极值化序列，有
\[
 \|Eg_n\|_q^q
 \leq\|Eg_n\|_\infty^{q-q_0}\cdot\|Eg_n\|_{q_0}^{q_0}
 \leq A^{q_0}\|Eg_n\|_\infty^{q-q_0}.
\]
左侧趋于正数 \(\mathcal C_q^q\)，因此丢掉有限项后，存在与 \(n\) 无关的 \(\kappa>0\)，可选 \(x_n\in\mathbb R^d\) 满足 \(|Eg_n(x_n)|\geq\kappa\)。选择 \(|c_n|=1\)，使 \(c_nEg_n(x_n)\) 为非负实数，并令
\[
 h_n(\xi)=c_ne^{2\pi\mathrm i x_n\cdot\xi}g_n(\xi).
\]
则
\[
 Eh_n(x)=c_nEg_n(x+x_n),\quad
 Eh_n(0)\geq\kappa,\quad \|h_n\|_2=1.
\]
它仍为同一问题的极值化序列。取弱收敛子列 \(h_n\rightharpoonup h\)，由于常数函数属于 \(L^2(S)\)，有
\[
 Eh(0)=\lim_n Eh_n(0)\geq\kappa.
\]
所以 \(h\neq0\)。\cref{b2:prop:23-nonzero-profile}给出 \(h_n\to h\) 的强收敛与极值性质。
\end{proof}

对球面，取 \(q_0=q_*=2(d+1)/(d-1)\)，便得到每个有限 \(q>q_*\) 的极值存在。证明同样适用于紧支撑的非零非负光滑曲面密度，因而一般紧曲率片也可使用\cref{b2:thm:23-curved-tomas}给出的 \(L^2\) 延拓界。上述论证的关键是正指数 \(q-q_0\)；在临界情形 \(q=q_0\)，寻找非零弱极限需要更精细的集中分析。

\ProseParagraphBreak
前面的归一化小帽在 \(q>q_*\) 时也显示出这一区别：其 \(L^\infty\) 输出至多为 \(\sigma(\theta_\delta)^{1/2}\to0\)，而 \(L^{q_*}\) 输出一致有界，故同一插值不等式给出 \(\|Eg_\delta\|_q\to0\)。因此这组单帽不能在较大指数处极值化；在临界指数处，该消失结论不再成立。

还有一个直接以全空间输出尾部表达的紧性判据。它的假设较强，但在具体问题中有时比寻找轮廓更易核验。

\Needspace{6\baselineskip}
\begin{proposition}[统一输出尾部判据]\label{b2:prop:23-extremizer-tail}
设 \(d\geq2\)，\(S\subset\mathbb R^d\) 为非空紧致光滑超曲面，取面积测度；\(E\) 为\cref{b2:def:23-extremizer}中的延拓算子。设 \(2<q<\infty\)，\(E:L^2(S)\to L^q(\mathbb R^d)\) 有界，并记
\(\mathcal C_q=\sup_{\|f\|_2=1}\|Ef\|_q\)。
假设 \(\|g_n\|_2=1\)、\(\|Eg_n\|_q\to\mathcal C_q\)。若在保持输入与输出范数的对称归一化后，
\[
 \lim_{R\to\infty}\sup_n\int_{|x|>R}|Eg_n(x)|^q\,\mathrm dx=0,
\]
则归一化后的序列具有强 \(L^2(S)\) 收敛子列，其极限为极值函数。
\end{proposition}
\begin{proof}
取弱收敛子列 \(g_n\rightharpoonup g\)。局部延拓紧性给出每个球上的一致收敛，因而也给出球内的 \(L^q\) 强收敛。Fatou 引理给出
\[
 \int_{|x|>R}|Eg(x)|^q\,\mathrm dx
 \leq\liminf_n\int_{|x|>R}|Eg_n(x)|^q\,\mathrm dx.
\]
把全空间分成球内与球外，再用
\(|u-v|^q\leq2^{q-1}(|u|^q+|v|^q)\)，先取大 \(R\)，再取大 \(n\)，即得 \(Eg_n\to Eg\) 于全空间 \(L^q\)。于是
\[
 \|Eg\|_q=\mathcal C_q,\quad
 \mathcal C_q\leq\mathcal C_q\|g\|_2,\quad
 \|g\|_2\leq1.
\]
由 \(\mathcal C_q>0\)，得到 \(\|g\|_2=1\)。Hilbert 空间中弱收敛加范数收敛便给出强收敛。
\end{proof}

每个固定 \(Eg_n\) 可积，并不保证尾部对全部 \(n\) 一致变小。调制造成的平移、帽集中造成的长管尺度，都可能破坏这个统一条件。因而尾部判据提供的是可检验的紧性目标，不能据此无条件断言临界曲面的极值存在。

\ProseParagraphBreak
\appref{b2:app:F}说明消失、分裂和紧性的区分；其中\cref{b2:sec:app-F-6}还说明限制问题的轮廓分析需要什么额外输入。对于临界限制问题，可能必须先证明一个精细延拓估计，以找到带有位置和尺度参数的非零轮廓，再证明余项在目标范数中消失。如果帽缩放产生抛物面轮廓，还需比较局部极限问题与原曲面问题的最佳常数。非零弱极限的强收敛机制由此成为轮廓分析的一环，位置、尺度和最佳常数的比较则决定怎样找到这样的极限。

% !TeX root = ../../Book2.tex
\OptionalSection{现代 Fourier 限制问题}{b2:sec:2306}

前面已经证明经典 \(L^2\) 输入范围，并从小帽与球面尾部得到必要条件。本节说明剩余困难怎样导向双线性估计、波包、解耦与多项式分割。先通过具体模型观察不同频率方向怎样相互作用，再介绍一般曲面上的深层限制与解耦定理。

\subsection{限制猜想}
\label{b2:subsec:230601}

以 \(S^{d-1}\subset\mathbb R^d\) 及面积测度为具体模型，经典限制猜想的延拓形式以如下范围为目标：
\[
 E:L^s(S^{d-1})\longrightarrow L^q(\mathbb R^d),\quad
 q>\frac{2d}{d-1},\quad
 \frac{d+1}{q}\leq\frac{d-1}{s'}.
\]
这里 \(1<s\leq\infty\)、\(q<\infty\)，\(1/s+1/s'=1\)；\(s=1\) 对应的平凡 \(L^\infty\) 输出界另行理解。第一个严格不等式来自球面测度 Fourier 变换的尾部，第二个来自小帽检验。它们描述经典目标范围，不是本章已经证明充分性的全部范围。具体维数以及边界上的强型、弱型等版本必须分别陈述。

\ProseParagraphBreak
当 \(s=2\) 时，小帽条件给出 \(q\geq q_*=2(d+1)/(d-1)\)，这正是 Stein--Tomas 已达到的范围。对 \(1<s<2\)，与平凡端点插值也给出相应的小帽边界。困难因而主要留在更大的输入指数与更小的输出指数。例如在 \(\mathbb R^3\) 的球面上，\(s=\infty\) 的经典目标是所有 \(q>3\)，而本章从 \(L^\infty(S^2)\subset L^2(S^2)\) 得到的只是 \(q\geq4\)。这个差距说明，仅改变输入的空间包含关系不能取得目标范围。

其他曲面也有相应的限制问题，但必须重新核对几何假设。紧非零曲率片具有同样形式的小帽约束；无界抛物面还涉及全局缩放，锥面则有退化的主曲率方向。球面、抛物面与锥面不能无说明地共用同一个全局命题。

\subsection{双线性归约}
\label{b2:subsec:230602}

取曲面上一个有限的光滑帽分解 \(\sum_\theta\chi_\theta=1\)，记 \(g_\theta=\chi_\theta g\)。则
\[
 Eg=\sum_\theta Eg_\theta,\quad
 |Eg|^2=\sum_{\theta,\theta'}
        Eg_\theta\,\overline{Eg_{\theta'}}.
\]
双线性归约把近邻帽与分离帽分开处理。相邻帽仍可能具有很强的相干性，需要在较小尺度重标度；分离帽的法向或切向形成横截关系，有机会得到比不加区别地应用 Hölder 更好的乘积估计。完整归约必须处理近对角项、分离项和尺度求和，不能把相邻帽直接删去。

\ProseParagraphBreak
乘积在 Fourier 侧对应测度的卷积，因而可以从频率相加来理解。设 \(\mu_\theta=g_\theta\,\mathrm d\sigma\)，则有限测度的 Fubini 公式给出
\[
 Eg_\theta\,Eg_{\theta'}
 =\mathcal F^{-1}(\mu_\theta*\mu_{\theta'}).
\]
这对应频率的和映射。若使用模平方中的共轭乘积，则应把第二个测度换成反射后的共轭测度；频率对应差映射。两种乘积可以用于不同归约，但不能在写公式时省掉共轭带来的变化。

一个平坦的二维模型已经显示横截常数。令 \(v_1,v_2\in\mathbb R^2\) 线性无关，\(f_1,f_2\) 为紧支撑 \(L^2(\mathbb R)\) 函数，并定义
\[
 E_jf_j(x)=\int_{\mathbb R}
       e^{2\pi\mathrm i t x\cdot v_j}f_j(t)\,\mathrm dt
       =\widehat f_j(-x\cdot v_j).
\]
作线性换元 \(y_j=x\cdot v_j\)，其 Jacobi 行列式绝对值为 \(|\det(v_1,v_2)|\)。Fubini 与一维 Plancherel 给出
\[
 \begin{aligned}
 \|E_1f_1\,E_2f_2\|_{L^2(\mathbb R^2)}^2
 &=|\det(v_1,v_2)|^{-1}
   \int_{\mathbb R^2}
      |\widehat f_1(-y_1)|^2\cdot|\widehat f_2(-y_2)|^2
      \,\mathrm dy_1\,\mathrm dy_2\\
 &=|\det(v_1,v_2)|^{-1}
      \cdot\|f_1\|_2^2\cdot\|f_2\|_2^2.
 \end{aligned}
\]
方向一旦靠近，常数就会恶化；若两方向是单位向量，范数常数恰为夹角正弦绝对值的 \(-1/2\) 次幂。这个模型不需要曲率，横截性已经足以使两个一维频率合成二维密度。

\ProseParagraphBreak
对于同一条曲线的不同片，曲率负责产生方向分离。例如 \(\gamma(t)=(t,t^2)\) 的和映射为
\[
 (t,s)\longmapsto(t+s,t^2+s^2),\quad
 |\det D\Psi(t,s)|=2|t-s|.
\]
在分离区间上，参数顺序固定，和映射一对一；因此能够用换元与 Plancherel 得到带 \(|t-s|^{-1}\) 权的精确乘积平方范数。\cref{b2:thm:25-parabola-bilinear}将推导这一公式，并通过逼近推广到一般输入。靠近对角线时，Jacobi 行列式退化，这正是分离片估计不能直接覆盖整个线性问题的原因。

高维和映射一般有正维纤维，不能照搬二维的一对一换元。进一步归约还需记录帽的重叠次数与每层尺度的系数；若更一般的多线性问题产生小于一的乘积指数，还必须按准范数的规则求和。横截模型中的增益能否在组合后保留，取决于这些重叠与求和估计。

\subsection{波包}
\label{b2:subsec:230603}

Knapp 管说明单个帽的主要空间尺度；波包还需要记录管的位置、振幅与相位。下面在抛物面上给出一条波包的精确公式。为使换元透明，先使用参数测度 \(\mathrm d\omega\)，写
\[
 n=d-1,\quad
 E_{\mathrm{par}}g(x',t)
 =\int_{\mathbb R^n}
       e^{2\pi\mathrm i(x'\cdot\omega+t|\omega|^2)}
       g(\omega)\,\mathrm d\omega.
\]
在固定紧参数片上，面积测度仅多出一个光滑正密度，可以吸收到振幅中；它与参数测度不应在公式里无说明地混用。

\Needspace{6\baselineskip}
\begin{proposition}[抛物面的单波包模型]\label{b2:prop:23-parabolic-packet}
设 \(n\geq1\)，抛物面采用参数测度，延拓算子定义为
\[
 E_{\mathrm{par}}g(x',t)
 =\int_{\mathbb R^n}e^{2\pi\mathrm i(x'\cdot\omega+t|\omega|^2)}
      g(\omega)\,\mathrm d\omega
 \quad(g\in L^1(\mathbb R^n)).
\]
取 \(a\in C_c^\infty(\mathbb R^n)\)、\(0<\delta\leq1\)、\(\omega_0,v\in\mathbb R^n\)，定义
\[
 g_{\delta,v}(\omega)
 =\delta^{-n/2}
   a\!\left(\frac{\omega-\omega_0}{\delta}\right)
   e^{-2\pi\mathrm i v\cdot\omega}.
\]
则 \(\|g_{\delta,v}\|_2=\|a\|_2\)。若
\[
 A(z,\tau)=\int_{\mathbb R^n}
       e^{2\pi\mathrm i(z\cdot u+\tau|u|^2)}a(u)\,\mathrm du,
\]
便有精确恒等式
\[
 \begin{aligned}
 E_{\mathrm{par}}g_{\delta,v}(x',t)
 ={}&\delta^{n/2}
 e^{2\pi\mathrm i((x'-v)\cdot\omega_0+t|\omega_0|^2)}\\
 &{}\cdot
 A\bigl(\delta(x'-v+2t\omega_0),\,\delta^2t\bigr).
 \end{aligned}
\]
对每个 \(N>0\)，存在只依赖 \(a,N\) 的 \(C_{a,N}\)，使在 \(|t|\leq\delta^{-2}\) 时，
\[
 |E_{\mathrm{par}}g_{\delta,v}(x',t)|
 \leq C_{a,N}\delta^{n/2}
       \bigl(1+\delta|x'-v+2t\omega_0|\bigr)^{-N}.
\]
\end{proposition}
\begin{proof}
令 \(\omega=\omega_0+\delta u\)。平方范数中的 Jacobian 为 \(\delta^n\)，恰好抵消振幅平方的 \(\delta^{-n}\)，得到输入范数等式。延拓相位展开为
\[
 (x'-v)\cdot\omega_0+t|\omega_0|^2
 +\delta(x'-v+2t\omega_0)\cdot u+\delta^2t|u|^2,
\]
便得到精确公式。

当 \(|\tau|\leq1\) 时，\(a(u)e^{2\pi\mathrm i\tau|u|^2}\) 具有固定紧支集，全部固定阶导数的 \(L^1\) 范数一致有界。在 \(A(z,\tau)\) 的线性 Fourier 积分中反复使用 \(1-\Delta_u\) 分部积分，得到
\[
 |A(z,\tau)|\leq C_{a,N}(1+|z|)^{-N}
 \quad(|\tau|\leq1).
\]
代入 \(z=\delta(x'-v+2t\omega_0)\)、\(\tau=\delta^2t\) 即得结论。
\end{proof}

令 \(R=\delta^{-2}\)，并把 \(\omega_0\) 限制在固定紧集。上式的主要区域为
\[
 |x'-v+2t\omega_0|\lesssim R^{1/2},\quad |t|\leq R.
\]
这是一条宽度约 \(R^{1/2}\)、长度约 \(R\) 的管，轴方向为抛物面的法向 \((-2\omega_0,1)\)。频率中心决定方向，调制参数 \(v\) 决定位置；一个帽内可以出现许多平行但位置不同的波包。公式中的快速衰减只在所指定的时间尺度内一致成立，不能把它读成一条在全时空中严格紧支撑的管。

\ProseParagraphBreak
一般输入的波包分解还需要一步：在每张帽内按位置调制展开，并控制系数平方和。\cref{b2:subsec:250502}将用 Fourier 级数实际构造这一步，证明平方系数由输入 \(L^2\) 范数控制，并估计远管尾部的可求和误差。不同帽之间的近似正交性还取决于频率支集的重叠。把波包替换成管的示性函数，会同时丢失相位、系数与尾部，所以只能用于提示几何，不能替代算子估计。

\subsection{解耦}
\label{b2:subsec:230604}

波包分解使单帽输出可理解，但不同帽仍可能相干叠加。解耦比较总延拓范数与各帽输出范数的平方和；它保留每张帽内部的相干性，再控制不同帽的组合。

具体取紧抛物面
\[
 S=\{(\omega,|\omega|^2):|\omega|\leq1\}\subset\mathbb R^d,
\]
在这里恢复使用面积测度，并以 \(E\) 记其延拓算子。设 \(R\geq1\)，将参数域切成直径约 \(R^{-1/2}\) 的有界重叠帽，用单位分解定义 \(E_\theta g\)。对任意半径 \(R\) 的球 \(B_R\)，记
\[
 w_{B_R}(x)=\left(1+\frac{|x-c(B_R)|}{R}\right)^{-10d},\quad
 \|F\|_{L^p(w_{B_R})}
 =\left(\int_{\mathbb R^d}|F|^p w_{B_R}\,\mathrm dx\right)^{1/p}.
\]
\cref{b2:thm:25-decoupling}所陈述的深层输入，在临界指数 \(p_c=2(d+1)/(d-1)\) 给出
\[
 \|Eg\|_{L^{p_c}(B_R)}
 \leq C_\varepsilon R^\varepsilon
       \left(\sum_\theta
          \|E_\theta g\|_{L^{p_c}(w_{B_R})}^2
       \right)^{1/2}
 \quad(\varepsilon>0).
\]
常数对 \(R\)、球心及 \(g\) 一致，但可以依赖 \(\varepsilon\)。抛物面的正式出版来源为 \cite{BourgainDemeter2015Proof}，系统讨论见 \cite{Demeter2020Fourier}；上述多尺度估计采用这两项文献的结论。

\ProseParagraphBreak
可以用一个初等比较看出这项输入的内容。帽数约为 \(R^{(d-1)/2}\)。单用三角不等式与离散 Cauchy--Schwarz，只能得到
\[
 \begin{aligned}
 \|Eg\|_{L^{p_c}(B_R)}
 &\leq\sum_\theta\|E_\theta g\|_{L^{p_c}(B_R)}\\
 &\lesssim R^{(d-1)/4}
   \left(\sum_\theta
      \|E_\theta g\|_{L^{p_c}(w_{B_R})}^2
   \right)^{1/2}.
 \end{aligned}
\]
第二步还用了 \(w_{B_R}\) 在球内有统一正下界。深层解耦把这个固定幂次损失改善到任意小的 \(R^\varepsilon\)。这里保留 \(R^\varepsilon\) 的尺度损失，右侧则是各片 \(L^{p_c}\) 范数的平方和。

加权空间球、帽的尺度与允许的幂次损失都是定理的一部分。在不同曲面或不同指数上，常数的尺度行为可能改变。尤其这里使用正定抛物面；不能不经修改地移到具有退化或不定第二基本形式的曲面。所谓“帽之间正交”只提供一种直觉，不能代替这个非 Hilbert 空间估计。

\subsection{多项式分割}
\label{b2:subsec:230605}

多项式分割试图把波包携带的空间质量分配到许多胞腔，再把贴近多项式零集的部分单独处理。其分配步骤需要一个外部存在性结论：对可积非负密度及给定整数 \(D\geq1\)，选择次数至多 \(D\) 的非零多项式，使补集有 \(O(D^d)\) 个胞腔，且每个胞腔的质量至多为常数乘 \(D^{-d}\) 倍总质量。这一分割定理及其证明见下述 Guth 文献。

一旦多项式选定，直线的穿越次数则有完全初等的控制。

\Needspace{6\baselineskip}
\begin{lemma}[直线的胞腔穿越次数]\label{b2:lem:23-polynomial-line-crossing}
设 \(d\geq1\)，\(D\) 为非负整数，\(P\) 为 \(\mathbb R^d\) 上次数至多 \(D\) 的非零实多项式。记 \(Z(P)=\{x:P(x)=0\}\)，称 \(\mathbb R^d\setminus Z(P)\) 的连通分支为胞腔。若仿射直线 \(\ell\) 不完全包含于 \(Z(P)\)，则 \(\ell\) 至多经过 \(D+1\) 个不同胞腔。
\end{lemma}
\begin{proof}
将直线写成 \(\ell(t)=a+tv\)、\(v\neq0\)。限制 \(P(a+tv)\) 为次数至多 \(D\) 且不恒为零的一元多项式，因此至多有 \(D\) 个不同实根。这些根把参数轴分成至多 \(D+1\) 个开区间。每个区间的像连通且避开 \(Z(P)\)，所以位于同一个胞腔中。不同区间可能落在同一胞腔，这只会使实际经过的胞腔数更少。
\end{proof}

对有厚度的管，必须同时说明零集邻域怎样删除。设 \(r>0\)，\(T\) 为直线 \(\ell\) 的半径 \(r\) 邻域，把每个胞腔缩为
\[
 \Omega^{(r)}
 =\{x\in\Omega:\operatorname{dist}(x,Z(P))>r\}.
\]
若 \(x\in T\cap\Omega^{(r)}\)，可取 \(y\in\ell\) 使 \(|x-y|\leq r\)。连接 \(x,y\) 的线段避开 \(Z(P)\)，所以 \(y\) 与 \(x\) 在同一胞腔内。因而管至多遇到 \(D+1\) 个这样的缩胞腔；若 \(\ell\subset Z(P)\)，整条管都位于被删除的零集邻域内。这一论证不能控制未经缩减的胞腔，也没有控制管在零集邻域内停留多久。

\ProseParagraphBreak
实际波包问题因此分成两部分。远离零集的胞腔部分可使用上述穿越预算，结合每个胞腔分到的质量做尺度递推；零集邻域中的部分还须区别横穿与近切向的管。长时间贴近零集的管可能沿低维几何聚集，单凭一元多项式的根数不能处理它们。波包又带有快速衰减尾部，所以把它们截到粗管时仍需保留并求和误差。

次数 \(D\) 同时增加可用胞腔数、穿越次数与零集邻域的复杂度，必须在递推中选择，不能把它当作免费的参数。正式出版的限制应用见 \cite{Guth2016Restriction}。线性限制估计需要把胞腔分配、切向管估计与尺度递推结合起来。

\ProseParagraphBreak
这些方法围绕同一个具体问题组织：单帽内的振荡可以保持相干，许多方向怎样组合才是难点。驻相描述单个局部相位，波包记录位置与方向，横截估计控制方向分离，解耦和分割则把这种控制用于许多尺度或许多空间区域。阅读一个具体定理时，应先核对曲面、输入与输出指数、空间半径及允许的尺度损失，再判断它是否足以处理手中的组合问题。


\begin{chaptersummary}
限制定理的对象是一个从测试函数连续延拓的算子。曲面测度、环境测度和带权输入范数承担不同职责；改变曲面尺度或密度时，相关因子必须重新计算。延拓对偶和正二次型使这些空间之间的关系可以核对，但二次型非负并不意味着卷积核处处非负，也不意味着算子保持函数的点态非负性。

\ProseParagraphBreak

帽集中给出一条可取等号的指数边界，球面尾部则要求另一个严格不等式。Stein--Tomas 在平方曲面输入线上达到帽边界；它的证明保留了曲面结构与核的振荡，不能由取绝对值后的衰减估计替代。径向输入、加权空间或不同频率测度可以改变问题，因此使用结论前应先确定这些附加结构。

\ProseParagraphBreak

最佳常数是否达到，需要有界性之外的紧性。调制、帽集中与多个轮廓的分离是不同损失方式，处理一种并不自动排除其余方式。现代方法用波包及方向几何组织大量帽的相互作用；管形描述还须配合相位、系数和尾部估计。单个模型的计算能够说明尺度，却不能独自证明一般限制猜想。
\end{chaptersummary}

\BookExercises

% 自拟；基础到中等：平面迹的最优加权补偿，求准确常数并证明临界权重失败。
\Needspace{6\baselineskip}
\begin{exercise}\label{b2:ex:23-new-1}
设整数 \(d\geq2\)、\(s\geq0\)，把 \(x\in\mathbb R^d\) 写成 \((x',t)\)。对 \(f\in\mathcal S(\mathbb R^d)\)，令
\[
 R_0f(\xi')=\widehat f(\xi',0),
 \quad
 \|f\|_{H_s}^2=\int_{\mathbb R^d}(1+t^2)^s\cdot|f(x',t)|^2\,\mathrm dx'\,\mathrm dt.
\]
这里 \(H_s\) 表示使上述范数有限的可测函数等价类空间。证明存在有限常数 \(C_s\)，使
\[
 \|R_0f\|_{L^2(\mathbb R^{d-1})}\leq C_s\|f\|_{H_s}
\]
当且仅当 \(s>1/2\)，并在成立时求最佳常数
\[
 C_s=\left(\int_{\mathbb R}(1+t^2)^{-s}\,\mathrm dt\right)^{1/2}.
\]
说明所得延拓怎样定义于全部 \(H_s\)，以及权重如何补偿平面法向的无限相干。
\end{exercise}
\begin{solution}
先对 Schwartz 输入用 Fubini，得到 \(R_0f\) 是函数
\(F(x')=\int_{\mathbb R}f(x',t)\,\mathrm dt\) 关于 \(x'\) 的 Fourier 变换。因此由切向 Plancherel 与加权 Cauchy--Schwarz，
\[
 \begin{aligned}
 \|R_0f\|_2^2
 &=\int\left|\int f(x',t)\,\mathrm dt\right|^2\,\mathrm dx'\\
 &\leq\left(\int(1+t^2)^{-s}\,\mathrm dt\right)
       \int(1+t^2)^s\cdot|f(x',t)|^2\,\mathrm dx'\,\mathrm dt.
 \end{aligned}
\]
第一个积分有限恰好要求 \(2s>1\)。在此条件下，同一个不等式保证一般 \(H_s\) 输入的法向积分对几乎每个 \(x'\) 绝对收敛，并定义一个 \(L^2(\mathbb R^{d-1})\) 函数。对它作平方意义的 Fourier 变换，就是所需延拓；Schwartz 类在加权平方范数中稠密，故该延拓唯一。

为证明常数和必要性，取 \(u\in\mathcal S(\mathbb R^{d-1})\)、\(\|u\|_2=1\)，考虑乘积输入 \(u(x')h(t)\)。记
\[
 I_R=\int_{-R}^R(1+t^2)^{-s}\,\mathrm dt,
 \quad h_R(t)=(1+t^2)^{-s}\cdot\mathbf1_{[-R,R]}(t).
\]
则 \(\int h_R=I_R\)，而 \(\int(1+t^2)^s|h_R|^2=I_R\)，所以相应的输出与输入范数之比为 \(I_R^{1/2}\)。虽然 \(h_R\) 未必光滑，但可在一个固定的稍大紧区间内用 \(C_c^\infty\) 函数逼近；该区间上权重有正上下界，平方逼近也使积分收敛。因此 Schwartz 输入可以任意接近同一比值。

若 \(s>1/2\)，令 \(R\to\infty\)，下界趋于所给 \(C_s\)，与上界一致。若 \(s\leq1/2\)，则 \(I_R\to\infty\)，不可能有统一的估计；其中 \(s=1/2\) 时 \(I_R\) 按 \(\log R\) 增长。这表明临界失败虽缓慢，却不能由增加一个固定常数修复。权重所提供的正是法向积分泛函的有界性，整个论证并未使用曲率。
\end{solution}

% 自拟；中等：可逆仿射变换下的延拓范数，区分推前测度与椭球面积测度的非恒定Jacobian。
\Needspace{6\baselineskip}
\begin{exercise}\label{b2:ex:23-new-2}
设整数 \(d\geq2\)，\(S=S^{d-1}\) 取面积测度 \(\sigma\)，固定
\(2(d+1)/(d-1)\leq q<\infty\)。记球面延拓 \(E:L^2(S)\to L^q(\mathbb R^d)\) 的最佳常数为 \(C_q\)。设 \(A\) 为可逆实矩阵、\(b\in\mathbb R^d\)，令
\[
 F(\xi)=A\xi+b,\quad \Sigma=F(S),\quad D=|\det A|.
\]
在 \(\Sigma\) 上先取推前测度 \(\mu=F_\#\sigma\)，即
\(\int H\,\mathrm d\mu=\int_S H(F(\xi))\,\mathrm d\sigma(\xi)\)。证明相应延拓的最佳 \(L^2(\mu)\to L^q\) 常数恰为 \(D^{-1/q}C_q\)。

\ProseParagraphBreak

再改用 \(\Sigma\) 的面积测度，证明面积因子为
\[
 J_A(\xi)=D|A^{-\mathsf T}\xi|.
\]
若 \(J_{\min}\)、\(J_{\max}\) 是它在球面上的最小值与最大值，证明此时的最佳常数 \(C_{\Sigma,q}\) 满足
\[
 D^{-1/q}J_{\min}^{1/2}C_q
 \leq C_{\Sigma,q}
 \leq D^{-1/q}J_{\max}^{1/2}C_q.
\]
对 \(A=\operatorname{diag}(a_1,\ldots,a_d)\)、\(a_j>0\)，明确写出两个面积因子。这说明只记录行列式是否足以处理原生面积测度。
\end{exercise}
\begin{solution}
令 \(h=g\circ F\)。推前测度的定义给出 \(\|h\|_{L^2(S)}=\|g\|_{L^2(\mu)}\)，且
\[
 E_\mu g(x)=e^{2\pi\mathrm i x\cdot b}\cdot Eh(A^{\mathsf T}x).
\]
换元 \(y=A^{\mathsf T}x\) 后，输出范数等于 \(D^{-1/q}\|Eh\|_q\)。输入的对应是满射等距映射，因而不仅有上界，两个最佳常数也恰好相差 \(D^{-1/q}\)。

在 \(\xi\in S\) 处取一组正交切向单位向量 \(v_1,\ldots,v_{d-1}\)。变换后面积是 \(Av_1,\ldots,Av_{d-1}\) 张成的平行多面体体积。其法向面积向量按余因子矩阵变换，等于 \((\det A)A^{-\mathsf T}\xi\)，至多差一个取向符号；取长度即得 \(J_A\)。也可将任意向量补为最后一列，比较两个行列式，得到同一余因子公式。因此
\[
 \int_\Sigma H(\eta)\,\mathrm d\sigma_\Sigma(\eta)
 =\int_S H(F(\xi))\cdot J_A(\xi)\,\mathrm d\sigma(\xi).
\]

令 \(h=(g\circ F)\sqrt{J_A}\)，则 \(\|h\|_2=\|g\|_{L^2(\Sigma)}\)，而面积延拓变为
\(e^{2\pi\mathrm i x\cdot b}\cdot E(\sqrt{J_A}\,h)(A^{\mathsf T}x)\)。记乘法算子 \(M h=\sqrt{J_A}\,h\)。于是
\[
 C_{\Sigma,q}=D^{-1/q}\|EM\|_{2\to q},
 \quad \|M\|_{2\to2}=J_{\max}^{1/2},
 \quad \|M^{-1}\|_{2\to2}=J_{\min}^{-1/2}.
\]
由 \(EM\) 的复合上界及 \(E=(EM)M^{-1}\)，分别得到题中上界和下界。

对正对角矩阵，\(|A^{-\mathsf T}\xi|\) 在单位球上的最小值为 \(1/\max_j a_j\)，最大值为 \(1/\min_j a_j\)。因此
\[
 J_{\min}=\frac D{\max_j a_j},
 \quad J_{\max}=\frac D{\min_j a_j}.
\]
各向异性变形使面积因子随法向改变，行列式本身没有记录这项变化。若所有 \(a_j=R\)，两个因子才都化为 \(R^{d-1}\)，并恢复球面缩放的准确幂次。
\end{solution}

% 自拟；综合应用：薄球壳频率投影的两种算子范数及其准确平方关系，连接限制、余面积与伴随。
\Needspace{6\baselineskip}
\begin{exercise}\label{b2:ex:23-new-3}
设整数 \(d\geq2\)、\(1\leq p\leq p_*=2(d+1)/(d+3)\)，对 \(0<\varepsilon<1/4\) 令
\[
 \Omega_\varepsilon=\{\,\xi\in\mathbb R^d:1-\varepsilon<|\xi|<1+\varepsilon\,\},
 \quad
 P_\varepsilon f=\mathcal F^{-1}(\mathbf1_{\Omega_\varepsilon}\widehat f).
\]
算子先在 Schwartz 类上定义。证明它有相应的连续延拓，并有
\[
 \|P_\varepsilon\|_{L^p\to L^2}\asymp\varepsilon^{1/2},
 \quad
 \|P_\varepsilon\|_{L^p\to L^{p'}}
   =\|P_\varepsilon\|_{L^p\to L^2}^2
   \asymp\varepsilon.
\]
隐含常数只依赖 \(d,p\)。这里不要求证明 \(P_\varepsilon\) 在同一个 \(L^p\) 空间上有界。
\end{exercise}
\begin{solution}
Plancherel、极坐标和半径缩放后的 Stein--Tomas 估计给出
\[
 \begin{aligned}
 \|P_\varepsilon f\|_2^2
 &=\int_{1-\varepsilon}^{1+\varepsilon}
       \|\widehat f\|_{L^2(rS^{d-1})}^2\,\mathrm dr\\
 &\leq C_{d,p}\|f\|_p^2
       \int_{1-\varepsilon}^{1+\varepsilon}
              r^{d-1-2d/p'}\,\mathrm dr
 \leq C'_{d,p}\varepsilon\|f\|_p^2.
 \end{aligned}
\]
半径始终位于 \((3/4,5/4)\)，所以幂因子一致有界。有限 \(p\) 下的稠密性给出 \(L^p\to L^2\) 延拓。为求下界，固定 \(\widehat f\in C_c^\infty\)，使其在 \(1/2\leq|\xi|\leq3/2\) 上恒为一。这个 \(f\) 是与 \(\varepsilon\) 无关的 Schwartz 函数，而
\[
 \|P_\varepsilon f\|_2^2=|\Omega_\varepsilon|\asymp\varepsilon.
\]
除以固定的 \(\|f\|_p^2\)，即得第一条范数等价。

把已证的 \(L^p\to L^2\) 算子记为 \(B\)。其伴随 \(B^*:L^2\to L^{p'}\) 有同一范数，包括 \(p=1\) 时的 \(L^\infty\) 值域。在测试函数上，实值幂等符号给出
\[
 \langle Bf,Bh\rangle_{L^2}
 =\int_{\Omega_\varepsilon}\widehat f(\xi)\cdot\overline{\widehat h(\xi)}\,\mathrm d\xi
 =\langle P_\varepsilon f,h\rangle.
\]
因此 \(B^*B\) 与原来 \(P_\varepsilon\) 的分布定义一致，提供 \(L^p\to L^{p'}\) 延拓，并满足 \(\|P_\varepsilon\|_{p\to p'}\leq\|B\|^2\)。反过来，对测试输入有
\[
 \|Bf\|_2^2
 =\langle P_\varepsilon f,f\rangle
 \leq\|P_\varepsilon\|_{p\to p'}\|f\|_p^2.
\]
对单位测试输入取上确界，得到反向不等式，故两范数恰有平方关系。这个论证利用了中间空间的 Hilbert 结构；第二条估计并非把第一条的输出指数直接换掉。
\end{solution}

% 自拟；基础核验与反例：显式球面核区分非负二次型与保持点态非负，并识别二次型的非平凡零空间。
\Needspace{6\baselineskip}
\begin{exercise}\label{b2:ex:23-new-4}
在 \(\mathbb R^3\) 中令 \(\sigma\) 为单位球面的面积测度，\(K=\check\sigma\)，并对 \(f\in\mathcal S(\mathbb R^3)\) 置 \(Tf=K*f\)。不用 Bessel 函数，直接证明
\[
 K(x)=\frac{2\sin(2\pi|x|)}{|x|}\quad(x\ne0),
 \quad K(0)=4\pi.
\]
证明 \(\int Tf\cdot\overline f\geq0\)，同时构造非负 \(f\in C_c^\infty\)，使 \(Tf\) 在某个开集内为负。再构造非零 Schwartz 函数满足 \(Tf=0\)。解释这三项性质为何可以同时成立。
\end{exercise}
\begin{solution}
由旋转对称性，可令 \(x=(0,0,r)\)。用球面极坐标并作换元 \(u=\cos\theta\)，得到
\[
 K(x)=2\pi\int_0^\pi e^{2\pi\mathrm i r\cos\theta}\sin\theta\,\mathrm d\theta
 =2\pi\int_{-1}^1e^{2\pi\mathrm i ru}\,\mathrm du
 =\frac{2\sin(2\pi r)}r.
\]
原点极限为 \(4\pi\)，也等于球面面积。对 Schwartz 输入，\(f\in L^1\) 且 \(K\) 有界，下面的积分换序绝对可行：
\[
 \int_{\mathbb R^3}Tf(x)\cdot\overline{f(x)}\,\mathrm dx
 =\int_{S^2}|\widehat f(\xi)|^2\,\mathrm d\sigma(\xi)\geq0.
\]
这里非负的是积分二次型。

取 \(x_0=(0,0,3/4)\)，则 \(K(x_0)<0\)。由连续性，存在 \(\delta>0\)，使 \(K\) 在 \(B(x_0,2\delta)\) 上为负。选非负 \(f\in C_c^\infty(B(0,\delta))\)，满足 \(\int f=1\)。对 \(x\in B(x_0,\delta)\) 和 \(y\in\operatorname{supp}f\)，有 \(x-y\in B(x_0,2\delta)\)，故
\(Tf(x)=\int K(x-y)\cdot f(y)\,\mathrm dy<0\)。所以这个算子不保持函数的点态非负性。

最后取非零 \(b\in C_c^\infty(B(0,1/2))\)，令 \(f=\check b\)。此时 \(f\in\mathcal S\setminus\{0\}\)，但 \(\widehat f|_{S^2}=0\)，并且
\[
 Tf(x)=\int_{S^2}e^{2\pi\mathrm i x\cdot\xi}\widehat f(\xi)\,\mathrm d\sigma(\xi)=0.
\]
二次型只测量 Fourier 变换在球面上的部分，既不检测所有非零输入，也不要求空间核取非负值。因而它的非负性、非平凡零空间和点态变号没有冲突。
\end{solution}

% 自拟；综合证明：径向输入的准确限制范围，由径向对偶求最佳常数并量化临界对数发散。
\Needspace{6\baselineskip}
\begin{exercise}\label{b2:ex:23-new-5}
设整数 \(d\geq2\)、\(1\leq p<\infty\)，只对径向 Schwartz 函数 \(f\) 考虑球面限制。记
\[
 A_d=\sigma(S^{d-1}),
 \quad j(x)=A_d^{-1}\int_{S^{d-1}}e^{2\pi\mathrm i x\cdot\xi}\,\mathrm d\sigma(\xi).
\]
证明存在统一的限制估计
\[
 \|\widehat f|_{S^{d-1}}\|_2\leq C_{d,p}\|f\|_p
\]
当且仅当 \(1\leq p<2d/(d+1)\)。成立时，证明径向类中的最佳常数是 \(A_d^{1/2}\|j\|_{p'}\)。在临界指数 \(p=2d/(d+1)\)，构造径向单位输入，其限制范数趋于无穷；可用截断对偶函数获得至少 \((\log R)^{(d-1)/(2d)}\) 的增长。
\end{exercise}
\begin{solution}
径向性使 \(\widehat f\) 在单位球面上为常数。对球面方向平均，并使用 Fubini，得到
\[
 \widehat f(e_1)=A_d^{-1}\int_{S^{d-1}}\widehat f(\xi)\,\mathrm d\sigma(\xi)
 =\int_{\mathbb R^d}f(x)\cdot j(x)\,\mathrm dx.
\]
球面对称使 \(j\) 为实值径向函数，因此最后一式没有反射或共轭差异。限制范数便是 \(A_d^{1/2}\) 乘这个积分的绝对值，Hölder 给出所述常数作为上界。

对 \(1<p<\infty\)，令 \(q=p'\)。在任意球 \(B(0,R)\) 上，径向函数
\(\operatorname{sgn}(j)|j|^{q-1}\cdot\mathbf1_{B(0,R)}\) 是这一积分泛函的对偶取等函数；零点处取零。归一化后，输入 \(L^p\) 范数为一，积分值为 \(\|j\|_{L^q(B(0,R))}\)。这些紧支撑径向函数可用径向 \(C_c^\infty\) 函数在 \(L^p\) 中逼近，例如用径向光滑核卷积。逼近支撑保持在一个稍大球内，那里 \(j\in L^q\)，所以积分值也收敛。再归一化并令 \(R\to\infty\)，就证明了最佳常数公式；若 \(j\notin L^q\)，同一构造也证明无界。

正文的球面驻相展开给出
\[
 |j(x)|\leq C_d(1+|x|)^{-(d-1)/2},
\]
并在每个充分大的单位径向区间中，找到占固定比例的一部分，使 \(|j(r)|\geq c_dr^{-(d-1)/2}\)。因此对有限 \(q\)，
\[
 j\in L^q(\mathbb R^d)
 \quad\Longleftrightarrow\quad q>\frac{2d}{d-1}.
\]
转回 \(p\) 就是 \(p<2d/(d+1)\)。当 \(p=1\) 时，\(|j|\leq1\) 且 \(j(0)=1\)，故 \(\|j\|_\infty=1\)；取向原点集中的非负径向光滑近似恒等族，积分趋于一，仍得到相同的最佳常数公式。

在临界 \(q_c=2d/(d-1)\) 处，上述上下界还给出
\[
 \int_{|x|<R}|j(x)|^{q_c}\,\mathrm dx\asymp\log R
 \quad(R\geq2).
\]
用前面的截断取等函数，选择其径向光滑逼近，使归一化后的积分至少保留一半，便得到所求径向 Schwartz 单位输入，限制范数至少为 \(c_d(\log R)^{1/q_c}\)。径向假设禁止任意单帽集中，所以必要条件由振荡尾部决定；这解释了本题范围为何能够超过一般输入的 Stein--Tomas 范围。
\end{solution}

% 自拟；中等：紧频率支集的整函数延拓与开集唯一性，给出不依赖PDE唯一延拓定理的证明。
\Needspace{6\baselineskip}
\begin{exercise}\label{b2:ex:23-new-6}
设整数 \(d\geq2\)，\(S\subset\mathbb R^d\) 为非空紧致光滑超曲面，取面积测度，并设 \(g\in L^1(S)\)。令 \(M=\max_{\xi\in S}|\xi|\)。证明延拓积分在 \(z\in\mathbb C^d\) 上仍有定义，并满足
\[
 |Eg(z)|\leq e^{2\pi M|\operatorname{Im}z|}\|g\|_1,
 \quad
 \|\partial^\alpha Eg\|_{L^\infty(\mathbb R^d)}
       \leq(2\pi M)^{|\alpha|}\|g\|_1.
\]
证明它在每一点有全局收敛的复幂级数。由此推出：若 \(Eg\) 在一个非空实开集上恒零，则 \(g=0\) 几乎处处；特别地，非零数据不能产生具有紧空间支集的延拓函数。不得直接调用偏微分方程的唯一延拓定理。
\end{exercise}
\begin{solution}
由于
\(\left|e^{2\pi\mathrm i z\cdot\xi}\right|
=e^{-2\pi\operatorname{Im}z\cdot\xi}\leq e^{2\pi M|\operatorname{Im}z|}\)，
积分绝对收敛，并立刻给出第一条估计。固定 \(z_0\in\mathbb C^d\)，把指数写成
\[
 e^{2\pi\mathrm i z\cdot\xi}
 =e^{2\pi\mathrm i z_0\cdot\xi}
   \sum_{k=0}^\infty\frac{\bigl(2\pi\mathrm i(z-z_0)\cdot\xi\bigr)^k}{k!}.
\]
当 \(|z-z_0|\leq R\) 时，级数绝对值之和至多为
\(e^{2\pi M|\operatorname{Im}z_0|}e^{2\pi MR}\)，乘 \(|g|\) 后可积。因此可以逐项积分，在每个复紧球上一致收敛；再把各次幂按多重指标展开，便得到在整个 \(\mathbb C^d\) 收敛的幂级数。逐项求导或在原积分中求导，得到
\[
 \partial^\alpha Eg(x)
 =\int_S(2\pi\mathrm i\xi)^\alpha
             e^{2\pi\mathrm i x\cdot\xi}g(\xi)\,\mathrm d\sigma(\xi).
\]
对实 \(x\) 取绝对值，就有第二条估计。

若 \(Eg\) 在非空实开集 \(U\) 上为零，取 \(x_0\in U\)。它在 \(x_0\) 的所有实偏导数都为零，因而上面以 \(z_0=x_0\) 展开的全部幂级数系数都为零。全局收敛说明 \(Eg\) 在整个实空间中恒零。

最后须由输出为零回到输入。对任意 \(\varphi\in\mathcal S(\mathbb R^d)\)，Fourier 反演与 Fubini 给出
\[
 \int_{\mathbb R^d}Eg(x)\cdot\widehat\varphi(x)\,\mathrm dx
 =\int_S\varphi(\xi)g(\xi)\,\mathrm d\sigma(\xi).
\]
绝对换序由 \(\|g\|_1\cdot\|\widehat\varphi\|_1<\infty\) 保证。左边为零，说明有限测度 \(g\,\mathrm d\sigma\) 作为分布为零，所以 \(g=0\) 几乎处处。具有紧支集的连续输出在某个非空外部开集为零，故也只能来自零数据。
\end{solution}

% 自拟；选做综合：L²到L∞延拓的全部极值与精确稳定性，补充有限输出指数下的紧性讨论。
\Needspace{6\baselineskip}
\begin{optionalexercise}\label{b2:ex:23-new-7}
设整数 \(d\geq2\)，\(S=S^{d-1}\) 取面积测度，\(A_d=\sigma(S)\)，\(E\) 为球面延拓。证明 \(E:L^2(S)\to L^\infty\) 的最佳常数是 \(A_d^{1/2}\)，且全部单位范数极值函数恰为
\[
 g_{a,\theta}(\xi)=A_d^{-1/2}e^{\mathrm i\theta}e^{-2\pi\mathrm i a\cdot\xi},
 \quad a\in\mathbb R^d,\quad\theta\in\mathbb R.
\]
进一步证明，对每个 \(\|g\|_2=1\)，
\[
 \inf_{a,\theta}\|g-g_{a,\theta}\|_2^2
 =2\left(1-\frac{\|Eg\|_\infty}{A_d^{1/2}}\right),
\]
且此下确界能取到。特别地，若 \(\|Eg\|_\infty\geq(1-\varepsilon)A_d^{1/2}\)，则 \(g\) 距离某个上述极值函数至多 \(\sqrt{2\varepsilon}\)。需要说明输出为何确实在某一点达到其最大模。
\end{optionalexercise}
\begin{solution}
Cauchy--Schwarz 给出 \(|Eg(x)|\leq A_d^{1/2}\|g\|_2\)。所列 \(g_{a,\theta}\) 的平方范数为一，并满足
\(Eg_{a,\theta}(a)=A_d^{1/2}e^{\mathrm i\theta}\)，所以最佳常数恰为 \(A_d^{1/2}\)。

对一般 \(g\in L^2(S)\)，先用光滑函数 \(g_n\) 在 \(L^2(S)\) 中逼近。光滑球面密度的驻相估计说明 \(Eg_n(x)\to0\) 当 \(|x|\to\infty\)，而
\[
 \|Eg_n-Eg\|_\infty\leq A_d^{1/2}\|g_n-g\|_2\longrightarrow0.
\]
因此 \(Eg\) 也是连续且在无穷远趋零的函数。若 \(g\ne0\)，Fourier 变换的单射性保证 \(Eg\) 不恒零；选一个模为正的值，充分大的球外模小于其一半，故最大模在一个紧球内达到。

若单位输入是极值函数，在某个最大点 \(a\) 有 \(|Eg(a)|=A_d^{1/2}\)。Cauchy--Schwarz 的取等条件迫使 \(g\) 与 \(e^{-2\pi\mathrm i a\cdot\xi}\) 成比例，单位范数再确定比例的模，恰得到 \(g_{a,\theta}\)。这证明了分类。

为求距离，令 \(h_a=A_d^{-1/2}e^{-2\pi\mathrm i a\cdot\xi}\)。复内积在第一变量线性，故
\[
 \langle g,h_a\rangle=A_d^{-1/2}Eg(a),
 \quad
 \|g-e^{\mathrm i\theta}h_a\|_2^2
 =2-2\operatorname{Re}\bigl(e^{-\mathrm i\theta}A_d^{-1/2}Eg(a)\bigr).
\]
固定 \(a\)，选相位使实部等于模，得到最小值
\(2-2A_d^{-1/2}|Eg(a)|\)。再选达到输出最大模的 \(a\)，就得到准确公式，并说明下确界取到。最后代入假设并开平方，即得稳定性界。这个极值问题由单点 Cauchy--Schwarz 取等完全决定，与有限输出指数时必须控制整个空间质量的情形有所不同。
\end{solution}

% 自拟；选做综合：Gaussian抛物面波包的全时空显式计算，核对管状尺度、长时间扩散与可积阈值。
\Needspace{6\baselineskip}
\begin{optionalexercise}\label{b2:ex:23-new-8}
设整数 \(n\geq1\)。在全抛物面 \(\{\,(\eta,|\eta|^2):\eta\in\mathbb R^n\,\}\) 上使用参数测度 \(\mathrm d\eta\)，定义
\[
 \mathcal Eg(x,t)=\int_{\mathbb R^n}
          e^{2\pi\mathrm i(x\cdot\eta+t|\eta|^2)}g(\eta)\,\mathrm d\eta.
\]
对 \(\delta>0\)、\(\eta_0\in\mathbb R^n\)，取单位平方范数数据
\[
 g_{\delta,\eta_0}(\eta)=2^{n/4}\delta^{-n/2}
              \cdot e^{-\pi|\eta-\eta_0|^2/\delta^2}.
\]
求其延拓的显式公式，并证明模的空间宽度随时间为
\[
 W_\delta(t)=\delta^{-1}\sqrt{1+4t^2\delta^4},
\]
中心位于 \(x=-2t\eta_0\)。对 \(1\leq q<\infty\)，证明
\[
 \mathcal Eg_{\delta,\eta_0}\in L^q(\mathbb R^{n+1})
 \quad\Longleftrightarrow\quad q>\frac{2(n+1)}n,
\]
且在成立时
\[
 \|\mathcal Eg_{\delta,\eta_0}\|_q
 =C_{n,q}\delta^{n/2-(n+2)/q},\quad0<C_{n,q}<\infty.
\]
解释 \(|t|\lesssim\delta^{-2}\) 时的管状描述，以及更长时间为何不能保持同一管宽。这里的模型不是带面积测度的紧球面。
\end{optionalexercise}
\begin{solution}
换元 \(\eta=\eta_0+\delta u\)，首先验证
\[
 \|g_{\delta,\eta_0}\|_2^2
 =2^{n/2}\int_{\mathbb R^n}e^{-2\pi|u|^2}\,\mathrm du=1.
\]
同一换元将延拓化为一个复 Gaussian 积分。对于 \(\operatorname{Re}a>0\)，有
\[
 \int_{\mathbb R^n}e^{-\pi a|u|^2}e^{2\pi\mathrm i v\cdot u}\,\mathrm du
 =a^{-n/2}e^{-\pi|v|^2/a}\quad(v\in\mathbb R^n).
\]
这里 \(a^{-n/2}\) 取右半平面上在正实轴为正的解析分支。该式先由实 Gaussian 的 Fourier 公式对 \(a>0\) 得到；参数位于右半平面紧集时，积分及参数导数均被一个多项式乘固定 Gaussian 控制，故两侧全纯，恒等定理把等式推广到整个右半平面。

取 \(a=1-2\mathrm i t\delta^2\)、\(v=\delta(x+2t\eta_0)\)，得到
\[
 \begin{aligned}
 \mathcal Eg_{\delta,\eta_0}(x,t)
 &=2^{n/4}\delta^{n/2}
   \cdot e^{2\pi\mathrm i(x\cdot\eta_0+t|\eta_0|^2)}\\
 &\quad\cdot(1-2\mathrm i t\delta^2)^{-n/2}
   \cdot\exp\left(-\frac{\pi\delta^2|x+2t\eta_0|^2}
                       {1-2\mathrm i t\delta^2}\right).
 \end{aligned}
\]
记 \(B(t)=1+4t^2\delta^4\)。取绝对值得
\[
 |\mathcal Eg_{\delta,\eta_0}(x,t)|
 =2^{n/4}\delta^{n/2}B(t)^{-n/4}
       \exp\left(-\frac{\pi\delta^2|x+2t\eta_0|^2}{B(t)}\right).
\]
指数中的二次项给出题述中心和宽度。当 \(|t|\leq c\delta^{-2}\) 时，\(B(t)\) 有依赖 \(c\) 的统一上下界，宽度与 \(\delta^{-1}\) 可比，中心模与 \(\delta^{n/2}\) 可比；这就是有限时间内沿倾斜中心线的一根管。若 \(|t|\gg\delta^{-2}\)，宽度则与 \(\delta|t|\) 可比，不能继续用固定宽度描述整个包。Gaussian 在任何空间点都不严格为零，管外描述的是衰减而非支集。

最后直接积分，避免把单个时刻的界误当成全时空界。对每个 \(t\)，平移中心并作 Gaussian 换元，得到
\[
 \int_{\mathbb R^n}|\mathcal Eg_{\delta,\eta_0}(x,t)|^q\,\mathrm dx
 =2^{nq/4}q^{-n/2}\delta^{nq/2-n}
       (1+4t^2\delta^4)^{-n(q-2)/4}.
\]
再令 \(u=t\delta^2\)，Tonelli 给出
\[
 \|\mathcal Eg_{\delta,\eta_0}\|_q^q
 =2^{nq/4}q^{-n/2}\delta^{nq/2-n-2}
       \int_{\mathbb R}(1+4u^2)^{-n(q-2)/4}\,\mathrm du.
\]
末项有限当且仅当 \(n(q-2)/2>1\)，即 \(q>2(n+1)/n\)；等号时仍有对数发散。积分有限时开 \(q\) 次方，就是所给准确尺度。特别在 \(q=2(n+2)/n\) 时，\(\delta\) 的幂为零，与平方输入的抛物面尺度相符；这一单个 Gaussian 家族的计算本身没有证明对全部输入的延拓估计。
\end{solution}

% !TeX root = ../../Book2.tex
\chapter{Bochner–Riesz 平均}
\label{b2:ch:24}
\BookChapterNumber{b2:ch:24}
\begin{chapterabstract}
Bochner--Riesz 平均通过球边界的消失阶调整频率截断。本章从符号缩放、阶数平均与核的振荡渐近出发，证明可积核范围、尾部必要条件、解析插值充分范围以及带正则性的逼近估计。球面限制提供另一类混合指数控制，Fefferman 定理则揭示高维硬球截断的几何障碍。半径、平滑阶数与输入空间由此呈现出不同的作用。
\end{chapterabstract}

\begin{learninggoals}\begin{itemize}
\item\label{b2:item:24-goal-1}核对半径缩放、阶数平均、Bessel 核常数及两参数极限，分清频率边界平滑与低频逼近阶数。
\item\label{b2:item:24-goal-2}用振荡主项和可控余项判断核的可积性与必要条件，理解端点为何需要下界而不只需要衰减上界。
\item\label{b2:item:24-goal-3}区分可积核范围、已证插值范围、经典目标与深层否定定理，并核对一致有界性在范数收敛中的作用。
\item\label{b2:item:24-goal-4}将限制估计用于径向壳层，保留混合指数和尺度因子；在组合平均或几何逼近中检查常数是否一致。
\end{itemize}\end{learninggoals}

矩形硬截断可以按坐标分解，球形硬截断却同时包含连续变化的法向。给球边界增加消失因子，能够减弱频率截断造成的远场振荡，但需要多少平滑仍取决于所问的空间。平方范数只看符号大小，端点空间要求核有足够的绝对可积性，中间指数则需要更细的抵消与几何估计。

\ProseParagraphBreak

学习本章时，可先分别观察两个参数。半径增大使更多频率进入平均，阶数增大则改变球内部的权重；两者同时变化甚至可以产生热算子或零算子，而不一定逼近恒等。\secref{b2:sec:2401}确定这些算子的定义与缩放，\secref{b2:sec:2402}再说明边界消失阶如何转化为空间核的衰减。

\ProseParagraphBreak

\secref{b2:sec:2403}将必要条件与充分估计放在同一组参数中比较。一个没有被尾部反例排除的指标，并不因此获得有界性；反过来，可积核方法或简单插值无法覆盖它，也不能据此宣布无界。逼近速度还需单独检查符号在原点的 Taylor 展开。组合不同半径可以消去低阶误差，却不会自动消除各条球边界的乘子障碍。

\ProseParagraphBreak

\appref{b2:app:A}的 Bessel 公式将给出径向核的具体表达，核的衰减又决定哪些指数可以直接使用 Young 不等式。最后一节将球内分解成许多薄层，把问题与\chapref{b2:ch:23}的限制估计相联系。Fefferman 定理说明球形硬截断的特殊障碍；其深层几何证明在此仅作概要介绍。

% !TeX root = ../../Book2.tex
\section{Bochner–Riesz 乘子}
\label{b2:sec:2401}

球形频率截断保留低频、删除高频，是一种自然的 Fourier 求和方法。它的困难集中在球面边界：边界两侧的符号突然跳变，逆变换便产生长距离振荡。Bochner–Riesz 平均在球内加入一个向边界逐渐消失的权，用参数 \(\delta\) 记录消失阶数。需要确定的是，对给定的 \(L^p\) 空间，多少平滑足以同时保证算子有界和半径增大时的逼近。

\subsection{定义}
\label{b2:subsec:240101}

本章采用
\(\widehat f(\xi)=\int_{\mathbb R^d}f(x)e^{-2\pi\mathrm i x\cdot\xi}\,dx\)。
先在有界符号能够直接作用的平方空间上定义算子，再讨论其他指数。

\Needspace{6\baselineskip}
\begin{definition}[Bochner--Riesz 平均]\label{b2:def:24-bochner-riesz}
设整数 \(d\geq1\)、\(\delta\geq0\)、\(R>0\)。定义频率符号
\[
 m_{\delta,R}(\xi)\coloneqq
 \begin{cases}
 (1-|\xi|^2/R^2)^\delta,&|\xi|<R,\\
 0,&|\xi|\geq R.
 \end{cases}
\]
相应的\term[bochner-riesz-pingjun]{Bochner--Riesz 平均}[Bochner--Riesz means]是
\[
 B_R^\delta f\coloneqq
       \mathcal F^{-1}(m_{\delta,R}\cdot\widehat f),
 \quad f\in L^2(\mathbb R^d).
\]
单位半径核记为
\[
 K_\delta\coloneqq\mathcal F^{-1}m_{\delta,1}.
\]
这里 \(\delta=0\) 时的符号是开球示性函数；球面上的取值对这个
Lebesgue 乘子没有影响。
\end{definition}
\symidx{\(B_R^\delta f\)：Bochner--Riesz 平均}
\symidx{\(K_\delta\)：单位半径 Bochner--Riesz 核}

由于 \(0\leq m_{\delta,R}\leq1\)，该定义在 \(L^2\) 中总有意义。
对 Schwartz 输入，还有普通积分表达式
\[
 B_R^\delta f(x)
 =\int_{|\xi|<R}e^{2\pi\mathrm i x\cdot\xi}
           (1-|\xi|^2/R^2)^\delta\widehat f(\xi)\,d\xi.
\]
它处处绝对收敛，输出光滑。若要在 \(1\leq p<\infty\) 上定义同一算子，
则须先证明
\[
 \|B_R^\delta f\|_p\leq C\|f\|_p
       \quad(f\in\mathcal S).
\]
Schwartz 类的稠密性随后给出唯一连续延拓，并在 \(L^p\cap L^2\)
上与平方定义一致。无穷指数采用\chapref{b2:ch:22}的伴随约定：
从相同实符号的 \(L^1\) 算子取 Banach 伴随；核可积时，这一约定
就是对有界输入作卷积。不能改用 Schwartz 类在 \(L^\infty\) 中的范数稠密性。

\ProseParagraphBreak
正参数使符号在球面处连续，但不使它在边界任意光滑。
非整数 \(\delta\) 的法向导数在超过相应阶数后出现奇性；
正整数阶也会在更高导数处遇到跨边界的不匹配。
因此，不能把“边界值连续”直接替换成\chapref{b2:ch:22} Mihlin 判据所需的全部导数控制。
稍后将用实际核估计确定有界范围。

\ProseParagraphBreak
这个算子整合球内部的频率，并不是只在球面上积分的延拓算子。
两者的联系来自逐半径分层。先取
\(\chi\in C_c^\infty(\mathbb R^d)\)，使 \(\chi=1\) 于原点附近，且支撑
严格位于单位球内。符号 \(\chi\cdot m_{\delta,1}\) 光滑紧支撑，其逆变换
为 Schwartz 函数，故相应卷积在所有 \(L^p\) 上有界。剩余部分才包含
球面边界的困难。这一区分会在核渐近、解析插值和薄壳分解中反复出现。

\subsection{缩放}
\label{b2:subsec:240102}

半径变化首先是一个准确的共轭关系。为了不混淆正、反缩放，明确规定
\[
 D_Rf(x)\coloneqq f(Rx),\quad R>0.
\]
它满足
\(\widehat{D_Rf}(\xi)=R^{-d}\widehat f(\xi/R)\) 和
\(\|D_Rf\|_p=R^{-d/p}\|f\|_p\)，无穷指数按 \(1/p=0\) 理解。

\Needspace{6\baselineskip}
\begin{proposition}[半径共轭与一致常数]\label{b2:prop:24-dilation}
设 \(d\geq1\)、\(\delta\geq0\)，\(B_R^\delta\)、\(K_\delta\)
如\cref{b2:def:24-bochner-riesz}。则对 Schwartz 函数有
\[
 B_R^\delta=D_R B_1^\delta D_{1/R},\quad
 K_{\delta,R}(x)=R^dK_\delta(Rx).
\]
若 \(B_1^\delta\) 在 \(L^p(\mathbb R^d)\)、\(1\leq p<\infty\)
上有界，则所有 \(R>0\) 的算子都有唯一相容延拓，且
\[
 \|B_R^\delta\|_{L^p\to L^p}
       =\|B_1^\delta\|_{L^p\to L^p}.
\]
在 \(L^\infty\) 上使用 \(L^1\) 伴随延拓时，同一等式成立。
\end{proposition}
\begin{proof}
先对 Schwartz 输入作 Fourier 变换：
\[
 \begin{aligned}
 \mathcal F(D_R B_1^\delta D_{1/R}f)(\xi)
 &=R^{-d}m_{\delta,1}(\xi/R)\,
             R^d\widehat f(\xi)\\
 &=m_{\delta,R}(\xi)\widehat f(\xi).
 \end{aligned}
\]
核公式则直接由逆变换中的 \(\xi=R\eta\) 得到。
两个拉回的范数因子互为倒数，所以共轭给出
\(\|B_R^\delta\|\leq\|B_1^\delta\|\)。反向用
\(B_1^\delta=D_{1/R}B_R^\delta D_R\)，得到反向不等式。
有限指数上的稠密延拓保留恒等式。无穷指数从 \(L^1\) 的共轭恒等式
取伴随，或在核可积时直接换元，得到相同结论。
\end{proof}

所以固定 \(\delta,p\) 的问题不是分别控制每个半径，而是证明一个
单位半径界。这个界一旦成立，就给出整个族的统一常数；缩放本身不能
证明单位半径有界。

\Needspace{6\baselineskip}
\begin{proposition}[有界性推出范数逼近]\label{b2:prop:24-strong-approximation}
设 \(d\geq1\)、\(\delta\geq0\)、\(1\leq p<\infty\)，并假设
\cref{b2:def:24-bochner-riesz}的 \(B_1^\delta\) 有界于
\(L^p(\mathbb R^d)\)。则对每个 \(f\in L^p(\mathbb R^d)\)，
\[
 \lim_{R\to\infty}\|B_R^\delta f-f\|_p=0.
\]
\end{proposition}
\begin{proof}
先取 \(\widehat h\in C_c^\infty\)，支撑包含在固定球中。
当 \(R\) 充分大时，整个支撑远离 \(|\xi|=R\)，在这片固定频率上，
\[
 (m_{\delta,R}-1)\widehat h\longrightarrow0
       \quad\text{于 }C_c^\infty
\]
的所有导数半范数中。事实上，正阶时对
\((1-|\xi|^2/R^2)^\delta\) 求导，分母始终远离零，每个非恒定项
含有趋零的 \(R^{-2}\) 因子；零阶时误差在这个支撑上最终恒零。
固定紧支撑上的导数收敛经 Fourier 反演和分部积分给出
\(B_R^\delta h-h\to0\) 于 Schwartz 拓扑，特别在 \(L^p\) 中收敛。

这样的 \(h\) 在有限指数 \(L^p\) 中稠密。为看清这一点，先以
\(C_c^\infty\) 函数逼近任意输入，再对这个 Schwartz 函数 \(u\)
使用频率截断
\(\widehat h_N=\chi(\xi/N)\widehat u(\xi)\)，其中固定光滑 \(\chi\)
紧支撑并在原点附近等于一。Schwartz 衰减控制截去的频率尾部及各阶导数，
故 \(h_N\to u\) 于 Schwartz 拓扑。

记统一算子界为 \(C\)。给定 \(\varepsilon>0\)，选上述 \(h\)
使 \(\|f-h\|_p<\varepsilon\)。于是
\[
 \|B_R^\delta f-f\|_p
 \leq (C+1)\varepsilon+\|B_R^\delta h-h\|_p.
\]
先令 \(R\to\infty\)，再令 \(\varepsilon\downarrow0\)，得到结论。
\end{proof}

证明中真正把稠密类推广到全部输入的是统一算子界。仅知道光滑输入
最终逼近良好，并不能反推普通 \(L^p\) 的有界性。
另一方面，这是逐输入的范数收敛，也不是算子范数收敛：
在 \(L^2\) 中，\(I-B_R^\delta\) 的符号在球外恒为一、处处模不超过一，
所以
\[
 \|I-B_R^\delta\|_{2\to2}=1\quad\text{对每个 }R>0.
\]
把单位输入的频率放到当前球外，总能留下全部误差。
固定输入的频率尾部会消失，随半径移动的输入却没有这种统一性。

\ProseParagraphBreak
阶数方向则有另一种结构：较高阶可以写成较低阶按半径的正权平均。
记 beta 积分
\[
 \mathrm B(a,b)=\int_0^1t^{a-1}(1-t)^{b-1}\,dt
       =\frac{\Gamma(a)\Gamma(b)}{\Gamma(a+b)},\quad a,b>0.
\]

\Needspace{6\baselineskip}
\begin{proposition}[阶数的单调传递]\label{b2:prop:24-order-monotonicity}
设 \(d\geq1\)、\(0\leq\delta_1<\delta_2\)，置
\(\eta=\delta_2-\delta_1\)。对\cref{b2:def:24-bochner-riesz}
的平均和任意 Schwartz 输入，有
\[
 B_R^{\delta_2}f=
 \frac1{\mathrm B(\delta_1+1,\eta)}
 \int_0^1 B_{R\sqrt t}^{\delta_1}f\,
          t^{\delta_1}(1-t)^{\eta-1}\,dt,\quad R>0.
\]
若 \(1\leq p<\infty\) 且
\(\sup_{R>0}\|B_R^{\delta_1}\|_{p\to p}\leq C\)，则
\(\sup_{R>0}\|B_R^{\delta_2}\|_{p\to p}\leq C\)。
在共同的 Schwartz 测试类上还有点态不等式
\[
 \sup_{R>0}|B_R^{\delta_2}f(x)|
       \leq\sup_{R>0}|B_R^{\delta_1}f(x)|.
\]
\end{proposition}
\begin{proof}
令 \(a=|\xi|^2/R^2\)。在 \(a<1\) 时，右侧的符号为
\[
 \begin{aligned}
 &\frac1{\mathrm B(\delta_1+1,\eta)}
     \int_a^1(1-a/t)^{\delta_1}t^{\delta_1}
                         (1-t)^{\eta-1}\,dt\\
 &\quad=\frac1{\mathrm B(\delta_1+1,\eta)}
     \int_a^1(t-a)^{\delta_1}(1-t)^{\eta-1}\,dt
       =(1-a)^{\delta_2}.
 \end{aligned}
\]
最后用 \(t=a+(1-a)u\) 换元。若 \(a\geq1\)，两侧均为零。
对 Schwartz 输入，可以将半径积分移过逆 Fourier 变换：
符号的模至多为一，\(\widehat f\) 可积，而归一化 beta 权的总质量为一。
这证明积分恒等式及其点态意义。

若较低阶的 \(L^p\) 界已知，\(t\mapsto B_{R\sqrt t}^{\delta_1}f\)
在 \(t>0\) 强连续。由缩放恒等式，这归结为有限指数 \(L^p\)
上拉回 \(D_r\) 的强连续性，后者先对光滑紧支撑函数验证，再用稠密性。
因此加权积分也可在 \(L^p\) 中解释，且其范数由
\(C\|f\|_p\) 控制。Minkowski 与测试函数的稠密延拓给出较高阶的界。
对点态式取绝对值，再用总质量为一的正权积分，得到最后的不等式。
\end{proof}

正权只出现在半径平均中，不说明实空间核处处非负。
例如低阶核本来有正负振荡，按半径作正权平均以后仍可能变号。
这个公式只能把有界性向更高阶传递，不能从大阶数的可积核反推出
小阶数的结论。若要把极大不等式从测试类推广到一般输入，也仍须
使用已建立的极大算子延拓，不能任意选择各半径的代表后取上确界。

\ProseParagraphBreak
固定阶数时 \(R\to\infty\) 逼近恒等；阶数同时变化，则可能得到热流。
这个例子提醒我们，两个参数不能只观察其中一个。

\Needspace{6\baselineskip}
\begin{proposition}[阶数与半径的联合平方极限]\label{b2:prop:24-heat-limit}
设 \(d\geq1\)、\(R_n\to\infty\)、\(\delta_n\geq0\)，并假设
\(\delta_n/R_n^2\to a\in[0,\infty]\)。对每个
\(f\in L^2(\mathbb R^d)\)，\cref{b2:def:24-bochner-riesz}
中的平均满足
\[
 B_{R_n}^{\delta_n}f\longrightarrow
 \begin{cases}
 f,&a=0,\\
 e^{\,a\Delta/(4\pi^2)}f,&0<a<\infty,\\
 0,&a=\infty
 \end{cases}
 \quad\text{于 }L^2.
\]
中间的算子按 Fourier 符号 \(e^{-a|\xi|^2}\) 定义。
\end{proposition}
\begin{proof}
固定 \(\xi\)，当 \(n\) 充分大时 \(|\xi|<R_n\)。
若 \(a<\infty\)，则
\[
 \delta_n\log(1-|\xi|^2/R_n^2)
 =-\frac{\delta_n}{R_n^2}|\xi|^2
       +O\!\left(\frac{\delta_n}{R_n^4}|\xi|^4\right)
       \longrightarrow-a|\xi|^2.
\]
误差趋零，因为 \(\delta_n/R_n^2\) 有界而 \(R_n^{-2}\to0\)。
若 \(a=\infty\) 且 \(\xi\neq0\)，则
\[
 0\leq m_{\delta_n,R_n}(\xi)
       \leq\exp(-\delta_n|\xi|^2/R_n^2)\longrightarrow0.
\]
零频单点不影响 Lebesgue 积分。所有符号均介于零与一之间，
故平方误差由 \(4|\widehat f|^2\) 控制。Plancherel 和支配收敛
给出所述三种强极限。热算子时间是 \(a/(4\pi^2)\)，来自本书
Fourier 变换中 Laplace 算子的符号 \(-4\pi^2|\xi|^2\)。
\end{proof}

特别地，\(R_n=\sqrt{\delta_n}\to\infty\) 产生符号 \(e^{-|\xi|^2}\)。
这是 \(L^2\) 中的强算子极限；有限阶核是否非负、其他
\(L^p\) 上是否存在统一算子界，还取决于近似核本身的性质。

\subsection{核表示}
\label{b2:subsec:240103}

为了读出边界平滑的作用，现在计算单位半径核。
核积分可以比普通平均的参数范围更广：只要 \(\delta>-1\)，
球内幂权仍可积；但负参数的符号在边界无界，不能当作平方有界乘子。

\Needspace{6\baselineskip}
\begin{proposition}[Bochner--Riesz 核]\label{b2:prop:24-kernel}
设 \(d\geq1\)、\(\delta>-1\)，在开单位球内令
\(m_\delta(\xi)=(1-|\xi|^2)^\delta\)，球外取零。
零阶按球示性函数解释。则
\[
 K_\delta(x)\coloneqq\int_{|\xi|<1}
             e^{2\pi\mathrm i x\cdot\xi}m_\delta(\xi)\,d\xi
\]
是连续有界的实值径向函数，对 \(x\neq0\) 满足
\[
 K_\delta(x)=\Gamma(\delta+1)\pi^{-\delta}
       |x|^{-d/2-\delta}J_{d/2+\delta}(2\pi|x|).
\]
在原点，
\[
 K_\delta(0)=
       \frac{\pi^{d/2}\Gamma(\delta+1)}
            {\Gamma(\delta+d/2+1)}.
\]
\end{proposition}
\begin{proof}
极坐标下，球内符号在边界的可积性恰由
\(\int_{1/2}^1(1-r)^\delta\,dr<\infty\) 保证。
因此逆变换连续有界；旋转不变性给出径向性，频率偶对称使正弦部分抵消，
给出实值性。

Bessel 公式使用\cref{b2:thm:A2-bochner-riesz}，其所需条件正是
\(d\geq1,\delta>-1\)。回忆该处计算的关键一步：将径向变换中的
Bessel 级数逐项积分后，出现
\[
 \int_0^1(1-s^2)^\delta s^{2k+2\nu+1}\,ds
 =\frac{\Gamma(k+\nu+1)\Gamma(\delta+1)}
        {2\Gamma(k+\nu+\delta+2)},\quad \nu=d/2-1.
\]
它消去原级数中的一个 Gamma 分母，剩下阶数
\(\nu+\delta+1=d/2+\delta\) 的 Bessel 级数。固定空间半径时，
绝对值级数由可积球内权控制，故这次逐项积分合法。
代回 Fourier 归一化中的 \(2\pi\)，得到所写系数。

原点值也可不借 Bessel 极限直接检查：
\[
 \begin{aligned}
 K_\delta(0)
 &=\frac{2\pi^{d/2}}{\Gamma(d/2)}
             \int_0^1(1-r^2)^\delta r^{d-1}\,dr\\
 &=\frac{\pi^{d/2}}{\Gamma(d/2)}
             \mathrm B(d/2,\delta+1)
 =\frac{\pi^{d/2}\Gamma(\delta+1)}
             {\Gamma(\delta+d/2+1)}.
 \end{aligned}
\]
\end{proof}

例如 \(d=2,\delta=0\) 时
\(K_0(x)=J_1(2\pi|x|)/|x|\)，原点值为 \(\pi\)，正是单位圆盘面积。
核的 Bessel 阶数与球面测度不同：球面公式出现 \(d/2-1\)，
球内加权公式出现 \(d/2+\delta\)。球内部的径向积分及其边界消失阶
正是这一差别的来源。下一节把公式化为可用于绝对积分和反例的远场渐近。

% !TeX root = ../../Book2.tex
\section{核的渐近}
\label{b2:sec:2402}

核公式给出了一个特殊函数表达式，要判断可积性，还必须知道无穷远处的主项及误差。上界能够证明收敛，排除端点则需要下界。本节把球内积分化为两个一维端点的贡献，读出边界消失阶怎样增加空间衰减，再由同一渐近分别处理核范数、逼近恒等和极大控制。

\subsection{Bessel 表示}
\label{b2:subsec:240201}

除 Bessel 级数外，还有一个便于看清几何的切片表示。
设 \(d\geq2,\delta>-1\)，把空间点旋转为 \(x=(0,\ldots,0,r)\)。
频率写为 \((\eta,t)\in\mathbb R^{d-1}\times\mathbb R\)，则
\[
 \begin{aligned}
 K_\delta(r)
 &=\int_{-1}^1e^{2\pi\mathrm i rt}
       \int_{|\eta|^2<1-t^2}
          (1-t^2-|\eta|^2)^\delta\,d\eta\,dt.
 \end{aligned}
\]
所有换序都由球内幂权的绝对可积性保证。
内层令 \(\eta=(1-t^2)^{1/2}u\)，得到
\[
 \int_{|\eta|^2<1-t^2}(1-t^2-|\eta|^2)^\delta\,d\eta
 =c_{d,\delta}(1-t^2)^{\delta+(d-1)/2},
\]
其中用\cref{b2:prop:24-kernel}的原点积分在维数 \(d-1\) 中计算，
\[
 c_{d,\delta}
 =\frac{\pi^{(d-1)/2}\Gamma(\delta+1)}
        {\Gamma(\delta+(d+1)/2)}.
\]
故令 \(\alpha=\delta+(d-1)/2\)，有
\[
 K_\delta(r)=c_{d,\delta}
       \int_{-1}^1e^{2\pi\mathrm i rt}(1-t^2)^\alpha\,dt.
\]
当 \(d=1\) 时，直接从定义得到同一公式，此时 \(c_{1,\delta}=1\)。

\ProseParagraphBreak
横截面的尺度贡献 \((1-t^2)^{(d-1)/2}\)，原来的边界权再贡献
\((1-t^2)^\delta\)。两者相加，给出端点 \(t=\pm1\) 的消失阶
\(\alpha\)。区间内部的振幅光滑，且相位导数不为零，局部化后可
反复分部积分；真正决定主项的是两个端点。

\ProseParagraphBreak
低维模型可以同时检查常数与符号。先取 \(d=1\)，置 \(u=2\pi x\)。
零阶时直接积分得到
\[
 K_0(x)=\int_{-1}^1e^{\mathrm i u\xi}\,d\xi
       =\frac{2\sin u}{u},\quad K_0(0)=2.
\]
对一阶权 \(1-\xi^2\)，在有限区间上对 \(u\) 求导合法，因而
\[
 K_1(x)=\left(1+\frac{d^2}{du^2}\right)
              \frac{2\sin u}{u}
       =\frac{4(\sin u-u\cos u)}{u^3},\quad K_1(0)=\frac43.
\]
后一个原点值可由 Taylor 展开取得，也等于
\(\int_{-1}^1(1-\xi^2)\,d\xi\)。
当 \(u=2\pi\) 时，\(K_1(x)<0\)，所以即使正阶平均的核已可积，
它也不是一个非负的概率核。

\ProseParagraphBreak
在三维，零阶球的横截面积为 \(\pi(1-t^2)\)，于是
\[
 K_0(x)=4\pi\,\frac{\sin u-u\cos u}{u^3},
 \quad u=2\pi|x|,\quad K_0(0)=\frac{4\pi}{3}.
\]
它与一维一阶核只差常数，但绝对积分的径向体积元已经由 \(dr\)
变为 \(4\pi r^2\,dr\)。同一个函数形式在不同底空间有不同的
可积性，判断时不能省去体积增长。

\subsection{端点振荡渐近}
\label{b2:subsec:240202}

所需的一维工具是\cref{b2:lem:A2-endpoint-asymptotic}：
若 \(\beta>-1\)，\(\chi\) 在 \([0,\infty)\) 光滑且紧支撑，则
\[
 \int_0^\infty e^{-2\pi\mathrm i rs}s^\beta\chi(s)\,ds
 =\Gamma(\beta+1)(2\pi r)^{-\beta-1}
       e^{-\mathrm i\pi(\beta+1)/2}\chi(0)
       +O(r^{-\beta-2}).
\]
这里 \(r\to+\infty\)，复相位对应负号 Fourier 指数。
该结论不要求 \(\beta\) 为整数，因而能处理任意实消失阶。

\ProseParagraphBreak
回忆这个工具的证明机制，可以避免把端点公式误当成一次普通分部积分。
模型积分带上 \(e^{-s}\) 后为
\[
 \int_0^\infty e^{-2\pi\mathrm i rs}s^\beta e^{-s}\,ds
       =\Gamma(\beta+1)(1+2\pi\mathrm i r)^{-\beta-1}.
\]
将 \(\chi(s)\) 写成
\(\chi(0)e^{-s}+s\,c(s)\)，余项在零点多消失一阶。
在 \(0\leq s\lesssim r^{-1}\) 上直接积分，得到 \(O(r^{-\beta-2})\)；
在 \(s\asymp2^j/r\) 的剩余小环带上，取整数 \(N>\beta+2\)
次分部积分，得到
\[
 C r^{-\beta-2}2^{j(\beta+2-N)}.
\]
这是可求和的几何级数，远离零点的部分则有任意阶衰减。
这个分解既利用了端点的尺度，又没有对奇异振幅在零点作非法求导。
附录中的证明同时控制了截断导数与无穷远的指数尾部。

\ProseParagraphBreak
现在固定 \(d\geq1,\delta\geq0\)，将切片积分在两个端点与区间内部
作光滑分割。在 \(t=1\) 附近令 \(s=1-t\)，便有
\[
 (1-t^2)^\alpha=s^\alpha(2-s)^\alpha,\quad
 e^{2\pi\mathrm i rt}=e^{2\pi\mathrm i r}e^{-2\pi\mathrm i rs}.
\]
端点振幅在 \(s=0\) 的值为 \(2^\alpha\)。上述公式给出的正端点主项为
\[
 c_{d,\delta}\,2^\alpha\Gamma(\alpha+1)
       (2\pi r)^{-\alpha-1}
       e^{\,2\pi\mathrm i r-\mathrm i\pi(\alpha+1)/2}.
\]
负端点贡献是它的共轭，内部项经分部积分可忽略到所需阶数。
代入 \(\alpha=\delta+(d-1)/2\)，整理常数得到
\[
 \begin{aligned}
 K_\delta(r)
 &=a_\delta r^{-\delta-(d+1)/2}
       \cos\!\left(2\pi r-\frac{\pi(d+1+2\delta)}4\right)
       +O(r^{-\delta-(d+3)/2}),\\
 a_\delta&=\frac{\Gamma(\delta+1)}{\pi^{\delta+1}}>0.
 \end{aligned}
\]
这也与\cref{b2:prop:24-kernel}中 Bessel 阶数
\(\nu=d/2+\delta\) 的大半径渐近一致。

\ProseParagraphBreak
每增加一阶 \(\delta\)，空间主项便多衰减一个幂次。
在固定维数下，若 \(\delta\) 限制在 \([0,\infty)\) 的一个固定紧区间，
这里的余项常数还可以一致选取：端点分割固定，光滑因子
\((2-s)^\alpha\) 的有限阶导数一致有界；环带估计中的 \(N\) 取为
大于全部 \(\alpha+2\) 的同一个整数，幂函数积分的分母也远离零。
Gamma 系数在该参数紧集连续有界。这种局部参数一致性可用于研究
可积性阈值附近的范数变化，却不是允许阶数无界增长的渐近结论。

\ProseParagraphBreak
余弦零点并不会使主项在大部分半径上消失。令
\[
 \beta_\delta=\delta+\frac{d+1}{2},\quad
 \phi_\delta=\frac{\pi(d+1+2\delta)}4.
\]
在每个充分大的单位区间 \([n,n+1]\) 内，集合
\[
 A_n=\{r\in[n,n+1]:|\cos(2\pi r-\phi_\delta)|\geq3/4\}
\]
都有同一个正长度。余项比主振幅小一个 \(r^{-1}\) 因子，
因此当 \(n\) 足够大时，
\[
 |K_\delta(r)|\geq c_{d,\delta}r^{-\beta_\delta}
       \quad(r\in A_n).
\]
上界处理整个空间，下界只需占固定比例的这些径向区域。
下面将这两种用途分开应用。

\subsection{空间衰减}
\label{b2:subsec:240203}

\Needspace{6\baselineskip}
\begin{proposition}[核的有限指数可积性]\label{b2:prop:24-kernel-tail}
设 \(d\geq1,\delta\geq0\)，\(K_\delta\) 为
\cref{b2:def:24-bochner-riesz}的单位半径核。
对每个 \(1\leq q<\infty\)，有
\[
 K_\delta\in L^q(\mathbb R^d)
 \quad\Longleftrightarrow\quad
 q\left(\delta+\frac{d+1}{2}\right)>d.
\]
核在全空间有界；上述有限指数条件在等号时不成立。
\end{proposition}
\begin{proof}
核在有界集合上连续有界。远场上界给出
\[
 \int_{|x|>1}|K_\delta(x)|^q\,dx
 \leq C\int_1^\infty r^{d-1-q\beta_\delta}\,dr,
 \quad \beta_\delta=\delta+(d+1)/2.
\]
当 \(q\beta_\delta>d\) 时收敛。
反过来，在前述 \(A_n\) 上使用下界及极坐标，
\[
 \int_{|x|>1}|K_\delta(x)|^q\,dx
       \geq c\sum_{n\geq n_0}n^{d-1-q\beta_\delta}.
\]
这里 \(r\asymp n\)，且每个 \(A_n\) 的长度都有固定正下界。
若 \(q\beta_\delta\leq d\)，这个级数发散；临界等号时就是调和级数。
\end{proof}

这只判断一个核函数的 \(L^q\) 性质。乘子在 \(L^p\) 上有界是对
全部输入的要求，两者不能互换。核的 \(L^1\) 可积性能够通过 Young
提供所有指数的充分条件，但不是每个中间指数的必要条件。

\Needspace{6\baselineskip}
\begin{theorem}[可积核范围]\label{b2:thm:24-integrable}
设 \(d\geq1,\delta\geq0\)，\(B_R^\delta,K_\delta\)
如\cref{b2:def:24-bochner-riesz}。则
\[
 K_\delta\in L^1(\mathbb R^d)
       \quad\Longleftrightarrow\quad \delta>\frac{d-1}{2}.
\]
在这一范围，算子由 \(K_{\delta,R}*f\) 定义于所有
\(L^p(\mathbb R^d)\)、\(1\leq p\leq\infty\)，并有
\[
 \sup_{R>0}\|B_R^\delta f\|_p
       \leq\|K_\delta\|_1\cdot\|f\|_p.
\]
对 \(p<\infty\)，\(B_R^\delta f\to f\) 于 \(L^p\)；
对有界一致连续函数，收敛在全空间一致范数中成立。
\end{theorem}
\begin{proof}
在上一命题中取 \(q=1\)，得到准确的核可积阈值。
缩放不改变核的 \(L^1\) 范数，Young 不等式遂给出算子界。
这个卷积与 Schwartz 类上的 Fourier 定义一致，有限指数上的延拓
因而相容；无穷指数也与 \(L^1\) 伴随约定一致，因为核实值且为偶函数。

还须核对逼近恒等的总质量。符号与核均可积，Fourier 反演表明
\(\widehat K_\delta=m_{\delta,1}\)。二者均在原点连续，故
\[
 \int_{\mathbb R^d}K_\delta(y)\,dy=m_{\delta,1}(0)=1.
\]
对任何 \(\rho>0\)，缩放后的绝对尾质量满足
\[
 \int_{|x|>\rho}|K_{\delta,R}(x)|\,dx
       =\int_{|y|>R\rho}|K_\delta(y)|\,dy\longrightarrow0.
\]
这些是可以变号的逼近恒等条件。

更直接地，在 \(L^p\) 中有
\[
 B_R^\delta f-f
       =\int K_\delta(y)\,[f(\,\cdot-y/R)-f]\,dy.
\]
有限指数的平移连续性使每个固定 \(y\) 对应的差范数趋零，
且它至多为 \(2\|f\|_p\)。用 \(|K_\delta|\) 作可积控制，
Minkowski 与支配收敛即得范数极限。对有界一致连续输入，
同样使用平移的一致范数连续性，得到一致极限。
\end{proof}

在 \(L^\infty\) 上只陈述有界一致连续输入的收敛，不能把它扩张为
所有有界函数的一致逼近。核卷积可以平滑一个跳跃，却不能在一致范数
中把跳跃的两侧同时逼近到任意精度。后面的习题将把这个边界层算出来。
总积分为一给出平移差公式，绝对值可积性控制误差，因此核允许变号。

\ProseParagraphBreak
可积阈值以上的径向衰减还控制全部半径的极大族。
设 \(\varepsilon=\delta-(d-1)/2>0\)，则
\[
 |K_\delta(x)|\leq C(1+|x|)^{-d-\varepsilon}.
\]
因而本问题满足\cref{b2:prop:17-approximate-identity}的具体核假设。
为看清严格正余量的用途，可以直接将卷积绝对值拆成小球与二进环带：
对 \(R>0\)，小球 \(|y|\leq R^{-1}\) 的贡献至多为 \(CMf(x)\)；
在
\(2^{j-1}R^{-1}<|y|\leq2^jR^{-1}\)、\(j\geq1\)，核大小与球体积相乘后
给出 \(C2^{-j\varepsilon}\)。于是
\[
 |B_R^\delta f(x)|
 \leq C\left(1+\sum_{j\geq1}2^{-j\varepsilon}\right)Mf(x).
\]

\Needspace{6\baselineskip}
\begin{corollary}[极大控制与 Lebesgue 点求和]\label{b2:cor:24-maximal}
设 \(d\geq1\)、\(\delta>(d-1)/2\)、\(1\leq p<\infty\)、
\(f\in L^p(\mathbb R^d)\)。采用核积分的代表，令
\[
 B_*^\delta f(x)\coloneqq\sup_{R>0}|B_R^\delta f(x)|.
\]
则 \(B_*^\delta f\leq C_{d,\delta}Mf\) 几乎处处。
它在 \(p>1\) 时强型有界，在 \(p=1\) 时为弱 \((1,1)\) 型。
在 \(f\) 的每个 Lebesgue 点 \(x\)，有
\[
 B_R^\delta f(x)\longrightarrow f(x)\quad(R\to\infty).
\]
\end{corollary}
\begin{proof}
前面的环带估计在 \(Mf(x)<\infty\) 处同时保证绝对收敛和极大上界。
对固定这样的 \(x\)，半径在任意紧区间内变化时，核函数连续，
并有与区间有关的共同可积控制；可用相同的环带和来作支配收敛。
所以 \(R\mapsto B_R^\delta f(x)\) 连续，上确界可取遍正有理半径，
从而是可测函数。极大函数的强、弱型界给出相应算子界。

若 \(x\) 为 Lebesgue 点，则 \(Mf(x)<\infty\)：小球平均由
Lebesgue 点性质控制，大球平均由 \(L^p\) 范数和 Hölder 控制，
其余半径留在紧区间。将卷积减去 \(f(x)\)，在每个环带上改用
\(|f(x-y)-f(x)|\) 的球平均。固定 \(j\) 时，球半径 \(2^j/R\to0\)，
该平均趋零；所有这些平均又被 \(Mf(x)+|f(x)|\) 控制。
由于 \(\sum_j2^{-j\varepsilon}<\infty\)，先控制级数尾，再令 \(R\to\infty\)，
就得到点态极限。
\end{proof}

如果把 \(\delta\) 降到核可积阈值，环带权不再可和，上述证明停止。
某个固定半径算子仍可能在中间指数上有界，但那并不自动控制极大族。
范数收敛、某个子列的几乎处处收敛与所有半径的几乎处处收敛，
需要各自的依据。

\ProseParagraphBreak
最后把两个端点的常数算清。对任意 \(K\in L^1(\mathbb R^d)\)，
卷积算子在 \(L^1\) 与 \(L^\infty\) 上的准确范数都是 \(\|K\|_1\)。
Young 已给出上界。在 \(L^1\) 上，取非负积分一的光滑近似恒等族
\(\rho_n\)，则
\[
 \|\rho_n\|_1=1,\quad K*\rho_n\to K\text{ 于 }L^1,
\]
这给出反向范数不等式。在 \(L^\infty\) 上，取
\[
 f(-y)=
 \begin{cases}
 \overline{K(y)}/|K(y)|,&K(y)\neq0,\\
 0,&K(y)=0.
 \end{cases}
\]
于是 \(\|f\|_\infty\leq1\) 且 \((K*f)(0)=\|K\|_1\)。
因为 \(K\) 的 \(L^1\) 平移连续，
\[
 |(K*f)(x+h)-(K*f)(x)|
       \leq\|K(\,\cdot+h)-K\|_1\cdot\|f\|_\infty,
\]
输出是连续函数，零点的值也给出本性上确界的下界。

\ProseParagraphBreak
应用于本章的可积核，得到
\[
 \|B_R^\delta\|_{1\to1}
       =\|B_R^\delta\|_{\infty\to\infty}
       =\|K_\delta\|_1>1.
\]
最后一个严格不等式来自振荡渐近：在无穷多径向区间内，
余弦主项为负且支配余项，所以核在一个正测度集合上为负。
核实值且总积分为一，因此
\(\int|K_\delta|=1+2\int(K_\delta)_- >1\)。
热核和 Poisson 核的非负性会给端点压缩；Bochner–Riesz 的总质量
归一化不具有同样的结论。

\ProseParagraphBreak
一维的普通指数范围至此可以完整辨认。零阶区间投影在
\(1<p<\infty\) 上由 Hilbert 变换有界，正阶再由
\cref{b2:prop:24-order-monotonicity}取得同一范围；
\(p=1,\infty\) 则要求正阶，核可积性给出充分性，下一节的尾部检验
给出必要性。高维中，仅靠这些事实仍留下大量中间指数，
这就是 Bochner–Riesz 问题的内容。

% !TeX root = ../../Book2.tex
\section{Bochner–Riesz 问题}
\label{b2:sec:2403}

可积核条件给出同时适用于所有 \(p\) 的范围，却未必给出固定 \(p\) 所需的最小平滑量。中心指数 \(p=2\) 只要求符号有界；离开中心后，核的尾部和不同方向的集中分别产生障碍。本节先给出可以直接检验的必要条件，再证明一个解析插值充分范围，最后在已知乘子有界性的前提下量化逼近速度。必要条件、充分条件和逼近阶数承担不同职责，不能用其中一个代替另外两个。

\subsection{\texorpdfstring{\(L^2\)}{L2} 理论}
\label{b2:subsec:240301}

设 \(d\geq1\)，固定 \(\delta\geq0\)。对每个 \(R>0\)，\cref{b2:def:24-bochner-riesz}中的符号
\(m_{\delta,R}(\xi)=(1-|\xi|^2/R^2)_+^\delta\)
在球内介于零和一之间，在原点邻域任意接近一；零阶时它是球的示性函数。因此符号的本性上确界恰为 \(1\)，\cref{b2:prop:22-l2}给出
\[
 \|B_R^\delta\|_{2\to2}=1,\quad
 \lim_{R\to\infty}\|B_R^\delta f-f\|_2=0
 \quad(f\in L^2).
\]
收敛的具体理由是 \(m_{\delta,R}(\xi)\to1\)，而
\(|m_{\delta,R}-1|^2\cdot|\widehat f|^2\leq|\widehat f|^2\in L^1\)；
先在 Fourier 一侧用支配收敛，再用 Plancherel 返回空间侧。
但 \(L^2\) 范数界不控制频率边界上的点态取值，也不自动控制 \(L^p\)、\(p\ne2\)。与限制问题的联系需要额外结构。

\ProseParagraphBreak
对于一维 \(\delta=0\)，球就是区间，Hilbert 变换给出所有 \(1<p<\infty\) 上的界。高维球不再是有限个半轴投影的乘积。矩形投影可以逐坐标应用 Hilbert 变换并用 Fubini 得到有界性；球边界的法向连续变化，出现了无限方向的组合。

\ProseParagraphBreak
平方理论还允许量化平均逼近的速度，而且显示有无边界消失阶会改变速度结构。
\Needspace{6\baselineskip}
\begin{proposition}\label{b2:prop:24-l2-rate}
设整数 \(d\geq1\)，\(B_R^\delta\) 是 \(L^2(\mathbb R^d;\mathbb C)\) 上符号为 \((1-|\xi|^2/R^2)_+^\delta\) 的 Fourier 乘子，\(R>0\)；零阶符号取开球的示性函数。设 \(\delta>0\)、\(0<s\leq2\)，且
\(f\in L^2(\mathbb R^d;\mathbb C)\) 满足
\((2\pi|\xi|)^s\cdot\widehat f(\xi)\in L^2\)。
以这一符号定义 \((-\Delta)^{s/2}f\)，则
\[
 \|B_R^\delta f-f\|_2
 \leq C_\delta(2\pi R)^{-s}\cdot\|(-\Delta)^{s/2}f\|_2.
\]
若 \(\delta=0\)，则对任意 \(s>0\) 同一结论成立，常数可以取一。
\end{proposition}
\begin{proof}
对 \(a\geq0\)，初等均值估计给出
\[
 |1-(1-a)_+^\delta|\leq\max\{1,\delta\}\min\{1,a\}.
\]
对 \(0<\delta<1\)，在 \(0\leq a\leq1\) 上有
\((1-a)^\delta\geq1-a\)，所以误差不超过 \(a\)；对 \(\delta\geq1\)，同一区间上的导数上界为 \(\delta\)。当 \(a\geq1\)，误差等于一。令 \(a=|\xi|^2/R^2\)。因为 \(0<s\leq2\)，有
\(\min\{1,a\}\leq a^{s/2}\)，所以
\[
 |m_{\delta,R}(\xi)-1|
 \leq C_\delta R^{-s}|\xi|^s.
\]
将两边乘以 \(|\widehat f|\)，取平方积分并用 Plancherel，即得第一项；因 \((-\Delta)^{s/2}\) 的符号是 \((2\pi|\xi|)^s\)，常数中出现 \((2\pi R)^{-s}\)。零阶硬截断的误差只在 \(|\xi|>R\)，在那里
\(\mathbf1_{\{|\xi|>R\}}\leq(|\xi|/R)^s\)
对任意 \(s>0\) 成立，故得到第二项。
\end{proof}
对于 \(\delta>0\)，在一个固定紧频带内有
\(R^2[1-(1-|\xi|^2/R^2)^\delta]\to\delta|\xi|^2\)。
若非零 \(\widehat f\) 支于该频带，支配收敛于是给出
\[
 R^2(B_R^\delta f-f)\longrightarrow
 -\frac{\delta}{4\pi^2}(-\Delta)f
 \quad\text{于 }L^2.
\]
因此通常不能仅靠输入具有更多光滑性把第一项误差提高到 \(R^{-s}\)、\(s>2\)。要提高阶数，应组合不同平均以消去 Taylor 首项，或者采用在低频附近恰好等于一的另一种截断。频率边界的平滑与低频的逼近阶数是两个不同因素。

\ProseParagraphBreak

这个例子也说明“越平滑，逼近必然越快”不是无条件原则。硬截断在固定紧频带上最终误差为零；正阶平均虽然在球边界更光滑，却在球内部偏离一，因而留下低阶 Taylor 误差。硬截断的困难主要是其他 \(L^p\) 的一致有界性和逐点求和；平方逼近本身则十分直接。

\subsection{临界指标}
\label{b2:subsec:240302}

\Needspace{6\baselineskip}
\begin{proposition}[尾部必要条件]\label{b2:prop:24-necessary}
设整数 \(d\geq1\)、\(\delta\geq0\)、\(1\leq p\leq\infty\)，\(B_1^\delta\) 为符号 \((1-|\xi|^2)_+^\delta\) 的 \(L^2\) Fourier 乘子，零阶符号取单位开球的示性函数。若它具有与 Schwartz 输入上的 Fourier 定义一致的有界延拓 \(L^p(\mathbb R^d;\mathbb C)\to L^p(\mathbb R^d;\mathbb C)\)，则在 \(p\ne2\) 时必须满足
\[
 \delta>d\left|\frac1p-\frac12\right|-\frac12.
\]
当右侧为负数时，该条件不提供比 \(\delta\ge0\) 更强的限制。
\end{proposition}
\begin{proof}
选 \(\widehat f\in C_c^\infty\)，在闭单位球附近等于 \(1\)。则 \(f\) 为 Schwartz 函数，且
\[
 B_1^\delta f=K_\delta
 =\mathcal F^{-1}\bigl[(1-|\xi|^2)_+^\delta\bigr].
\]
所以对 \(p<\infty\)，算子有界性迫使 \(K_\delta\in L^p\)。
\cref{b2:prop:24-kernel-tail}给出的精确核可积性条件为
\[
 p\bigl(\delta+(d+1)/2\bigr)>d,
 \quad\text{即}\quad
 \delta>\frac dp-\frac{d+1}{2}.
\]
恢复严格不等式的来源：核主项的余弦在每个足够大的单位径向区间内都有一段固定比例集合保持模下界，余项比主项低一阶。故径向绝对积分至少为下列级数的某个正常数倍：
\[
 \sum_{n\geq n_0}n^{d-1-p(\delta+(d+1)/2)}.
\]
阈值等号时正是发散的调和级数。

符号为实值，在 Schwartz 类上有
\(\langle B_1^\delta f,g\rangle=\langle f,B_1^\delta g\rangle\)。
若 \(1<p<\infty\)，由此和 \(L^p\) 对偶性，\(p\) 上的界也给出 \(p'\) 上相同的界。分别应用刚得的尾部条件，再取较强者，得到
\[
 \delta>\max\left\{\frac dp,\frac d{p'}\right\}
 -\frac{d+1}{2}
 =d\left|\frac1p-\frac12\right|-\frac12.
\]

若 \(p=1\)，直接得到 \(\delta>(d-1)/2\)。若 \(p=\infty\)，对 Schwartz 函数 \(f,g\) 使用实符号配对和假设的 \(L^\infty\) 界，有
\[
 |\langle B_1^\delta f,g\rangle|
 =|\langle f,B_1^\delta g\rangle|
 \leq C\|f\|_1\cdot\|g\|_\infty.
\]
对紧支撑光滑 \(g\) 的单位球取上确界，便得到
\(\|B_1^\delta f\|_1\leq C\|f\|_1\)，于是也可应用 \(p=1\) 的必要条件。给定延拓在 Schwartz 类上的配对关系与 \(L^1\) 范数的对偶公式，便把无穷指数的问题转到可积端点。对于 \(p=2\)，尾部检验给出 \(\delta>-1/2\)，所有 \(\delta\geq0\) 都满足。
\end{proof}
严格性来自径向调和级数；没有对尾部作精确下界，就不能宣称阈值等号也失败。这个检验并未排除临界函数
\(\max\{d|1/p-1/2|-1/2,0\}\) 中为零的内部指数上的硬截断；下一小节的深层几何定理才排除高维硬球。

\subsection{Bochner–Riesz 猜想}
\label{b2:subsec:240303}

经典 Bochner–Riesz 目标是：对 \(d\ge2\)、\(1<p<\infty\)，在
\[
 \delta>\max\left\{d\left|\frac1p-\frac12\right|-\frac12,\ 0\right\}
\]
时取得 \(L^p\) 有界性。这里严格的正平滑量将硬球情形排除在目标外。维数、端点形式、径向输入限制以及局部版本各是不同问题。

\ProseParagraphBreak

下面采用\term[fefferman-qiuchengzi-dingli]{Fefferman 球乘子定理}[Fefferman ball multiplier theorem]的下述版本。

\Needspace{6\baselineskip}
\begin{theorem}[深层结构定理：Fefferman 球乘子定理]\label{b2:thm:24-fefferman}
设整数 \(d\geq2\)，\(1<p<\infty\)。在 \(L^2(\mathbb R^d;\mathbb C)\) 上以符号
\(\mathbf1_{\{\,|\xi|<1\,\}}\) 定义单位球 Fourier 乘子。它具有与 Schwartz 类上一致的有界延拓
\(L^p(\mathbb R^d;\mathbb C)\to L^p(\mathbb R^d;\mathbb C)\)，当且仅当 \(p=2\)。
\end{theorem}
\begin{proofsketch}
假设球乘子在某个非二指数有界。对球作放大、频移并取极限，配合随机符号，可推出各方向半空间投影的一致向量值界。Besicovitch 型细矩形构造使输入的平方函数质量分散，而输出在一个小测度集合中累积；对 \(p>2\) 这与向量值界矛盾。对偶性处理 \(p<2\)，切片传递处理高维。
几何矩形构造及传递论证的详细证明见原文
\cite[Theorem 1; Section 2, pp.~330--336]{Fefferman1971Ball}。
\end{proofsketch}
这个定理并不说任何正的 \(\delta\) 都足够。球乘子的失败与正平滑后的范围之间，仍隔着一个依赖维数与 \(p\) 的估计问题。

\subsection{解析插值与 \texorpdfstring{\(L^p\)}{Lᵖ} 逼近估计}
\label{b2:subsec:240304}

可积核给出了 \(L^1\) 和 \(L^\infty\) 两个强端点，平方理论给出了中心指数。它们对应不同的平滑阶数；要连接它们，必须同时改变算子中的参数，而不能只对一个未知有界的固定算子插值。先证明复参数竖线上的核估计，再把参数与输入指数联系起来。

\Needspace{6\baselineskip}
\begin{lemma}[固定实部的复阶 Bochner--Riesz 核]\label{b2:lem:24-complex-kernel}
设整数 \(d\geq1\)，固定实数 \(\sigma>-1\)。对 \(z=\sigma+\mathrm it\)、\(t\in\mathbb R\)，令
\[
 m_z(\xi)=
 \begin{cases}
  \exp\bigl(z\log(1-|\xi|^2)\bigr),&|\xi|<1,\\
  0,&|\xi|\geq1,
 \end{cases}
 \quad K_z=\mathcal F^{-1}m_z,
\]
其中球内使用实对数。符号 \(m_z\) 属于 \(L^1\)，所以 \(K_z\) 是有界连续函数。存在只依赖 \(d,\sigma\) 的常数 \(C\)，使
\[
 |K_{\sigma+\mathrm it}(x)|
 \leq C e^{C(1+t^2)}
 (1+|x|)^{-\sigma-(d+1)/2},
 \quad x\in\mathbb R^d,\ t\in\mathbb R.
\]
特别地，若 \(\sigma>(d-1)/2\)，则
\(\|K_{\sigma+\mathrm it}\|_1\leq C e^{C(1+t^2)}\)。
\end{lemma}
\begin{proof}
记 \(\nu=d/2+z=u+\mathrm it\)，其中 \(u=d/2+\sigma>-1/2\) 固定。\subsecref{b2:subsec:230402}已经以整函数 \(P=1/\Gamma\) 构造复阶 Bessel 级数，并证明归一化球内分布的核恒等式。这里 \(\operatorname{Re}z>-1\)，是普通可积符号，乘回 Gamma 因子即得
\[
 K_z(x)=\Gamma(z+1)\pi^{-z}
 |x|^{-\nu}J_\nu(2\pi|x|),\quad x\ne0.
\]
这是未归一化平均的核，因子 \(\Gamma(z+1)\) 必须保留。复 Gamma 的积分定义给出
\[
 |\Gamma(\sigma+1+\mathrm it)|
 \leq\int_0^\infty e^{-v}v^\sigma\,\mathrm dv
 =\Gamma(\sigma+1).
\]
因此分子 Gamma 因子在这条竖线上甚至一致有界。

小半径与初值由级数控制。对 \(0<r\leq1\)，
\[
 r^{-\nu}J_\nu(2\pi r)
 =\pi^\nu\sum_{k=0}^\infty
 \frac{(-\pi^2r^2)^k}{k!\Gamma(k+\nu+1)}.
\]
因为 \(u+1>0\)，Gamma 递推给出
\[
 |\Gamma(k+\nu+1)|^{-1}
 \leq
 \frac{|P(u+1+\mathrm it)|}
 {\prod_{\ell=0}^{k-1}(u+1+\ell)}.
\]
分母是增长的正实乘积；配合 \(k!\)，级数及其任意固定阶 \(r\) 导数均可在 \(0\leq r\leq1\) 一致求和。前章从 Gamma 无限乘积得到
\[
 |P(u+1+\mathrm it)|\leq C_u e^{C_u(1+t^2)}.
\]
于是 \(|r^{-\nu}J_\nu(2\pi r)|\leq C e^{C(1+t^2)}\)，并且
\[
 |J_\nu(2\pi)|+
 \left|\frac{\mathrm d}{\mathrm dr}J_\nu(2\pi r)
 \bigg|_{r=1}\right|
 \leq C e^{C(1+t^2)}.
\]
对 \(r^\nu\) 求导带来的 \(1+|\nu|\) 因子也可并入此界。前面的 Gamma 分子界因而给出 \(0\leq|x|\leq1\) 的目标估计，原点按连续极限解释。

对 \(r\geq1\)，令
\[
 w(r)=r^{1/2}J_\nu(2\pi r),\quad
 E(r)=|w'(r)|^2+4\pi^2|w(r)|^2.
\]
复阶 Bessel 方程仍给出
\[
 w''+4\pi^2w=(\nu^2-1/4)r^{-2}w.
\]
能量 \(E\) 是实的非负量，不需要系数 \(\nu^2\) 为实数。微分并消去自由振荡项后，
\[
 |E'(r)|
 \leq C|\nu^2-1/4|r^{-2}\cdot E(r)
 \leq C_u(1+t^2)r^{-2}\cdot E(r).
\]
积分 \(\int_1^\infty r^{-2}\,\mathrm dr=1\)，结合已证的 \(r=1\) 初值界，得到
\[
 E(r)\leq E(1)e^{C_u(1+t^2)}
 \leq C e^{C(1+t^2)},\quad r\geq1.
\]
因此
\[
 |J_\nu(2\pi r)|
 \leq C e^{C(1+t^2)}r^{-1/2}.
\]
核公式中另外的幂因子模为 \(\pi^{-\sigma}r^{-u}\)，所以大半径衰减为
\(r^{-u-1/2}=r^{-\sigma-(d+1)/2}\)。
与小半径估计合并即得结论。最后，当
\(\sigma+(d+1)/2>d\) 时，径向积分可积，便得到统一的 \(L^1\) 核界。
\end{proof}

与球面限制中的竖线不同，此处 \(\nu\) 的实部是 \(d/2+\sigma\)。能量法仍然适用，因为对固定实部有
\(|\nu^2-1/4|\leq C(1+t^2)\)；不能直接把只对固定实阶陈述的 Bessel 渐近当作这一统一估计。另一个区别在插值的目标：本节要得到 \(L^p\to L^p\)，而\cref{b2:lem:23-compact-analytic-interpolation}的陈述给出的是 \(L^p\to L^{p'}\)。下面保留其共同定义域的做法，重新选择测试函数指数。

\Needspace{6\baselineskip}
\begin{theorem}[可直接证明的插值范围]\label{b2:thm:24-interpolation-range}
设整数 \(d\geq1\)、\(1<p<\infty\)，实数 \(\delta\) 满足
\[
 \delta>(d-1)\left|\frac1p-\frac12\right|,
\]
并在 Schwartz 类上以
\[
 B_R^\delta f=\mathcal F^{-1}
 \bigl[(1-|\xi|^2/R^2)_+^\delta\cdot\widehat f(\xi)\bigr],
 \quad R>0,
\]
定义 Bochner--Riesz 平均。则这些算子具有与 \(L^2\) 定义相容的唯一的 \(L^p(\mathbb R^d;\mathbb C)\) 有界延拓，且
\[
 \sup_{R>0}\|B_R^\delta\|_{p\to p}\leq C_{d,p,\delta},
 \quad
 \lim_{R\to\infty}\|B_R^\delta f-f\|_p=0
 \quad(f\in L^p).
\]
\end{theorem}
\begin{proof}
当 \(p=2\) 时，结论已由本节开头的 Plancherel 论证得到，甚至允许 \(\delta=0\)。先设 \(1<p<2\)，令
\[
 \theta=\frac2p-1\in(0,1),\quad
 \sigma=\frac\delta\theta.
\]
题设恰好使 \(\sigma>(d-1)/2\)，而目标参数满足 \(\delta=\theta\sigma\)。在闭条带 \(0\leq\operatorname{Re}w\leq1\) 上，考察
\[
 S_wf=\mathcal F^{-1}(m_{\sigma w}\cdot\widehat f),
 \quad
 m_{\sigma w}(\xi)=(1-|\xi|^2)_+^{\sigma w}.
\]
球内使用实对数，球外为零，球面零测集的取值无关。在整条带上
\(|m_{\sigma w}|\leq1\)，且符号支撑于闭单位球。
对内点参数的任意固定阶导数，只多出
\(\sigma^k\log^k(1-|\xi|^2)\)；
在参数紧集上，其绝对积分由球面法向变量中可积的
\(u^\varepsilon|\log u|^k\) 控制。因此符号与光滑频率测试函数的配对全纯，并连续到两条边界。

先明确一套可随参数变化的共同测试空间。令 \(f,g\) 是有界紧支撑简单函数。它们的 Fourier 变换有界，所以
\[
 B_w(f,g)=\int_{|\xi|<1}
 m_{\sigma w}(\xi)\cdot\widehat f(\xi)
 \cdot\widehat g(-\xi)\,\mathrm d\xi
\]
是绝对收敛的双线性配对。它也等于
\(\int(S_wf)(x)\cdot g(x)\,\mathrm dx\)：符号可积使核有界，两个输入紧支撑，故两种写法由 Fubini 相等。这个共同定义域包括稍后使用的所有变化测试函数。

取 \(\|f\|_p=\|g\|_{p'}=1\)，把它们的非零值分别写成
\(f=\sum_j a_j\mathbf1_{E_j}\)、
\(g=\sum_k b_k\mathbf1_{F_k}\)，同一族内的集合两两不交且有界。定义
\[
 \begin{aligned}
 f_w&=\sum_j\frac{a_j}{|a_j|}
 |a_j|^{p(1+w)/2}\cdot\mathbf1_{E_j},\\
 g_w&=\sum_k\frac{b_k}{|b_k|}
 |b_k|^{p'(1-w)/2}\cdot\mathbf1_{F_k}.
 \end{aligned}
\]
在 \(w=\theta\) 时，两个复幂指数都等于一，所以
\(f_\theta=f\)、\(g_\theta=g\)。在竖边上则有
\[
 \|f_{\mathrm it}\|_2=\|g_{\mathrm it}\|_2=1,\quad
 \|f_{1+\mathrm it}\|_1=1,\quad
 \|g_{1+\mathrm it}\|_\infty=1.
\]
测试函数 \(g_w\) 的指数与 \(f_w\) 方向相反，正是为了在右边界配对 \(L^1\) 输出与 \(L^\infty\) 测试。

函数 \(F(w)=B_w(f_w,g_w)\) 在条带内全纯并连续到边界，因为 \(f_w,g_w\) 都是有限解析组合。并且各系数的模仅依赖 \(\operatorname{Re}w\)，于是
\[
 |F(w)|\leq |B(0,1)|\cdot
 \|f_w\|_1\cdot\|g_w\|_1\leq C_{f,g}
\]
在整个闭条带成立。有限解析组合及这项统一上界共同保证了条带插值所需的解析性与增长控制。

在 \(w=\mathrm it\) 上，符号模不超过一，Plancherel 与 Cauchy--Schwarz 给出
\(|F(\mathrm it)|\leq1\)。
在 \(w=1+\mathrm it\) 上，\cref{b2:lem:24-complex-kernel}及 Young 不等式给出
\[
 \|S_{1+\mathrm it}\|_{1\to1}
 \leq C e^{C(1+\sigma^2t^2)},
 \quad
 |F(1+\mathrm it)|\leq C e^{C(1+\sigma^2t^2)}.
\]
取固定 \(\kappa>C\sigma^2\)，对
\(e^{\kappa(w-\theta)^2}F(w)\)
应用\cref{b2:thm:17-three-lines}。这个因子在竖直方向压住右边界增长，在插值点等于一，边界的实部只带来固定常数。由此得到
\[
 |B_\theta(f,g)|\leq C_{d,p,\delta}.
\]
齐次性与 \(L^p\) 对偶性说明
\(\|S_\theta f\|_p\leq C_{d,p,\delta}\|f\|_p\)
在紧支撑简单函数上成立；稠密性给出唯一的 \(L^p\) 延拓。
它与原 Fourier 定义相容：对 Schwartz 输入作有界紧支撑简单函数逼近，可同时在 \(L^1,L^2,L^p\) 中收敛，频率配对与 \(L^2\) 乘子极限均保持一致。因此 \(S_\theta=B_1^\delta\)。

若 \(2<p<\infty\)，则 \(1<p'<2\)，并且
\(|1/p'-1/2|=|1/p-1/2|\)。
实值符号 \(m_\delta\) 使 \(B_1^\delta\) 在测试类上自伴随，已证的 \(L^{p'}\) 界通过对偶性给出 \(L^p\) 界。最后，\cref{b2:prop:24-dilation}的同指数共轭公式使各半径范数相同，\cref{b2:prop:24-strong-approximation}用统一算子界与稠密测试类给出 \(R\to\infty\) 时的强收敛。
\end{proof}
这是一个明确已证范围，一般窄于经典目标。它的证明只结合 \(L^2\) 与可积核端点，没有利用球面限制的额外改进。比较两条临界直线，能看清深层理论应在哪个区间增加信息。

\ProseParagraphBreak
已经证明的一致乘子界也能与输入正则性配合，给出其他 \(L^p\) 中的逼近速度。
\Needspace{6\baselineskip}
\begin{proposition}\label{b2:prop:24-lp-rate}
设整数 \(d\geq1\)，\(1<p<\infty\)，\(\delta>0\)，\(0<s\leq2\)。设符号 \((1-|\xi|^2/R^2)_+^\delta\) 的 Fourier 乘子具有与 \(L^2\) 定义相容的延拓 \(B_R^\delta:L^p\to L^p\)，并满足
\[
 M\coloneqq\sup_{R>0}\|B_R^\delta\|_{p\to p}<\infty.
\]
令 \(H^{s,p}(\mathbb R^d)\) 按\cref{b2:def:22-bessel-space}定义，即
\(\|f\|_{H^{s,p}}=\|J^sf\|_p\)，其中 \(J^s\) 的符号为
\((1+4\pi^2|\xi|^2)^{s/2}\)。则对每个 \(f\in H^{s,p}(\mathbb R^d)\) 和 \(R\geq1\)，
\[
 \|B_R^\delta f-f\|_p
 \leq C_{d,p,s,\delta,M}R^{-s}\cdot\|f\|_{H^{s,p}}.
\]
\end{proposition}
\begin{proof}
先对 Schwartz 输入证明；该类在 \(H^{s,p}\) 中稠密，且 \(J^{-s}\) 在 \(L^p\) 上有界，最终可把估计传到一般输入。取固定径向
\(\chi\in C_c^\infty(\mathbb R^d)\)，在 \(|\xi|\leq1/4\) 上为一，支撑于 \(|\xi|<1/2\)。令
\(P_Rf=\mathcal F^{-1}[\chi(\xi/R)\cdot\widehat f]\)，并记 \(g=J^sf\)。分解
\[
 (I-B_R^\delta)f
 =(I-B_R^\delta)P_Rf+
 (I-B_R^\delta)(I-P_R)f.
\]
后一个输入只含 \(|\xi|\geq R/4\) 的高频。把其 \(R^{-s}\) 因子提出后，作用于 \(g\) 的乘子为
\[
 a_R(\xi)=R^s(1-\chi(\xi/R))
 (1+4\pi^2|\xi|^2)^{-s/2}.
\]
在其支撑上 \(R\leq4|\xi|\)。每次微分 Bessel 因子带来相应的 \((1+|\xi|)^{-1}\)，微分截断的项只出现在 \(|\xi|\asymp R\) 上。因此对 \(|\alpha|\leq d+2\)，
\[
 |\xi|^{|\alpha|}\cdot|\partial^\alpha a_R(\xi)|
 \leq C_\alpha R^s(1+|\xi|)^{-s}\leq C_\alpha
\]
一致成立。\cref{b2:thm:22-mihlin}给出
\(\|(I-P_R)f\|_p\leq C R^{-s}\|g\|_p\)，再应用
\(\|I-B_R^\delta\|_{p\to p}\leq1+M\)
便控制高频误差。

低频误差的归一化乘子为
\[
 \ell_R(\xi)=R^s[1-(1-|\xi|^2/R^2)^\delta]\cdot
 \chi(\xi/R)(1+4\pi^2|\xi|^2)^{-s/2}.
\]
在截断支撑上 \(v=|\xi|^2/R^2\leq1/4\)，远离符号的球边界。由微积分基本定理，
\[
 1-(1-v)^\delta=v\,b_\delta(v),\quad
 b_\delta(v)=\delta\int_0^1(1-tv)^{\delta-1}\,\mathrm dt.
\]
所以 \(b_\delta\) 及每个所需阶数的导数在 \(0\leq v\leq1/4\) 有界。
把低频乘子写成
\[
 \ell_R(\xi)=R^{s-2}|\xi|^2\cdot
 b_\delta(|\xi|^2/R^2)\cdot\chi(\xi/R)
 (1+4\pi^2|\xi|^2)^{-s/2}.
\]
链式与乘积求导可分成两个区域记录：对所有 \(|\alpha|\leq d+2\)，
\[
 |\xi|^{|\alpha|}\cdot|\partial^\alpha\ell_R(\xi)|
 \leq C_\alpha
 \begin{cases}
 R^{s-2}|\xi|^2,&0<|\xi|\leq1,\\
 (|\xi|/R)^{2-s},&1\leq|\xi|\leq R/2.
 \end{cases}
\]
第一行利用了 \(|\xi|^2\) 的零点和其余因子在单位球内的光滑有界性；高阶导数即使消去这个零点，外面的 \(|\xi|^{|\alpha|}\) 也补足所需幂次。第二行中，微分 \(b_\delta\) 或截断产生的尺度因子由 \(R^{-1}\lesssim|\xi|^{-1}\) 控制，微分 Bessel 因子带来 \(|\xi|^{-1}\)。这解释了每次求导后仍保持同一零阶界的原因。

因为 \(R\geq1\)、\(s\leq2\)，两行都一致有界；支撑外乘子为零。再用 Mihlin 定理，得到低频误差不超过 \(C R^{-s}\|g\|_p\)。合并两部分后，以 \(H^{s,p}\) 中的 Schwartz 逼近传到一般 \(f\)，即得结论。
\end{proof}
在一致乘子界成立的指数范围内，这个论证把两种误差分开：高频尾部由输入的光滑性控制，低频误差由平均符号在原点的 Taylor 阶数控制。若进一步构造消掉低阶 Taylor 系数的组合平均，后一限制可以改变，但需要重新核对组合的统一算子界。

\ProseParagraphBreak

例如，对固定 \(d,p\) 选择处于已证插值范围的 \(\delta\)，由\cref{b2:prop:22-integer-sobolev}的整数阶范数等价，一个 \(W^{1,p}\) 输入有 \(R^{-1}\) 的范数逼近速度，一个 \(W^{2,p}\) 输入有 \(R^{-2}\) 的速度。这里只是范数误差，不能把它直接变成全参数逐点误差率；后者需要位置上的极大控制或更强的正则性与嵌入。

\ProseParagraphBreak

必要直线与插值直线的距离可以用一个数值模型说明。在 \(d=3\)、\(p=4\) 中，尾部条件要求 \(\delta>1/4\)，本章直接证明的插值范围要求 \(\delta>1/2\)。例如 \(\delta=3/4\) 在已证范围内，\(\delta=1/8\) 被必要条件排除，而 \(\delta=3/8\) 不能仅靠这两项结果判断。说它“没有被反例排除”与说它“已证明有界”有明确区别。

\ProseParagraphBreak

指数对偶把 \(p=4\) 与 \(p=4/3\) 的范围对应起来，因为球符号为实且算子在测试类上自伴随。缩放又使每个已证同指数界自动对所有半径具有同一常数。因此进一步扩大范围的关键，是获得更精细的平滑量与输入指数关系。若一篇文献只给固定空间球内的局部输出或径向输入，其命题还不能直接填入这个全空间、一般输入的空缺。

\ProseParagraphBreak

在 \(p\) 接近二时，尾部必要条件的右侧可能为负，此时它只要求原来已经约定的 \(\delta\geq0\)，看起来没有门槛。然而高维 \(\delta=0\) 已由球乘子定理排除所有非二指数。这个额外的几何障碍来自多方向叠加，并不是单个核尾部的一条更仔细积分能看出来的。正阶问题是否能在接近二的范围内降到任意小的平滑量，则要依赖更深的曲率和限制估计。

\ProseParagraphBreak

对角线上的 \(p=2\) 则没有这些困难：符号有界已经给出范数一，且所有 \(\delta\geq0\) 都随固定阶数的半径增加强收敛到恒等。这个特殊点是各种插值证明的可靠起点，但它不控制其他指数上的几何组合。把一个证明完全归结为“平方理论加插值”时，应能够说清另一个端点来自哪里，以及为什么它有足够强的范围和统一常数。
\ProseParagraphBreak
消去低频误差的组合可以直接构造，而无需增加单个平均的边界消失阶。半径放大到 \(\sqrt2R\) 时，低频二次项的系数减半；用权 \(2,-1\) 相减，就能保留常数项并消去该误差。

\Needspace{6\baselineskip}
\begin{proposition}[两半径组合的逼近阶数]\label{b2:prop:24-two-radius-rate}
设整数 \(d\geq1\)，\(1<p<\infty\)，\(\delta>0\)，且符号
\((1-|\xi|^2/R^2)_+^\delta\) 的相容 \(L^p\) Fourier 乘子满足
\(\sup_{R>0}\|B_R^\delta\|_{p\to p}\leq M<\infty\)。
定义
\[
 C_R=2B_{\sqrt2R}^\delta-B_R^\delta,\quad R>0.
\]
则 \(\|C_R\|_{p\to p}\leq3M\)，并且对 \(0<s\leq4\)、
\(R\geq1\) 和 \(f\in H^{s,p}(\mathbb R^d)\)，有
\[
 \|C_Rf-f\|_p
 \leq C_{d,p,s,\delta,M}R^{-s}\cdot\|f\|_{H^{s,p}}.
\]
这里 \(H^{s,p}\) 及其范数按\cref{b2:def:22-bessel-space}定义。
若 \(\delta=1\)，上述速率对每个 \(s>0\) 都成立，常数可依赖 \(s\)。
\end{proposition}
\begin{proof}
范数界由两个平均的界直接相加得到。对低频区域 \(|\xi|<R/2\)，令 \(v=|\xi|^2/R^2\)，则组合符号为
\[
 c_\delta(v)=2(1-v/2)^\delta-(1-v)^\delta
 =1-\frac{\delta(\delta-1)}4v^2+O(v^3).
\]
更精确地，\(c_\delta(0)=1\)、\(c_\delta'(0)=0\)，所以 Taylor 积分余项把 \(1-c_\delta(v)\) 写成
\(v^2b(v)\)，其中 \(b\) 在 \([0,1/4]\) 连同所有所需导数有界。

使用\cref{b2:prop:24-lp-rate}证明中的同一低频截断 \(P_R\)，并记 \(g=J^sf\)。高频输入估计
\[
 \|(I-P_R)f\|_p\leq C R^{-s}\cdot\|g\|_p
\]
实际上对任意 \(s>0\) 成立，因为该处仅使用 \(|\xi|\geq R/4\)，不使用 \(s\leq2\)。对它应用 \(I-C_R\)，代价至多为 \(1+3M\)。

低频归一化误差乘子现在是
\[
 R^{s-4}|\xi|^4\cdot b(|\xi|^2/R^2)\cdot
 \chi(\xi/R)(1+4\pi^2|\xi|^2)^{-s/2}.
\]
对 \(|\alpha|\leq d+2\)，其第 \(\alpha\) 阶导数乘以 \(|\xi|^{|\alpha|}\) 后，在 \(0<|\xi|\leq1\) 上由
\(C_\alpha R^{s-4}|\xi|^4\) 控制，在
\(1\leq|\xi|\leq R/2\) 上由
\(C_\alpha(|\xi|/R)^{4-s}\) 控制。理由与前一命题相同：原点的四阶零点控制小频率，归一化链式求导控制大频率。因为 \(0<s\leq4\)、\(R\geq1\)，这些界与 \(R\) 无关。Mihlin 定理给出低频误差的 \(R^{-s}\) 界，再合并高频部并由稠密性推广。

若 \(\delta=1\)，则
\[
 2(1-v/2)-(1-v)=1
\]
在整个低频区域精确成立，低频误差为零；只剩已经对任意 \(s>0\) 成立的高频估计。这证明最后一个断言，但仍使用了一阶平均原先具有统一 \(L^p\) 界的假设。
\end{proof}

组合权有负数，所以它不保留正性，也不能由正权阶数单调公式读出。它改变的是球内部的逼近阶数，而不是每条球边界的曲率障碍。特别是 \(\delta=1\) 时，精确低频抵消改善逼近阶数的作用，仍限于该乘子有界的指数范围。

\ProseParagraphBreak
通常四阶速率也会饱和。取非零 Schwartz 函数 \(f\)，使 \(\widehat f\) 光滑紧支撑，并固定 \(\delta>0\)、\(\delta\ne1\)。当 \(R\) 足够大时，整个频率支撑都在低频区；上面的 Taylor 展开连同频率导数一致成立，所以逆变换在 Schwartz 拓扑中给出
\[
 R^4(C_R-I)f
 \longrightarrow
 -\frac{\delta(\delta-1)}{64\pi^4}(-\Delta)^2f.
\]
系数来自 \(v^2=|\xi|^4/R^4\) 与
\((-\Delta)^2\) 的符号 \(16\pi^4|\xi|^4\)。该极限非零，因此不能只要求输入再光滑一些，就把一般两半径组合的阶数提高到 \(s>4\)。若要继续消去误差，必须增加组合条件或改变平均；稳定的算子界与低频逼近阶数应分别设计。

% !TeX root = ../../Book2.tex
\section{与限制问题的关系}
\label{b2:sec:2404}

球面限制控制薄壳上的 Fourier 质量，Bochner--Riesz 符号则把不同球面层加权求和。两者共享曲率与频率集中，但限制估计首先产生输入、输出指数不同的界。本节把这条联系中的测度与缩放写清，再用空间局部化和多面体逼近说明：从混合指数估计到同指数乘子估计，还缺少哪些控制。

\subsection{球面 Fourier 变换}
\label{b2:subsec:240401}

先取 \(d\geq2\)，单位球面 \(S^{d-1}\) 的面积测度记为 \(\sigma\)，半径 \(r>0\) 的球面面积测度记为 \(\sigma_r\)。它们满足
\[
 \int_{rS^{d-1}}h(\xi)\,\mathrm d\sigma_r(\xi)
 =r^{d-1}\int_{S^{d-1}}h(r\omega)\,\mathrm d\sigma(\omega).
\]
右侧的 \(r^{d-1}\) 来自曲面面积的缩放，与径向积分中的 \(\mathrm dr\) 分别保留。

对 \(f\in\mathcal S(\mathbb R^d)\)，定义单位参数下的球面层算子
\[
 T_rf(x)\coloneqq
 \int_{S^{d-1}}e^{2\pi\mathrm i r x\cdot\omega}
       \cdot\widehat f(r\omega)\,\mathrm d\sigma(\omega).
\]
以真实面积测度积分的限制再延拓算子则是 \(r^{d-1}T_r\)。因此极坐标给出
\begin{equation}\label{b2:eq:24-layer-measure}
 \begin{aligned}
 B_1^\delta f(x)
 &=\int_0^1(1-r^2)^\delta
       \int_{rS^{d-1}}e^{2\pi\mathrm i x\cdot\xi}
          \cdot\widehat f(\xi)\,\mathrm d\sigma_r(\xi)\,\mathrm dr\\
 &=\int_0^1(1-r^2)^\delta r^{d-1}T_rf(x)\,\mathrm dr,
 \quad \delta>-1.
 \end{aligned}
\end{equation}
对 \(\delta\geq0\)，这就是\cref{b2:def:24-bochner-riesz}中的平均；对 \(-1<\delta<0\)，此处先按球内可积符号的逆 Fourier 积分解释。测试函数的 Fourier 变换有界，而
\(\int_0^1(1-r^2)^\delta r^{d-1}\,\mathrm dr<\infty\)，所以换序绝对合法。这个公式中若使用 \(\sigma_r\)，就不应再另乘一次 \(r^{d-1}\)。

\ProseParagraphBreak
内层算子只保留一个频率半径，径向积分才恢复球内部的频率。球面 Fourier 变换的振荡尾部提供 \((d-1)/2\) 阶衰减；径向权在 \(r=1\) 的消失阶又在端点振荡积分中提供 \(\delta+1\) 阶增益，合起来与本章核的 \(\delta+(d+1)/2\) 阶衰减一致。若先对每一层取绝对值，再对半径求和，不同半径之间的抵消便不再可用。下面先看这种直接估计究竟能够证明什么。

\subsection{限制估计}
\label{b2:subsec:240402}

本处需要\cref{b2:thm:23-stein-tomas}的球面端点版本。设
\[
 p_*=\frac{2(d+1)}{d+3},\quad
 q_*=p_*'=\frac{2(d+1)}{d-1}.
\]
对单位球面及面积测度，限制算子 \(Rf=\widehat f|_{S^{d-1}}\) 满足
\[
 \|Rf\|_{L^2(S^{d-1})}\leq A_d\|f\|_{p_*}.
\]
其有界伴随为
\[
 Eg(x)=\int_{S^{d-1}}e^{2\pi\mathrm i x\cdot\omega}
        \cdot g(\omega)\,\mathrm d\sigma(\omega),\quad
 \|Eg\|_{q_*}\leq A_d\|g\|_2.
\]
这里对一般 \(L^{p_*}\) 输入的限制按连续延拓定义，并不是任意 Fourier 代表的逐点取值。球面为有限正测度，故\cref{b2:prop:23-ttstar}也给出 \(ER:L^{p_*}\to L^{q_*}\) 及其范数上界 \(A_d^2\)。

\ProseParagraphBreak
半径改变时，需要同时缩放输入和输出。令 \(f_r(y)=f(y/r)\)，则
\[
 \widehat{f_r}(\omega)=r^d\widehat f(r\omega),\quad
 \|f_r\|_{p_*}=r^{d/p_*}\|f\|_{p_*}.
\]
因此
\[
 \|\widehat f(r\,\cdot)\|_{L^2(S^{d-1})}
 \leq A_dr^{-d/q_*}\|f\|_{p_*},
\]
\[
 T_rf(x)=r^{-d}(ERf_r)(rx),\quad
 \|T_rf\|_{q_*}
 \leq A_d^2r^{-2d/q_*}\|f\|_{p_*}.
\]
最后一个幂也等于 \(d/p_*-d/q_*-d\)。在\eqref{b2:eq:24-layer-measure}中乘以真实面积因子 \(r^{d-1}\) 后，层算子的界便带有
\[
 r^{d-1-2d/q_*}=r^\alpha,\quad
 \alpha\coloneqq\frac{d-1}{d+1}.
\]
这一幂次在零半径处仍然可积，稍后将用于全平均。

\ProseParagraphBreak
把逐半径估计先用于非负权，可以同时得到频率平方界与混合指数界。

\Needspace{6\baselineskip}
\begin{proposition}[非负径向权的限制界]\label{b2:prop:24-radial-weight}
设整数 \(d\geq2\)，
\[
 p_*=\frac{2(d+1)}{d+3},\quad
 q_*=\frac{2(d+1)}{d-1},\quad
 \alpha=\frac{d-1}{d+1}.
\]
设 \(w:(0,\infty)\to[0,\infty)\) 可测且
\[
 I_w\coloneqq\int_0^\infty w(r)r^\alpha\,\mathrm dr<\infty.
\]
对 \(f\in\mathcal S(\mathbb R^d)\)，定义
\[
 A_wf(\xi)=w(|\xi|)^{1/2}\cdot\widehat f(\xi),\quad
 S_wf=\mathcal F^{-1}\!\left[w(|\xi|)\cdot\widehat f(\xi)\right].
\]
则这两个算子分别唯一连续延拓为
\[
 A_w:L^{p_*}(\mathbb R^d)\longrightarrow L^2(\mathbb R^d),\quad
 S_w:L^{p_*}(\mathbb R^d)\longrightarrow L^{q_*}(\mathbb R^d),
\]
并有
\[
 \|A_wf\|_2\leq A_d I_w^{1/2}\|f\|_{p_*},\quad
 \|S_wf\|_{q_*}\leq A_d^2 I_w\|f\|_{p_*}.
\]
常数 \(A_d\) 可取单位球面 Stein--Tomas 限制界的常数，与 \(w,f\) 无关。
\end{proposition}
\begin{proof}
先取 Schwartz 输入。极坐标、非负性和刚才的半径限制估计给出
\[
 \begin{aligned}
 \|A_wf\|_2^2
 &=\int_0^\infty w(r)r^{d-1}
       \|\widehat f(r\,\cdot)\|_{L^2(S^{d-1})}^2\,\mathrm dr\\
 &\leq A_d^2\|f\|_{p_*}^2
       \int_0^\infty w(r)r^{d-1-2d/q_*}\,\mathrm dr
 =A_d^2 I_w\|f\|_{p_*}^2.
 \end{aligned}
\]
这就给出第一项连续延拓。第二项的测试函数定义也有意义：在 \(r\leq1\) 时 \(r^{d-1}\leq r^\alpha\)；在 \(r\geq1\) 时，\(\widehat f\) 的任意阶衰减控制 \(r^{d-1-\alpha}\)。所以 \(w(|\xi|)\cdot\widehat f(\xi)\) 绝对可积。

将 \(S_wf\) 与另一 Schwartz 函数 \(h\) 配对，Fubini 给出
\[
 \langle S_wf,h\rangle
 =\int_{\mathbb R^d}w(|\xi|)\cdot
       \widehat f(\xi)\overline{\widehat h(\xi)}\,\mathrm d\xi
 =\langle A_wf,A_wh\rangle.
\]
于是
\[
 |\langle S_wf,h\rangle|
 \leq A_d^2 I_w\|f\|_{p_*}\cdot\|h\|_{p_*}.
\]
\(L^{p_*}\) 的对偶为 \(L^{q_*}\)，所以这个分布由一个 \(L^{q_*}\) 函数表示，并得到第二项估计。最后由 Schwartz 类在 \(L^{p_*}\) 中稠密，完成唯一连续延拓。
\end{proof}

这里的平方配对就是 \(TT^*\) 方法，但已经在实际使用的径向权上重新给出了论证，未要求这个加权测度具有紧支集。权的非负性保证 \(\langle S_wf,f\rangle\geq0\)；它不保证逆 Fourier 核处处非负。对复值径向权，只要以 \(|w|\) 定义的 \(I_{|w|}\) 有限，同样把配对中的相位因子放在两个平方因子之间，仍得到第二项上界。

\Needspace{6\baselineskip}
\begin{corollary}[薄球壳的两种上界]\label{b2:cor:24-thin-shell}
设整数 \(d\geq2\)，\(p_*=2(d+1)/(d+3)\)、\(q_*=p_*'\)，且 \(0<\varepsilon<1/2\)。记
\[
 \Omega_\varepsilon
 =\{\,\xi\in\mathbb R^d:1-\varepsilon<|\xi|<1+\varepsilon\,\},\quad
 \Pi_\varepsilon f
 =\mathcal F^{-1}\!\left[\mathbf1_{\Omega_\varepsilon}\cdot\widehat f\right]
\]
先定义于 Schwartz 类。则
\[
 \|\mathbf1_{\Omega_\varepsilon}\cdot\widehat f\|_2
 \leq C_d\varepsilon^{1/2}\|f\|_{p_*},\quad
 \|\Pi_\varepsilon f\|_{q_*}
 \leq C_d\varepsilon\|f\|_{p_*}.
\]
两项分别按上述范数具有唯一连续延拓。
\end{corollary}
\begin{proof}
在\cref{b2:prop:24-radial-weight}中取
\(w=\mathbf1_{(1-\varepsilon,1+\varepsilon)}\)。
由于 \(1/2<r<3/2\)，有
\[
 I_w=\int_{1-\varepsilon}^{1+\varepsilon}r^{(d-1)/(d+1)}\,\mathrm dr
 \leq C_d\varepsilon.
\]
分别代入频率平方界与混合指数界，即得两项估计及延拓。
\end{proof}

第一式中的 \(\varepsilon^{1/2}\) 来自径向积分后开平方；第二式中的 \(\varepsilon\) 来自限制与伴随各提供一次平方根。它们控制的输出空间不同，不能仅因第二项也是频率投影，就把它读成 \(L^{p_*}\) 上的投影范数。

\ProseParagraphBreak
更一般地，若 \(w\geq0\) 支于固定半径区间 \([1/2,3/2]\)，则混合算子范数至多为
\(C_d\int w(r)\,\mathrm dr\)。在离球边界约 \(2^{-j}\) 的壳层上，\((1-r^2)^\delta\) 大小约为 \(2^{-j\delta}\)，壳层厚度约为 \(2^{-j}\)，因此该层的混合指数贡献至多为
\(C_{d,\delta}2^{-j(\delta+1)}\)。只要 \(\delta>-1\)，这些边界层就可求和。要把它用于整个球，还必须处理半径趋零的部分。

\Needspace{6\baselineskip}
\begin{proposition}[全平均的混合指数界]\label{b2:prop:24-mixed-bound}
设整数 \(d\geq2\)，\(p_*=2(d+1)/(d+3)\)、\(q_*=p_*'\)，\(\delta>-1\)、\(R>0\)。对复值 Schwartz 函数，定义
\[
 B_R^\delta f(x)
 =\int_{|\xi|<R}e^{2\pi\mathrm i x\cdot\xi}
       (1-|\xi|^2/R^2)^\delta
       \cdot\widehat f(\xi)\,\mathrm d\xi.
\]
则它有唯一连续延拓
\[
 B_R^\delta:L^{p_*}(\mathbb R^d)\longrightarrow L^{q_*}(\mathbb R^d),
\]
并满足
\[
 \|B_R^\delta f\|_{q_*}
 \leq\frac{A_d^2}{2}
       \mathrm B\!\left(\frac d{d+1},\delta+1\right)
       R^{2d/(d+1)}\|f\|_{p_*}.
\]
其中 \(A_d\) 为单位球面面积测度的 \(L^{p_*}(\mathbb R^d)\to L^2(S^{d-1})\) 限制常数，\(\mathrm B\) 为 beta 函数。
\end{proposition}
\begin{proof}
先令 \(R=1\)，在径向权命题中令 \(w(r)=(1-r^2)^\delta\) 对 \(0<r<1\) 成立，在其余半径取零。
相应的积分为
\[
 I_w=\int_0^1(1-r^2)^\delta r^{(d-1)/(d+1)}\,\mathrm dr
 =\frac12\,\mathrm B\!\left(\frac d{d+1},\delta+1\right).
\]
最后一步令 \(t=r^2\)。零端的参数 \(d/(d+1)>0\)，球面边界端的参数 \(\delta+1>0\)，所以两端均可积。\(T_r\) 的范数在 \(r\downarrow0\) 时按 \(r^{-2d/q_*}\) 增大，而体积权 \(r^{d-1}\) 将总幂次变成 \((d-1)/(d+1)>-1\)，保证了零端的可积性。

对任意 \(R>0\)，直接在径向权中取
\[
 w_R(r)=
 \begin{cases}
  (1-r^2/R^2)^\delta,&0<r<R,\\
  0,&r\geq R.
 \end{cases}
\]
换元 \(r=Rs\) 后，
\[
 I_{w_R}=R^{\,1+(d-1)/(d+1)}I_w
        =R^{2d/(d+1)}I_w.
\]
再次应用径向权命题即得所述界。测试函数稠密性给出唯一延拓。
\end{proof}

对 \(\delta\geq0\)，这里的算子与本章原定义一致；负参数只在本命题明确规定的积分与连续延拓意义下使用。若 \(-1<\delta<0\)，其符号在球面附近本性无界，因而不是有界的 \(L^2\) 同指数乘子。这与命题不矛盾，因为输入、输出空间已经改变。

\ProseParagraphBreak
半径因子也能用一般缩放核对：
\[
 d\left(\frac1{p_*}-\frac1{q_*}\right)=\frac{2d}{d+1}.
\]
同指数时两个范数的尺度因子相消；这里并不相消。对单位半径证明的混合估计，不能在推广到任意半径时删掉这个幂次。

\subsection{乘子估计}
\label{b2:subsec:240403}

在无限测度空间 \(\mathbb R^d\) 上，\(L^{q_*}\) 不包含 \(L^{p_*}\)。例如选
\[
 \frac d{q_*}<a\leq\frac d{p_*},
\]
则 \((1+|x|)^{-a}\) 属于 \(L^{q_*}\)，却不属于 \(L^{p_*}\)。因而前面的混合界没有一个免费的输出空间嵌入，可以把它改成同指数界。

\ProseParagraphBreak
把输出局部化到空间球，可以明确看见缺失的常数。对半径 \(L>0\) 的球 \(B(0,L)\)，Hölder 给出
\[
 \|F\|_{L^{p_*}(B(0,L))}
 \leq |B(0,L)|^{1/p_*-1/q_*}
       \|F\|_{L^{q_*}(B(0,L))}
 =c_d L^{2d/(d+1)}\|F\|_{L^{q_*}(B(0,L))}.
\]
随着空间球增大，这个体积因子也增大。薄壳投影因此只能直接给出
\[
 \|\Pi_\varepsilon f\|_{L^{p_*}(B(0,L))}
 \leq C_d\varepsilon L^{2d/(d+1)}\|f\|_{p_*}.
\]
例如取 \(L=\varepsilon^{-1}\)，右侧的系数是
\(\varepsilon^{-(d-1)/(d+1)}\)，并未随壳厚趋零而减小。这只是球内估计；球外输出尚未处理，不能把它宣布为全空间乘子界。

\ProseParagraphBreak
有限个壳层也能呈现同一问题。设
\[
 \Lambda_j
 =\{\,\xi:1-2^{-j}<|\xi|<1-2^{-j-1}\,\},\quad
 Q_jf=\mathcal F^{-1}(\mathbf1_{\Lambda_j}\cdot\widehat f),
 \quad j\geq1.
\]
由径向权估计，对任意有限 \(N\) 和复系数 \(a_1,\ldots,a_N\)，有
\[
 \left\|\sum_{j=1}^N a_jQ_jf\right\|_{q_*}
 \leq C_d\left(\sum_{j=1}^N|a_j|2^{-j}\right)\|f\|_{p_*}.
\]
即使 \(a_j=1\)，右侧也对 \(N\) 一致；但换成空间球上的 \(p_*\) 范数时，仍须乘 \(L^{2d/(d+1)}\)。当 \(a_j\) 的大小为 \(2^{-j\delta}\) 时，混合范数中的可和条件是 \(\delta>-1\)。这个求和没有提供足以取消体积因子的额外信息。

\ProseParagraphBreak
频率局限于同一固定薄壳也不能恢复所缺的嵌入。固定 \(0<\varepsilon<1/2\)，取非零
\(\widehat\phi\in C_c^\infty(\Omega_\varepsilon)\)，于是 \(\phi\in\mathcal S\) 且 \(\Pi_\varepsilon\phi=\phi\)。对正整数 \(N\) 与距离参数 \(L\)，令
\[
 F_{N,L}(x)=\sum_{k=1}^N\phi(x-kLe_1).
\]
它的 Fourier 支集仍在同一薄壳内。固定 \(N\)，当 \(L\to\infty\) 时，对 \(s=p_*\) 与 \(s=q_*\) 分别有
\[
 \|F_{N,L}\|_s\longrightarrow N^{1/s}\|\phi\|_s.
\]
证明只需先把 \(\phi\) 在 \(L^s\) 中逼近为紧支撑函数；距离充分大时，各平移的紧支撑互不相交，范数的 \(s\) 次方便准确相加。原函数与近似函数的和相差至多 \(N\) 倍单项误差；对固定 \(N\) 再令误差趋零，即得极限。

\ProseParagraphBreak
所以在这些薄壳频率函数上，
\[
 \frac{\|F_{N,L}\|_{p_*}}{\|F_{N,L}\|_{q_*}}
 \longrightarrow
 N^{1/p_*-1/q_*}\frac{\|\phi\|_{p_*}}{\|\phi\|_{q_*}}
 =N^{2/(d+1)}\frac{\|\phi\|_{p_*}}{\|\phi\|_{q_*}}.
\]
先取大 \(N\)，再取足够大的有限距离，就能使比值任意大。这是一个有限平移模型，说明即使只考察同一频率壳层内的输出，也没有统一的 \(L^{q_*}\to L^{p_*}\) 范数转换。它不是薄壳投影同指数不有界性的独立证明；它排除的是从已有混合估计直接转换指数的那一步。

\ProseParagraphBreak
要研究同指数 \(L^p\) 界，通常还需按空间尺度分解核，把限制估计与核尾部估计结合，再控制各壳层的组合。\chapref{b2:ch:19}的 Littlewood--Paley 分解把以原点为中心、大小按倍数变化的频率环带组织起来；这里的各层都贴近同一单位球面，径向宽度不断缩小，不能不加说明地套用同一平方函数结论。不交频率在 Plancherel 中确实正交，但这只给 \(p=2\) 的平方范数，其他指数还要处理方向与空间位置上的相干。

\subsection{曲率}
\label{b2:subsec:240404}

球、椭球和多面体有助于分清局部几何与全局组合。先看线性换元：它改变球面在欧氏坐标中的形状，却不改变同指数乘子范数。

\Needspace{6\baselineskip}
\begin{proposition}[可逆线性换元下的乘子共轭]\label{b2:prop:24-linear-conjugacy}
设整数 \(d\geq1\)，\(A\) 为实可逆 \(d\times d\) 矩阵，\(m\in L^\infty(\mathbb R^d)\)，并令
\[
 T_mf=\mathcal F^{-1}(m\cdot\widehat f),\quad
 U_Af(x)=f(A^{\mathsf T}x).
\]
若 \(1<p,q<\infty\)，则 \(T_m:L^p\to L^q\) 有界，当且仅当 \(T_{m\circ A}:L^p\to L^q\) 有界；在有界情形下，
\[
 T_{m\circ A}=U_A^{-1}T_mU_A,\quad
 \|T_{m\circ A}\|_{p\to q}
 =|\det A|^{1/q-1/p}\cdot\|T_m\|_{p\to q}.
\]
这里算子均指从 Schwartz 类的连续延拓。
\end{proposition}
\begin{proof}
Fourier 换元公式与空间范数换元分别为
\[
 \widehat{U_Af}(\xi)=|\det A|^{-1}\cdot\widehat f(A^{-1}\xi),\quad
 \|U_Af\|_s=|\det A|^{-1/s}\cdot\|f\|_s
 \quad(1\leq s<\infty).
\]
因此 \(U_A^{-1}T_mU_A\) 的符号为 \(m(A\xi)\)，得到共轭恒等式。若 \(T_m\) 有界，两个拉回范数给出
\[
 \|T_{m\circ A}\|_{p\to q}
 \leq|\det A|^{1/q-1/p}\cdot\|T_m\|_{p\to q}.
\]
反向对 \(A^{-1}\) 和符号 \(m\circ A\) 应用同一结论，得到逆向不等式及有界性的等价。测试函数稠密性使共轭关系保持到两个 Banach 空间之间。
\end{proof}

取 \(m(\xi)=(1-|\xi|^2)_+^\delta\)、\(\delta\geq0\)，就得到椭球符号
\((1-|A\xi|^2)_+^\delta\) 的平均。在同指数 \(p=q\) 时，行列式因子完全相消；因此固定椭球与球具有相同的有界性范围和算子范数。在不同指数上则保留这个因子，与前面球半径的缩放一致。若另外估计各欧氏方向导数，矩阵与其逆矩阵会进入导数公式；这种正则性常数的依赖不能与纯 \(L^p\) 乘子范数混为一谈。

\ProseParagraphBreak
平面边界提供另一种可直接构造的模型。需要恢复的前章工具是一维 Hilbert 变换：
\[
 \widehat{Hf}(\xi)=-\mathrm i\,\operatorname{sgn}(\xi)
                  \cdot\widehat f(\xi),\quad
 \|Hf\|_{L^p(\mathbb R)}\leq C_{H,p}\|f\|_{L^p(\mathbb R)}
 \quad(1<p<\infty).
\]
其平方定义见\cref{b2:thm:18-hilbert-l2}；核 \(1/(\pi x)\) 的大小与差分满足 Calderón--Zygmund 条件，所以\cref{b2:thm:18-cz-lp}给出中间指数有界性及与平方定义的兼容。由于 \(\mathrm iH\) 的符号为 \(\operatorname{sgn}\xi\)，负半轴投影为 \((I-\mathrm iH)/2\)。在高维中沿固定方向应用这个一维算子，Fubini 保持同一个 \(L^p\) 上界。

\Needspace{6\baselineskip}
\begin{proposition}[有限半空间交的投影]\label{b2:prop:24-polytope-projection}
设整数 \(d\geq1\)，\(1<p<\infty\)，正整数 \(N\)，单位向量
\(\nu_1,\ldots,\nu_N\in S^{d-1}\) 及实数 \(b_1,\ldots,b_N\) 给出
\[
 P=\bigcap_{j=1}^N\{\,\xi\in\mathbb R^d:\xi\cdot\nu_j<b_j\,\}.
\]
令 \(T_Pf=\mathcal F^{-1}(\mathbf1_P\cdot\widehat f)\) 先在 Schwartz 类上定义。则 \(T_P\) 唯一连续延拓为 \(L^p(\mathbb R^d)\) 上的有界算子，并且
\[
 \|T_P\|_{p\to p}\leq
 \left(\frac{1+C_{H,p}}2\right)^N.
\]
这里 \(C_{H,p}\) 是一维 Hilbert 变换的 \(L^p\) 范数上界；常数不依赖法向与各平面的位置。集合 \(P\) 可以无界，空集情形也包括在内。
\end{proposition}
\begin{proof}
对单位向量 \(\nu\)，用旋转把一维 Hilbert 变换放到该方向，记作 \(H_\nu\)。它在平方空间的符号为
\(-\mathrm i\,\operatorname{sgn}(\xi\cdot\nu)\)，且 Fubini 给出
\(\|H_\nu\|_{p\to p}\leq C_{H,p}\)。于是
\[
 Q_{\nu,0}=\frac{I-\mathrm iH_\nu}{2}
\]
在 \(L^2\) 上是 \(\{\xi\cdot\nu<0\}\) 的投影，在 \(L^p\) 上有界至多 \((1+C_{H,p})/2\)。令
\[
 M_{b\nu}f(x)=e^{2\pi\mathrm i b\nu\cdot x}\cdot f(x).
\]
调制在 \(L^p\) 上等距，且
\(Q_{\nu,b}=M_{b\nu}Q_{\nu,0}M_{-b\nu}\)
的符号为 \(\mathbf1_{\{\xi\cdot\nu<b\}}\)。

在 \(L^2\) 上，这些算子都是 Fourier 乘法，因此彼此交换，有限复合的符号为
\[
 \prod_{j=1}^N\mathbf1_{\{\xi\cdot\nu_j<b_j\}}=\mathbf1_P.
\]
每项的 \(L^p\) 延拓与其 \(L^2\) 定义在交集上一致，所以对 Schwartz 输入，复合既等于 \(T_P\)，又有上述 \(N\) 个常数的乘积界。由 \(L^p\) 稠密性完成延拓。边界包含在有限个超平面中，Lebesgue 测度为零，故使用开或闭半空间不改变乘子。
\end{proof}

矩形是这个结论的一个特例，也可以直接逐坐标使用区间投影。一般有限凸多面体的各个平面方向可以不同，但有限次复合仍有界。这个证明给出的是随面数积累的上界，尚未说明它在越来越多面时是否真的必须变大。下面用逼近球的有限模型回答后一问题。

\Needspace{6\baselineskip}
\begin{proposition}[逼近球的多面体范数发散]\label{b2:prop:24-polytope-divergence}
设整数 \(d\geq2\)，\(1<p<\infty\)、\(p\neq2\)。取 \(0<\eta_k<1\)、\(\eta_k\to0\)，并取单位球面的有限 \(\eta_k\)-网 \(V_k\subset S^{d-1}\)，即每个单位向量到 \(V_k\) 某点的欧氏距离小于 \(\eta_k\)。定义外切凸多面体
\[
 P_k=\bigcap_{\nu\in V_k}\{\,\xi:\xi\cdot\nu<1\,\},\quad
 T_{P_k}f=\mathcal F^{-1}(\mathbf1_{P_k}\cdot\widehat f).
\]
则每个 \(T_{P_k}\) 都在 \(L^p(\mathbb R^d)\) 上有界，但
\[
 \|T_{P_k}\|_{p\to p}\longrightarrow\infty.
\]
\end{proposition}
\begin{proof}
有限网由单位球面的紧致性取得。单位开球包含于 \(P_k\)。若 \(\xi\in P_k\)、\(\xi\neq0\)，令 \(u=\xi/|\xi|\)，并选 \(\nu\in V_k\) 满足 \(|u-\nu|<\eta_k\)。单位向量的恒等式给出
\[
 u\cdot\nu=1-\frac{|u-\nu|^2}{2}
 >1-\frac{\eta_k^2}{2}.
\]
结合 \(\xi\cdot\nu<1\)，得到
\[
 B(0,1)\subset P_k
 \subset B\!\left(0,(1-\eta_k^2/2)^{-1}\right).
\]
所以各 \(P_k\) 有界，且其示性函数在单位球面之外逐点趋于 \(\mathbf1_{B(0,1)}\)。球面是 Lebesgue 零测集，故符号几乎处处收敛。

每个多面体投影的有界性已由\cref{b2:prop:24-polytope-projection}证明。若所述范数不趋于无穷，则存在一个子列，其范数统一不超过某个有限 \(M\)。对任意固定 \(f\in\mathcal S\)，符号差的模至多为一，且 \(\widehat f\in L^1\)，故支配收敛给出
\[
 \|( \mathbf1_{P_k}-\mathbf1_{B(0,1)})\cdot\widehat f\|_1\longrightarrow0.
\]
因而这一子列的 \(T_{P_k}f\) 一致收敛到硬球投影 \(B_1^0f\)。Fatou 引理于是给出
\[
 \|B_1^0f\|_p^p
 \leq\liminf_k\|T_{P_k}f\|_p^p
 \leq M^p\|f\|_p^p.
\]
稠密性将硬球投影延拓为有界 \(L^p\) 算子，与\cref{b2:thm:24-fefferman}在 \(d\geq2\)、\(p\neq2\) 时的结论矛盾。
\end{proof}

这个发散结论使用了本章明确标出的深层球乘子定理，但由它到有限多面体模型的极限步骤都已写出。它没有提供关于面数的定量发散速度；在 \(p=2\) 时，各个非空多面体投影的范数反而恒为一。几何区域逐点逼近球，并不能提供非二指数上一致的算子常数。

\ProseParagraphBreak
曲率的两种作用因此需要分开理解。对一张固定的驻相片，非零曲率能改善振荡核衰减；当越来越多方向共同出现时，输出又可能形成多方向集中。前一种效应服务于正平滑平均的核估计，后一种效应参与硬球投影的失败。分析新的频率区域时，应同时考察单片衰减、片的数量以及各片怎样组合；只有局部衰减还不够。

\ProseParagraphBreak
本章各项结果也各有明确的适用范围：尾部检验给出必要条件，可积核与解析插值给出同指数充分范围，球面限制给出本节的混合指数界，球乘子定理排除高维硬截断的非二指数。应用文献中的进一步估计时，还应核对输入是否径向、空间是否局部、半径是否统一及是否允许尺度损失，不能把改变了其中条件的结论直接填入普通全空间 \(L^p\) 乘子问题。


\begin{chaptersummary}
Bochner--Riesz 符号的球边界决定核的远场振荡，原点附近的展开则决定光滑输入的低频误差。前者控制可积性与乘子障碍，后者可以通过半径组合改变。两种机制虽然出现在同一个符号中，却需要不同的估计；更高的边界消失阶并不无条件意味着更快的逼近。

\ProseParagraphBreak

可积核给出全部普通指数的稳定性、有限指数下的范数逼近以及由极大函数控制的逐点结论。临界核失去绝对可积性以后，这套证明停止，但平方理论仍然成立。尾部检验、解析插值与高维硬球否定定理各自提供不同信息，留下的正阶区间不能靠形式代入或指数包含补齐。

\ProseParagraphBreak

限制理论通过壳层给出混合指数估计。将它转成同指数乘子界，必须另外处理空间局部化、尺度求和和方向组合。固定多面体的有界性也不能在缺少一致常数时传递给球。无论改变半径、阶数还是几何形状，都应同时追踪算子常数和收敛方式，才能判断一个极限真正保留了什么。
\end{chaptersummary}

\BookExercises

% 自拟；中等：固定频带上的阶数半群、生成元定义域与准确范数差，区分子空间恒等和全空间投影。
\Needspace{6\baselineskip}
\begin{exercise}\label{b2:ex:24-new-1}
设整数 \(d\geq1\)、\(R>0\)，令
\[
 \mathcal H_R=\{\,f\in L^2(\mathbb R^d):
          \widehat f=0\ \text{几乎处处于 }\{|\xi|\geq R\}\,\}.
\]
把 \(B_R^\delta\) 限制到这个 Hilbert 子空间，证明 \(\delta\geq0\) 构成强连续的压缩半群。若生成元按
\(Gf=\lim_{h\downarrow0}(B_R^hf-f)/h\) 定义，证明
\[
 \widehat{Gf}(\xi)=\log(1-|\xi|^2/R^2)\cdot\widehat f(\xi)
       \quad(|\xi|<R),
\]
球外取零；其定义域恰为使右侧在球内属于 \(L^2\) 的 \(f\in\mathcal H_R\)。再证明对 \(0<\delta<\eta\)，
\[
 \|B_R^\delta-B_R^\eta\|_{\mathcal H_R\to\mathcal H_R}
 =\left(1-\frac\delta\eta\right)
       \left(\frac\delta\eta\right)^{\delta/(\eta-\delta)},
\]
而 \(\|B_R^\delta-I\|_{\mathcal H_R\to\mathcal H_R}=1\) 对每个 \(\delta>0\) 成立。说明为什么同一个族在全空间上的零阶算子不能被称为恒等。
\end{exercise}
\begin{solution}
球内令 \(a(\xi)=1-|\xi|^2/R^2\)，则 \(0<a\leq1\)。符号相乘给出
\(B_R^\delta B_R^\eta=B_R^{\delta+\eta}\)，其模不超过一；零阶符号在 \(\mathcal H_R\) 上为一。对每个固定输入，\(a^h\to1\) 且 \(|a^h-1|\leq1\)，Plancherel 与支配收敛证明零点强连续。半群性质和压缩界把连续性移到其余参数。

记 \(\ell(\xi)=-\log a(\xi)\geq0\)。当 \(\ell\widehat f\in L^2\) 时，
\[
 \frac{e^{-h\ell}-1}{h}\widehat f\longrightarrow-\ell\widehat f,
 \quad
 \left|\frac{e^{-h\ell}-1}{h}\widehat f\right|
       \leq\ell|\widehat f|.
\]
平方支配收敛给出生成元极限。反过来，若差商在 \(L^2\) 中有极限，取一个几乎处处收敛的子列；其点态极限只能是 \(-\ell\widehat f\)，所以这个函数必在 \(L^2\)，也必是原来的平方极限。定义域的两向包含均已得到。

算子范数由符号的本性上确界给出。函数 \(u^\delta-u^\eta\) 在 \(0\leq u\leq1\) 非负，内部最大点满足
\(u^{\eta-\delta}=\delta/\eta\)。代入便得题中准确公式；相应的 \(u\) 值对应正测度邻域的球内频率，所以也是本性上确界。当比较正阶与零阶时，\(1-u^\delta\) 在 \(u\downarrow0\) 趋于一，故范数差恒为一。这不妨碍每个固定输入强连续。

在整个 \(L^2\) 上，\(B_R^0\) 是频率球的正交投影，消去了球外数据，因而不是恒等。只有先明确限制到 \(\mathcal H_R\)，上述族才具有半群在零参数处所需的恒等元。
\end{solution}

% 自拟；中等：半径趋零时的质量扩散轮廓，求L1极限范数并与中间指数的强趋零比较。
\Needspace{6\baselineskip}
\begin{exercise}\label{b2:ex:24-new-2}
设整数 \(d\geq1\)、\(\delta>(d-1)/2\)，记 \(K_\delta\) 为单位半径核。对 \(f\in L^1(\mathbb R^d)\)，令 \(M_f=\int f\)。证明当 \(R\downarrow0\) 时，对每个 \(1\leq q\leq\infty\)，
\[
 R^{-d}(B_R^\delta f)(\,\cdot\,/R)
       \longrightarrow M_fK_\delta\quad\text{于 }L^q,
\]
其中 \(q=\infty\) 指一致范数。由此证明
\[
 \lim_{R\downarrow0}R^{-d(1-1/q)}\|B_R^\delta f\|_q
       =|M_f|\cdot\|K_\delta\|_q.
\]
说明 \(L^1\) 输入的输出何时强趋于零，并证明对每个 \(1<p<\infty\) 及每个 \(f\in L^p\)，都有 \(B_R^\delta f\to0\) 于 \(L^p\)。
\end{exercise}
\begin{solution}
可积核定理和频率符号可积性说明 \(K_\delta\in L^1\cap C_0\)，因而属于全部 \(L^q\)、\(1\leq q\leq\infty\)。它的平移在有限 \(q\) 范数中连续；在一致范数中也连续，因为 \(C_0\) 函数一致连续。缩放核公式给出
\[
 R^{-d}(B_R^\delta f)(z/R)=\int K_\delta(z-Ry)\cdot f(y)\,\mathrm dy.
\]
减去 \(M_fK_\delta(z)\)，再用 Minkowski 或一致范数的积分上界，得到
\[
 \left\|R^{-d}(B_R^\delta f)(\,\cdot\,/R)-M_fK_\delta\right\|_q
 \leq\int|f(y)|\cdot
          \|K_\delta(\,\cdot-Ry)-K_\delta\|_q\,\mathrm dy.
\]
对固定 \(y\)，范数趋于零；它至多为 \(2\|K_\delta\|_q\)。以 \(|f|\) 为可积控制，支配收敛就给出第一条结论。空间换元把左侧未减去轮廓的范数写成
\(R^{-d(1-1/q)}\|B_R^\delta f\|_q\)，故也得到第二条结论。

特别地，\(\|B_R^\delta f\|_1\to|M_f|\cdot\|K_\delta\|_1\)，所以强趋于零恰好要求 \(M_f=0\)。当 \(M_f\ne0\) 时甚至不存在任何 \(L^1\) 强极限：一方面
\(\|B_R^\delta f\|_\infty\leq R^d\|K_\delta\|_\infty\cdot\|f\|_1\to0\)，
另一方面其 \(L^1\) 范数不趋零。若有强极限，几乎处处收敛子列会迫使极限为零，形成矛盾。

若 \(1<p<\infty\) 且 \(f\in L^1\cap L^p\)，上面第二式的 \(q=p\) 情形与 \(d(1-1/p)>0\) 给出 \(L^p\) 强趋零。一般 \(L^p\) 输入可用 \(L^1\cap L^p\) 逼近，再用 \(\sup_R\|B_R^\delta\|_{p\to p}\leq\|K_\delta\|_1\) 推广。这里的小半径把质量摊到越来越大的空间尺度；总质量仍可保留，而中间指数范数已经趋零。
\end{solution}

% 自拟；综合计算：热模型的L²算子范数收敛及一阶修正，用一致尾部控制升级点态Taylor展开。
\Needspace{6\baselineskip}
\begin{exercise}\label{b2:ex:24-new-3}
设整数 \(d\geq1\)，\(H\) 为符号 \(e^{-|\xi|^2}\) 的热算子，\(Q\) 为符号 \(\tfrac12|\xi|^4e^{-|\xi|^2}\) 的平方乘子。证明
\[
 \left\|n\bigl(B_{\sqrt n}^n-H\bigr)+Q\right\|_{2\to2}
       \longrightarrow0\quad(n\to\infty),
\]
并求出准确的一阶范数极限
\[
 \lim_{n\to\infty}n\|B_{\sqrt n}^n-H\|_{2\to2}=2e^{-2}.
\]
需要证明对全部频率的一致控制，仅在固定紧频带内展开并不足够。
\end{exercise}
\begin{solution}
令 \(a=|\xi|^2\)，并记
\[
 D_n(a)=n\left[e^{-a}-(1-a/n)_+^n\right]\geq0.
\]
在任意固定的紧 \(a\) 区间内，Taylor 展开一致给出
\[
 n\log(1-a/n)=-a-\frac{a^2}{2n}+O(n^{-2}),
 \quad D_n(a)\longrightarrow\frac{a^2}2e^{-a}.
\]
余项常数在这一步允许依赖所固定的紧区间。

为控制其外部，先在 \(0\leq a\leq n/2\) 置
\(u_n=-n\log(1-a/n)-a\)。对数级数给出
\[
 0\leq u_n
 =n\sum_{k=2}^\infty\frac{(a/n)^k}{k}
 \leq\frac{a^2}{2n(1-a/n)}\leq\frac{a^2}n.
\]
所以
\(D_n(a)=ne^{-a}(1-e^{-u_n})\leq a^2e^{-a}\)。若 \(a>n/2\)，则直接有
\(D_n(a)\leq ne^{-a}<2ae^{-a}\)，包括 \(a\geq n\) 的截断区域。因此整个半轴上的尾部都由
\((a^2+2a)e^{-a}\) 一致控制，而这个函数在无穷远趋零。先选大紧区间使尾部小，再在紧区间内取极限，即得
\[
 \sup_{a\geq0}\left|D_n(a)-\frac{a^2}2e^{-a}\right|\longrightarrow0.
\]
Plancherel 将符号的一致收敛转为第一条算子范数收敛。最后，\(a^2e^{-a}/2\) 在 \(a=2\) 处取得最大值 \(2e^{-2}\)；由一致收敛，其上确界也是 \(D_n\) 上确界的极限，得到所求准确常数。
\end{solution}

% 自拟；反例训练：临界核的弱L1性质不自动给卷积弱(1,1)，以正核和扩张输入展示端点Young失效。
\Needspace{6\baselineskip}
\begin{exercise}\label{b2:ex:24-new-4}
设整数 \(d\geq1\)、\(\delta_c=(d-1)/2\)。对可测函数 \(u\)，记
\[
 \|u\|_{L^{1,\infty}}
 =\sup_{\lambda>0}\lambda\cdot
       |\{\,x:|u(x)|>\lambda\,\}|.
\]
证明临界 Bochner--Riesz 核满足
\(K_{\delta_c}\in L^{1,\infty}\setminus L^1\)。再考虑正核
\(k(x)=(1+|x|)^{-d}\)，证明它也属于弱 \(L^1\)，但不存在常数 \(C\) 使
\[
 \|k*f\|_{L^{1,\infty}}\leq C\|f\|_1
\]
对所有非负光滑紧支撑输入成立。说明这个反例对临界 Bochner--Riesz 算子本身能够、不能够推出什么。
\end{exercise}
\begin{solution}
临界核渐近和原点处的有界性给出
\(|K_{\delta_c}(x)|\leq C(1+|x|)^{-d}\)。因此其超过 \(\lambda\) 的集合包含于一个体积至多为 \(C'/\lambda\) 的球，弱 \(L^1\) 准范数有限。另一方面，振荡主项在每个充分大的单位径向区间中都有固定比例的非零下界，径向绝对积分至少贡献 \(c/r\)；调和级数发散，所以它不在 \(L^1\)。

同一个分布函数上界说明 \(k\in L^{1,\infty}\)。取非负 \(\psi\in C_c^\infty(\mathbb R^d)\)，使 \(\int\psi=1\)，且在单位球上 \(\psi\geq c_0>0\)，令
\[
 f_R(x)=R^{-d}\cdot\psi(x/R),\quad R\geq4.
\]
这些输入的 \(L^1\) 范数为一。若 \(|x|<R/2\)，在 \(|y|<R\) 上 \(f_R(y)\geq c_0R^{-d}\)，而平移后的积分区域包含 \(B(0,R/2)\)，故
\[
 (k*f_R)(x)
 \geq c_0R^{-d}\int_{|z|<R/2}(1+|z|)^{-d}\,\mathrm dz
 \geq c_dR^{-d}\log R.
\]
最后一个下界来自 \(1\leq|z|\leq R/2\) 上与 \(\mathrm dr/r\) 可比的径向积分。取 \(\lambda\) 为右侧的一半，超水平集的测度至少为 \(|B(0,R/2)|\)，所以
\[
 \|k*f_R\|_{L^{1,\infty}}\geq c'_d\log R\longrightarrow\infty.
\]
这排除了所声称的弱型界。

反例说明不能把 Young 的 \(L^1\) 核假设机械改成弱 \(L^1\) 核假设，并在端点保留同样的结论。它没有证明临界 Bochner--Riesz 算子无弱型界，因为 \(k\) 是另一个没有振荡抵消的核；对原算子的端点判断仍需利用它自身的结构。
\end{solution}

% 自拟；综合证明：保留三半径消项并推广有限半径外推，证明导数界、饱和阶及整数参数例外。
\Needspace{6\baselineskip}
\begin{exercise}\label{b2:ex:24-new-5}
设整数 \(d\geq1\)、\(1<p<\infty\)、\(\delta>0\)，并假定
\(\sup_{R>0}\|B_R^\delta\|_{p\to p}<\infty\)。先求一组不依赖 \(\delta\) 的三个权，使
\[
 C_R=w_1B_R^\delta+w_2B_{\sqrt2R}^\delta+w_3B_{\sqrt3R}^\delta
\]
的低频符号对每个 \(\delta>0\) 都等于 \(1+O(|\xi|^6/R^6)\)，并证明满足这一统一要求的权唯一。进一步，对整数 \(m\geq2\)，证明
\[
 w_j=\frac{(-1)^{m-j}j^{m-1}}{(j-1)!(m-j)!},
 \quad C_R^{(m)}=\sum_{j=1}^m w_j B_{\sqrt jR}^\delta
\]
满足
\[
 \|C_R^{(m)}f-f\|_p\leq C R^{-s}\|f\|_{H^{s,p}},
 \quad R\geq1,\quad0<s\leq2m.
\]
这里 \(H^{s,p}\) 使用 \(J^s=(I-\Delta)^{s/2}\) 的范数；常数可依赖 \(d,p,\delta,m,s\) 和已有的一致算子界，不依赖 \(R,f\)。求第一个可能非零的 Taylor 系数，证明当 \(\delta\notin\{1,\ldots,m-1\}\) 时通常不能超过 \(2m\) 阶；当 \(\delta\) 属于这个有限整数集时，证明上述界反而允许任意 \(s>0\)。
\end{exercise}
\begin{solution}
令 \(v=|\xi|^2/R^2\)。要保留常数并消去 \(v,v^2\)，相应的插值权应满足
\[
 w_1+w_2+w_3=1,
 \quad w_1+\frac{w_2}2+\frac{w_3}3=0,
 \quad w_1+\frac{w_2}4+\frac{w_3}9=0.
\]
方程的唯一解为 \((1/2,-4,9/2)\)。若 \(\delta=1\)，原符号本无二次项；这里的三条条件仍选出同一组通用于所有 \(\delta\) 的插值权，单独要求这一特殊参数的误差阶则不会唯一确定三个权。

一般情形取节点 \(t_j=1/j\)。次数至多 \(m-1\) 的多项式 \(P\) 满足 Lagrange 插值恒等式
\[
 P(0)=\sum_{j=1}^mP(t_j)
          \prod_{k\ne j}\frac{-t_k}{t_j-t_k}.
\]
直接约分所得乘积，就是题中的 \(w_j\)。因此
\[
 \sum_jw_j=1,
 \quad \sum_jw_jj^{-\ell}=0\quad(1\leq\ell<m).
\]
这些条件唯一确定权重，因为节点互异。对 \(P(t)=t^m\)，其插值误差是 \(\prod_j(t-t_j)\)，在零点取值得
\[
 \sum_jw_jj^{-m}=\frac{(-1)^{m-1}}{m!}.
\]
在 \(0\leq v<1\) 内逐项展开，因而有
\[
 \sum_{j=1}^m w_j(1-v/j)^\delta
 =1-\frac{\delta(\delta-1)\cdots(\delta-m+1)}{(m!)^2}v^m
       +O(v^{m+1}).
\]
对固定 \(m,\delta\)，余项在 \(0\leq v\leq1/4\) 上连同任意有限阶导数受控。

下面证明整个 \(L^p\) 误差界。首先，组合算子的统一范数至多为
\(\sum_j|w_j|\) 乘原来的统一界。取光滑截断 \(\chi\)，在 \(|\xi|\leq1/4\) 为一，支撑于 \(|\xi|<1/2\)，并令 \(g=J^sf\)。高频输入可写成 \(R^{-s}T_{a_R}g\)，其中
\[
 a_R(\xi)=R^s\bigl(1-\chi(\xi/R)\bigr)
                         (1+4\pi^2|\xi|^2)^{-s/2}.
\]
在其非零区域有 \(|\xi|\geq R/4\)。对每个多重指标 \(\alpha\)，乘积与链式法则给出
\(|\xi|^{|\alpha|}|\partial^\alpha a_R(\xi)|\leq C_{\alpha,s}\)：
落在截断上的导数只出现于 \(|\xi|\asymp R\)，落在 Bessel 权上的导数则每次提供一个由 \(|\xi|^{-1}\) 控制的因子。故 Mihlin 定理给出一致界，再应用 \(C_R^{(m)}-I\) 的一致界，得到高频误差。

低频误差符号可以写成 \(v^m b(v)\)，其中 \(b\) 在上述紧区间光滑。除去所求的 \(R^{-s}\) 后，对 \(g\) 的符号为
\[
 b_R(\xi)=R^{s-2m}|\xi|^{2m}
    \cdot b(|\xi|^2/R^2)\chi(\xi/R)
                     (1+4\pi^2|\xi|^2)^{-s/2}.
\]
给出其尺度导数界可以明确看见 \(2m\) 的作用。多项式 \(|\xi|^{2m}\) 的每个非零 \(k\) 阶导数至多为 \(C_k|\xi|^{2m-k}\)；归一化光滑因子的导数至多为 \(C_\alpha R^{-|\alpha|}\)，在 \(0<|\xi|\leq R/2\) 上可进一步用 \(C_\alpha|\xi|^{-|\alpha|}\) 控制。Bessel 权的导数也满足相应的尺度界。于是乘积求导后有
\[
 |\xi|^{|\alpha|}\cdot|\partial^\alpha b_R(\xi)|
 \leq C_{\alpha,m,\delta,s}
        R^{s-2m}|\xi|^{2m}(1+|\xi|)^{-s}.
\]
当 \(|\xi|\leq1\) 时，右侧由 \(CR^{s-2m}\) 控制；当 \(1\leq|\xi|\leq R/2\) 时，由 \(C(|\xi|/R)^{2m-s}\) 控制。若 \(s\leq2m\)、\(R\geq1\)，两者一致有界。取到 Mihlin 定理需要的阶数，低频误差也得到 \(CR^{-s}\|g\|_p\) 上界，完成证明。

若 \(\delta\notin\{1,\ldots,m-1\}\)，上述首项系数非零。固定一个非零 Schwartz 输入，使 \(\widehat f\) 光滑紧支撑于远离零点的环带。固定频带上的展开连同全部有限阶导数一致成立，所以逆变换甚至在 Schwartz 拓扑中收敛，并有
\[
 R^{2m}(C_R^{(m)}-I)f\longrightarrow
 -\frac{\delta(\delta-1)\cdots(\delta-m+1)}
          {(m!)^2(2\pi)^{2m}}\,(-\Delta)^mf.
\]
右端非零，排除了对所有这类光滑输入都成立的更高幂次速度。若 \(\delta\in\{1,\ldots,m-1\}\)，则每个 \((1-v/j)^\delta\) 在 \(v<1\) 是次数低于 \(m\) 的多项式，插值条件使组合精确等于一。低频误差于是完全消失；前面的高频论证对每个 \(s>0\) 都适用，给出最后的例外结论。
\end{solution}

% 自拟；综合计算：可积核阈值附近端点范数的准确爆破常数，使用一致渐近和周期均值而非只判发散。
\Needspace{6\baselineskip}
\begin{exercise}\label{b2:ex:24-new-6}
设整数 \(d\geq1\)，令 \(\delta_c=(d-1)/2\)。利用正文中对紧阶数区间一致成立的核渐近，证明
\[
 \lim_{\varepsilon\downarrow0}
      \varepsilon\|K_{\delta_c+\varepsilon}\|_1
 =\frac{4\Gamma((d+1)/2)}{\pi^{3/2}\Gamma(d/2)}.
\]
因此两个端点 \(L^1\)、\(L^\infty\) 的算子范数在接近可积性阈值时都按 \(\varepsilon^{-1}\) 增长。证明中应说明为什么余弦的相位随 \(\varepsilon\) 变化不影响极限。
\end{exercise}
\begin{solution}
记 \(\omega_{d-1}=2\pi^{d/2}/\Gamma(d/2)\)，并令
\[
 a_\delta=\frac{\Gamma(\delta+1)}{\pi^{\delta+1}},
 \quad \theta_\delta=\frac{\pi(d+1+2\delta)}4.
\]
对 \(0\leq\varepsilon\leq1\)、\(r\geq1\)，一致渐近为
\[
 K_{\delta_c+\varepsilon}(r)
 =a_{\delta_c+\varepsilon}r^{-d-\varepsilon}
       \cos(2\pi r-\theta_{\delta_c+\varepsilon})
       +O(r^{-d-\varepsilon-1}).
\]
利用 \(\bigl||u+v|-|u|\bigr|\leq|v|\)，余项对径向绝对积分的贡献至多为
\(C\int_1^\infty r^{-2-\varepsilon}\,\mathrm dr\leq C\)。在单位球内，核由其频率符号的 \(L^1\) 范数一致控制，贡献也为 \(O(1)\)。所以
\[
 \|K_{\delta_c+\varepsilon}\|_1
 =\omega_{d-1}a_{\delta_c+\varepsilon}
   \int_1^\infty r^{-1-\varepsilon}
          |\cos(2\pi r-\theta_{\delta_c+\varepsilon})|\,\mathrm dr+O(1).
\]

函数 \(|\cos(2\pi r-\theta)|\) 的周期均值为 \(2/\pi\)，与 \(\theta\) 无关。减去这个均值后的函数有一个一致有界的原函数
\[
 H_\theta(r)=\int_1^r
       \left(|\cos(2\pi u-\theta)|-\frac2\pi\right)\mathrm du.
\]
有界性由整周期积分为零、剩余长度小于一直接得到，且界与相位无关。对权 \(r^{-1-\varepsilon}\) 分部积分，边界项及导数积分均一致有界，故
\[
 \int_1^\infty r^{-1-\varepsilon}|\cos(2\pi r-\theta)|\,\mathrm dr
 =\frac2{\pi\varepsilon}+O(1)
\]
对所有相位及 \(0<\varepsilon\leq1\) 一致成立。乘 \(\varepsilon\) 并令其趋零，得到极限
\[
 \frac2\pi\omega_{d-1}a_{\delta_c}
 =\frac{4\Gamma((d+1)/2)}{\pi^{3/2}\Gamma(d/2)}.
\]
最后，正文已证明可积卷积核的两个端点算子范数都恰为核的 \(L^1\) 范数；缩放不改变此范数。因此上述极限同时给出了全部半径的两个端点常数的准确增长。
\end{solution}

% 自拟；反例与泛函分析：正阶趋零时的扩展乘子范数必爆破，区分单个符号平滑、统一界及单输入失控。
\Needspace{6\baselineskip}
\begin{exercise}\label{b2:ex:24-new-7}
设整数 \(d\geq2\)、\(1<p<\infty\)、\(p\ne2\)。对 \(\delta>0\)，令
\[
 N_p(\delta)=
 \begin{cases}
 \|B_1^\delta\|_{p\to p},&B_1^\delta\text{ 有有界的 }L^p\text{ 延拓},\\
 +\infty,&\text{否则}.
 \end{cases}
\]
以正文的 Fefferman 球乘子定理为前提，证明
\(N_p(\delta)\to+\infty\) 当 \(\delta\downarrow0\)。进一步，若 \(\delta_n\downarrow0\) 是一个满足每个 \(N_p(\delta_n)<\infty\) 的序列，证明存在同一个 \(f\in L^p\)，使
\[
 \limsup_{n\to\infty}\|B_1^{\delta_n}f\|_p=+\infty.
\]
本题不假定任意正阶都已有界，也不要求构造适用于所有输入的失败结论。
\end{exercise}
\begin{solution}
若第一个结论不成立，则存在 \(\delta_n\downarrow0\) 的子序列与有限常数 \(C\)，使 \(N_p(\delta_n)\leq C\)。其符号在球内部趋于一，在球外为零，且模始终不超过一。因此对任意 Schwartz 输入，Plancherel 与支配收敛给出
\[
 B_1^{\delta_n}f\longrightarrow B_1^0f\quad\text{于 }L^2.
\]
选取几乎处处收敛的子列，再用 Fatou，得到
\[
 \|B_1^0f\|_p^p
 \leq\liminf_n\|B_1^{\delta_n}f\|_p^p
 \leq C^p\|f\|_p^p.
\]
Schwartz 类在有限指数 \(L^p\) 中稠密，因而这会给出硬球乘子的有界延拓，与\cref{b2:thm:24-fefferman}矛盾。此论证排除了所有趋零而范数有界的阶数列，正是扩展意义下的 \(N_p(\delta)\to+\infty\)。

现在假定题中全由有界算子组成的序列确实存在。若每个 \(f\in L^p\) 都满足 \(\sup_n\|B_1^{\delta_n}f\|_p<\infty\)，一致有界原理将给出
\(\sup_nN_p(\delta_n)<\infty\)，与第一步矛盾。因此存在一个输入使输出范数的上确界无穷。每个固定 \(n\) 的输出范数有限，删除任意有限前缀仍然无界，所以其上极限也为无穷。

这是一条依赖深层球乘子定理的退化结论，并没有提供爆破的定量速度。它也不证明某个给定正阶有界；若小阶数已经无界，第一条结论按定义仍然成立，第二条的序列假设则可能无法满足。量词中的统一常数和同一个输入都不能省略。
\end{solution}

% 自拟；综合应用：一维跳跃的完整边界轮廓与过冲，利用累积振荡尾部证明准确Lp误差尺度。
\Needspace{6\baselineskip}
\begin{exercise}\label{b2:ex:24-new-8}
在 \(\mathbb R\) 上固定有限实数 \(\delta>0\)，令 \(K_\delta\) 为单位半径核，并置
\[
 F_\delta(t)=\int_{-\infty}^tK_\delta(y)\,\mathrm dy,
 \quad H(t)=\mathbf1_{(0,\infty)}(t),
 \quad G_\delta=F_\delta-H.
\]
对 \(f=\mathbf1_{(0,1)}\)，证明 \(B_R^\delta f(t/R)\to F_\delta(t)\) 对每个固定 \(t\) 成立。证明 \(F_\delta\) 在某些点大于一，在某些点小于零；不能只从 \(K_\delta\) 变号推出这件事。最后证明对每个 \(1\leq p<\infty\)，
\[
 \lim_{R\to\infty}R^{1/p}\|B_R^\delta f-f\|_p
       =2^{1/p}\|G_\delta\|_p\in(0,\infty),
\]
而一致范数误差始终不小于 \(1/2\)。
\end{exercise}
\begin{solution}
一维中 \(\delta>0\) 已在可积核范围，核为偶实函数且总积分为一。因此 \(F_\delta\) 连续，满足
\[
 F_\delta(-t)=1-F_\delta(t),
 \quad F_\delta(0)=\frac12,
 \quad F_\delta(-\infty)=0,\quad F_\delta(+\infty)=1.
\]
直接对区间积分，得到
\[
 B_R^\delta f(t/R)=\int_{t-R}^tK_\delta(y)\,\mathrm dy
       \longrightarrow F_\delta(t).
\]
这给出左端跳跃在宽度 \(R^{-1}\) 内的完整轮廓。

要证明过冲，必须计算累积尾部。记
\(a_\delta=\Gamma(\delta+1)/\pi^{\delta+1}\)、
\(\theta_\delta=\pi(\delta+1)/2\)。将正文的端点展开多保留一项，得到
\[
 K_\delta(r)=a_\delta r^{-\delta-1}
 \left[\cos(2\pi r-\theta_\delta)
       -\frac{\delta(\delta+1)}{4\pi r}
           \sin(2\pi r-\theta_\delta)\right]
       +O(r^{-\delta-3}).
\]
具体地，在频率端点 \(\xi=1\) 令 \(u=1-\xi\)，振幅为
\(2^\delta u^\delta(1-\delta u/2+O(u^2))\)。对前两项使用端点 Fourier 公式，另一端取共轭；Gamma 递推产生显示的 \(\delta(\delta+1)/(4\pi r)\)，余项由下一阶端点估计控制。这说明这里需要的余项确实比首项低两个幂次。

对首个余弦项分部积分两次，对显示的正弦项分部积分一次，其余项按绝对值积分，得到
\[
 \int_t^\infty K_\delta(r)\,\mathrm dr
 =-\frac{a_\delta}{2\pi}t^{-\delta-1}
       \sin(2\pi t-\theta_\delta)+O(t^{-\delta-2}).
\]
所以
\[
 F_\delta(t)=1+\frac{a_\delta}{2\pi}t^{-\delta-1}
       \sin(2\pi t-\theta_\delta)+O(t^{-\delta-2})
       \quad(t\to+\infty).
\]
选择趋于无穷的一列 \(t\)，使正弦恒为一，主项最终严格大于余项，故 \(F_\delta(t)>1\)。再用 \(F_\delta(-t)=1-F_\delta(t)\)，得到负侧的严格负值。对这样一个固定 \(t\)，前面的轮廓极限还说明足够大的有限半径平均也在对应位置出现过冲。

尾积分公式使 \(|G_\delta(t)|\leq C(1+|t|)^{-\delta-1}\) 在两端成立；在有界区间它有界。因此 \(G_\delta\in L^p\) 对全部有限 \(p\geq1\) 成立，且由零点两侧趋于 \(1/2\) 和 \(-1/2\) 可知它不是零函数。将区间示性函数写成 \(H(x)-H(x-1)\)，卷积恒等式给出
\[
 B_R^\delta f(x)-f(x)
       =G_\delta(Rx)-G_\delta(Rx-R)
       \quad\text{几乎处处}.
\]
于是
\[
 R^{1/p}\|B_R^\delta f-f\|_p
       =\|G_\delta-G_\delta(\,\cdot-R)\|_p.
\]
对任意 \(G\in L^p\)，两个相距趋于无穷的平移副本满足
\(\|G-G(\,\cdot-R)\|_p\to2^{1/p}\|G\|_p\)。证明先取紧支撑截断 \(G_A=G\mathbf1_{[-A,A]}\)：当 \(R>2A\)，两支撑不交，范数的 \(p\) 次方恰为 \(2\|G_A\|_p^p\)。再由三角不等式，替换 \(G_A\) 为 \(G\) 所造成的范数误差至多为 \(2\|G-G_A\|_p\)，令 \(A\to\infty\) 即得。应用于 \(G_\delta\)，完成准确误差极限。

最后，\(B_R^\delta f\) 连续，而 \(f\) 在零点有高度一的跳跃。任意连续函数在这个跳跃两侧的本性一致误差，至少是其零点值到零和一的距离的较大者，因而不小于 \(1/2\)。这不是单个点值的误差，而是在两侧正测度邻域内都无法同时消除的误差；所以不可能有一致范数收敛。
\end{solution}

% !TeX root = ../../Book2.tex
\OptionalChapter{多线性振荡积分与多线性 Fourier 限制}{b2:ch:25}
\BookChapterNumber{b2:ch:25}
\begin{chapterabstract}
多线性估计保留多个输入及其方向关系，使相位抵消与横截几何可以分别检验。从相位和投影族的相容性出发，Loomis--Whitney 不等式、平直横截延拓和曲线双线性模型将显示方向分离带来的增益。一般曲面上，波包把频率方向变为空间中的细长管，多线性限制与解耦则组织这些管在不同尺度上的相互作用。
\end{chapterabstract}

\begin{learninggoals}\begin{itemize}
\item\label{b2:item:25-goal-1}相对于给定投影族判断相位能否被输入吸收，区分标量非退化、相位相互作用和尺度必要条件。
\item\label{b2:item:25-goal-2}计算法向横截量及测度换元因子，说明横截性、曲率与参数归一化各自控制什么。
\item\label{b2:item:25-goal-3}由频率和映射的行列式与逆像重数建立双线性估计，核对对角退化和乘积频率的重叠。
\item\label{b2:item:25-goal-4}准确使用小于一的输出准范数、局部尺度递推、波包尾部和解耦中的平方组合，保留数量损失与尺度损失。
\end{itemize}\end{learninggoals}

线性限制把许多频率方向汇入一个输出，多线性理论则先保留不同方向的乘积。几族细长管若具有独立方向，它们共同相遇的区域会缩小；若输入能够分别消去相位，再大的振荡参数也可能毫无作用。这两种现象说明，不能只根据一个相位的 Hesse 矩阵或一张曲面的曲率判断多输入估计。

\ProseParagraphBreak

本章先把可以直接计算的结构写清楚。\secref{b2:sec:2501}讨论相位与投影、缩放及频率和；\secref{b2:sec:2502}用 Loomis--Whitney 证明平直横截模型；\secref{b2:sec:2503}将曲线乘积变成频率和映射的密度。这些模型分别显示投影维数、横截行列式和逆像重数怎样进入估计常数。

\ProseParagraphBreak

一般曲面的困难在于不同位置、方向和尺度的相互作用。\secref{b2:sec:2504}陈述局部多线性限制的深层输入，并说明把它用于帽分解时还需怎样分配方向。\secref{b2:sec:2505}进一步讨论尺度归纳和波包，保留级数收敛、格点尾部与局部半径；解耦的平方组合则要与点态平方函数分开理解。一个模型中的无损常数，不能自动取消一般定理允许的 \(R^\varepsilon\) 因子。

\ProseParagraphBreak

初次阅读可先通过前三节的具体模型理解横截性：换元公式中的行列式衡量方向分离，Plancherel 等式则把乘积的平方范数化为频率密度。后两节讨论曲面随位置变化后出现的困难，其中一般限制与解耦定理以证明概要介绍。九道选做题围绕等号、退化、集中、重数及求和损失展开，第九题把振荡卷积的内部指数与强端点放在同一个例子中比较。

% !TeX root = ../../Book2.tex
\section{多线性积分形式}
\label{b2:sec:2501}

多线性估计保留多个输入之间的相互作用。同样一个相位，在不同的投影族下可能表现完全不同：有些振荡能够分别吸收到输入中，有些则把不同变量耦合起来。讨论衰减之前，需要先辨认积分的定义域、输入范数以及允许输入怎样随参数变化。

\subsection{多线性型}
\label{b2:subsec:250101}

设 \(d,k,d_1,\ldots,d_k\) 为正整数，\(\pi_j:\mathbb R^d\to\mathbb R^{d_j}\) 为满射线性映射，\(P:\mathbb R^d\to\mathbb R\) 为实相位，\(a\in C_c^\infty(\mathbb R^d;\mathbb C)\)。先对有界可测输入定义
\[
 \Lambda_\lambda^P(f_1,\ldots,f_k)
 \coloneqq\int_{\mathbb R^d}e^{2\pi\mathrm i\lambda P(x)}
             a(x)\prod_{j=1}^k f_j(\pi_jx)\,\mathrm dx.
\]
紧支撑振幅保证此时绝对收敛。满射性也有一个测度上的用途：输入在零测集上改变代表，只会在其逆像上改变被积函数，而该逆像由线性换元与 Fubini 可知仍是零测集。因而这个积分首先是输入等价类上的多线性标量型。

\ProseParagraphBreak

要把定义扩展到 \(L^{p_j}\) 输入，不能只因振幅紧支撑就认为乘积可积。例如在二维取两个相同投影 \(\pi_1(x,y)=\pi_2(x,y)=x\)，两个局部可积函数的乘积仍可能不可积。实际需要先建立
\[
 \int |a(x)|\prod_j|f_j(\pi_jx)|\,\mathrm dx
 \leq C\prod_j\|f_j\|_{p_j},
\]
或者先在稠密测试类上证明振荡型的有界估计，再按该估计定义延拓。后一个延拓未必还能解释为每个输入的绝对收敛积分。若某个输入指数为无穷，也不能使用有限指数的范数稠密论证。

\ProseParagraphBreak

投影几何在最简单的例子中已经可见。对坐标投影 \(\pi_1(x,y)=x\)、\(\pi_2(x,y)=y\)，若振幅支撑于 \(I\times J\)，则
\[
 |\Lambda_\lambda^P(f,g)|
 \leq \|a\|_\infty
       \left(\int_I|f|\right)\left(\int_J|g|\right).
\]
有限区间上的 Hölder 把它转成任意 \(1\leq p,q\leq\infty\) 的局部界。若两个投影都等于 \(x\)，则需要直接控制 \(f(x)g(x)\)，指数关系随之改变。相位完全相同，并不能消除这种输入几何的差别。

\subsection{相位相互作用}
\label{b2:subsec:250102}

相位中的一元项可以藏到某一个输入中。这里的“一元”是相对于投影而言，并不限于实数变量上的函数。

\Needspace{6\baselineskip}
\begin{definition}[相对于投影族的退化相位]\label{b2:def:25-degenerate-phase}
设 \(d,k,d_1,\ldots,d_k\) 为正整数，\(\pi_j:\mathbb R^d\to\mathbb R^{d_j}\) 为给定满射线性映射，\(\Omega\subset\mathbb R^d\) 为开集，\(P\) 为其上的实可测函数。若存在实可测函数 \(P_j\)，定义于 \(\pi_j(\Omega)\)，使
\[
 P(x)=\sum_{j=1}^k P_j(\pi_jx)
 \quad\text{对几乎所有 }x\in\Omega,
\]
则称 \(P\) 在 \(\Omega\) 上是相对于该投影族的\term[tuihua-xiangwei]{退化相位}[degenerate phase]。若限定研究多项式相位及多项式分解，应另行把 \(P,P_j\) 都限定为多项式。
\end{definition}

\Needspace{6\baselineskip}
\begin{proposition}[可吸收相位与统一衰减]\label{b2:prop:25-phase-absorption}
设 \(d,k,d_1,\ldots,d_k\) 为正整数，\(\Omega\subset\mathbb R^d\) 为开集，\(a\in C_c^\infty(\Omega;\mathbb C)\)，\(\pi_j:\mathbb R^d\to\mathbb R^{d_j}\) 为满射线性映射。设 \(P,Q:\Omega\to\mathbb R\)、\(P_j:\pi_j(\Omega)\to\mathbb R\) 为可测函数，且在 \(\Omega\) 上几乎处处有
\[
 P=Q+\sum_jP_j\circ\pi_j.
\]
对实数 \(\lambda\) 和有界紧支撑复输入，定义
\[
 \Lambda_\lambda^P(f_1,\ldots,f_k)
 =\int_\Omega e^{2\pi\mathrm i\lambda P(x)}
                  a(x)\prod_jf_j(\pi_jx)\,\mathrm dx,
\]
\(\Lambda_\lambda^Q\) 同样定义。令
\(M_{j,\lambda}f_j=e^{2\pi\mathrm i\lambda P_j}\cdot f_j\)，并在 \(\pi_j(\Omega)\) 外任意延拓相位。则
\[
 \Lambda_\lambda^P(f_1,\ldots,f_k)
 =\Lambda_\lambda^Q(M_{1,\lambda}f_1,\ldots,M_{k,\lambda}f_k).
\]
各 \(M_{j,\lambda}\) 在每个 \(L^{p_j}\)、\(1\leq p_j\leq\infty\) 上都是可逆等距映射。因此两型在该测试类上具有相同的最佳范数常数。特别地，若 \(Q=0\)，且存在固定有界紧支撑输入 \(g_j\) 使
\[
 \int a(x)\prod_jg_j(\pi_jx)\,\mathrm dx\ne0,
\]
便不可能有 \(C|\lambda|^{-\epsilon}\prod_j\|f_j\|_{p_j}\) 型的统一上界，其中 \(C\) 与输入、\(|\lambda|\geq1\) 无关，\(\epsilon>0\)。
\end{proposition}
\begin{proof}
把指数中相位之和拆成乘积，就得到型的恒等式。乘子模为一，其逆为相反相位，所以输入范数及可取输入集合不变，最佳常数相同。

在最后的情形取 \(f_{j,\lambda}=e^{-2\pi\mathrm i\lambda P_j}\cdot g_j\)。所得输出等于题设非零积分，输入范数则与 \(\lambda\) 无关。若有正幂次衰减，令 \(|\lambda|\to\infty\) 就会使这个非零数的模趋于零，产生矛盾。这里使用的是随参数变化的一族输入；统一算子界本来就必须同时控制它们。
\end{proof}

在坐标投影 \((x,y)\mapsto x,y\) 下，\(x^2+y^2\) 是退化相位，尽管它作为二维标量相位的 Hesse 矩阵可逆。反之，\(xy\) 不能写成纯 \(x\) 项与纯 \(y\) 项之和。若 \(P\in C^2(I\times J)\)，其中 \(I,J\) 为开区间，则这种分解恰好对应 \(\partial_x\partial_yP=0\)。充分性可由矩形上的积分恒等式看出：
\[
 \begin{aligned}
 P(x,y)-P(x,y_0)-P(x_0,y)+P(x_0,y_0)
 =\int_{x_0}^x\int_{y_0}^y
       \partial_x\partial_yP(u,v)\,\mathrm dv\,\mathrm du.
 \end{aligned}
\]
必要性不要求先知道分解项可微：分解使矩形差分几乎处处为零，再由 \(P\) 的连续性推广到所有矩形；对差分取极限即可。

\Needspace{6\baselineskip}
\begin{proposition}[混合二次相位的平方估计]\label{b2:prop:25-mixed-phase}
设 \(\lambda\in\mathbb R\setminus\{0\}\)，对复 Schwartz 函数定义
\[
 \mathcal L_\lambda(f,g)
 =\int_{\mathbb R^2}e^{2\pi\mathrm i\lambda xy}
                       f(x)g(y)\,\mathrm dx\,\mathrm dy.
\]
它唯一连续延拓为 \(L^2(\mathbb R)\times L^2(\mathbb R)\) 上的双线性型，且最佳常数为
\[
 |\mathcal L_\lambda(f,g)|
 \leq|\lambda|^{-1/2}\|f\|_2\cdot\|g\|_2.
\]
\end{proposition}
\begin{proof}
先在 Schwartz 类上使用 Fubini，得到
\[
 \mathcal L_\lambda(f,g)=\int f(x)\widehat g(-\lambda x)\,\mathrm dx.
\]
Cauchy--Schwarz、换元与 Plancherel 给出
\(\|\widehat g(-\lambda\,\cdot)\|_2=|\lambda|^{-1/2}\|g\|_2\)，从而得到估计。以 Schwartz 函数分别逼近两个 \(L^2\) 输入，双线性差分及此界保证型的极限存在且与逼近无关。

任取非零 Schwartz 函数 \(g\)，再取
\(f(x)=\overline{\widehat g(-\lambda x)}\)，便在 Cauchy--Schwarz 中取等。因此显示的常数不能减小。一般平方输入上的值由一维配对定义；不要求原来的二维积分绝对收敛。
\end{proof}

这项计算给出了真正的参数增益。若另乘分离振幅 \(a_1(x)a_2(y)\)，则先把它们吸收到 \(f,g\) 中，上界再乘 \(\|a_1\|_\infty\cdot\|a_2\|_\infty\)。给相位添加任意纯 \(x\) 项、纯 \(y\) 项不会改变相位型的最佳常数。一般多投影问题需要相应的非退化条件；单凭“不能分解”并没有在此建立适用于所有维数、相位和指数的统一定理。

\subsection{缩放}
\label{b2:subsec:250103}

缩放先约束无振荡的绝对积分，也帮助辨认局部化后的尺度因子。对于固定紧支撑振幅，全空间的缩放等式通常会被截断打破，不能不加区别地搬用。

\Needspace{6\baselineskip}
\begin{proposition}[投影积分的维数必要条件]\label{b2:prop:25-dimension-necessary}
设 \(d,d_j,k\) 为正整数，\(\pi_j:\mathbb R^d\to\mathbb R^{d_j}\) 为满射线性映射，\(1\leq p_j\leq\infty\)，约定 \(1/\infty=0\)。若存在有限常数 \(C\)，使对全部非负有界紧支撑输入都有
\[
 \int_{\mathbb R^d}\prod_{j=1}^k f_j(\pi_jx)\,\mathrm dx
 \leq C\prod_j\|f_j\|_{p_j},
\]
则
\[
 d=\sum_j\frac{d_j}{p_j},
 \quad
 \dim V\leq\sum_j\frac{\dim\pi_j(V)}{p_j}
 \quad\text{对每个线性子空间 }V\subset\mathbb R^d.
\]
\end{proposition}
\begin{proof}
取各输入为相应空间中单位球的示性函数。所有投影的核的公共交若非零，左侧在该方向不变且在原点附近为一，积分即无穷，与假设矛盾。公共核为零时，映射 \(x\mapsto(\pi_1x,\ldots,\pi_kx)\) 单射，有限维范数比较保证共同逆像有界。因此这组输入的积分有限且严格为正。

将它们换成 \(f_j(r\,\cdot)\)。左侧乘 \(r^{-d}\)，右侧乘 \(r^{-\sum_jd_j/p_j}\)。同一个常数须对所有 \(r>0\) 有效；分别令 \(r\downarrow0\) 和 \(r\to\infty\)，得到维数等式。

再固定子空间 \(V\)，记 \(m=\dim V\)、\(V_j=\pi_j(V)\)、\(m_j=\dim V_j\)。对 \(R\geq1\)，考虑
\[
 \mathcal T_R=\{v+w:v\in V,\ |v|\leq R,\ 
                      w\in V^\perp,\ |w|\leq1\}.
\]
其体积为 \(c_VR^m\)。映射 \(\pi_j\) 的线性有界性说明，\(\pi_j(\mathcal T_R)\) 包含于某个以 \(V_j\) 为长方向的盒状区域：沿 \(V_j\) 长度至多 \(C_jR\)，沿 \(V_j^\perp\) 厚度至多 \(C_j\)。取该区域示性函数为 \(f_j\)，其 \(L^{p_j}\) 范数至多为 \(C'_jR^{m_j/p_j}\)，无穷指数时为一。原积分在 \(\mathcal T_R\) 上被积函数为一，故
\[
 c_VR^m\leq C'\,R^{\sum_jm_j/p_j}.
\]
令 \(R\to\infty\)，得到子空间不等式。零维分量按体积因子一解释，不影响论证。
\end{proof}

只检查总维数是不够的。例如在 \(\mathbb R^2\) 上取两个相同投影到 \(x\) 轴，并取 \(p_1=p_2=1\)，虽然总维数等式成立，竖直直线却违反子空间条件；沿该直线没有任何输入提供衰减。子空间条件因而记录了总维数等式无法发现的方向退化。

\ProseParagraphBreak

曲面延拓还具有非各向同性缩放。抛物面的切向频率变化为 \(\delta\)，法向变化为 \(\delta^2\)，相应空间尺度必须分别记录。

\Needspace{6\baselineskip}
\begin{proposition}[抛物面的重标度]\label{b2:prop:25-parabolic-rescaling}
设 \(n\geq1\)，对紧支撑的 \(g\in L^1(\mathbb R^n;\mathbb C)\) 定义参数测度下的延拓
\[
 E_{\mathrm{par}}g(x',t)
 =\int_{\mathbb R^n}e^{2\pi\mathrm i(x'\cdot\xi+t|\xi|^2)}
                       g(\xi)\,\mathrm d\xi.
\]
取 \(\xi_0\in\mathbb R^n\)、\(\delta>0\)，令
\[
 h(u)=g(\xi_0+\delta u),
 \quad
 T_{\delta,\xi_0}(x',t)
 =\bigl(\delta(x'+2t\xi_0),\,\delta^2t\bigr).
\]
则
\[
 E_{\mathrm{par}}g(x',t)
 =\delta^n e^{2\pi\mathrm i(x'\cdot\xi_0+t|\xi_0|^2)}
       E_{\mathrm{par}}h\bigl(T_{\delta,\xi_0}(x',t)\bigr).
\]
对任意可测区域 \(\Omega\subset\mathbb R^{n+1}\) 和 \(0<q\leq\infty\)，有
\[
 \|E_{\mathrm{par}}g\|_{L^q(\Omega)}
 =\delta^{\,n-(n+2)/q}
       \|E_{\mathrm{par}}h\|_{L^q(T_{\delta,\xi_0}\Omega)}.
\]
无穷指数按 \(1/q=0\) 解释；有限指数等式也允许两侧同为无穷。若 \(g\in L^p\)、\(1\leq p\leq\infty\)，则另有
\(\|g\|_p=\delta^{n/p}\|h\|_p\)。
\end{proposition}
\begin{proof}
在延拓积分中令 \(\xi=\xi_0+\delta u\)，把相位展开为常数项与
\(\delta(x'+2t\xi_0)\cdot u+\delta^2t|u|^2\)，便得第一个等式。线性映射 \(T_{\delta,\xi_0}\) 是剪切与两个不同膨胀的复合，其行列式为 \(\delta^{n+2}\)。对绝对值的 \(q\) 次方换元并取 \(q\) 次根，得到有限指数公式；无穷指数由可逆线性映射保持零测集得到。输入范数等式是 \(n\) 维的普通换元。
\end{proof}

例如一张 \(\delta\) 帽经过重标度变成固定大小的参数片，但半径 \(R\) 的空间球变成一个剪切后的椭球：切向尺度约为 \(\delta R\)，时间尺度为 \(\delta^2R\)。不能把两者都写成 \(\delta^2R\)。若要在重标度后继续使用球上的定理，尚需覆盖这些区域，并计入求和因子。\cref{b2:prop:23-parabolic-packet}进一步处理了光滑单帽的衰减；本章后面用 Fourier 级数把一般输入分成这样的局部波包。

\subsection{近似正交性}
\label{b2:subsec:250104}

平方理论中的正交性发生在输出频率上。下面把乘积情形直接写入陈述，避免只检查原输入而遗漏相互作用。

\Needspace{6\baselineskip}
\begin{proposition}[输出频率和的有限重叠]\label{b2:prop:25-output-overlap}
设有限族 \(u_\alpha,v_\alpha\) 为 \(\mathbb R^d\) 上的 Schwartz 函数，且
\(\operatorname{supp}\widehat u_\alpha\subset A_\alpha\)、
\(\operatorname{supp}\widehat v_\alpha\subset B_\alpha\)，其中 \(A_\alpha,B_\alpha\) 为紧集。若几乎处处有
\[
 \sum_\alpha\mathbf1_{A_\alpha+B_\alpha}(\xi)\leq L,
\]
则
\[
 \left\|\sum_\alpha u_\alpha v_\alpha\right\|_2^2
 \leq L\sum_\alpha\|u_\alpha v_\alpha\|_2^2.
\]
更一般地，任意有限族 \(L^2\) 函数的 Fourier 支撑若至多 \(L\) 重叠，也有同一个平方和估计。
\end{proposition}
\begin{proof}
Fourier 乘积公式给出
\(\widehat{u_\alpha v_\alpha}=\widehat u_\alpha*\widehat v_\alpha\)，其支撑包含于 \(A_\alpha+B_\alpha\)。在几乎每个频率点，至多有 \(L\) 个非零项，有限维 Cauchy--Schwarz 因而给出
\[
 \left|\sum_\alpha\widehat{u_\alpha v_\alpha}(\xi)\right|^2
 \leq L\sum_\alpha|\widehat{u_\alpha v_\alpha}(\xi)|^2.
\]
积分并使用 Plancherel，即得结论。这个最后的论证只使用平方可积性和支撑重叠，因此也证明了第二个断言。
\end{proof}

输入频率片互不相交，不保证它们的和集合互不相交。例如四片分别靠近 \(0,1,2,3\)，两组和 \(0+3\) 与 \(1+2\) 仍靠近同一个频率。乘积输出因而可能同相叠加。对许多项求和时，重叠数 \(L\) 是必须计算的量。

\ProseParagraphBreak

空间局部化也会改变这个数。若选 \(\psi\in\mathcal S(\mathbb R^d)\) 且 \(\operatorname{supp}\widehat\psi\subset B(0,c)\)，令 \(\psi_R(x)=\psi(x/R)\)，则 \(\psi_Ru_\alpha v_\alpha\) 的 Fourier 支撑包含于
\[
 A_\alpha+B_\alpha+B(0,c/R).
\]
因此前面的证明给出局部化平方估计，但所需的是这些扩大集合的重叠数。可以选 \(\psi\) 使其在单位球上模有正下界，从而控制空间球上的范数。若改用空间紧支撑的光滑截断，它的 Fourier 变换通常没有紧支撑，只具有快速衰减尾部；此时应按尾部求和，不能把“稍微扩大支撑”当作精确等式。

% !TeX root = ../../Book2.tex
\section{横截性}
\label{b2:sec:2502}

横截性把方向分离写成一个定量下界。对 \(d\) 个超曲面片，最完整的情形是各取一个法向后，它们始终张成整个环境空间。这个条件控制不同方向共同形成的体积；即使各片完全平直，也能产生单个因子所没有的乘积估计。

\subsection{横截性}
\label{b2:subsec:250201}

\Needspace{6\baselineskip}
\begin{definition}[完整横截的超曲面片]\label{b2:def:25-transversality}
设整数 \(d\geq2\)，\(S_1,\ldots,S_d\subset\mathbb R^d\) 为 \(C^1\) 超曲面片，在各片上选单位法向 \(n_j\)。若存在 \(\nu>0\)，使对任意 \(\xi_j\in S_j\) 都有
\[
 \bigl|\det(n_1(\xi_1),\ldots,n_d(\xi_d))\bigr|\geq\nu,
\]
则称这些片满足一致的\term[hengjie-xing]{横截性}[transversality]条件，或称它们完整横截。法向的列次序和各自取向不影响这个绝对值条件。
\end{definition}

行列式是法向所张平行多面体的体积。在二维，它等于两法向夹角的正弦的绝对值；在三维，仅有两两夹角远离零还不足以排除三个方向共面。条件中的量词也不能只落在每片的中心：估计面对的是整片输入，需要同时控制所有可能的法向组合。

\ProseParagraphBreak

若曲面由 \(\Sigma_j:U_j\subset\mathbb R^{d-1}\to\mathbb R^d\) 参数化，切向偏导的外积可用其法向向量 \(Y_j\) 表示，方向由参数次序选定，而
\[
 |Y_j(u)|=
 \sqrt{\det\bigl(D\Sigma_j(u)^{\mathsf T}D\Sigma_j(u)\bigr)}.
\]
这个量是参数测度变成面积测度的 Jacobi 因子，满足
\[
 \mathrm d\sigma_j(\Sigma_j(u))=|Y_j(u)|\,\mathrm du,
 \quad
 Y_j(u)=|Y_j(u)|\,n_j(\Sigma_j(u)).
\]
因而
\[
 \begin{aligned}
 &\bigl|\det(Y_1(u_1),\ldots,Y_d(u_d))\bigr|\\
 &\quad=\left(\prod_j|Y_j(u_j)|\right)
    \bigl|\det(n_1(\Sigma_1(u_1)),\ldots,n_d(\Sigma_d(u_d)))\bigr|.
 \end{aligned}
\]
当各 Jacobi 因子有固定正上下界时，这两种横截条件可以互换；常数必须带上这些界。参数被压缩时，原始向量 \(Y_j\) 的长度可能变小，即使单位法向完全不变。因此采用参数测度的定理与采用面积测度的定理应先完成这个换算。

\ProseParagraphBreak

抛物面 \(\Sigma(\xi)=(\xi,|\xi|^2)\)、\(\xi\in\mathbb R^{d-1}\)，提供了可直接计算的例子。其一个未归一化法向是 \(Y(\xi)=(-2\xi,1)\)。在三维，三个参数点 \(\xi_1,\xi_2,\xi_3\in\mathbb R^2\) 满足
\[
 \bigl|\det(Y(\xi_1),Y(\xi_2),Y(\xi_3))\bigr|
 =4\bigl|\det(\xi_2-\xi_1,\xi_3-\xi_1)\bigr|.
\]
从后两列减去第一列，再沿最后一行展开便得到这个公式。因此完整横截要求参数三角形具有定量面积，而不只是三个顶点两两分离。曲率让同一曲面的法向发生变化，但取出的方向是否张成空间，仍须单独判断。

\subsection{Loomis–Whitney 不等式}
\label{b2:subsec:250202}

先把法向取为坐标轴。此时沿第 \(j\) 个法向不变的函数，只依赖删去 \(x_j\) 后的坐标。普通 Hölder 不能直接使用这些函数的全空间积分，因为每个因子都沿一个方向无限延伸；需要利用不同缺失方向互相补足的结构。

\Needspace{6\baselineskip}
\begin{theorem}[Loomis--Whitney 不等式]\label{b2:thm:25-loomis-whitney}
设整数 \(d\geq2\)，\(\pi_j:\mathbb R^d\to\mathbb R^{d-1}\) 删除第 \(j\) 个坐标。对非负可积函数 \(F_j\) 有\term[loomis-whitney-budengshi]{Loomis--Whitney 不等式}[Loomis--Whitney inequality]
\[
 \int_{\mathbb R^d}\prod_{j=1}^d
      F_j(\pi_jx)^{1/(d-1)}\,\mathrm dx
 \leq\prod_{j=1}^d
      \left(\int_{\mathbb R^{d-1}}F_j\right)^{1/(d-1)}.
\]
其中常数一是最佳常数。
\end{theorem}
\begin{proof}
对维数归纳。当 \(d=2\) 时，两个因子分别只依赖 \(x_2,x_1\)，Tonelli 使左侧恰好等于右侧。现设 \(d\geq3\)，并已知 \(d-1\) 维结论。先对 \(x_d\) 积分。第 \(d\) 个因子与它无关，对其余 \(d-1\) 个因子使用等指数 Hölder，令
\[
 G_j(\pi'_jx')=
 \int_{\mathbb R}F_j(\pi_j(x',x_d))\,\mathrm dx_d
 \quad(1\leq j<d),
\]
其中 \(x'=(x_1,\ldots,x_{d-1})\)，\(\pi'_j\) 从这 \(d-1\) 个坐标中删去第 \(j\) 个。每个 \(G_j\) 是 \(\mathbb R^{d-2}\) 上的非负可积函数。于是原积分至多为
\[
 \int_{\mathbb R^{d-1}}F_d(x')^{1/(d-1)}
                  \prod_{j<d}G_j(\pi'_jx')^{1/(d-1)}\,\mathrm dx'.
\]
对 \(F_d\) 因子与其余乘积使用指数 \(d-1\) 和 \((d-1)/(d-2)\) 的 Hölder，得到
\[
 \left(\int F_d\right)^{1/(d-1)}
 \left[\int_{\mathbb R^{d-1}}
       \prod_{j<d}G_j(\pi'_jx')^{1/(d-2)}\,\mathrm dx'\right]^{(d-2)/(d-1)}.
\]
方括号由归纳假设控制，随后用 Tonelli 所给的
\(\int G_j=\int F_j\)，正好得到所求右侧。也可先将输入截断为有界紧支撑函数，再用单调收敛去掉截断；这样每个中间积分都由相应的有限量控制。

为证明最佳性，取非零、非负可积函数 \(h_i:\mathbb R\to[0,\infty)\)，并令
\[
 F_j(\pi_jx)=\prod_{i\ne j}h_i(x_i).
\]
左侧乘积中每个 \(h_i\) 出现 \(d-1\) 次，所以被积函数为 \(\prod_i h_i(x_i)\)，积分等于 \(\prod_i\int h_i\)。右侧也是同一个数。因此常数不能小于一。
\end{proof}

在三维中，这个式子写成
\[
 \int_{\mathbb R^3}
   \bigl[F_1(x_2,x_3)F_2(x_1,x_3)F_3(x_1,x_2)\bigr]^{1/2}\,\mathrm dx
 \leq\prod_{j=1}^3\left(\int F_j\right)^{1/2}.
\]
每一对坐标都出现一次，正是缺失方向相互补足的方式。全空间尺度也完全平衡：把所有坐标放大 \(r\)，两侧都产生 \(r^d\) 的体积因子。

\ProseParagraphBreak

把输入取为截面的重数函数，还能得到一个最直接的管束模型。它只处理每族内部完全平行、各族方向恰为坐标轴的管，尚未允许方向在帽内变化。

\Needspace{6\baselineskip}
\begin{corollary}[坐标方向平行管的重数估计]\label{b2:cor:25-parallel-tubes}
设整数 \(d\geq2\)、\(r>0\)。对每个 \(j\)，取有限族
\[
 T_{j,\alpha}=\{x\in\mathbb R^d:
                \pi_jx\in z_{j,\alpha}+[-r,r]^{d-1}\},
 \quad a_{j,\alpha}\geq0,
\]
并令 \(N_j(x)=\sum_\alpha a_{j,\alpha}\mathbf1_{T_{j,\alpha}}(x)\)。则
\[
 \int_{\mathbb R^d}\prod_{j=1}^dN_j(x)^{1/(d-1)}\,\mathrm dx
 \leq(2r)^d\prod_{j=1}^d
                      \left(\sum_\alpha a_{j,\alpha}\right)^{1/(d-1)}.
\]
若把每条管再截短为任意有限线段所对应的子集，同一个上界仍然成立。
\end{corollary}
\begin{proof}
将 \(F_j\) 取为相应截面立方体示性函数的加权和，则
\(\int F_j=(2r)^{d-1}\sum_\alpha a_{j,\alpha}\)。代入 Loomis--Whitney 不等式即得结论。截短只会减小每个非负重数函数，故不改变上界的有效性。
\end{proof}

这里同族内的管可以完全重叠，输入的权重总和已经计算了这种重数。它说明横截估计控制的是不同族的共同聚集，而不是要求每一族自行具有不交性。一般多线性 Kakeya 理论还需要处理每族的方向变化，不能由这个固定方向模型直接取极限得到。

\subsection{平直横截延拓估计}
\label{b2:subsec:250203}

平面延拓沿法向不变，可以通过线性换元变成上面的删坐标结构。换元的体积因子与各平面上的面积因子不同，必须分别计算；准确的行列式幂次就来自这两项。

\Needspace{6\baselineskip}
\begin{proposition}[平直横截延拓]\label{b2:prop:25-flat-extension}
设整数 \(d\geq2\)，\(n_1,\ldots,n_d\in\mathbb R^d\) 为单位向量，
\(N=(n_1,\ldots,n_d)\) 可逆。令 \(H_j=n_j^\perp\)，\(\sigma_j\) 为其欧氏面积测度。先对
\(g_j\in L^1(H_j;\mathbb C)\cap L^2(H_j;\mathbb C)\) 定义
\[
 E_jg_j(x)=\int_{H_j}e^{2\pi\mathrm i x\cdot\xi}
                          g_j(\xi)\,\mathrm d\sigma_j(\xi).
\]
令 \(q=2/(d-1)\)，则乘积映射连续延拓到全部 \(L^2(H_j)\) 输入，并满足最佳常数估计
\[
 \left\|\prod_{j=1}^dE_jg_j\right\|_{L^q(\mathbb R^d)}
 \leq|\det N|^{-1/2}\prod_j\|g_j\|_{L^2(H_j)}.
\]
当 \(q<1\) 时，左侧指 \((\int|\prod_jE_jg_j|^q)^{1/q}\) 的准范数。一般平方输入的单个延拓由平面上的逆 Fourier 变换再沿法向保持常值来解释。
\end{proposition}
\begin{proof}
记 \(D=|\det N|>0\)，令 \(P_j\) 为到 \(H_j\) 的正交投影。平面上的逆 Fourier 变换 \(v_j=\mathcal F_{H_j}^{-1}g_j\) 满足
\[
 E_jg_j(x)=v_j(P_jx),
 \quad \|v_j\|_{L^2(H_j)}=\|g_j\|_{L^2(H_j)}.
\]
第一式在初始可积输入上是普通积分；平方输入由平面上的 Plancherel 定义 \(v_j\)。满射线性投影的零测集逆像仍为零测集，所以拉回函数的几乎处处代表不影响乘积。

令 \(x=Ny\)。因 \(P_jn_j=0\)，映射 \(P_jNy\) 不依赖 \(y_j\)。把剩下 \(d-1\) 个坐标组成的线性映射记为
\[
 M_j:\mathbb R^{d-1}\to H_j,
 \quad
 M_j(\pi_jy)=\sum_{\ell\ne j}P_jn_\ell\,y_\ell.
\]
它的面积 Jacobi 因子为 \(D\)。事实上，以这些投影列向量和单位法向 \(n_j\) 组成的体积，在将列 \(P_jn_\ell\) 换回 \(n_\ell\) 时不变，因为二者相差 \(n_j\) 的倍数；最后所得行列式就是原行列式，至多相差列置换的符号。这也证明 \(M_j\) 可逆。

令 \(F_j(\pi_jy)=|v_j(M_j(\pi_jy))|^2\)。面积换元与 Plancherel 给出
\[
 \int_{\mathbb R^{d-1}}F_j
 =D^{-1}\|v_j\|_2^2=D^{-1}\|g_j\|_2^2.
\]
空间换元 \(x=Ny\) 却贡献 \(D\)，因此 Loomis--Whitney 给出
\[
 \begin{aligned}
 \int_{\mathbb R^d}\prod_j|E_jg_j(x)|^{2/(d-1)}\,\mathrm dx
 &=D\int_{\mathbb R^d}\prod_jF_j(\pi_jy)^{1/(d-1)}\,\mathrm dy\\
 &\leq D^{-1/(d-1)}\prod_j\|g_j\|_2^{2/(d-1)}.
 \end{aligned}
\]
取 \((d-1)/2\) 次幂，就得到所求行列式常数。这个推导直接适用于一般平方输入的 \(v_j\)，无须预设延拓积分绝对收敛。

乘积映射的连续性也可从同一估计得出：两个输入组的乘积之差按逐个替换因子展开为 \(d\) 项，每一项应用已经证明的界。若 \(q\geq1\)，使用三角不等式；若 \(0<q<1\)，则先使用
\(\|\sum_\ell h_\ell\|_q^q\leq\sum_\ell\|h_\ell\|_q^q\)。稠密类上的逼近因而有唯一极限，且与上面用 \(v_j\) 定义的乘积一致。

最后，在 \(y\) 坐标中选
\[
 F_j(\pi_jy)=\prod_{i\ne j}e^{-2\pi y_i^2},
 \quad v_j(z)=F_j(M_j^{-1}z)^{1/2}.
\]
这些 \(v_j\) 为非零 Gaussian 函数，逆向取其 Fourier 变换得到可积且平方可积的 \(g_j\)。Loomis--Whitney 中恰好取等，其余步骤都是等式，所以 \(D^{-1/2}\) 确实为最佳常数。
\end{proof}

在二维，两个法向夹角为 \(\vartheta\) 时，常数为
\(|\sin\vartheta|^{-1/2}\)。方向靠近时的发散是准确的，不能通过改进证明消去。在三维，三个法向可以两两保持明显分离而逐渐共面，行列式仍会趋于零；这说明完整横截下界具有不可替代的作用。

\ProseParagraphBreak

各平面还可以平移为仿射平面 \(H_j+\xi_j^0\)。用平移后的密度定义延拓，只会在 \(E_jg_j\) 前乘
\(e^{2\pi\mathrm i x\cdot\xi_j^0}\)，所以乘积的模和最佳常数不变。这项估计只需要相互横截，不需要曲率；曲率在后续曲面问题中的作用，是让同一曲面的不同片提供不同方向，并在局部重标度之后控制法向的变化。

% !TeX root = ../../Book2.tex
\section{曲线的双线性延拓估计}
\label{b2:sec:2503}

两个延拓函数的乘积保留了两片频率怎样相遇的信息。对二维曲线，两个参数的总维数恰好等于输出频率空间的维数；若频率和映射可以换元，乘积的 Fourier 密度就能直接写出。这样得到的平方范数公式不仅证明一个上界，还能读出横截常数、参数测度和多重覆盖各自的作用。

\subsection{抛物线延拓的频率和映射}
\label{b2:subsec:250301}

令 \(\gamma(t)=(t,t^2)\)。对有限区间 \(I\) 和复函数
\(f\in L^2(I)\)，定义
\[
 E_If(x,y)=\int_I e^{2\pi\mathrm i(xt+yt^2)}f(t)\,\mathrm dt.
\]
这里采用参数测度 \(\mathrm dt\)，不是抛物线的弧长测度。
由于
\(\|f\|_{L^1(I)}\leq |I|^{1/2}\cdot\|f\|_{L^2(I)}\)，
积分绝对收敛，并给出一个有界连续函数；连续性来自支配收敛。
它也可以理解为有限复测度 \(\gamma_*(f\,\mathrm dt)\) 的逆 Fourier 变换。

取两个不交紧区间 \(I,J\)，先设 \(I\) 位于 \(J\) 左侧。Fubini 给出
\[
 E_If(x,y)\,E_Jg(x,y)
 =\int_{I\times J}
 e^{2\pi\mathrm i[(t+s)x+(t^2+s^2)y]}f(t)g(s)
 \,\mathrm dt\,\mathrm ds.
\]
因此要研究的是参数乘积测度经过映射
\[
 \Psi(t,s)=(t+s,t^2+s^2)
\]
以后是否具有平方可积密度。这一步把空间侧的乘积变成频率侧的和，尚未使用任何振荡估计。

\ProseParagraphBreak
写 \(u=t+s\)、\(v=t^2+s^2\)，则
\[
 2v-u^2=(s-t)^2,\quad
 t=\frac{u-\sqrt{2v-u^2}}2,\quad
 s=\frac{u+\sqrt{2v-u^2}}2.
\]
在 \(I\times J\) 上，正平方根的选择由 \(t<s\) 决定。
所以和与平方和先确定无序对 \(\{t,s\}\)，两个区间的先后次序再确定有序对。
这说明 \(\Psi\) 在整个输入矩形上一对一；它不是仅从局部逆函数定理得到的结论。
另一方面，
\[
 D\Psi(t,s)=
 \begin{pmatrix}1&1\\2t&2s\end{pmatrix},
 \quad
 |\det D\Psi(t,s)|=2(s-t).
\]
若两个区间至少相隔 \(\rho>0\)，这个行列式至少为 \(2\rho\)。
两类信息缺一不可：非零行列式负责局部换元，全局一对一负责避免把不同逆像的贡献遗漏。

\subsection{双线性估计}
\label{b2:subsec:250302}

\Needspace{6\baselineskip}
\begin{theorem}[分离抛物线片的双线性估计]\label{b2:thm:25-parabola-bilinear}
设 \(I,J\subset\mathbb R\) 为长度有限的紧区间，
\(\operatorname{dist}(I,J)\geq\rho>0\)，
\(f\in L^2(I;\mathbb C)\)、\(g\in L^2(J;\mathbb C)\)。
对 \((x,y)\in\mathbb R^2\)，令
\[
 \begin{aligned}
 E_If(x,y)&=\int_Ie^{2\pi\mathrm i(xt+yt^2)}f(t)\,\mathrm dt,\\
 E_Jg(x,y)&=\int_Je^{2\pi\mathrm i(xs+ys^2)}g(s)\,\mathrm ds.
 \end{aligned}
\]
则乘积属于 \(L^2(\mathbb R^2)\)，并且
\[
 \begin{split}
 \|E_If\,E_Jg\|_{L^2(\mathbb R^2)}^2
 &=\int_{I\times J}
 \frac{|f(t)|^2\cdot|g(s)|^2}{2|t-s|}
 \,\mathrm dt\,\mathrm ds\\
 &\leq(2\rho)^{-1}
 \cdot\|f\|_{L^2(I)}^2\cdot\|g\|_{L^2(J)}^2.
 \end{split}
\]
\end{theorem}
\begin{proof}
交换两个区间不改变结论，所以设 \(t<s\)。
先取光滑输入。上一小节的逆映射和换元公式说明，
\(\Psi_*(f(t)g(s)\,\mathrm dt\,\mathrm ds)\) 的密度为
\[
 h(u,v)=
 \begin{cases}
 \displaystyle\frac{f(t(u,v))g(s(u,v))}
 {2[s(u,v)-t(u,v)]},&(u,v)\in\Psi(I\times J),\\[6pt]
 0,&(u,v)\notin\Psi(I\times J).
 \end{cases}
\]
边界上的取值不影响积分。这个公式的含义可以用任意有界可测函数
\(\varphi\) 检验：
\begin{equation}\label{b2:eq:25-curve-pushforward}
 \int_{\mathbb R^2}\varphi(u,v)h(u,v)\,\mathrm du\,\mathrm dv
 =\int_{I\times J}\varphi(\Psi(t,s))f(t)g(s)
 \,\mathrm dt\,\mathrm ds.
\end{equation}
对绝对值同样换元，得到
\(\|h\|_1=\|f\|_1\cdot\|g\|_1\)，因此 \(h\) 的逆 Fourier 积分有意义。
在\cref{b2:eq:25-curve-pushforward}中取
\(\varphi(u,v)=e^{2\pi\mathrm i(xu+yv)}\)，便得到
\(\check h(x,y)=E_If(x,y)E_Jg(x,y)\)。
再作一次换元，
\[
 \|h\|_2^2
 =\int_{I\times J}
 \frac{|f(t)|^2\cdot|g(s)|^2}{4(s-t)^2}
 \,2(s-t)\,\mathrm dt\,\mathrm ds
 =\int_{I\times J}
 \frac{|f(t)|^2\cdot|g(s)|^2}{2(s-t)}
 \,\mathrm dt\,\mathrm ds.
\]
分母的下界证明 \(h\in L^2\)，Plancherel 给出所需等式和上界。

最后说明一般输入的逼近为何得到同一个乘积。
取 \(f_n\to f\) 于 \(L^2(I)\)、\(g_n\to g\) 于 \(L^2(J)\)，其中
\(f_n,g_n\) 光滑。已证估计给出
\[
 \begin{split}
 \|E_If_nE_Jg_n-E_If_mE_Jg_m\|_2
 \leq(2\rho)^{-1/2}\bigl(
 &\|f_n-f_m\|_2\cdot\|g_n\|_2\\
 &+\|f_m\|_2\cdot\|g_n-g_m\|_2\bigr).
 \end{split}
\]
所以乘积在 \(L^2\) 中收敛。同时有限区间上的 \(L^2\) 收敛蕴含
\(L^1\) 收敛，而
\(\sup_{x,y}|E_I(f_n-f)(x,y)|\leq\|f_n-f\|_1\)；
两个延拓分别一致收敛，且一致有界，故乘积也一致收敛到
\(E_IfE_Jg\)。从 \(L^2\) 收敛列取几乎处处收敛子列，便识别了这两个极限。
等式右侧也收敛：有
\[
 \bigl\||f_n|^2-|f|^2\bigr\|_1
 \leq\|f_n-f\|_2\cdot(\|f_n\|_2+\|f\|_2),
\]
对 \(g_n\) 同理，而权 \(1/(2|t-s|)\) 有界。
于是等式和不等式均传递到一般 \(L^2\) 输入。
\end{proof}

这个公式可以算到具体数值。取 \(a>0\)、\(\rho>0\)，令
\[
 I=[0,a],\quad J=[a+\rho,2a+\rho],\quad
 f=a^{-1/2}\mathbf1_I,\quad g=a^{-1/2}\mathbf1_J.
\]
两个输入的平方范数均为一。先对 \(t\) 积分，
\[
 \int_0^a\frac{\mathrm dt}{s-t}
 =\log s-\log(s-a).
\]
再对 \(s\in[a+\rho,2a+\rho]\) 积分，得到
\[
 \|E_If\,E_Jg\|_2^2
 =\frac{(2a+\rho)\log(2a+\rho)
 -2(a+\rho)\log(a+\rho)+\rho\log\rho}{2a^2}.
\]
其中各项的线性部分恰好抵消。令 \(\lambda=\rho/a\)，也可写成
\[
 \frac{(2+\lambda)\log(2+\lambda)
 -2(1+\lambda)\log(1+\lambda)+\lambda\log\lambda}{2a}.
\]
固定 \(a\) 而令 \(\rho\downarrow0\)，这个数趋于 \((\log2)/a\)。
因此方向分离的统一上界随 \(\rho\) 恶化，并不表示每一对固定形状的输入都以相同速度发散；
统一估计必须兼顾可随分离尺度改变、向最近端点集中的输入。

\ProseParagraphBreak
定理控制的是两个分离片的乘积。它没有断言同一片的延拓属于全空间
\(L^4\)：同片乘积会包含 \(t=s\) 附近的参数对，那里
\(2|t-s|\) 不能用正数统一控制。\chapref{b2:ch:23}预告的横截增益在这里已经成为精确公式，
但把分离乘积重新组合成一个线性估计，还需要处理对角附近各个尺度。

\subsection{切向横截与频率和映射}
\label{b2:subsec:250303}

对两条参数曲线，抛物线的特殊公式可以换成
\[
 \Psi(t,s)=\gamma_1(t)+\gamma_2(s),\quad
 \mathcal J(t,s)=
 |\det(\gamma_1'(t),\gamma_2'(s))|.
\]
分母现在衡量两条切向的横截程度。曲线不一定有曲率，也不必属于同一条曲线。
若参数域由多个片组成，同一个输出频率还可能有多个逆像；这时应先把各支的复密度相加，再取平方。

\Needspace{6\baselineskip}
\begin{proposition}[有限重数的曲线双线性估计]\label{b2:prop:25-finite-multiplicity}
设 \(K_1,K_2\subset\mathbb R\) 各为有限个紧区间之并，
\(\gamma_j\) 是 \(K_j\) 的开邻域上取值于 \(\mathbb R^2\) 的 \(C^1\) 映射。
令 \(\Psi(t,s)=\gamma_1(t)+\gamma_2(s)\)，并假设
\[
 \mathcal J(t,s)=|\det(\gamma_1'(t),\gamma_2'(s))|
 \geq\nu>0\quad((t,s)\in K_1\times K_2).
\]
设整数 \(M\geq1\)，且几乎每个 \(z\in\mathbb R^2\) 在
\(K_1\times K_2\) 中至多有 \(M\) 个 \(\Psi\)-逆像。
对复函数 \(f\in L^2(K_1)\)、\(g\in L^2(K_2)\)，定义
\[
 E_1f(x)=\int_{K_1}e^{2\pi\mathrm ix\cdot\gamma_1(t)}f(t)\,\mathrm dt,
 \quad
 E_2g(x)=\int_{K_2}e^{2\pi\mathrm ix\cdot\gamma_2(s)}g(s)\,\mathrm ds.
\]
则乘积的 Fourier 密度为
\[
 h(z)=\sum_{\substack{(t,s)\in K_1\times K_2\\\Psi(t,s)=z}}
 \frac{f(t)g(s)}{\mathcal J(t,s)}
 \quad\text{几乎处处},
\]
且
\[
 \begin{split}
 \|E_1f\,E_2g\|_{L^2(\mathbb R^2)}^2
 &=\|h\|_2^2\\
 &\leq M\int_{K_1\times K_2}
 \frac{|f(t)|^2\cdot|g(s)|^2}{\mathcal J(t,s)}
 \,\mathrm dt\,\mathrm ds\\
 &\leq\frac M\nu
 \cdot\|f\|_{L^2(K_1)}^2\cdot\|g\|_{L^2(K_2)}^2.
 \end{split}
\]
若和映射在输入域上一对一，则第一项不等式成为不带 \(M\) 的等式。
\end{proposition}
\begin{proof}
逆函数定理给输入域中每一点一个开邻域，使 \(\Psi\) 在该邻域上为
\(C^1\) 微分同胚。输入域紧，所以只需有限个这样的邻域
\(V_1,\ldots,V_L\)。令
\[
 A_1=(K_1\times K_2)\cap V_1,\quad
 A_\ell=((K_1\times K_2)\cap V_\ell)
 \setminus\bigcup_{k<\ell}V_k.
\]
这些不交可测块覆盖输入域，每一块都落在一个逆函数坐标内。
在 \(A_\ell\) 上作普通换元，得到一项由该支逆像定义的密度。
把有限项相加，就得到陈述中的 \(h\)；同一个参数对只属于一块，因此不会重复计算。
参数函数取不同的几乎处处代表也不改变输出密度的几乎处处等价类，
因为这些 \(C^1\) 坐标变换及其逆在局部紧块上把零测集送到零测集。

绝对值积分满足
\[
 \int_{\mathbb R^2}|h(z)|\,\mathrm dz
 \leq\int_{K_1\times K_2}|f(t)|\cdot|g(s)|\,\mathrm dt\,\mathrm ds
 =\|f\|_1\cdot\|g\|_1<\infty.
\]
因此逐块换元与 Fubini 对一般 \(L^2\) 输入已经合法，并直接给出
\(\check h=E_1fE_2g\)。
在几乎每个输出点，对至多 \(M\) 项用 Cauchy--Schwarz，
\[
 |h(z)|^2
 \leq M\sum_{\Psi(t,s)=z}
 \frac{|f(t)|^2\cdot|g(s)|^2}{\mathcal J(t,s)^2}.
\]
随后对每支换回参数变量，每项的体积因子为 \(\mathcal J(t,s)\)，故
\[
 \|h\|_2^2
 \leq M\int_{K_1\times K_2}
 \frac{|f(t)|^2\cdot|g(s)|^2}{\mathcal J(t,s)}
 \,\mathrm dt\,\mathrm ds.
\]
横截下界说明右侧有限；Plancherel 和 \(\mathcal J\geq\nu\) 完成证明。
一对一时逆像和只有一项，不需要 Cauchy--Schwarz，因而得到精确等式。
\end{proof}

重数因子可能实际出现。取
\[
 K_1=\bigcup_{k=0}^{M-1}[2k,2k+1],\quad K_2=[0,1],
\]
在第 \(k\) 个参数片上令
\(\gamma_1(t)=(t-2k,0)\)，并令 \(\gamma_2(s)=(0,s)\)。
这些公式各自在参数片的开邻域内光滑，且 \(\mathcal J=1\)。
单位正方形内部的每个点恰有 \(M\) 个逆像。
取 \(f=M^{-1/2}\mathbf1_{K_1}\)、\(g=\mathbf1_{K_2}\)，则
\[
 \|f\|_2=\|g\|_2=1,\quad
 h=M^{1/2}\mathbf1_{[0,1]^2},\quad
 \|E_1f\,E_2g\|_2^2=M.
\]
各支的相位在同一个输出频率处完全相同，所以密度先叠加，再产生平方。
若将平方范数误写成各支平方的直接和，就会漏掉这个 \(M\)。

\ProseParagraphBreak
另一个例子把横截常数与曲率分开。令
\[
 \gamma_1(t)=(t,0),\quad
 \gamma_2(s)=(s\cos\vartheta,s\sin\vartheta),
 \quad 0<\vartheta<\pi.
\]
在任意两个有限参数区间上，和映射是可逆线性映射，
\(\mathcal J=|\sin\vartheta|\)。因此
\[
 \|E_1f\,E_2g\|_2
 =|\sin\vartheta|^{-1/2}
 \cdot\|f\|_2\cdot\|g\|_2.
\]
两条曲线的曲率都为零，但夹角仍给出精确的乘积界；
夹角趋零时常数真实发散。抛物线中的曲率则保证相隔的参数点能够产生不同切向。

\ProseParagraphBreak
若输入按曲线弧长给定，测度因子也必须进入密度。
具体地，设 \(K_j\) 为有限个紧区间之并，
\(\gamma_j\) 在其开邻域上为 \(C^1\) 单射浸入，
两片单位切向的行列式绝对值有一致正下界。
令 \(\mathrm d\sigma_j=|\gamma_j'|\,\mathrm dt\)，取
\(a_j\in L^2(\gamma_j(K_j),\sigma_j;\mathbb C)\)，并定义
\[
 E_{\sigma_j}a_j(x)=\int_{\gamma_j(K_j)}
 e^{2\pi\mathrm ix\cdot\xi}a_j(\xi)\,\mathrm d\sigma_j(\xi).
\]
对应的参数输入为
\[
 f(t)=a_1(\gamma_1(t))|\gamma_1'(t)|,\quad
 g(s)=a_2(\gamma_2(s))|\gamma_2'(s)|.
\]
当和映射一对一时，令
\(\tau_j=\gamma_j'/|\gamma_j'|\) 为单位切向，精确公式变成
\[
 \|E_{\sigma_1}a_1\,E_{\sigma_2}a_2\|_2^2
 =\int_{\gamma_1(K_1)\times\gamma_2(K_2)}
 \frac{|a_1(\xi)|^2\cdot|a_2(\eta)|^2}
 {|\det(\tau_1(\xi),\tau_2(\eta))|}
 \,\mathrm d\sigma_1(\xi)\,\mathrm d\sigma_2(\eta).
\]
这是因为
\(\mathcal J=|\gamma_1'|\cdot|\gamma_2'|\cdot|\det(\tau_1,\tau_2)|\)，
约去一次速度因子后，剩下两次弧长微分。
参数式与几何式是同一个计算，分母中使用的横截量却有不同归一化。

二维计算由有限个逆像组成。若在 \(\mathbb R^d\)、\(d>2\) 中改用两个超曲面，
两个参数的总维数是 \(2(d-1)>d\)；即使和映射具有满秩微分，其典型纤维仍有维数 \(d-2\)。
此时必须沿纤维积分，不能照抄上面的有限逆像求和。
下一节用 \(d\) 个独立法向说明高维乘积的另一种横截结构。

% !TeX root = ../../Book2.tex
\section{多线性 Fourier 限制}
\label{b2:sec:2504}

平直横截估计把 \(d\) 个法向方向的互补关系变成乘积的可积性。
若曲面弯曲，同一个片内的法向会随参数变化，不能再用一个线性坐标变换把所有延拓同时化为删除坐标的函数。
一般多线性限制定理处理的正是这种变化；它保留局部空间尺度，并允许任意小的幂次损失。
理解这个结论时，应把定理本身、测度换算和反馈到线性问题的组合步骤分开。

\subsection{多线性延拓估计}
\label{b2:subsec:250401}

下面采用\term[bennett-carbery-tao-duoxianxing-xianzhi]{Bennett--Carbery--Tao 多线性限制定理}[Bennett--Carbery--Tao multilinear restriction theorem]的局部版本。
参数化需要定义在开集上，才能讨论导数；实际积分则限制在其中的紧片上。
这个区分也使边界附近的光滑控制有明确含义。

\Needspace{6\baselineskip}
\begin{theorem}[深层结构定理：Bennett--Carbery--Tao 局部多线性限制]\label{b2:thm:25-bct}
设整数 \(d\geq2\)。对 \(1\leq j\leq d\)，设
\(U_j\subset\mathbb R^{d-1}\) 为开集，
\(K_j\subset U_j\) 为非空紧集，
\(\Sigma_j:U_j\to\mathbb R^d\) 为固定的光滑浸入。
以
\[
 \langle Y_j(u),v\rangle
 =\det\bigl(\partial_1\Sigma_j(u),\ldots,
 \partial_{d-1}\Sigma_j(u),v\bigr)
 \quad(v\in\mathbb R^d)
\]
定义切向楔积所对应的法向量 \(Y_j(u)\)。
假设存在 \(\nu>0\)，使
\[
 |\det(Y_1(u_1),\ldots,Y_d(u_d))|\geq\nu
 \quad\text{对所有 }(u_1,\ldots,u_d)\in K_1\times\cdots\times K_d.
\]
对 \(g_j\in L^2(K_j;\mathbb C)\)，定义参数测度延拓
\[
 E_jg_j(x)=\int_{K_j}
 e^{2\pi\mathrm ix\cdot\Sigma_j(u)}g_j(u)\,\mathrm du.
\]
令 \(q=2/(d-1)\)。则对每个 \(\epsilon>0\)，存在常数 \(C_\epsilon\)，
使任意 \(R\geq1\)、\(a\in\mathbb R^d\) 及上述输入均满足
\[
 \left\|\prod_{j=1}^dE_jg_j\right\|_{L^q(B(a,R))}
 \leq C_\epsilon R^\epsilon\prod_{j=1}^d\|g_j\|_{L^2(K_j)}.
\]
这里对任意 \(q>0\) 均按
\(\|F\|_{L^q(B)}=(\int_B|F|^q)^{1/q}\) 解释左侧。
常数可依赖 \(d,\epsilon,\nu\)、固定参数片及其开邻域上的光滑参数化，
不依赖 \(a,R\) 或输入。
\end{theorem}
\begin{proofsketch}
把参数片分成与空间半径相适应的小帽，每个小帽的延拓再写成沿相应法向传播的波包。
横截下界保证来自 \(d\) 组的管束方向始终独立，但每组内部仍有许多位置和方向。
多线性 Kakeya 估计控制这些横截管束的重叠，尺度归纳再将这种几何控制反馈到延拓乘积。
原证明以 Gaussian 族的单调性处理相应的多线性管束估计；
局部化误差与各层尺度的常数必须在归纳中统一控制，最后保留允许的 \(R^\epsilon\) 损失。

完整管束和尺度归纳证明在此略去。上述版本是 Jonathan Bennett、Anthony Carbery 与陶哲轩的《On the Multilinear Restriction and Kakeya Conjectures》%
\cite[Theorem 1.16, p.~267; Sections 2--4]{Bennett2006Multilinear}
的 \(L^2\) 端点特例。对固定紧参数片，横截条件由连续性延伸到稍大的邻域，
可在其中加入固定光滑截断；必要时再作有限次参数片细分。
这些固定选择进入常数，不随空间半径改变。
\end{proofsketch}

任意球心不需要另一个定理。写 \(x=a+y\)，并令
\[
 g_{j,a}(u)=e^{2\pi\mathrm ia\cdot\Sigma_j(u)}g_j(u).
\]
则
\[
 E_jg_j(a+y)=E_jg_{j,a}(y),\quad
 \|g_{j,a}\|_2=\|g_j\|_2.
\]
所以以原点为中心的结论通过输入调制直接传到任意球心，常数保持不变。
这里的输入允许依赖球心，定理本来就对全部 \(L^2\) 输入一致成立。

\ProseParagraphBreak
令 \(p=dq=2d/(d-1)\)。定理也常用几何平均形式书写：
\begin{equation}\label{b2:eq:25-bct-geometric-mean}
 \left\|\left(\prod_{j=1}^d|E_jg_j|\right)^{1/d}
 \right\|_{L^p(B(a,R))}
 \leq C_\epsilon^{1/d}R^{\epsilon/d}
 \prod_{j=1}^d\|g_j\|_2^{1/d}.
\end{equation}
这是原不等式取 \(d\) 次根，因为左侧的 \(d\) 次幂恰为
\(\|\prod_jE_jg_j\|_{L^q}\)。
若希望右侧仍记成 \(R^\epsilon\)，先在原定理中使用 \(d\epsilon\)，再重命名常数即可。
乘积指数是 \(q\)，几何平均指数是 \(p\)，二者不能混用。

局部球和小幂次损失都是结论的一部分。固定 \(\epsilon>0\) 后，
\(R^\epsilon\) 随 \(R\to\infty\) 增大，故不能直接用单调收敛取得无损全空间界；
先让 \(\epsilon\downarrow0\) 也不合法，因为定理没有给出
\(C_\epsilon\) 在零附近的一致上界。
平直模型的全空间无损结论来自其特殊坐标结构，不能由这里的局部曲面定理直接替代。

\subsection{多线性限制}
\label{b2:subsec:250402}

先说明参数测度与面测度。若 \(\Sigma_j\) 在当前参数片上单射，
写 \(S_j=\Sigma_j(K_j)\)、\(J_j(u)=|Y_j(u)|\)，则面积公式为
\[
 \mathrm d\sigma_j(\Sigma_j(u))=J_j(u)\,\mathrm du.
\]
若面测度输入为 \(h_j\)，为了使延拓积分相同，应取
\[
 g_j(u)=h_j(\Sigma_j(u))J_j(u).
\]
相应的平方范数不是直接相等，而是
\[
 \|g_j\|_{L^2(K_j)}^2
 =\int_{S_j}|h_j(\xi)|^2
 \cdot J_j(\Sigma_j^{-1}(\xi))\,\mathrm d\sigma_j(\xi).
\]
若 \(0<c_j\leq J_j\leq C_j\)，这给出
\[
 c_j^{1/2}\cdot\|h_j\|_{L^2(S_j)}
 \leq\|g_j\|_{L^2(K_j)}
 \leq C_j^{1/2}\cdot\|h_j\|_{L^2(S_j)}.
\]
固定紧片上的浸入确有这样的界；若参数化发生变化，不能将这些常数无条件省去。
参数化不单射时，应先分成单射片，或在面积公式中保留覆盖重数。

\ProseParagraphBreak
单位法向 \(n_j=Y_j/J_j\) 与参数法向满足
\[
 |\det(Y_1,\ldots,Y_d)|
 =\left(\prod_{j=1}^dJ_j\right)
 |\det(n_1,\ldots,n_d)|.
\]
所以使用单位法向横截下界时，必须同时记录各 \(J_j\) 的正下界。
例如平片重参数化 \(\Sigma(u)=(bu,0)\)、\(b>0\)，使
\(J=b^{d-1}\)；虽然几何曲面没有改变，参数输入的平方范数却等于相应面测度输入平方范数的 \(b^{d-1}\) 倍。
几何横截与参数尺度分别出现，正是上述换算要保留的内容。

\ProseParagraphBreak
“多线性限制”通常以延拓乘积表述。固定其余输入以后，当然可以把其中一个输入看作线性变量；
但由此产生的伴随含有其他延拓函数，不能等同于一个不带权的曲面限制算子。
例如在有限球 \(B\) 内，固定
\(H=\prod_{j=2}^dE_jg_j\)，并取 \(\varphi\in L^1(B)\)。
由于各延拓有界，Fubini 给出
\[
 \int_B E_1g_1(x)H(x)\overline{\varphi(x)}\,\mathrm dx
 =\int_{K_1}g_1(u)
 \overline{\widehat{\overline H\,\varphi\mathbf1_B}
 (\Sigma_1(u))}\,\mathrm du.
\]
右侧限制的是含有 \(\overline H\) 的乘积 Fourier 变换。
这个配对恒等式解释了名称之间的联系，却不自动给出每个单独延拓的同等可积性。

输出指数还决定怎样求和。\(d=2\) 时 \(q=2\)，
上一节的曲线模型可以在受控重数下给出全空间平方估计；
\(d=3\) 时 \(q=1\)；
\(d\geq4\) 时 \(0<q<1\)，左侧仅为准范数。
对任意测度空间上的有限族可测复函数 \(F_\ell\)，此时成立的是
\[
 \left\|\sum_{\ell=1}^L F_\ell\right\|_q^q
 \leq\sum_{\ell=1}^L\|F_\ell\|_q^q,
\]
因为点态有
\(|\sum_\ell F_\ell|^q\leq(\sum_\ell|F_\ell|)^q
\leq\sum_\ell|F_\ell|^q\)。
若每项的准范数至多为 \(A\)，这一估计只给 \(L^{1/q}A\)。
例如 \(d=4\)、\(q=2/3\)，两个不交单位立方体的示性函数各有准范数一，
其和的准范数却为 \(2^{3/2}>2\)。
因此不能把普通三角不等式写入这个指数。
几何平均形式的 \(p=2d/(d-1)>1\) 是通常的范数，
但采用它时必须同时将乘积和输入范数都取 \(d\) 次根。

\subsection{对线性限制的反馈}
\label{b2:subsec:250403}

若把一个曲面输入分成有限个帽，线性延拓是各帽延拓之和。
在某一点，大和可能由许多横截方向共同贡献，也可能主要来自接近一个低维空间的方向。
下面的有限维选择说明这个区分怎样实际进行；它只用数值大小和方向，不预先假定一个无限尺度归约。

\Needspace{6\baselineskip}
\begin{lemma}[有限帽的方向分配]\label{b2:lem:25-finite-cap-allocation}
设整数 \(d\geq2\)、\(N\geq1\)，
\(n_1,\ldots,n_N\in\mathbb R^d\) 为单位向量，
\(z_1,\ldots,z_N\in\mathbb C\)，并设 \(0<\eta<1\)。
记 \(S=\sum_{\theta=1}^Nz_\theta\)。至少有下列一种情况成立：
\begin{enumerate}
\item\label{b2:item:25-cap-transverse}存在 \(d\) 个不同指标 \(\theta_1,\ldots,\theta_d\)，满足
\[
 |\det(n_{\theta_1},\ldots,n_{\theta_d})|
 \geq\eta^{d-1},\quad
 |S|\leq2N
 \left(\prod_{j=1}^d|z_{\theta_j}|\right)^{1/d}.
\]
\item\label{b2:item:25-cap-cluster}存在由至多 \(d-1\) 个所给向量张成的非零子空间 \(V\)，
\(\dim V\leq d-1\)，使
\[
 |S|\leq
 2\sum_{\{\theta:\operatorname{dist}(n_\theta,V)\leq\eta\}}
 |z_\theta|.
\]
\end{enumerate}
第二种情况中的候选子空间可以预先从有限族中选取，该族的大小至多为
\(\sum_{r=1}^{d-1}\binom Nr\)，其中 \(r>N\) 的项取零。
\end{lemma}
\begin{proof}
若 \(A=\sum_\theta|z_\theta|=0\)，第二种情况对
\(V=\operatorname{span}(n_1)\) 成立。以下设 \(A>0\)，并令
\[
 \mathcal H=\{\theta:|z_\theta|\geq A/(2N)\}.
\]
其补集的总质量小于 \(A/2\)，所以
\(\sum_{\theta\in\mathcal H}|z_\theta|\geq A/2\)，特别地
\(\mathcal H\ne\varnothing\)。
从 \(\mathcal H\) 中任选一个向量作为第一列。
若尚有大项方向到已选列张成的空间的距离大于 \(\eta\)，就将它加入；
否则停止。

若选出 \(d\) 列，Gram--Schmidt 体积公式说明其行列式绝对值等于第一列的长度一，
乘以后续各列到前列张成空间的距离，故大于 \(\eta^{d-1}\)。
每个被选输入的绝对值至少为 \(A/(2N)\)，从而
\[
 |S|\leq A
 \leq2N\left(\prod_{j=1}^d|z_{\theta_j}|\right)^{1/d}.
\]
若不足 \(d\) 列就停止，所选列张成一个维数不超过 \(d-1\) 的空间 \(V\)，
全部 \(\mathcal H\) 方向都在其 \(\eta\) 邻域内。因此
\[
 |S|\leq A\leq
 2\sum_{\theta\in\mathcal H}|z_\theta|
 \leq2\sum_{\operatorname{dist}(n_\theta,V)\leq\eta}|z_\theta|.
\]
最后，所选列来自至多 \(d-1\) 个指标；枚举这些有限子集便得到候选族的数量上界。
\end{proof}

把 \(z_\theta\) 取为第 \(\theta\) 个帽的延拓 \(u_\theta(x)\)，
引理就在每个空间点给出一个分配。
为将第一种情况接到多线性定理，还要从帽中心的法向下界过渡到整帽下界。
对单位向量列，行列式的多线性展开给出
\[
 |\det(v_1,\ldots,v_d)-\det(w_1,\ldots,w_d)|
 \leq\sum_{j=1}^d|v_j-w_j|.
\]
因而若每帽内的单位法向偏离参考方向不超过
\(\eta^{d-1}/(2d)\)，参考行列式至少为 \(\eta^{d-1}\) 的帽组在所有参数点上仍有
\(\eta^{d-1}/2\) 的下界。再用面积因子的正下界，就满足
\cref{b2:thm:25-bct}的参数法向条件。

\ProseParagraphBreak
有限帽的横截部分可以直接求和。设上述 \(N\) 个帽及其参数化固定，
\(\mathcal T\) 为满足参考行列式下界的有序 \(d\) 元帽组，
并记
\[
 G(x)=\max_{(\theta_1,\ldots,\theta_d)\in\mathcal T}
 \left(\prod_{j=1}^d|u_{\theta_j}(x)|\right)^{1/d}.
\]
若 \(\mathcal T\) 为空，令 \(G=0\)。
写 \(u_\theta=E_\theta g_\theta\)、\(a_\theta=\|g_\theta\|_2\)，并设
\(p=2d/(d-1)\)、\(q=2/(d-1)\)。
对有限帽组取各个 BCT 常数的最大值，得到
\[
 \begin{split}
 \|G\|_{L^p(B(a,R))}^p
 &\leq\sum_{\mathcal T}\int_{B(a,R)}
 \prod_{j=1}^d|u_{\theta_j}|^q\\
 &\leq C_\epsilon^qR^{\epsilon q}
 \sum_{\mathcal T}\prod_{j=1}^da_{\theta_j}^q\\
 &\leq C_\epsilon^qR^{\epsilon q}
 \left(\sum_{\theta=1}^Na_\theta^q\right)^d.
 \end{split}
\]
先对非负 \(p\) 次幂积分求和，就能将各帽组的估计合并。
由于 \(q\leq2\)，有限和的 Hölder 不等式又给出
\[
 \sum_\theta a_\theta^q
 \leq N^{1-q/2}
 \left(\sum_\theta a_\theta^2\right)^{q/2}.
\]
合并并取 \(p\) 次根，得到可核算的横截部分估计
\begin{equation}\label{b2:eq:25-cap-broad-bound}
 \|G\|_{L^p(B(a,R))}
 \leq C_\epsilon^{1/d}R^{\epsilon/d}
 N^{(d-2)/2}
 \left(\sum_{\theta=1}^N\|g_\theta\|_2^2\right)^{1/2}.
\end{equation}
若各 \(g_\theta\) 是同一参数输入在不交帽上的限制，最后的平方和就是原输入的平方范数。
引理在横截点上控制的线性和至多为 \(2NG\)，所以还会多出一个 \(2N\)。
在固定帽数下，这些都是有限常数；若帽数随空间尺度增长，就必须把其幂次计入递推，不能视作无代价求和。

\ProseParagraphBreak
剩余的方向簇未必只是一个小帽。以三维抛物面
\[
 \Sigma(t,s)=(t,s,t^2+s^2),\quad
 n(t,s)=\frac{(-2t,-2s,1)}{\sqrt{1+4t^2+4s^2}}
\]
为例，在三个中心 \((-1,0),(0,0),(1,0)\) 处，法向分别为
\[
 \frac{(2,0,1)}{\sqrt5},\quad (0,0,1),\quad
 \frac{(-2,0,1)}{\sqrt5}.
\]
它们两两夹角的正弦分别为 \(2/\sqrt5,2/\sqrt5,4/5\)，均远离零，
但三个方向都在 \(\operatorname{span}(e_1,e_3)\) 内，三列行列式为零。
取中心附近三个互不相交的等面积参数帽，在各帽上放相同的正常数输入，
则在 \(x=0\) 处三个延拓值相同且为正，线性和同时受到三帽贡献。
这些帽在参数平面上相距固定距离，不能合并成一个直径趋零的小帽。
若各帽半径为 \(\delta\)，其法向到上述平面的距离至多为 \(C\delta\)，
因此三帽行列式也至多为 \(C'\delta\)；细分帽不会凭空给出一个与 \(\delta\) 无关的完整横截下界。

这个例子解释了低维方向簇为何需要单独处理。
有限帽引理允许逐点选择子空间，但每点选出的帽组可能不同；
后续必须构造可求和的空间分配，控制各方向簇的范数，并在缩放后记录参数域、空间区域和常数。
若集中在单个小帽，可以尝试使用前文的抛物缩放；若像上例那样沿一条频率方向分布，
则还需要较低维或切向的估计。
从有限分配与横截部分的界走向线性估计，还需控制各方向簇在不同尺度上的组合。下一节将用尺度归纳、波包和分割研究这些问题。

% !TeX root = ../../Book2.tex
\section{多线性限制的几何与尺度方法}
\label{b2:sec:2505}

前面的横截模型解释了不同方向怎样改善乘积估计。把这种增益用于许多帽和许多空间区域，还需要尺度递推、波包分解以及对几何组合的控制。尺度递推追踪估计在放大与缩小之间的变化，单帽波包显示频率如何决定空间中的管，多帽解耦则量化不同方向合成时的损失。

\subsection{尺度归纳}
\label{b2:subsec:250501}

设 \(A(R)\) 表示半径 \(R\) 球上的某个规范化最佳常数；若定理允许任意球心或一族曲面，便应先对这些参数取上确界。尺度归纳希望把 \(A(R)\) 与较小尺度的常数联系起来。一种常见模型是
\[
 A(R)\leq C_\eta R^\eta A(R^{1/2}).
\]
右侧不只有幂次损失：每迭代一次，还会多出一个 \(C_\eta\)。下面把两者都算出。

\Needspace{6\baselineskip}
\begin{lemma}[平方根尺度递推]\label{b2:lem:25-scale-iteration}
设 \(A:[1,\infty)\to[0,\infty)\) 为有限值函数，满足
\[
 M\coloneqq\sup_{1\leq R\leq2}A(R)<\infty.
\]
假设对每个 \(\eta>0\)，存在与 \(R\) 无关的有限常数 \(C_\eta\geq1\)，使
\[
 A(R)\leq C_\eta R^\eta A(R^{1/2})\quad(R>2).
\]
则对每个 \(\varepsilon>0\)，存在有限常数 \(K_\varepsilon\)，使
\[
 A(R)\leq K_\varepsilon R^\varepsilon\quad(R\geq1).
\]
常数 \(K_\varepsilon\) 可以依赖 \(M\) 及相应的 \(C_\eta\)，不要求它在 \(\varepsilon\downarrow0\) 时保持有界。
\end{lemma}
\begin{proof}
固定 \(\eta>0\) 及 \(R>2\)，令 \(m\) 为满足 \(R^{2^{-m}}\leq2\) 的最小正整数。前 \(m\) 步均可应用递推，故
\[
 \begin{aligned}
 A(R)
 &\leq C_\eta^m
       R^{\eta(1+2^{-1}+\cdots+2^{-(m-1)})}
       A(R^{2^{-m}})\\
 &\leq M C_\eta^m R^{2\eta}.
 \end{aligned}
\]
另一方面，
\[
 m\leq1+\log_2\!\left(\frac{\log R}{\log2}\right).
\]
记 \(b_\eta=\log_2 C_\eta\geq0\)，则
\[
 C_\eta^m
 \leq C_\eta\left(\frac{\log R}{\log2}\right)^{b_\eta}.
\]
所以固定 \(\eta\) 后，全部常数累积只是一个固定的对数幂。

现在给定 \(\varepsilon>0\)，先取 \(\eta=\varepsilon/4\)。这样 \(C_\eta\) 和 \(b_\eta\) 已经固定。函数
\((\log R)^{b_\eta}R^{-\varepsilon/2}\)
在 \(R\geq2\) 上有界：令 \(t=\log R\)，它就是 \(t^{b_\eta}e^{-\varepsilon t/2}\)，指数衰减控制固定幂次。因此
\[
 A(R)\leq K_\varepsilon
          R^{2\eta+\varepsilon/2}
       =K_\varepsilon R^\varepsilon.
\]
初始区间 \(1\leq R\leq2\) 由 \(M\) 控制，增大常数即可。
\end{proof}

顺序在这里很重要：先选 \(\eta\)，再吸收由 \(C_\eta\) 决定的对数幂。若让 \(\eta\) 随 \(R\) 改变，却不知道 \(C_\eta\) 怎样增长，上述吸收便没有统一常数。这个引理也不负责证明具体限制问题所需的递推式；它只说明一旦递推和初始尺度界成立，怎样正确结束迭代。

\ProseParagraphBreak
每一步还必须留在同一允许的曲面类中，包括参数化的导数界、支集大小和横截下界。\cref{b2:prop:25-parabolic-rescaling}给出抛物面的精确重标度；一般曲面的小片重标度则要控制偏离二次模型的误差。如果归纳后曲率或横截常数不断恶化，就不能把下一步仍写成同一个 \(A(R^{1/2})\)，更不能只凭幂次几何级数断言得到任意小损失。

\subsection{波包}
\label{b2:subsec:250502}

在空间半径 \(R\) 上，一个宽度 \(R^{-1/2}\) 的频率帽对应宽约 \(R^{1/2}\)、长约 \(R\) 的管。这与\cref{b2:fig:23-cap-tube}及\cref{b2:prop:23-parabolic-packet}的尺度一致，但单个光滑帽函数还不代表一般 \(L^2\) 输入。下面用帽内 Fourier 级数把一般输入拆成位置不同的波包，并保留系数的平方能量。

\ProseParagraphBreak
为使归一化可直接检查，这一构造使用抛物面的参数测度；面积测度的光滑密度在固定参数片上可以吸收到输入中，但相应的范数比较应随之注明。

\Needspace{6\baselineskip}
\begin{proposition}[单帽的波包展开]\label{b2:prop:25-single-cap-packets}
设整数 \(n\geq1\)，\(R\geq1\)、\(\delta=R^{-1/2}\)，且 \(\xi_0\in\mathbb R^n\)、\(|\xi_0|\leq1\)。取
\[
 g\in L^2(\mathbb R^n),\quad
 \operatorname{supp}g\subset
       \xi_0+\delta[-1/2,1/2]^n,
\]
并以参数测度定义
\[
 E_{\mathrm{par}}g(x',t)
 =\int_{\mathbb R^n}
       e^{2\pi\mathrm i(x'\cdot\xi+t|\xi|^2)}
       \cdot g(\xi)\,\mathrm d\xi.
\]
固定 \(\rho\in C_c^\infty((-1,1)^n)\)，在 \([-1/2,1/2]^n\) 上恒为一。则存在系数 \(c_k\)、\(k\in\mathbb Z^n\)，及只依赖 \(\xi_0,\delta,\rho,k\) 的频率函数 \(\psi_k\)，使
\[
 g=\sum_{k\in\mathbb Z^n}c_k\psi_k
       \quad\text{于 }L^2(\mathbb R^n),\quad
 \sum_k|c_k|^2=\|g\|_2^2.
\]
若 \(W_k=E_{\mathrm{par}}\psi_k\)，则按有限立方格点部分和取极限时，
\[
 E_{\mathrm{par}}g=\sum_kc_kW_k
\]
在整个 \(\mathbb R^{n+1}\) 上一致收敛。对每个 \(N>0\)，在 \(|t|\leq R\) 内还有
\[
 |W_k(x',t)|
 \leq C_{N,n,\rho}\delta^{n/2}
  \left(1+\left|\delta(x'+2t\xi_0)+k/2\right|\right)^{-N}.
\]
可选一个只依赖维数的固定 \(A_0\)，使对任意 \(L>0\)，
\[
 \sup_{|(x',t)|\leq R}
 \left|E_{\mathrm{par}}g(x',t)
       -\sum_{|k|\leq A_0\sqrt R}c_kW_k(x',t)\right|
 \leq C_{L,n,\rho}R^{-L}\|g\|_2.
\]
所有常数均与 \(R,\xi_0,g\) 无关。
\end{proposition}
\begin{proof}
将帽拉回单位盒，令 \(h(u)=g(\xi_0+\delta u)\)。它在
\(Q=[-1,1]^n\) 上支于较小的盒 \([-1/2,1/2]^n\)。正交基
\[
 e_k(u)=2^{-n/2}e^{\pi\mathrm i k\cdot u},\quad k\in\mathbb Z^n,
\]
给出 Fourier 级数
\[
 h=\sum_ka_ke_k\quad\text{于 }L^2(Q),\quad
 \sum_k|a_k|^2=\|h\|_{L^2(Q)}^2
             =\delta^{-n}\|g\|_2^2.
\]
令
\[
 c_k=\delta^{n/2}a_k,\quad
 \psi_k(\xi)=2^{-n/2}\delta^{-n/2}
       \rho\!\left(\frac{\xi-\xi_0}{\delta}\right)
       \cdot e^{\pi\mathrm i k\cdot(\xi-\xi_0)/\delta}.
\]
由于 \(\rho\cdot h=h\)，把 Fourier 级数乘以有界函数 \(\rho\)，再作换元，就得到所述 \(L^2\) 重构与精确系数能量。每个频率函数支于放大一倍的帽，且
\[
 \|\psi_k\|_2=2^{-n/2}\|\rho\|_2.
\]
它们有固定的 \(L^2\) 大小，但并不因这个归一化就彼此正交；系数能量等式来自乘 \(\rho\) 之前的正交 Fourier 基。

换元 \(\xi=\xi_0+\delta u\) 给出原波包的具体公式
\[
 \begin{aligned}
 W_k(x',t)
 ={}&2^{-n/2}\delta^{n/2}
       e^{2\pi\mathrm i(x'\cdot\xi_0+t|\xi_0|^2)}\\
 &{}\cdot
 \int_{\mathbb R^n}\rho(u)\cdot
       e^{2\pi\mathrm i[
          \delta(x'+2t\xi_0)\cdot u
          +t\delta^2|u|^2+k\cdot u/2]}\,\mathrm du.
 \end{aligned}
\]
频率部分和的误差支于固定的放大帽，其 \(L^2\) 范数趋零，故其 \(L^1\) 范数也趋零。延拓积分核的模为一，因此输出误差的一致范数至多为这个 \(L^1\) 误差。频率侧的 \(L^2\) 逼近因而转化为延拓函数在全空间上的一致收敛。

若 \(|t|\leq R\)，则 \(|t\delta^2|\leq1\)。振幅
\(\rho(u)\cdot e^{2\pi\mathrm i t\delta^2|u|^2}\)
的支集固定，各固定阶导数一致有界。对剩余的线性频率
\[
 z+k/2,\quad z=\delta(x'+2t\xi_0),
\]
反复分部积分，即得所述任意阶衰减。

还需把这些尾部实际求和。对任意 \(z\in\mathbb R^n\)，格点间距固定，半径 \(2^j\) 的球内至多有 \(C_n2^{jn}\) 个点。因此当 \(2N>n\) 时，
\[
 \sup_z\sum_{k\in\mathbb Z^n}
       (1+|z+k/2|)^{-2N}<\infty.
\]
离散 Cauchy--Schwarz 与系数能量给出
\[
 \sum_k|c_k|\cdot|W_k(x',t)|
 \leq C_{N,n,\rho}\delta^{n/2}\|g\|_2
 \quad(|t|\leq R).
\]
所以该时间片内的波包和还逐点绝对收敛。

在球 \(|(x',t)|\leq R\) 内，因 \(|\xi_0|\leq1\)，有
\(|z|\leq3\sqrt R\)。取固定 \(A_0\geq12\)，若 \(|k|>A_0\sqrt R\)，则
\[
 |z+k/2|\geq|k|/4.
\]
同样按格点环带求和，得到
\[
 \begin{aligned}
 \sum_{|k|>A_0\sqrt R}|c_k|\cdot|W_k(x',t)|
 &\leq C_{N,n,\rho}\delta^{n/2}\|g\|_2
        R^{n/4-N/2}\\
 &=C_{N,n,\rho}R^{-N/2}\|g\|_2.
 \end{aligned}
\]
取 \(N>\max\{n/2,2L\}\)，便得到命题的局部误差界。
\end{proof}

在这个构造中，主要管区域为
\[
 |x'+2t\xi_0+k/(2\delta)|\lesssim\delta^{-1},\quad
 |t|\leq R.
\]
其宽度约为 \(R^{1/2}\)，长度约为 \(R\)，长轴沿抛物面的法向 \((-2\xi_0,1)\)。帽的中心决定方向，格点 \(k\) 决定位置；同一帽中的管彼此平行。空间球内只需保留 \(O(R^{n/2})\) 个位置指标，其余项可以按任意预先指定的负幂控制。

\ProseParagraphBreak
这里有三个不同的结论。频率重构在 \(L^2\) 中成立；紧频率支集又使延拓部分和在全空间一致收敛；管外快速衰减及上述有限格点截断则限于 \(|t|\leq R\) 或 \(|(x',t)|\leq R\)。最后一项不能无条件推广到任意大时间。波包并非管的示性函数：公式中的相位与系数仍须保留，截到粗管时产生的尾部也须按需要的阶数求和。

\ProseParagraphBreak
这只是单张帽的展开。若把许多帽相加，各帽之间有不同的法向，不能从一张帽的系数能量等式推出 BCT 或解耦。为了估计这些组合，接下来需要同时看管与空间分区的关系，以及不同帽输出的范数怎样相加。

\subsection{多项式分割}
\label{b2:subsec:250503}

多项式分割选择一个多项式，把大部分空间分成若干胞腔，并将零集的邻域单独保留。它希望让各胞腔承担可控的质量，同时利用管的穿越次数限制重复计算。选择分割多项式及控制各胞腔质量属于额外的存在性结论，不能由直线的根数直接推出。

\ProseParagraphBreak
\cref{b2:lem:23-polynomial-line-crossing}证明：次数至多 \(D\) 的非零多项式，其零集之外的胞腔，至多有 \(D+1\) 个能被同一条不全含于零集的直线经过。对于半径为 \(\rho\) 的粗管，要把胞腔缩到与零集距离大于 \(\rho\) 的部分，才能把管的相遇转回中心直线的穿越。原因是管中这样的点到中心线的连接线段避开零集，故其中心线点仍在同一胞腔。未经缩减的胞腔与零集邻域，不由这个次数界控制。

即使多项式只有一次，零集邻域中的管也会表现出两种不同情形。平面模型足以把角度和停留长度算清。

\Needspace{6\baselineskip}
\begin{proposition}[管与平面厚邻域]\label{b2:prop:25-plane-wall}
设整数 \(d\geq2\)，\(\nu,v\in S^{d-1}\)，\(a\in\mathbb R^d\)、\(b\in\mathbb R\)，且 \(L,r,\rho>0\)。令
\[
 H=\{\,x:x\cdot\nu=b\,\},\quad
 W_r(H)=\{\,x:|x\cdot\nu-b|\leq r\,\},
\]
\[
 \gamma(s)=a+sv\quad(0\leq s\leq L),\quad
 T_\rho=\bigcup_{0\leq s\leq L}\overline{B(\gamma(s),\rho)}.
\]
记可能与平面厚邻域相遇的轴参数集合为
\[
 I=\{\,s\in[0,L]:
       \overline{B(\gamma(s),\rho)}\cap W_r(H)\neq\varnothing\,\}.
\]
以一维 Lebesgue 测度记 \(|I|\)。若 \(v\cdot\nu\neq0\)，则
\[
 |I|\leq\min\left\{L,\frac{2(r+\rho)}{|v\cdot\nu|}\right\}.
\]
若 \(v\cdot\nu=0\)，则 \(I=[0,L]\) 或 \(I=\varnothing\)，前一种情形恰好对应 \(|a\cdot\nu-b|\leq r+\rho\)。特别地，若轴线在 \(H\) 内且 \(\rho\leq r\)，则整条管 \(T_\rho\) 都包含于 \(W_r(H)\)。
\end{proposition}
\begin{proof}
球 \(\overline{B(\gamma(s),\rho)}\) 与平面厚邻域相遇，当且仅当其球心到平面的距离至多为 \(r+\rho\)。于是
\[
 I=\{\,s\in[0,L]:
       |a\cdot\nu-b+s(v\cdot\nu)|\leq r+\rho\,\}.
\]
若斜率 \(v\cdot\nu\neq0\)，不等式在整个实轴上给出长度
\(2(r+\rho)/|v\cdot\nu|\) 的区间，与 \([0,L]\) 相交后得到所述上界。若斜率为零，条件与 \(s\) 无关，所以参数集合全有或全无。轴线位于平面内时，管中每点到该平面的距离不超过 \(\rho\)，最后一项随即成立。
\end{proof}

若 \(\vartheta\) 是管轴与平面的夹角，则 \(|v\cdot\nu|=|\sin\vartheta|\)。取波包尺度 \(L=R\)、\(r\) 与 \(\rho\) 都约为 \(R^{1/2}\)，固定正角度给出的停留长度只有 \(O(R^{1/2})\)；当角度小到 \(R^{-1/2}\) 的量级时，上界可以达到整条管的长度 \(R\)。平行而贴墙的管更可能一直留在墙内，不能把它作为普通横穿管处理。

\ProseParagraphBreak
这说明零集邻域中的分析不能只问“中心线与零集交了几次”。横穿部分可利用角度下界缩短停留区域；近切向部分则可能沿低维几何聚集，需要不同的估计。对高次零集，局部切平面会随位置改变，还会有奇点及多个局部片，平面模型本身不处理这些困难。

\ProseParagraphBreak
Guth 的《A Restriction Estimate Using Polynomial Partitioning》\cite{Guth2016Restriction}给出了这种方法的正式限制应用。多项式次数一方面增加胞腔数，另一方面也增加穿越预算和零集邻域的复杂度，必须放进尺度递推中选择。这篇文献进一步结合分割定理、切向管分析和尺度递推，把次数与角度的控制转化为线性限制估计。

\subsection{解耦}
\label{b2:subsec:250504}

单帽波包保留帽内的相干性；解耦则控制许多帽叠加后的输出。下面采用\term[jieou-dingli]{解耦定理}[decoupling theorem]的抛物面临界版本。帽的尺度、分割函数和空间权都属于陈述的一部分。

\Needspace{6\baselineskip}
\begin{theorem}[深层结构定理：抛物面临界解耦]\label{b2:thm:25-decoupling}
设整数 \(d\geq2\)、\(R\geq1\)，令
\[
 S=\{\,(\xi',|\xi'|^2):|\xi'|\leq1\,\}\subset\mathbb R^d,\quad
 p_c=\frac{2(d+1)}{d-1},\quad \delta=R^{-1/2}.
\]
曲面取面积测度 \(\sigma\)，并定义
\[
 Eg(x)=\int_Se^{2\pi\mathrm i x\cdot\xi}
       \cdot g(\xi)\,\mathrm d\sigma(\xi)
 \quad(g\in L^2(S,\sigma)).
\]
在参数域上取尺度 \(\delta\) 的规则方格帽覆盖，并固定非负光滑单位分解 \((\chi_\theta)_\theta\)：每个 \(\chi_\theta\) 支于相应方格的固定倍数放大，其支集重叠次数一致有界，
\[
 \sum_\theta\chi_\theta=1\quad\text{于 }S,\quad
 |\partial^\alpha\chi_\theta|\leq C_\alpha\delta^{-|\alpha|}.
\]
导数在参数坐标中计算，\(C_\alpha\) 与 \(R,\theta\) 无关。记
\[
 E_\theta g\coloneqq E(\chi_\theta\cdot g).
\]
若 \(B_R=B(c,R)\subset\mathbb R^d\)，令
\[
 w_{B_R}(x)=\left(1+\frac{|x-c|}{R}\right)^{-10d},\quad
 \|F\|_{L^{p_c}(w_{B_R})}
 =\left(\int_{\mathbb R^d}|F(x)|^{p_c}
          w_{B_R}(x)\,\mathrm dx\right)^{1/p_c}.
\]
则对每个 \(\varepsilon>0\)，存在 \(C_\varepsilon<\infty\)，使
\[
 \|Eg\|_{L^{p_c}(B_R)}
 \leq C_\varepsilon R^\varepsilon
       \left(\sum_\theta
          \|E_\theta g\|_{L^{p_c}(w_{B_R})}^2
       \right)^{1/2}.
\]
常数可以依赖 \(d,\varepsilon\) 及固定分割的重叠、导数界，与 \(R,c,g\) 无关。
\end{theorem}
\begin{proofsketch}
频率薄邻域中的全空间解耦经空间局部化得到上述延拓形式；固定的光滑帽分割用于定义各项，有限重叠只改变常数。完整证明结合多线性限制、抛物重标度和反复尺度迭代，最后保留任意小幂次损失。\Cref{b2:lem:25-scale-iteration}说明了递推中的幂次平衡；定理所需的具体多尺度构造见下述原论文。

正式来源为 Bourgain 与 Demeter 的《The Proof of the \(\ell^2\) Decoupling Conjecture》%
\cite[Theorem 1.1; equation (7), pp.~352--355]{BourgainDemeter2015Proof}。
该版本要求正定第二基本形式，临界指数如上。标准抛物面满足这些条件；不能把同一结论未经修改套到第二基本形式不定或退化的曲面。
\end{proofsketch}

定义中的 \(E_\theta g\) 始终使用 \(\chi_\theta\cdot g\)，所以
\(Eg=\sum_\theta E_\theta g\) 精确成立。若改用互不相交的可测帽示性函数，应从头采用那种分割的版本；不能既把帽当作重叠覆盖，又让它们的示性函数相加充当一。权 \(w_{B_R}\) 在球内有统一正下界，在球外保留快速衰减尾部，这与波包没有严格管支集相适应。换成双边可比的权只改变常数，不能任意删去球外部分而保持同一陈述。

\ProseParagraphBreak
还应区分两个容易混写的右端。对有限个函数 \(F_\theta\) 及任意正测度 \(\mu\)，当 \(2\leq p<\infty\) 且各项属于 \(L^p(\mu)\) 时，Minkowski 不等式给出
\[
 \begin{aligned}
 \left\|\left(\sum_\theta|F_\theta|^2\right)^{1/2}\right\|_{L^p(\mu)}^2
 &=\left\|\sum_\theta|F_\theta|^2\right\|_{L^{p/2}(\mu)}\\
 &\leq\sum_\theta
       \bigl\||F_\theta|^2\bigr\|_{L^{p/2}(\mu)}
 =\sum_\theta\|F_\theta\|_{L^p(\mu)}^2.
 \end{aligned}
\]
因此点态平方函数的范数至多为各函数范数的 \(\ell^2\) 和。把解耦定理右端换成前者，会提出一个通常更强的估计，不能由显示的解耦不等式直接推出。

\ProseParagraphBreak
一个有限模型可把差别算出。取 \(\mathbb R^d\) 中 \(N\) 个互不相交、体积为一的方块 \(Q_1,\ldots,Q_N\)，令 \(F_j=\mathbf1_{Q_j}\)。则
\[
 \left\|\left(\sum_{j=1}^N|F_j|^2\right)^{1/2}\right\|_p=N^{1/p},
 \quad
 \left(\sum_{j=1}^N\|F_j\|_p^2\right)^{1/2}=N^{1/2}.
\]
在临界 \(p=p_c\) 处，两者的比值为 \(N^{1/(d+1)}\)。这个例子只比较两种范数结构，没有把方块示性函数宣称为曲面帽的延拓。若各个 \(F_j\) 完全相同，两者又相等；差别取决于各项在空间中的分布，不能用一句“正交”概括。

\ProseParagraphBreak
在 \(p=2\) 时，上述两种范数由积分的线性性相等；非平方指数处则存在实际差距。抛物面临界解耦的内容是以 \(R^\varepsilon\) 控制许多频率方向的叠加，并保留加权球和指定帽尺度。单帽的系数能量、平面穿越计算及尺度递推引理各解决其中一部分问题，不能单独代替这项多帽估计。


\begin{chaptersummary}
多线性非退化必须连同投影族解释。增加一个输入投影，可能使原来有相互作用的相位变得可吸收；输入尺度与投影子空间也会限制可能的指数。横截性则把独立方向转成可量化的行列式下界，平直模型已经能够显示这种增益，因此它与单张曲面的曲率不是同一条件。

\ProseParagraphBreak

曲线的频率和映射把双线性估计转成换元问题：行列式控制局部密度，逆像重数控制不同来源的叠加。允许同片相互作用时，对角退化可能破坏线性推论；输入频带互不相交，也不能保证它们的和集合互不相交。每次使用近似正交，都应重新检查乘积输出的频率支集。

\ProseParagraphBreak

一般曲面的多线性限制与解耦需要更深的几何和迭代。输出指数小于一时，应先在准范数的幂次层面求和，再核对项数造成的损失。波包的平方系数、管外衰减和帽之间的组合是不同步骤；有限个模型的计算能够解释这些步骤的尺度，但不能独自完成一般定理。局部半径、权重、横截常数和任意小幂次损失都属于结论本身。
\end{chaptersummary}

\BookExercises

% 自拟；中等：三个投影下可吸收相位的充要微分判据，明确额外输入如何改变退化性。
\Needspace{6\baselineskip}
\begin{optionalexercise}\label{b2:ex:25-new-1}
设 \(I,J\subset\mathbb R\) 为非空开区间，\(P\in C^3(I\times J;\mathbb R)\)，考虑投影 \(x,y,x+y\)。证明存在实函数
\(A\in C^3(I)\)、\(B\in C^3(J)\)、\(C\in C^3(I+J)\)，使
\[
 P(x,y)=A(x)+B(y)+C(x+y),
\]
当且仅当
\[
 \partial_x\partial_y(\partial_x-\partial_y)P=0.
\]
用此判据比较 \(xy\)、\(x^2y+xy^2\) 和 \(x^2y-xy^2\)。最后说明：对非零非负 \(a\in C_c^\infty(I\times J)\)，若相位可如此分解，则
\[
 \int e^{2\pi\mathrm i\lambda P(x,y)}a(x,y)
            f_1(x)f_2(y)f_3(x+y)\,\mathrm dx\,\mathrm dy
\]
不可能对所有输入具有统一的正幂次 \(\lambda\) 衰减，即使原相位对只有 \(x,y\) 的两个投影不能吸收。
\end{optionalexercise}
\begin{solution}
若分解成立，则 \(\partial_x\partial_yP=C''(x+y)\)，再作用 \(\partial_x-\partial_y\) 就为零。反过来，令 \(h=P_{xy}\)。假设说明 \(h\) 沿方向 \((1,-1)\) 的导数为零。矩形与每条直线 \(x+y=t\) 的交集是连通线段，因此 \(h\) 在这条线段上为常数，可以写成 \(h_0(t)\)。用局部横截线取值可知 \(h_0\in C^1(I+J)\)；各局部定义因沿线段恒定而一致。

在区间 \(I+J\) 上取二次原函数 \(C\)，使 \(C''=h_0\)，便有 \(C\in C^3\)。令 \(Q=P-C(x+y)\)，则 \(Q_{xy}=0\)。固定 \((x_0,y_0)\in I\times J\)，两次使用微积分基本定理给出
\[
 Q(x,y)=Q(x,y_0)+Q(x_0,y)-Q(x_0,y_0).
\]
右侧即为所需的 \(A(x)+B(y)\)，完成全矩形上的分解。

前两个例子有显式表示
\[
 xy=\frac{(x+y)^2-x^2-y^2}{2},
 \quad x^2y+xy^2=\frac{(x+y)^3-x^3-y^3}{3}.
\]
而对 \(P=x^2y-xy^2\)，题中的三阶微分算子给出常数四，所以不能分解。尤其 \(xy\) 对两个投影时具有混合相互作用，增加 \(x+y\) 投影后却成为可吸收相位。

最后在三个投影像上选固定紧支撑截断，均在 \(\operatorname{supp}a\) 的相应投影上为一。将它们分别乘以
\(e^{-2\pi\mathrm i\lambda A}\)、\(e^{-2\pi\mathrm i\lambda B}\)、\(e^{-2\pi\mathrm i\lambda C}\)，得到三个输入。各个固定 \(L^p\) 范数不随参数改变，相位却完全抵消，积分恒为正数 \(\int a\)。这排除了任何趋于零的统一衰减界。第三个例子不满足分解条件，其衰减问题因而需要其他估计。
\end{solution}

% 自拟；综合：三维Loomis–Whitney的归一化等号分类，并由缺陷控制一个平方根距离。
\Needspace{6\baselineskip}
\begin{optionalexercise}\label{b2:ex:25-new-2}
设 \(F_1,F_2,F_3\geq0\) 为 \(\mathbb R^2\) 上可积函数，各自积分为一，令
\[
 \mathcal I=\int_{\mathbb R^3}
       [F_1(y,z)F_2(x,z)F_3(x,y)]^{1/2}\,\mathrm dx\,\mathrm dy\,\mathrm dz.
\]
证明 \(\mathcal I=1\) 当且仅当存在 \(\mathbb R\) 上的概率密度 \(a,b,c\)，使几乎处处有
\[
 F_1(y,z)=b(y)c(z),
 \quad F_2(x,z)=a(x)c(z),
 \quad F_3(x,y)=a(x)b(y).
\]
再令 \(a(x)=\int F_2(x,z)\,\mathrm dz\)、\(b(y)=\int F_1(y,z)\,\mathrm dz\)，证明对一般输入都有
\[
 \left\|\sqrt{F_3}-\sqrt{a(x)b(y)}\right\|_{L^2(\mathbb R^2)}^2
       \leq2(1-\mathcal I).
\]
零边际处的函数值按几乎处处等价类处理，不假设密度严格为正。
\end{optionalexercise}
\begin{solution}
在 \(b(y)>0\)、\(a(x)>0\) 的地方分别置
\(p_y(z)=F_1(y,z)/b(y)\)、\(q_x(z)=F_2(x,z)/a(x)\)。它们是概率密度；其余参数处可以任选一个固定概率密度，因为相应权重为零。非负性与 Tonelli 保证零边际对应的原函数也几乎处处为零。记
\[
 A(x,y)=\int\sqrt{p_y(z)q_x(z)}\,\mathrm dz\leq1.
\]
先沿 \(z\) 积分，随后对 \((x,y)\) 用 Cauchy--Schwarz，得到
\[
 \mathcal I=\int\sqrt{F_3(x,y)a(x)b(y)}\cdot A(x,y)\,\mathrm dx\,\mathrm dy
 \leq\mathcal J\leq1,
\]
其中 \(\mathcal J=\int\sqrt{F_3ab}\)。最后一步中两个平方根的平方积分都等于一。

若 \(\mathcal I=1\)，则两步均取等。第二步的取等条件给出 \(F_3=ab\) 几乎处处；第一步于是使 \(A(x,y)=1\) 对 \(a(x)b(y)\,\mathrm dx\,\mathrm dy\) 几乎处处成立。概率密度的 Cauchy--Schwarz 取等说明 \(p_y=q_x\) 几乎处处于 \(z\)。由 Fubini 可固定一个适当的 \(x_0\)，使对 \(b(y)\,\mathrm dy\) 几乎所有 \(y\)，均有 \(p_y=q_{x_0}\)。令 \(c=q_{x_0}\)，再固定一个适当的 \(y\)，得到 \(q_x=c\) 对 \(a(x)\,\mathrm dx\) 几乎所有 \(x\) 成立。因而三个密度恰有题述形式，零边际处也由前面的零函数说明覆盖。反过来，代入乘积形式后，被积函数为 \(a(x)b(y)c(z)\)，Fubini 给出积分一。

最后，无需取等也有
\[
 \left\|\sqrt{F_3}-\sqrt{ab}\right\|_2^2
       =2-2\mathcal J\leq2-2\mathcal I.
\]
这条定量结论把接近等号的程度转化为该边际乘积距离的控制。
\end{solution}

% 自拟；基础到中等：中心横截量向整片传播的定量尺度，并构造两两分离却以该尺度退化的法向族。
\Needspace{6\baselineskip}
\begin{optionalexercise}\label{b2:ex:25-new-3}
设整数 \(d\geq2\)，\(n_j,\widetilde n_j\in S^{d-1}\)。证明
\[
 \left|\det(\widetilde n_1,\ldots,\widetilde n_d)
             -\det(n_1,\ldots,n_d)\right|
       \leq\sum_{j=1}^d|\widetilde n_j-n_j|.
\]
由此给出一个充分条件，把帽中心处的横截下界 \(\nu>0\) 传到整张帽：若第 \(j\) 个单位法向映射为 Lipschitz、常数 \(L_j\)，帽参数距中心至多为 \(r\)，求确保下界仍至少为 \(\nu/2\) 的条件。

\ProseParagraphBreak

再用三维法向
\[
 n_1=e_1,\quad n_2=e_2,\quad
 n_3=\frac{(1,1,\varepsilon)}{\sqrt{2+\varepsilon^2}},
 \quad0<\varepsilon\leq1,
\]
证明即使两两夹角的正弦有统一正下界，为保证完整横截，所允许的法向扰动仍须记录 \(\nu\) 的量级。
\end{optionalexercise}
\begin{solution}
逐列把 \(n_j\) 换成 \(\widetilde n_j\)，写出行列式差的望远镜和。第 \(j\) 项只有一列是 \(\widetilde n_j-n_j\)，其余列都是单位向量；由行列式绝对值不超过列长之积，该项至多为 \(|\widetilde n_j-n_j|\)。求和即得第一式。

若中心处的行列式绝对值至少为 \(\nu\)，各帽上法向变化至多为 \(rL_j\)，三角不等式给出任意取点的下界
\[
 |\det(n_1(u_1),\ldots,n_d(u_d))|
       \geq\nu-r\sum_{j=1}^dL_j.
\]
因此 \(r\sum_jL_j\leq\nu/2\) 足够。若只知各法向扰动不超过 \(h\)，则 \(h\leq\nu/(2d)\) 足够。

给定三维模型的行列式为
\(\nu_\varepsilon=\varepsilon/\sqrt{2+\varepsilon^2}\)。前两个法向正交，第三个与它们的夹角正弦均至少为 \(1/\sqrt2\)。然而把第三个换成
\(n_3^0=(1,1,0)/\sqrt2\)，完整行列式立刻变为零，两两角度却仍保持这一正下界。其变化量满足
\[
 |n_3-n_3^0|^2
 =2\left(1-\sqrt{1-\nu_\varepsilon^2}\right)
 =\frac{2\nu_\varepsilon^2}{1+\sqrt{1-\nu_\varepsilon^2}},
\]
因而介于 \(\nu_\varepsilon^2\) 与 \(2\nu_\varepsilon^2\) 之间。这展示了量级为 \(\nu\) 的扰动就能破坏横截性。若希望把它实现为帽上的变化，可取
\(n_3(u)=(1,1,\varepsilon-u)/\sqrt{2+(\varepsilon-u)^2}\)、\(0\leq u\leq\varepsilon\)；其 Lipschitz 常数一致有界，却在距中心 \(O(\nu_\varepsilon)\) 处退化。因此不能只用两两角度选取一个与完整横截量无关的帽尺度。
\end{solution}

% 自拟；中等：分离抛物线片的最好常数、全部极值化序列的端点集中刻画及弱极限。
\Needspace{6\baselineskip}
\begin{optionalexercise}\label{b2:ex:25-new-4}
令 \(I=[0,1]\)、\(J=[2,3]\)，采用抛物线延拓
\(E_If(x,y)=\int_I e^{2\pi\mathrm i(xt+yt^2)}f(t)\,\mathrm dt\)。证明双线性 \(L^2\) 界的最佳常数为 \(1/\sqrt2\)，但没有非零取等输入。对单位范数输入，定义缺陷
\[
 \Delta(f,g)=\frac12-\|E_If\cdot E_Jg\|_2^2.
\]
把到两端点的距离的一阶矩记为
\[
 \begin{aligned}
 m_I(f)&\coloneqq\int_I(1-t)|f(t)|^2\,\mathrm dt,\\
 m_J(g)&\coloneqq\int_J(s-2)|g(s)|^2\,\mathrm ds.
 \end{aligned}
\]
证明
\[
 \frac16\bigl(m_I(f)+m_J(g)\bigr)
 \leq\Delta(f,g)
 \leq\frac12\bigl(m_I(f)+m_J(g)\bigr).
\]
据此证明：一列单位输入极值化，当且仅当其平方质量分别集中到 \(t=1\)、\(s=2\)。把这里的“集中”写成准确的积分条件，并说明这样的输入列在各自 \(L^2\) 中都弱趋于零。
\end{optionalexercise}
\begin{solution}
由\cref{b2:thm:25-parabola-bilinear}，乘积平方范数为
\(\int_{I\times J}|f(t)|^2\cdot|g(s)|^2/[2(s-t)]\,\mathrm dt\,\mathrm ds\)。因 \(s-t\geq1\)，得到常数 \(1/\sqrt2\)。取
\(f_\eta=\eta^{-1/2}\mathbf1_{[1-\eta,1]}\)、
\(g_\eta=\eta^{-1/2}\mathbf1_{[2,2+\eta]}\)，其平方范数都为一，且输出平方范数介于 \(1/[2(1+2\eta)]\) 与 \(1/2\) 之间；令 \(\eta\downarrow0\) 就证明最佳性。取等要求乘积密度集中在 \(s-t=1\)，该集合在矩形中测度为零，故非零输入不能取等。

对单位输入，把 \(|f|^2\)、\(|g|^2\) 看成概率密度。置 \(a=1-t\)、\(b=s-2\)，则 \(0\leq a+b\leq2\)，并有
\[
 \frac12-\frac1{2(s-t)}=\frac{a+b}{2(1+a+b)}.
\]
这个量介于 \((a+b)/6\) 与 \((a+b)/2\) 之间；对乘积概率测度积分，即得两个缺陷界。

因此 \(\Delta(f_n,g_n)\to0\) 当且仅当题中的两个非负一阶矩都趋零。它又等价于对每个 \(0<r<1\)，
\[
 \int_0^{1-r}|f_n(t)|^2\,\mathrm dt\longrightarrow0,
 \quad \int_{2+r}^3|g_n(s)|^2\,\mathrm ds\longrightarrow0.
\]
一方向由矩至少为 \(r\) 乘远端质量得到；反方向把矩分成距端点小于 \(r\) 的部分和其余部分，前者至多为 \(r\)，后者趋零，再令 \(r\downarrow0\)。

对任意 \(h\in L^2(I)\)，在端点邻域上的配对至多为该邻域内的 \(\|h\|_2\)，在其外部则由 \(f_n\) 的平方质量趋零控制。先固定小邻域令 \(n\to\infty\)，再缩小邻域，得到 \(\langle f_n,h\rangle\to0\)。\(g_n\) 同理。因此极值化序列弱趋零，且因范数始终为一，不可能有强收敛子列。
\end{solution}

% 自拟；综合：同片频率和的双重逆像公式，并用平移连续性处理任意非零L²输入的对角发散。
\Needspace{6\baselineskip}
\begin{optionalexercise}\label{b2:ex:25-new-5}
设 \(I\subset\mathbb R\) 为长度为正的紧区间，\(f,g\in L^2(I)\)，令
\[
 E_If(x,y)=\int_I e^{2\pi\mathrm i(xt+yt^2)}f(t)\,\mathrm dt.
\]
证明下式成立，允许两端同时为 \(+\infty\)：
\[
 \|E_If\cdot E_Ig\|_{L^2(\mathbb R^2)}^2
 =\int_{\substack{t,s\in I\\t<s}}
       \frac{|f(t)g(s)+f(s)g(t)|^2}{2(s-t)}\,\mathrm dt\,\mathrm ds.
\]
这里频率和映射在有序三角形上为一一映射；说明为什么不能直接沿用两片分离时的正 Jacobi 行列式下界。进一步证明，对每个非零 \(f\in L^2(I)\)，都有
\[
 \|E_If\|_{L^4(\mathbb R^2)}=\infty.
\]
不得事先假定频率密度属于 \(L^2\)，也不得把最后一问限制为光滑或处处为正的输入。
\end{optionalexercise}
\begin{solution}
记有序三角形为 \(D=\{(t,s)\in I^2:t<s\}\)，边界按零测集处理。映射
\[
 \Psi(t,s)=(t+s,t^2+s^2)
\]
在 \(D\) 上一一：由两个分量可确定 \(t+s\) 及
\(ts=((t+s)^2-(t^2+s^2))/2\)，再按 \(t<s\) 排列二次方程的两个根。其 Jacobi 行列式为 \(2(s-t)>0\)，但在靠近对角线时没有统一正下界。

因为有限区间上的 \(L^2\) 输入也属于 \(L^1\)，将乘积积分分成 \(t<s\)、\(s<t\) 两部分并交换后者的变量，得到
\[
 E_If\cdot E_Ig
 =\int_D e^{2\pi\mathrm i[(t+s)x+(t^2+s^2)y]}
       [f(t)g(s)+f(s)g(t)]\,\mathrm dt\,\mathrm ds.
\]
对角线在乘积 Lebesgue 测度下为零。换元定义频率密度
\[
 h(\Psi(t,s))
 =\frac{f(t)g(s)+f(s)g(t)}{2(s-t)}
 \quad((t,s)\in D),
\]
并在像集外置零。换元先给出
\[
 \|h\|_1
 \leq\int_D\bigl(|f(t)|\cdot|g(s)|+|f(s)|\cdot|g(t)|\bigr)
          \,\mathrm dt\,\mathrm ds
 \leq\|f\|_1\cdot\|g\|_1.
\]
故 \(h\in L^1(\mathbb R^2)\)，且乘积作为连续函数及缓增分布均等于 \(\check h\)。再次按非负积分换元可知，题中右端正是 \(\int|h|^2\)，此时它可能无穷。

若 \(h\in L^2\)，Plancherel 给出所需等式。反过来，若 \(\check h\in L^2\)，其 \(L^2\) Fourier 变换也在 \(L^2\) 中，并且作为缓增分布等于 \(h\)。分布唯一性说明二者几乎处处相等，因而 \(h\in L^2\)。这证明一端有限当且仅当另一端有限，也就建立了允许无穷值的公式。

取 \(g=f\)，公式化为
\[
 \|E_If\|_4^4
 =2\int_D\frac{|f(t)|^2\cdot|f(s)|^2}{s-t}\,\mathrm dt\,\mathrm ds.
\]
将 \(f\) 延拓为区间外的零函数。取任意 \(M>0\)，令
\(\rho(t)=\min\{|f(t)|^2,M\}\)。若 \(f\ne0\)，则 \(\rho\) 非零、有紧支集且属于 \(L^2(\mathbb R)\)。由 \(L^2\) 平移连续性，
\[
 \int_{\mathbb R}\rho(t)\rho(t+h)\,\mathrm dt
       \longrightarrow\|\rho\|_2^2>0
       \quad(h\downarrow0).
\]
所以存在 \(\eta>0\)，使上述积分在 \(0<h<\eta\) 时至少为
\(\|\rho\|_2^2/2\)。Tonelli 与 \(s=t+h\) 给出
\[
 \int_D\frac{|f(t)|^2\cdot|f(s)|^2}{s-t}\,\mathrm dt\,\mathrm ds
 \geq\int_0^\eta\frac1h
       \left(\int_{\mathbb R}\rho(t)\rho(t+h)\,\mathrm dt\right)\mathrm dh
 =\infty.
\]
这里零延拓自动排除了区间外的点，故这一下界对原积分有效。这也说明，同片相互作用中的对角退化确实破坏了全空间的线性 \(L^2\) 到 \(L^4\) 推论。
\end{solution}

% 自拟；中等：准范数求和相对于一次和与平方和的最佳项数损失，以及固定区域上的尺度反例。
\Needspace{6\baselineskip}
\begin{optionalexercise}\label{b2:ex:25-new-6}
设 \(0<q<1\)，\(M\geq1\) 为整数，\(h_1,\ldots,h_M\in L^q\)。证明
\[
 \left\|\sum_{\ell=1}^M h_\ell\right\|_q
 \leq M^{1/q-1}\sum_{\ell=1}^M\|h_\ell\|_q,
\]
以及
\[
 \left\|\sum_{\ell=1}^M h_\ell\right\|_q
 \leq M^{1/q-1/2}
       \left(\sum_{\ell=1}^M\|h_\ell\|_q^2\right)^{1/2}.
\]
证明两个关于 \(M\) 的常数都是最佳的。对多线性输出指数
\(q=2/(d-1)\)、\(d\geq4\)，若项数 \(M(R)\asymp R^\kappa\)、\(\kappa>0\)，说明即使各项具有统一平方预算
\(\sum_\ell\|h_\ell\|_q^2\leq1\)，求和仍可能出现怎样的幂次损失。要求在固定区间 \((0,1)\) 上实现这个例子，并具体计算 \(d=4\) 时的两个损失。
\end{optionalexercise}
\begin{solution}
点态不等式 \(|\sum_\ell h_\ell|^q\leq\sum_\ell|h_\ell|^q\) 给出
\[
 \left\|\sum_\ell h_\ell\right\|_q
       \leq\left(\sum_\ell a_\ell^q\right)^{1/q},
 \quad a_\ell=\|h_\ell\|_q.
\]
有限序列的均值不等式分别给出
\[
 \left(\sum_\ell a_\ell^q\right)^{1/q}
       \leq M^{1/q-1}\sum_\ell a_\ell,
 \quad
 \left(\sum_\ell a_\ell^q\right)^{1/q}
       \leq M^{1/q-1/2}
             \left(\sum_\ell a_\ell^2\right)^{1/2}.
\]
例如可分别对凹函数 \(t^q\)、\(t^{q/2}\) 使用 Jensen 不等式得到。若各 \(h_\ell\) 支集互不相交且具有相同非零准范数，幂次求和及两个均值不等式同时取等，所以两个常数都不能减小。

具体地，把 \((0,1)\) 分成
\(E_\ell=((\ell-1)/M,\ell/M)\)，并取
\[
 h_\ell=M^{1/q-1/2}\mathbf1_{E_\ell}.
\]
每项的 \(L^q\) 准范数为 \(M^{-1/2}\)，因而平方预算恰为一；总和则在 \((0,1)\) 上恒为 \(M^{1/q-1/2}\)，其准范数也等于这个数。于是当 \(M(R)\asymp R^\kappa\) 时，可能的损失达到
\[
 R^{\kappa(1/q-1/2)}
       =R^{\kappa(d-2)/2}
\]
的量级。这不能被任意小的 \(R^\varepsilon\) 吸收：只要
\(0<\varepsilon<\kappa(d-2)/2\)，二者的比值就无界。若 \(d=4\)，则 \(q=2/3\)，相对于一次和、平方和的最佳损失分别为 \(M^{1/2}\)、\(M\)。准范数的幂次次可加性并不意味着在其他组合方式下求和没有代价。
\end{solution}

% 自拟；中等：等差频带的加法重数精确求和，给出乘积输出重叠因子不可去除的定量例子。
\Needspace{6\baselineskip}
\begin{optionalexercise}\label{b2:ex:25-new-7}
取非零 \(\phi\in\mathcal S(\mathbb R)\)，满足
\(\operatorname{supp}\widehat\phi\subset[-1/10,1/10]\)，令
\[
 u_j(x)=e^{2\pi\mathrm i jx}\phi(x),
 \quad1\leq j\leq N.
\]
说明这些输入的 Fourier 支集两两不交，并计算
\[
 \left\|\left(\sum_{j=1}^N u_j\right)^2\right\|_2^2
 \quad\text{与}\quad
 \sum_{j,k=1}^N\|u_ju_k\|_2^2.
\]
由精确比值证明，输入频带分离不能推出所有乘积项之间具有与 \(N\) 无关的平方正交界，并说明本例中输出频率重叠所需的因子量级。
\end{optionalexercise}
\begin{solution}
调制把 \(\widehat\phi\) 平移到区间
\([j-1/10,j+1/10]\)，这些区间两两不交。但乘积的调制频率只有 \(j+k\)，所以
\[
 \left(\sum_{j=1}^N u_j(x)\right)^2
       =\phi(x)^2\sum_{s=2}^{2N}r_N(s)e^{2\pi\mathrm i sx},
\]
其中
\[
 r_N(s)=
 \begin{cases}
 s-1,&2\leq s\leq N+1,\\
 2N+1-s,&N+1<s\leq2N.
 \end{cases}
\]
又因 \(\widehat{\phi^2}=\widehat\phi*\widehat\phi\) 支持在
\([-1/5,1/5]\)，不同 \(s\) 对应的整个频率组彼此正交。Plancherel 因而给出
\[
 \left\|\left(\sum_{j=1}^N u_j\right)^2\right\|_2^2
 =\left(N^2+2\sum_{r=1}^{N-1}r^2\right)\|\phi\|_4^4
 =\frac{2N^3+N}{3}\|\phi\|_4^4.
\]
另一方面，每个 \(u_ju_k\) 的绝对值均为 \(|\phi|^2\)，故
\[
 \sum_{j,k=1}^N\|u_ju_k\|_2^2=N^2\|\phi\|_4^4.
\]
两个平方量的比值为
\[
 \frac{2N^2+1}{3N},
\]
随 \(N\) 线性增长。因此，若要用右侧乘积项的平方和控制左侧平方范数，常数不可能独立于 \(N\)。在中心频率 \(s=N+1\) 处，共有 \(N\) 个有序乘积具有相同 Fourier 支集，最大重叠数正是 \(N\)；有限重叠估计中的平方范数因子 \(N\) 在量级上确实不能去掉。问题发生在相同和频率组内部的相干叠加，输入之间的正交并未消除它。
\end{solution}

% 自拟；中等：平方函数与范数平方和的锐利反向比较，并以严格频带分离的Schwartz函数实现最大差距。
\Needspace{6\baselineskip}
\begin{optionalexercise}\label{b2:ex:25-new-8}
设 \(2\leq p<\infty\)，\(u_1,\ldots,u_N\in L^p\)，记
\[
 S_p=\left\|\left(\sum_{j=1}^N|u_j|^2\right)^{1/2}\right\|_p,
 \quad A_p=\left(\sum_{j=1}^N\|u_j\|_p^2\right)^{1/2}.
\]
证明锐利的双边比较
\[
 S_p\leq A_p\leq N^{1/2-1/p}S_p.
\]
进一步证明，即使各 \(u_j\) 为 Schwartz 函数且 Fourier 支集紧致、两两不交，右侧因子也不能缩小。可使用上一题的 \(\phi\)，比较
\[
 u_j^{(L)}(x)=e^{2\pi\mathrm i jx}\phi(x-jL)
\]
在 \(L=0\) 及 \(L\to\infty\) 时的行为。说明这一比较对一般平方函数推理与解耦右侧的关系意味着什么。
\end{optionalexercise}
\begin{solution}
在 \(L^{p/2}\) 中用三角不等式，得到
\[
 S_p^2
 =\left\|\sum_j|u_j|^2\right\|_{p/2}
 \leq\sum_j\||u_j|^2\|_{p/2}=A_p^2.
\]
再对有限序列 \((\|u_j\|_p)_j\) 用 \(\ell^p\) 与 \(\ell^2\) 的比较，并在每一点用同样的幂次关系，得到
\[
 A_p
 \leq N^{1/2-1/p}\left(\sum_j\|u_j\|_p^p\right)^{1/p}
 \leq N^{1/2-1/p}
       \left(\int\left(\sum_j|u_j|^2\right)^{p/2}\right)^{1/p}.
\]
这就是第二个界。特别地，\(p=2\) 时两个量总相等。

现在固定 \(N\)，取题述 \(u_j^{(L)}\)。平移只给 Fourier 变换乘以模为一的因子，调制则把支集移到
\([j-1/10,j+1/10]\)，所以所有 \(L\) 下都保持 Fourier 支集两两不交。范数不受平移或调制影响，故
\[
 A_p=\sqrt N\,\|\phi\|_p.
\]
当 \(L=0\) 时，各函数具有相同绝对值，所以 \(S_p=A_p\)，第一边的常数一已经最佳。

当 \(L\to\infty\) 时，证明
\[
 S_p(u_1^{(L)},\ldots,u_N^{(L)})
       \longrightarrow N^{1/p}\|\phi\|_p.
\]
为此取截断 \(\phi_B=\phi\mathbf1_{[-B,B]}\)。若 \(L>2B\)，不同平移后的截断支集互不相交，相应平方函数范数恰为
\(N^{1/p}\|\phi_B\|_p\)。向量值 \(L^p(\ell^2)\) 范数的三角不等式及已经证明的第一边给出截断误差至多为
\[
 \left(\sum_{j=1}^N\|\phi-\phi_B\|_p^2\right)^{1/2}
       =\sqrt N\,\|\phi-\phi_B\|_p.
\]
先令 \(L\to\infty\)，再令 \(B\to\infty\)，即得所称极限。因此 \(A_p/S_p\) 趋于
\(N^{1/2-1/p}\)，第二边常数也最佳，而且这一最佳性已在所要求的 Schwartz 频带分离类中实现。

当 \(p>2\) 时，同样的频带分离既可对应 \(S_p=A_p\)，也可使两者之比达到随项数增长的最大量级。因此，只凭一般范数关系或频带正交，不能把解耦估计右侧的 \(A_p\) 无损换成较小的 \(S_p\)。曲面延拓的专门平方函数估计则还会利用额外的几何结构。
\end{solution}
% 自拟：综合可积核估计、二阶 van der Corput、Plancherel 与插值，并检验精确零均值仍不足以给出强端点衰减。
\Needspace{6\baselineskip}
\begin{optionalexercise}[振荡卷积的内部指数与端点障碍]\label{b2:ex:25-oscillatory-convolution-endpoints}
在 \(\mathbb R\) 上取复函数空间。固定非零实函数 \(a\in C_c^\infty(\mathbb R)\)、\(a\geq0\)，对 \(\lambda\geq1\) 令
\[
 k_\lambda(x)=a(x)e^{2\pi\mathrm i\lambda x^2},\quad
 T_\lambda f=k_\lambda*f,
 \quad
 \beta_p=\frac12-\left|\frac1p-\frac12\right|\quad(1\leq p\leq\infty),
\]
其中 \(1/\infty=0\)，Fourier 变换使用本卷 \(2\pi\) 约定。
\begin{enumerate}
\item\label{b2:item:25-oscillatory-interior}证明对每个 \(1\leq p\leq\infty\)，存在不依赖 \(\lambda,f\) 的常数 \(C_{a,p}\)，使
\[
 \|T_\lambda f\|_p\leq C_{a,p}\lambda^{-\beta_p}\|f\|_p.
\]
特别说明 van der Corput 估计的常数为何可对全部输出频率统一选取。
\item\label{b2:item:25-oscillatory-endpoints}证明端点算子范数实际上满足
\[
 \|T_\lambda\|_{L^1\to L^1}
 =\|T_\lambda\|_{L^\infty\to L^\infty}=\|a\|_1,
\]
并解释为何核的振荡不能强迫这两个范数随 \(\lambda\) 衰减。
\item\label{b2:item:25-oscillatory-cancellation}另固定 \(b\in C_c^\infty(\mathbb R)\)、\(b\geq0\)、\(\int b=1\)，令
\[
 c_\lambda=\int_{\mathbb R}k_\lambda(x)\,\mathrm dx,\quad
 \widetilde k_\lambda=k_\lambda-c_\lambda b,\quad
 \widetilde T_\lambda f=\widetilde k_\lambda*f.
\]
证明 \(\int\widetilde k_\lambda=0\)，且 \(\widetilde T_\lambda\) 仍满足第一问的指数估计（常数可依赖 \(a,b,p\)），但两个端点范数均不趋于零。
\end{enumerate}
\end{optionalexercise}
\begin{solution}
Young 不等式\cref{b2:thm:03-young}先给出所有指数上的统一界 \(\|T_\lambda\|_{p\to p}\leq\|a\|_1\)。另一方面，
\[
 \widehat k_\lambda(\xi)
 =\int e^{2\pi\mathrm i\lambda(x^2-(\xi/\lambda)x)}a(x)\,\mathrm dx.
\]
选一个固定有限区间包含 \(\operatorname{supp}a\)。其中的相位二阶导数恒为 \(2\)，与 \(\xi,\lambda\) 都无关；\cref{b2:thm:21-vdc}的二阶版本只用这个下界及振幅的 \(\|a\|_\infty+\|a'\|_1\)，故
\[
 \sup_{\xi\in\mathbb R}|\widehat k_\lambda(\xi)|\leq C_a\lambda^{-1/2}.
\]
先对 Schwartz 输入用卷积定理与 Plancherel，再由 \(L^2\) 稠密性延拓，得 \(\|T_\lambda\|_{2\to2}\leq C_a\lambda^{-1/2}\)。将它分别与 \(L^1\)、\(L^\infty\) 的统一界作\cref{b2:thm:17-riesz-thorin}中的插值，所得衰减指数在 \(1\leq p\leq2\) 时为 \(1-1/p\)，在 \(2\leq p\leq\infty\) 时为 \(1/p\)，合起来就是 \(\beta_p\)。

为核对端点，先证明任意 \(k\in L^1(\mathbb R)\) 的卷积算子 \(S_k\) 都有
\(\|S_k\|_{1\to1}=\|S_k\|_{\infty\to\infty}=\|k\|_1\)。上界是 Young 不等式。在 \(L^1\) 端取非负光滑逼近恒等 \(\rho_\varepsilon\)，其积分为 \(1\)；由 \(k*\rho_\varepsilon\to k\) 于 \(L^1\)，得到算子范数至少为 \(\|k\|_1\)。在 \(L^\infty\) 端，取
\[
 F(y)=\overline{\operatorname{phase}(k(-y))},\quad
 \operatorname{phase}(z)=z/|z|\ (z\ne0),\quad \operatorname{phase}(0)=0.
\]
它的本性上确界至多为 \(1\)，而 \((k*F)(0)=\int|k|\)。此外 \(k*F\) 有连续代表，因为它的平移差一致受 \(\|\tau_hk-k\|_1\cdot\|F\|_\infty\) 控制并趋零，故该点值也给出本性上确界的下界。应用于 \(k_\lambda\) 即得所求等式。更具体地，在这个非负振幅模型中，\(L^\infty\) 输入 \(F_\lambda(y)=e^{-2\pi\mathrm i\lambda y^2}\) 能在输出点 \(0\) 消去核的相位，且输入范数恒为 \(1\)；\(L^1\) 的集中输入则在输出中近似保留整个核的绝对值质量。

最后，\(c_\lambda=\widehat k_\lambda(0)\)，所以 \(|c_\lambda|\leq C_a\lambda^{-1/2}\)，且 \(\int\widetilde k_\lambda=c_\lambda-c_\lambda\int b=0\)。由 \(|\widehat b|\leq\|b\|_1=1\)，有
\[
 \|\widetilde k_\lambda\|_1\leq\|a\|_1+C_a\lambda^{-1/2},\quad
 \|\widehat{\widetilde k_\lambda}\|_\infty
 \leq2C_a\lambda^{-1/2}.
\]
再次用 Young、Plancherel 与同样的两段插值，得到所需的全部指数界。可是反向三角不等式和刚证的端点范数公式给出
\[
 \|\widetilde T_\lambda\|_{1\to1}
 =\|\widetilde T_\lambda\|_{\infty\to\infty}
 =\|\widetilde k_\lambda\|_1
 \geq\|a\|_1-C_a\lambda^{-1/2}.
\]
因 \(a\ne0\)，右端趋向正数 \(\|a\|_1\)。于是精确零均值可以与内部指数上的振荡衰减并存，却不能单独推出强 \(L^1\) 或强 \(L^\infty\) 的参数衰减。
\end{solution}


